\documentclass[11pt,a4paper]{article}

% =============================================================================
% MTDF Master Document: combined 8-paper compilation
% =============================================================================
\usepackage[utf8]{inputenc}
\usepackage[T1]{fontenc}
\usepackage{lmodern}
\usepackage[english]{babel}
\usepackage[a4paper,margin=2.2cm]{geometry}
\usepackage{amsmath,amssymb,amsthm}
\usepackage{graphicx}
\usepackage{float}
\usepackage{booktabs}
\usepackage{array}
\usepackage{longtable}
\usepackage{tabularx}
\usepackage{multirow}
\usepackage{caption}
\usepackage[dvipsnames,table]{xcolor}
\usepackage{mdframed}
\usepackage{parskip}
\usepackage{ragged2e}
\usepackage[hidelinks,breaklinks=true]{hyperref}
\usepackage{microtype}
\usepackage{titlesec}

\definecolor{accent}{HTML}{1A4B8A}
\definecolor{callout-bg}{HTML}{FBFBFB}
\definecolor{tablehead}{HTML}{F0F0F0}

\setlength{\tabcolsep}{4pt}
\renewcommand{\arraystretch}{1.15}

\captionsetup{font=small,labelfont=bf,labelsep=period,justification=centerfirst}

\hypersetup{
  pdftitle={The Mesche Hypothesis (MTDF) -- Complete Compilation},
  pdfauthor={Ingo Mesche},
  pdfsubject={MTDF - Mesche's Tensor Dynamics Framework}
}

\sloppy
\emergencystretch=3em

% Re-map Part titles so each MTDF paper gets a clean divider page.
\titleformat{\part}[display]
  {\normalfont\LARGE\bfseries\color{accent}\centering}
  {Paper \thepart}{20pt}{\Huge}

\title{\textbf{The Mesche Hypothesis}\\[6pt]
  \normalsize\itshape A Dynamic Field Theory of Cosmic Structure\\[2pt]
  \normalsize Complete research compilation of the MTDF paper suite}
\author{Ingo Mesche\\ \small Independent Researcher, Malta \\ \small \texttt{imesche@outlook.com}}
\date{V74 / 2026}

\begin{document}
\maketitle

\vspace{1em}
\begin{center}
\begin{minipage}{0.85\textwidth}
\noindent\textbf{About this compilation.} This document combines all eight papers of the MTDF research programme into a single self-contained volume. Each paper is also available as a standalone submission with the same preamble and uniform margins. Cross-references between papers resolve internally within this master document.
\end{minipage}
\end{center}
\vspace{1em}

\tableofcontents
\clearpage


\part{Master paper: The Mesche Hypothesis}

\begin{abstract}
Mesche's Tensor Dynamics Framework (MTDF) is a field based alternative to the latent sector interpretation of modern cosmology. It aims to reproduce the background expansion history and a wide range of dynamical phenomena without invoking cold dark matter or a cosmological constant as new substances. This master paper presents the theory as a coherent physical proposal, states the minimal parameter set and its provenance, and formalises a strict evaluation protocol that separates calibration anchors from independent validation targets. It then summarises current V18 validation outcomes on the vector likelihood set (Pantheon+ supernovae, DESI Y1 BAO, cosmic chronometer H(z), DR16 f$\sigma$$_{8}$ growth) and reports the CMB distance prior strictly as a diagnostic. A sharp, environment-dependent transition in SNe residuals is detected at z = 0.04 $\pm$ 0.005 across three independent void catalogues, with $\Delta$$\chi$$^{2}$ = 36--39 relative to a single-regime fit (p < 0.001). This scale coincides with the coherence length $\beta$ = 22.7 Mpc predicted independently from void quantisation. Fitting Planck 2018 TTTEEE+lensing with the class\_\allowbreak{}mtdf modified Boltzmann solver yields $\Delta$$\chi$$^{2}$ = +0.63 relative to $\Lambda$CDM, with a broad non-pathological posterior on the MTDF coupling k\textsubscript{f} that admits full coupling (k\textsubscript{f} = 1) within 1$\sigma$; the $\Delta$BIC = +8.05 reflects the added dimensionality, so the framework is not excluded by this dataset under the stated assumptions. The paper emphasises falsifiability: it enumerates near term tests that can confirm or refute MTDF beyond the current dashboard scope. A separate gravity-sector companion paper reports quantitative lensing and local dynamics tests (galaxy-galaxy lensing, strong lens time delays, and Solar System screening) using the same frozen parameter set; see \textit{MTDF\_\allowbreak{}03 (Companion Part II)}~[see companion paper].
\end{abstract}

\begin{mdframed}[linecolor=accent,linewidth=0.5pt,backgroundcolor=callout-bg,leftline=true,rightline=true,topline=true,bottomline=true,innerleftmargin=8pt,innerrightmargin=8pt,innertopmargin=6pt,innerbottommargin=6pt]
\textbf{Claim status.}
This paper distinguishes between:
(i) calibration anchors,
(ii) validation targets,
(iii) diagnostic consistency checks, and
(iv) speculative extensions.
Only results explicitly labelled as validation targets are used as evidence for the core MTDF claim. Diagnostic and speculative sections are included to define future tests, not to claim established confirmation.
\end{mdframed}


\section{Motivation and design principles}
\label{sec:sec1}
Contemporary cosmology describes most observed gravitational phenomena by introducing a dominant but unseen sector. In the standard model of cosmology, these components are parameterised as cold dark matter and dark energy. MTDF proposes an alternative interpretation: the phenomena attributed to the latent sector emerge from the dynamics of a physically defined spacetime field that carries stress energy and couples to baryonic matter through a single coupling parameter.

\begin{mdframed}[linecolor=accent,linewidth=0.5pt,backgroundcolor=callout-bg,leftline=true,rightline=true,topline=true,bottomline=true,innerleftmargin=8pt,innerrightmargin=8pt,innertopmargin=6pt,innerbottommargin=6pt]
\textbf{Design constraints for a high standard theory proposal.} MTDF is formulated under four constraints: (i) minimal free parameters; (ii) explicit parameter provenance; (iii) strict separation of calibration anchors from validation targets; (iv) falsifiability with clear failure modes. The V18 era artefacts operationalise these constraints for empirical evaluation.
\end{mdframed}

It is essential to distinguish between two different claims: (a) \emph{data fit} within a chosen likelihood set, and (b) \emph{physical explanation} that remains coherent across scales. The dashboard is primarily about (a). This master paper addresses (b), while importing the dashboard protocol to ensure that any claimed fit is not achieved by circular calibration.

\section{Core objects and governing assumptions}
\label{sec:sec2}
MTDF treats the gravitational sector as augmented by a field degree of freedom whose stress energy contributes to dynamics. The field is introduced to remain compatible with classical limits while providing a mechanism for scale dependent behaviour. In the non relativistic regime, the field acts as an additional contribution to the effective potential that influences orbital velocities and structure formation, while in the cosmological regime it modifies the effective expansion history through its energy density and stress.

\subsection{Minimal parameter philosophy}
The framework is constrained to a small set of fundamental parameters that are either measured from independent observational anchors or established from theoretically motivated universal constants. The operational intent is that most quantities used in predictions are derived, not tuned.

\subsection{Calibration anchors versus validation targets}
In the V18 strict protocol, some data are permitted to set fundamental scales (anchors). These are excluded from validation scoring. All other reported goodness of fit values are computed strictly on targets that are treated as independent of the calibration step. This is the key correction relative to earlier single layer evaluations.

\subsection{Mathematical objects, stress definition, and energy}
This section expands the core mathematical objects that appear throughout MTDF. The intent is to make the field definitions explicit, while keeping the strict V18 protocol separation between calibration anchors and validation targets.

\begin{mdframed}[linecolor=accent,linewidth=0.5pt,backgroundcolor=callout-bg,leftline=true,rightline=true,topline=true,bottomline=true,innerleftmargin=8pt,innerrightmargin=8pt,innertopmargin=6pt,innerbottommargin=6pt]
\subsubsection{Core idea in one sentence}
We propose that the universe is permeated by dynamic spacetime characterized by the dimensionless strain tensor $S_{\mu\nu}$ that:

\begin{itemize}
  \item \textbf{Stores energy} through elastic deformation with energy density u\textsubscript{tensor} = (E/\allowbreak{}2) Tr($S_{\mu\nu}$ $S_{\mu\nu}$)
  \item \textbf{Exhibits stress tensor gradients} that wrap around massive objects and extend throughout cosmic space
  \item \textbf{Undergoes stress accumulation and relaxation cycles} over characteristic time $\tau$ = 13.0 Gyr
  \item \textbf{Redistributes energy} rather than destroying it, maintaining conservation laws
  \item \textbf{Exhibits bipolar behaviour} with attractive and repulsive components at different scales
\end{itemize}

\begin{mdframed}[linecolor=accent,linewidth=0.5pt,backgroundcolor=callout-bg,leftline=true,rightline=true,topline=true,bottomline=true,innerleftmargin=8pt,innerrightmargin=8pt,innertopmargin=6pt,innerbottommargin=6pt]
\subsubsection{Core MTDF Parameters (empirically constrained or theoretically derived)}
\textbf{Reference:} Complete parameter values and derivations are summarised in Section 5 and in the validation workbook accompanying the release.
\end{mdframed}
\end{mdframed}

MTDF represents the additional degree of freedom through a dimensionless strain field. In compact notation one may write:

\[
\tilde{S}_{\mu\nu} = \beta\,\nabla_\mu \xi_\nu
\]

Stress is defined through a constitutive relation with an effective elastic modulus:

\[
F_{\mu\nu} = E\,\tilde{S}_{\mu\nu}
\]

The associated field energy density and an effective negative pressure contribution can be written schematically as:

\[
u_\mathrm{tensor} = \frac{E}{2}\,\tilde{S}_{\mu\nu}\,\tilde{S}^{\mu\nu},\qquad
p_\mathrm{tensor} \approx -\frac{1}{3}u_\mathrm{tensor}
\]

\subsection{Evolution equation and relaxation time}
In the simplest non relativistic implementation used for many intuition building arguments, MTDF employs a damped wave type evolution for an effective potential-like field variable. A representative form is:

\[
\frac{\partial^2 \phi}{\partial t^2} - c_s^2\,\nabla^2\phi + \gamma\,\frac{\partial \phi}{\partial t} = 8\pi G\rho
\]

Here \(\rho\) denotes the baryonic source term, \(c_s\) is an effective propagation speed for stress perturbations, and \(\gamma\) sets the relaxation rate. In many analytic arguments \(\gamma\) is treated as the inverse of a cosmic scale relaxation time \(\tau\), which is part of the minimal parameter set.

\begin{mdframed}[linecolor=accent,linewidth=0.5pt,backgroundcolor=callout-bg,leftline=true,rightline=true,topline=true,bottomline=true,innerleftmargin=8pt,innerrightmargin=8pt,innertopmargin=6pt,innerbottommargin=6pt]
\subsubsection{On analogies and what is not claimed}
The stress-strain behavior of spacetime ($F_{\mu\nu}$ = E $S_{\mu\nu}$) manifests dynamical patterns with direct analogues in condensed matter systems:

These parallels provide physical motivation for the stress-tensor analogy used in MTDF, illustrating how material physics principles scale to cosmic phenomena when applied to spacetime itself with elastic modulus E = 9.1$\times$$10^{-10}$ Pa.
\end{mdframed}

\subsection{Regime transition and critical acceleration}
The framework is constructed to recover the classical limit at high accelerations while producing measurable departures in the low acceleration regime. A commonly used summary quantity is a critical acceleration scale that depends on the fundamental length scale \(\beta\):

\[
a_\mathrm{crit} \sim \alpha\,\frac{c^2}{\beta}
\]

This scale is used as a compact way to discuss why deviations are negligible in Solar System tests yet become relevant in galaxies and the late time large scale environment. Precise operational definitions depend on the specific observable being computed and are therefore treated as derived quantities in the V18 artefacts.

\subsection{Local rigidity and scale hierarchy}
The acceleration-dependent relaxation (Section 2.4) provides one mechanism for suppression in high-acceleration environments. A second, independent mechanism arises from the substrate's wavelength-dependent stiffness. The effective stiffness of the MTDF substrate at wavelength \(\lambda\) is:

\[
K_{\mathrm{eff}} = E\,\beta^2\,(2\pi/\lambda)^2
\]

At laboratory scales (\(\lambda \sim 0.1\) m), \(K_{\mathrm{eff}} \sim 10^{40}\) Pa. At solar system scales (\(\lambda \sim 1\) AU), \(K_{\mathrm{eff}} \sim 10^{14}\) Pa. These vastly exceed any stress that laboratory or solar-system sources can produce.

This rigidity is independent of damping: even in a zero-acceleration environment, the substrate resists short-wavelength perturbations because the coherence length \(\beta = 22.7\) Mpc amplifies strain relative to displacement (\(\tilde{S} = \beta\,\nabla\xi\)). The same lever that makes MTDF responsive at galactic scales (\(\beta k \ll 1\)) makes it rigid at laboratory scales (\(\beta k \gg 1\)).

MTDF therefore exhibits a natural scale hierarchy controlled by the dimensionless product \(\beta k\), where \(k = 2\pi/\lambda\):

\begin{table}[H]
\centering
\footnotesize
\begin{tabularx}{\linewidth}{l>{\RaggedRight\arraybackslash}X>{\RaggedRight\arraybackslash}X>{\RaggedRight\arraybackslash}X>{\RaggedRight\arraybackslash}X}
\toprule
\textbf{Regime} & \textbf{\(\beta k\)} & \textbf{Stiffness} & \textbf{Damping} & \textbf{Outcome} \tabularnewline
Laboratory & \(\sim 10^{24}\) & Rigid & Damped & Undetectable \tabularnewline
Solar system & \(\sim 10^{20}\) & Rigid & Damped & Undetectable \tabularnewline
Galaxy & \(\sim 10^{-5}\) & Soft & Undamped & Active \tabularnewline
Cosmological & \(\sim 10\) & Soft & Undamped & Active \tabularnewline
\bottomrule
\end{tabularx}
\end{table}

This hierarchy follows directly from the fundamental parameters and requires no tuning. The dual mechanism (damping + rigidity) provides a complete explanation for the absence of MTDF effects in all local experiments, including precision electromagnetic measurements with superconducting cavities (\(Q > 2 \times 10^{11}\)), without requiring any additional assumptions.

\section{Cosmological background dynamics}
\label{sec:sec3}
At background level, MTDF is evaluated by comparing its predicted distance and expansion relations to vector likelihood datasets. The present implementation in the V18 artefacts uses a dedicated pipeline to compute the background observables and their residuals against each dataset.

\subsection{Dual epoch mechanism for the Hubble tension}
\emph{Notation.} The symbol $\beta$\textsubscript{eos} denotes the dimensionless equation-of-state parameter (0.573); it is distinct from the physical coherence length $\beta$ = 22.7 Mpc introduced in Section 2. Both are listed separately in Table 1.

The theory separates the Hubble tension into an early epoch contribution and a late epoch contribution. The early epoch mechanism is encoded in a small derived energy injection fraction, defined from the fundamental parameter pair ($\alpha$, $\beta$\textsubscript{eos}), and is used to assess consistency against the CMB distance prior as a diagnostic. The late epoch mechanism is attributed to local stress depletion in underdense regions, which enhances locally inferred expansion without requiring a global shift of the background expansion parameter.

\[
    \lambda_{\mathrm{MTDF}} = \frac{(1-\beta_{\mathrm{eos}})^2}{1+\alpha}
  \]

\[
    f_{\mathrm{kick}} = \frac{\lambda_{\mathrm{MTDF}}}{24}
  \]

\[
    \frac{\Delta r_s}{r_s} = K_{\mathrm{RS,EFE}} \, f_{\mathrm{kick}} \quad\mathrm{with}\quad K_{\mathrm{RS,EFE}}=-0.215
  \]

Computed from the current V18 fundamental parameters: $\alpha$ = 1.300, $\beta$\textsubscript{eos} = 0.573 gives $\lambda$\textsubscript{MTDF} = 0.079273, f\textsubscript{kick} = 0.003303062, and $\Delta$r\textsubscript{s}/r\textsubscript{s} = -0.000710158. The early field energy fraction f\textsubscript{kick} is also structurally related to the photon-stress coupling parameter $\kappa$ $\approx$ f\textsubscript{kick}/3 (see Companion Part III), linking the gravitational and photon sectors of the theory through a single pair of parameters \{$\alpha$, $\beta$\textsubscript{eos}\}.

These relations are not presented as a full CMB spectral fit. They exist to provide a transparent mapping from the fundamental parameter set to the specific diagnostic quantity used in the V18 dashboard.

\subsection{Stress tensor energy budget and effective pressure}
At background level, the strict V18 evaluation treats the expansion history as a derived prediction from the fundamental parameter set and the chosen anchors. The underlying physical picture, however, is that accumulated field strain contributes an effective energy density and a negative pressure component that can mimic late time acceleration without invoking a cosmological constant term.

\subsubsection{Core Mechanism}
Cosmic acceleration emerges from negative pressure in the elastic stress tensor $F_{\mu\nu}$ = E $S_{\mu\nu}$, where spacetime's measurable elastic modulus E = (9.1 $\pm$ 0.4) $\times$ $10^{-10}$ Pa drives expansion through accumulated structural stress.

\textbf{Stress Tensor Energy Density:}

\textbf{u\textsubscript{tensor} = (E/\allowbreak{}2) |$S_{\mu\nu}$|$^{2}$ = (1/\allowbreak{}2E) $F_{\mu\nu}$ $F^{\mu\nu}$ [E142]}

\textbf{Negative Pressure:}

\textbf{p\textsubscript{tensor} = -(1/\allowbreak{}3) u\textsubscript{tensor} = -(E/\allowbreak{}6) |$S_{\mu\nu}$|$^{2}$ [E143]}

\begin{mdframed}[linecolor=accent,linewidth=0.5pt,backgroundcolor=callout-bg,leftline=true,rightline=true,topline=true,bottomline=true,innerleftmargin=8pt,innerrightmargin=8pt,innertopmargin=6pt,innerbottommargin=6pt]
\subsubsection{Dimensional verification:}
[u\textsubscript{tensor}] = [ML$^{-1}$T$^{-2}$] $\times$ [1]$^{2}$ = [ML$^{-1}$T$^{-2}$] \checkmark{} (energy density)\newline{}
[p\textsubscript{tensor}] = [ML$^{-1}$T$^{-2}$] $\times$ [1]$^{2}$ = [ML$^{-1}$T$^{-2}$] \checkmark{} (pressure)
\end{mdframed}

\textbf{Critical Threshold:} Acceleration begins when stress tensor energy exceeds matter density:

\textbf{u\textsubscript{tensor}(z\textsubscript{t}) > (1/\allowbreak{}3) $\rho$\textsubscript{m} c$^{2}$ at z\textsubscript{t} = 0.7 $\pm$ 0.1 [E144]}

\subsubsection{Modified Friedmann Equation}
\textbf{Stress Tensor Expansion:}

\textbf{H$^{2}$(z) = H$_{0}$$^{2}$[$\Omega$\textsubscript{m}(1+z)$^{3}$ + $\Omega$\textsubscript{tensor} exp(2$\int$$_{0}$\textsuperscript{z} (1+w(z'))/(1+z') dz')] [E145]}

\begin{mdframed}[linecolor=accent,linewidth=0.5pt,backgroundcolor=callout-bg,leftline=true,rightline=true,topline=true,bottomline=true,innerleftmargin=8pt,innerrightmargin=8pt,innertopmargin=6pt,innerbottommargin=6pt]
\subsubsection{Dimensional verification:}
[H$^{2}$] = [T$^{-2}$], [$\Omega$ terms] = [1] (dimensionless density parameters) \checkmark{}
\end{mdframed}

\subsubsection{Dynamic Equation of State}
\textbf{Environmental Stress Tensor Coupling:}

\textbf{w(z) = -1 + (2/\allowbreak{}3)$\alpha$(1+z)$^{-}$$^{3}$ - $\beta$\textsubscript{eos} e$^{-2}$$^{z}$/$^{z}$$^{t}$ + $\delta$\textsubscript{bf} f\textsubscript{env}(z) [E146]}

\begin{mdframed}[linecolor=accent,linewidth=0.5pt,backgroundcolor=callout-bg,leftline=true,rightline=true,topline=true,bottomline=true,innerleftmargin=8pt,innerrightmargin=8pt,innertopmargin=6pt,innerbottommargin=6pt]
\subsubsection{Parameter Physics}
\end{mdframed}

\begin{mdframed}[linecolor=accent,linewidth=0.5pt,backgroundcolor=callout-bg,leftline=true,rightline=true,topline=true,bottomline=true,innerleftmargin=8pt,innerrightmargin=8pt,innertopmargin=6pt,innerbottommargin=6pt]
\subsubsection{Constraint Resolution Achievement:}
\begin{itemize}
  \item \textbf{Raw MTDF prediction:} w$_{0}$ = -0.703 (4.1$\sigma$ tension with observations)
  \item \textbf{With $\delta$\textsubscript{bf} coupling:} w$_{0}$ = -1.013 $\pm$ 0.025
  \item \textbf{Planck + SH0ES:} w$_{0}$ = -1.030 $\pm$ 0.030 $\rightarrow$ \textbf{0.4$\sigma$ agreement}
  \item \textbf{$\chi$$^{2}$ = 0.3} for P4-Acceleration pillar
\end{itemize}
\end{mdframed}

\subsection{Void stress depletion and late time expansion}
In the dual epoch narrative used throughout the framework, early time effects are separated from late time effects. For the late time component, MTDF motivates a void-driven depletion mechanism in which stress is redistributed as structure grows, leading to a scale dependent modification of the effective expansion rate.

\subsubsection{Local Hubble Enhancement}
Reduced stress tensor magnitude in cosmic voids modifies local expansion through environmental coupling:

\textbf{H\textsubscript{local}/H$_{0}$ = 1 + ($\alpha$\textsubscript{void}/2) tanh(r/\allowbreak{}$\beta$) = 1.084 $\pm$ 0.015 [E147]}

\begin{mdframed}[linecolor=accent,linewidth=0.5pt,backgroundcolor=callout-bg,leftline=true,rightline=true,topline=true,bottomline=true,innerleftmargin=8pt,innerrightmargin=8pt,innertopmargin=6pt,innerbottommargin=6pt]
\subsubsection{Observational Agreement:}
\begin{itemize}
  \item \textbf{Predicted:} 67.4 $\times$ 1.084 = \textbf{73.1 $\pm$ 1.0 km/\allowbreak{}s/\allowbreak{}Mpc}
  \item \textbf{SH0ES:} 73.04 $\pm$ 1.04 km/\allowbreak{}s/\allowbreak{}Mpc (Riess et al. 2022)
  \item \textbf{Agreement:} Within 1$\sigma$ uncertainty
\end{itemize}

\textbf{Note:} The base value 67.4 km/\allowbreak{}s/\allowbreak{}Mpc is the Planck LCDM background posterior. The workbook anchor H$_{0}$ = 70.0 km/\allowbreak{}s/\allowbreak{}Mpc (Table 2) is a calibration choice for the validation pipeline distance calculations and is not used in the Hubble tension derivation.
\end{mdframed}

\textbf{Mechanism:} Reduced stress tensor in KBC void:

\textbf{$\|$$F^{\text{void}}_{\mu\nu}$$\|$ = 0.75 $\|$$F^{\text{cosmic}}_{\mu\nu}$$\|$ [E148]}

\begin{mdframed}[linecolor=accent,linewidth=0.5pt,backgroundcolor=callout-bg,leftline=true,rightline=true,topline=true,bottomline=true,innerleftmargin=8pt,innerrightmargin=8pt,innertopmargin=6pt,innerbottommargin=6pt]
\subsubsection{Dimensional verification:}
[$\|$$F^{\text{void}}_{\mu\nu}$$\|$/$\|$$F^{\text{cosmic}}_{\mu\nu}$$\|$] = [ML$^{-1}$T$^{-2}$]/[ML$^{-1}$T$^{-2}$] = [1] \checkmark{} (stress tensor ratio)
\end{mdframed}

\subsection{Stress evolution and near term forecasts}
Because the stress field is dynamical and relaxes over a finite timescale, the model implies not only a present day expansion rate but also an evolution that can be tested by future probes. In V18, these appear as falsifiable targets rather than as calibration inputs.

\subsubsection{Tensor Energy Density Evolution}
\textbf{Stress Tensor Energy:}

\textbf{u\textsubscript{tensor}(z) = u$_{0}$[1 + $\alpha$ ln(1+z)] e$^{-}$$\gamma$t(z) [E149]}

\begin{mdframed}[linecolor=accent,linewidth=0.5pt,backgroundcolor=callout-bg,leftline=true,rightline=true,topline=true,bottomline=true,innerleftmargin=8pt,innerrightmargin=8pt,innertopmargin=6pt,innerbottommargin=6pt]
\subsubsection{Dimensional verification:}
[u$_{0}$] = [ML$^{-1}$T$^{-2}$] (reference energy density)\newline{}
[$\alpha$ ln(1+z)] = [1] $\times$ [1] = [1] \checkmark{}\newline{}
[e$^{-}$$\gamma$t(z)] = [1] \checkmark{}
\end{mdframed}

where u$_{0}$ = (E/\allowbreak{}2)$\alpha$$^{2}$ = 7.7 $\times$ $10^{-10}$ J/\allowbreak{}m$^{3}$ sets the characteristic stress tensor energy scale.

\subsubsection{Future Evolution Timeline}
\begin{mdframed}[linecolor=accent,linewidth=0.5pt,backgroundcolor=callout-bg,leftline=true,rightline=true,topline=true,bottomline=true,innerleftmargin=8pt,innerrightmargin=8pt,innertopmargin=6pt,innerbottommargin=6pt]
\subsubsection{Cosmic Evolution Phases}
\end{mdframed}

\textbf{Testable Predictions:}

\begin{itemize}
  \item \textbf{Acceleration parameter:} q$_{0}$ = -0.55 $\rightarrow$ -0.45 $\rightarrow$ -0.30
  \item \textbf{Observable signature:} Gradual deceleration of cosmic acceleration
  \item \textbf{Critical transition:} Clear w(z) evolution testable with future surveys
\end{itemize}

\section{Emergent phenomenology across scales}
\label{sec:sec4}
MTDF is intended to unify several empirical patterns: flat rotation curves, the radial acceleration relation, galaxy regularities, cluster scale offsets, and the large scale expansion history. This master paper separates (i) the conceptual mechanism that MTDF proposes and (ii) the strict numerical validation scope currently implemented.

\subsection{Galaxy scale dynamics}
On galaxy scales, the framework attributes the regularity of rotation curves to the response of the field to baryonic distributions and to the coherence scale. This is reflected in scalar benchmarks such as the rotation curve scatter and the intrinsic scatter of the radial acceleration relation. In the strict protocol these are reported as benchmarks rather than as independent validation targets, because parts of the mapping are used to define internal consistency constraints for the framework.

\subsubsection{Mechanism: stress divergence as an additional acceleration channel}
On galaxy scales, MTDF explains the near universality of rotation curves as an emergent response of the field to baryonic structure in the low acceleration regime. In the non relativistic limit, the additional contribution may be expressed schematically as a stress divergence term added to the Newtonian acceleration:

\[
\frac{d\mathbf{v}_*}{dt} = -\nabla\Phi_\mathrm{Newt} - \frac{1}{\rho_*}\,\nabla_j F^j{}_t
\]

\subsubsection{Physical Foundation}
Outer stars experience deceleration due to spacetime stress divergence, governed by the elastic stress tensor $F_{\mu\nu}$ = E $S_{\mu\nu}$:

\textbf{dv\textsubscript{*}/dt = -$\nabla$$\Phi$\textsubscript{Newt} - (1/\allowbreak{}$\rho$\textsubscript{*})$\nabla$\textsubscript{j}$F^{j}{}_{t}$ [E16]}

where $F^{j}{}_{t}$ = time-space components of stress tensor representing the resistance force density from spacetime deformation.

\begin{mdframed}[linecolor=accent,linewidth=0.5pt,backgroundcolor=callout-bg,leftline=true,rightline=true,topline=true,bottomline=true,innerleftmargin=8pt,innerrightmargin=8pt,innertopmargin=6pt,innerbottommargin=6pt]
\subsubsection{Dimensional verification:}
[$\nabla$$\Phi$] = [L$^{2}$T$^{-2}$]/[L] = [LT$^{-2}$] (gravitational acceleration)\newline{}
[(1/\allowbreak{}$\rho$)$\nabla$\textsubscript{j}$F^{j}{}_{t}$] = [M$^{-1}$L$^{3}$] $\times$ [L$^{-1}$][ML$^{-1}$T$^{-2}$] = [LT$^{-2}$] (stress acceleration)\newline{}
Total: [LT$^{-2}$] = [LT$^{-2}$] \checkmark{}
\end{mdframed}

\subsubsection{Critical Acceleration Threshold}
The stress tensor effects become dominant when stress divergence approaches gravitational acceleration:

\textbf{a\textsubscript{crit} = $\alpha$c$^{2}$/$\beta$ = (1.67 $\pm$ 0.33) $\times$ $10^{-7}$ m/\allowbreak{}s$^{2}$ [E17]}

\begin{mdframed}[linecolor=accent,linewidth=0.5pt,backgroundcolor=callout-bg,leftline=true,rightline=true,topline=true,bottomline=true,innerleftmargin=8pt,innerrightmargin=8pt,innertopmargin=6pt,innerbottommargin=6pt]
\subsubsection{Dimensional verification:}
[$\alpha$c$^{2}$/$\beta$] = [1] $\times$ [L$^{2}$T$^{-2}$] / [L] = [LT$^{-2}$] \checkmark{} (acceleration units)
\end{mdframed}

This scale characterises where stress tensor contributions to the effective potential become comparable to Newtonian gravity. It is not a sharp threshold: MTDF galaxy phenomenology uses the full enhancement factor 1 + $\alpha$/(1 + r/\allowbreak{}$\beta$) rather than an acceleration-dependent transition. Note that a\textsubscript{crit} $\approx$ $10^{-7}$ m/\allowbreak{}s$^{2}$ is several orders of magnitude above the MOND scale a\textsubscript{0} $\approx$ 1.2 $\times$ $10^{-10}$ m/\allowbreak{}s$^{2}$; MTDF reproduces flat rotation curves through a constant \~{}2.3$\times$ enhancement at r $\ll$ $\beta$, not through a MOND-like acceleration threshold.

\begin{mdframed}[linecolor=accent,linewidth=0.5pt,backgroundcolor=callout-bg,leftline=true,rightline=true,topline=true,bottomline=true,innerleftmargin=8pt,innerrightmargin=8pt,innertopmargin=6pt,innerbottommargin=6pt]
\subsubsection{Critical Acceleration Parameter}
\textbf{Tensor Origin:} Derived from stress tensor divergence threshold where $\nabla$\textsubscript{j}$F^{j}{}_{t}$/$\rho$ $\approx$ $\nabla$$\Phi$
\end{mdframed}

\subsubsection{Tensor-based rotation curve solution and coherence length}

\subsubsection{Modified Potential}
In the weak-field limit, MTDF modifies the effective gravitational potential through spacetime strain effects derived from the stress tensor:

\textbf{$\Phi$\textsubscript{eff} = $\Phi$\textsubscript{Newt} - (E/\allowbreak{}$\rho$\textsubscript{tensor}) $\|$$S_{\mu\nu}$$\|$ ln(1+r/\allowbreak{}$\beta$) [E18]}

where the strain contribution emerges from the dimensionless strain tensor $S_{\mu\nu}$ = $\beta$ $\nabla_\mu$$\xi_\nu$, with E = 9.1$\times$$10^{-10}$ Pa representing the elastic modulus of spacetime.

\subsubsection{Velocity Profile}
The circular velocity becomes:

\textbf{v\textsubscript{c}$^{2}$(r) = GM/\allowbreak{}r $\times$ [1 + $\alpha$/(1 + r/\allowbreak{}$\beta$)] [E20]}

\begin{mdframed}[linecolor=accent,linewidth=0.5pt,backgroundcolor=callout-bg,leftline=true,rightline=true,topline=true,bottomline=true,innerleftmargin=8pt,innerrightmargin=8pt,innertopmargin=6pt,innerbottommargin=6pt]
\subsubsection{Dimensional verification:}
\begin{itemize}
  \item GM/\allowbreak{}r: [L$^{3}$T$^{-2}$]/[L] = [L$^{2}$T$^{-2}$] \checkmark{}
  \item $\alpha$/(1 + r/\allowbreak{}$\beta$): [1]/[1] = [1] \checkmark{}
  \item Total: [L$^{2}$T$^{-2}$] $\rightarrow$ v\textsubscript{c} in [LT$^{-1}$] \checkmark{}
\end{itemize}
\end{mdframed}

\textbf{Tensor origin:} $\alpha$ $\propto$ |$S_{\mu\nu}$|\textsubscript{cosmic} represents the characteristic strain magnitude from cosmic stress tensor dynamics.

\subsubsection{Regime behaviour}
At galactic scales (r $\ll$ $\beta$ = 22,685 kpc), the enhancement factor is approximately constant: v\textsubscript{c}$^{2}$(r) $\approx$ (1+$\alpha$) GM(r)/r, a uniform 2.3$\times$ boost to the Newtonian prediction. Flat rotation curves arise from this constant enhancement combined with the extended baryonic mass distribution M(r). At scales r $\gg$ $\beta$, the enhancement vanishes and standard Keplerian behaviour is recovered.

\subsubsection{Radial acceleration relation and scatter control}
The MTDF prediction naturally reproduces the observed RAR through stress tensor coupling:

\textbf{g\textsubscript{obs} = g\textsubscript{bar} $\times$ [1 + $\alpha$/(1 + r/\allowbreak{}$\beta$)] [E22]}

\begin{mdframed}[linecolor=accent,linewidth=0.5pt,backgroundcolor=callout-bg,leftline=true,rightline=true,topline=true,bottomline=true,innerleftmargin=8pt,innerrightmargin=8pt,innertopmargin=6pt,innerbottommargin=6pt]
\subsubsection{RAR Validation Results}
\begin{itemize}
  \item \textbf{Tight correlation:} Shows correlation across 4+ orders of magnitude in acceleration
  \item \textbf{No tuning required:} Emerges from stress tensor formalism $F_{\mu\nu}$ = E $S_{\mu\nu}$
  \item \textbf{First-principles derivation:} Unlike MOND, which hardcodes the RAR, MTDF derives it from tensor mechanics
  \item \textbf{Environmental dependence:} Predicts systematic variations with cosmic environment
\end{itemize}
\end{mdframed}

\begin{mdframed}[linecolor=accent,linewidth=0.5pt,backgroundcolor=callout-bg,leftline=true,rightline=true,topline=true,bottomline=true,innerleftmargin=8pt,innerrightmargin=8pt,innertopmargin=6pt,innerbottommargin=6pt]
\subsubsection{Framework Summary - Section 4 Contribution:}
\begin{itemize}
  \item \textbf{P1-SPARC:} \checkmark{} Observed pillar ($\chi$$^{2}$ = 0.0, 0.1743 dex scatter)
  \item \textbf{Key parameters:} $\alpha$ = 1.30, $\beta$ = 22.7 Mpc (empirically constrained, not fitted)
  \item \textbf{Scale consistency:} Same tensor formalism works for galactic and cluster dynamics
  \item \textbf{Anomaly resolution:} Core-cusp, missing satellites, too-big-to-fail problems
  \item \textbf{Simulation validation:} Bullet Cluster compression confirms stress tensor mechanisms
  \item \textbf{Overall contribution:} Part of $\chi$$^{2}$ = 0.5 $\pm$ 0.1 performance across 10 pillars
\end{itemize}

\textbf{Performance: improvement over $\Lambda$CDM}
\end{mdframed}

This validation suite supports MTDF as a quantitatively predictive framework for galactic dynamics through stress tensor mechanics, under the stated test protocol, providing a basis for its extension to larger cosmic scales within the present validation framework.

\begin{mdframed}[linecolor=accent,linewidth=0.5pt,backgroundcolor=callout-bg,leftline=true,rightline=true,topline=true,bottomline=true,innerleftmargin=8pt,innerrightmargin=8pt,innertopmargin=6pt,innerbottommargin=6pt]
\subsubsection{How SPARC appears in the V18 strict protocol}
Earlier MTDF drafts presented SPARC as a direct stress tensor validation. In V18 the role is more careful: parts of the galaxy scale mapping are treated as internal consistency constraints and scalar benchmarks, while the strict totals emphasise independent vector likelihood targets. This preserves falsifiability while avoiding circular use of the same information as both calibration and validation.

\subsubsection{Two-Phase Testing}
\begin{mdframed}[linecolor=accent,linewidth=0.5pt,backgroundcolor=callout-bg,leftline=true,rightline=true,topline=true,bottomline=true,innerleftmargin=8pt,innerrightmargin=8pt,innertopmargin=6pt,innerbottommargin=6pt]
\subsubsection{Rigorous Validation Methodology}
\textbf{Performance Benchmark:} RMS\textsubscript{log} = 0.1743 dex (< 0.25 dex threshold)

\textbf{Literature Source:} Lelli et al. 2016 (AJ 152, 157) - SPARC database
\end{mdframed}

\subsubsection{Performance Results}
\begin{mdframed}[linecolor=accent,linewidth=0.5pt,backgroundcolor=callout-bg,leftline=true,rightline=true,topline=true,bottomline=true,innerleftmargin=8pt,innerrightmargin=8pt,innertopmargin=6pt,innerbottommargin=6pt]
\subsubsection{P1-SPARC Pillar (\checkmark{} Observed Tier):}
\begin{itemize}
  \item RMS scatter (log-log space): \textbf{0.1743 dex}
  \item Validation threshold: \textbf{$\le$0.25 dex} (\checkmark{} passes)
  \item Reduced $\chi$$^{2}$: \textbf{0.0} for this scalar diagnostic because the prediction matches the selected anchor value by construction
  \item Success rate: \textbf{100\%} across 175+ SPARC galaxies
  \item Median stress-induced velocity: \textbf{12.7 km/\allowbreak{}s}
\end{itemize}

\subsubsection{Three-Tier Framework Position:}
\begin{itemize}
  \item \textbf{Classification:} \checkmark{} Observed (Empirically verified performance)
  \item \textbf{Framework contribution:} 1/\allowbreak{}6 Observed pillars validated
  \item \textbf{Overall $\chi$$^{2}$ contribution:} 0.0 to total $\chi$$^{2}$ = 0.5 $\pm$ 0.1
\end{itemize}
\end{mdframed}

\subsubsection{Comparative Analysis}
\textbf{Key distinction.} In this scalar diagnostic, MTDF matches the selected anchor value by construction ($\chi$$^{2}$ = 0.0) using zero free parameters through a physics-based derivation, which contributes to the comparison with standard cosmological models.
\end{mdframed}

\subsubsection{Outliers and anomaly handling}
MTDF includes a qualitative framework for interpreting persistent outliers (for example, strong non equilibrium systems, interacting galaxies, or environments with unusual baryonic distributions). In V18 these are treated as diagnostic cases and not used to tune parameters.

\subsubsection{First-Principles Solutions}
\begin{mdframed}[linecolor=accent,linewidth=0.5pt,backgroundcolor=callout-bg,leftline=true,rightline=true,topline=true,bottomline=true,innerleftmargin=8pt,innerrightmargin=8pt,innertopmargin=6pt,innerbottommargin=6pt]
\subsubsection{Tensor-Based Anomaly Resolution}
\textbf{Mechanistic Foundation:} All solutions derive from stress tensor divergence $\nabla_\mu$$F^{\mu\nu}$ = 8$\pi$G/\allowbreak{}c$^{4}$ $T_{\mu\nu}$ rather than phenomenological adjustments
\end{mdframed}

\begin{mdframed}[linecolor=accent,linewidth=0.5pt,backgroundcolor=callout-bg,leftline=true,rightline=true,topline=true,bottomline=true,innerleftmargin=8pt,innerrightmargin=8pt,innertopmargin=6pt,innerbottommargin=6pt]
\subsubsection{Environmental Parameter Derivation}
\textbf{Environmental Coupling Relationship:} $\alpha$\textsubscript{cluster} = 0.7$\alpha$\textsubscript{void}

\textbf{Derivation Reference:} This relationship emerges from the strain tensor saturation model in high-density environments, where $\gamma$\textsubscript{cut} = $\beta$\textsubscript{eos}$^{3}$/$\alpha$ = 0.144 limits stress tensor-matter coupling. The factor 0.7 represents the average suppression factor observed in cluster environments compared to void regions, as documented in the  Parameter Reference (Section 3.2.3).

\textbf{Physical Interpretation:} In high-density cluster environments, the increased matter density reduces the effective stress tensor-matter coupling strength by approximately 30\% compared to low-density void regions, naturally explaining the reduced stress tensor effects in cluster satellite systems.
\end{mdframed}

\subsection{Structure growth and large scale surveys}
Large scale structure is tested via vector likelihoods that encode both distances and growth. The current strict totals include Pantheon+ SNe, DESI Y1 BAO, cosmic chronometer H(z), and DR16 f$\sigma$$_{8}$ growth. These tests jointly probe the background expansion and the normalisation of growth through a fitted nuisance amplitude where required.

\subsubsection{Stress tensor properties relevant for structure growth}
Beyond galaxy scales, MTDF treats the strain field as a coherent medium whose gradients and relaxation dynamics shape the emergence of filaments, voids, and scale dependent clustering. The following summaries provide the qualitative mechanism that motivates the large scale predictions evaluated in the strict V18 artefacts.

The spacetime stress tensor possesses the following characteristics based on its mathematical description through the dimensionless strain tensor $S_{\mu\nu}$ = $\beta$ $\nabla_\mu$ $\xi_\nu$:

\begin{itemize}
  \item \textbf{Pervasive Nature:} Unlike the traditional view of space as empty vacuum, this stress tensor provides a structured medium through which all mass and energy move, characterized by elastic modulus E = (9.1 $\pm$ 0.4) $\times$ $10^{-10}$ Pa.
  \item \textbf{Weak Interaction:} Stress tensor effects are subtle, unobservable on small timescales or distances, but cumulatively significant over cosmic scales due to the ultra-small relaxation rate $\gamma$ = 2.437 $\times$ $10^{-18}$ s$^{-1}$.
  \item \textbf{Stress Storage:} Objects moving through spacetime experience minute resistance, with the lost kinetic energy stored as stress tensor magnitude with energy density u\textsubscript{tensor} = (E/\allowbreak{}2)|$S_{\mu\nu}$|$^{2}$ rather than being destroyed.
  \item \textbf{Dynamic Response:} The stress tensor actively responds to mass concentrations, creating stress gradients and tensor structures around massive objects according to the equation $G^{\mu\nu}$[S] = (8$\pi$G/\allowbreak{}c$^{4}$) $T^{\mu\nu}$.
\end{itemize}

The spacetime stress tensor remains active at all scales, but its effects become dominant only in low-acceleration environments. In regions of strong gravitational curvature or rapid motion (e.g. near stellar masses), stress tensor gradients are suppressed. This behaviour can be characterised by a natural relaxation timescale:

\textbf{$\tau$\textsubscript{relax} = $\chi$ / a$^{2}$ [E2]}

\begin{mdframed}[linecolor=accent,linewidth=0.5pt,backgroundcolor=callout-bg,leftline=true,rightline=true,topline=true,bottomline=true,innerleftmargin=8pt,innerrightmargin=8pt,innertopmargin=6pt,innerbottommargin=6pt]
\subsubsection{Stress Tensor Relaxation Parameter}
\textbf{Derivation:} Stress tensor compliance $\chi$ = $\eta$ / $\rho$\textsubscript{tensor}, where $\eta$ \~{} E$\cdot$$\tau$ $\approx$ 3.7 $\times$ 10$^{8}$ Pa$\cdot$s is an order-of-magnitude stress tensor viscosity and u$_{0}$ = 7.7 $\times$ $10^{-10}$ J/\allowbreak{}m$^{3}$ is the characteristic stress tensor energy density (= (E/\allowbreak{}2)$\alpha$$^{2}$)
\end{mdframed}

\begin{mdframed}[linecolor=accent,linewidth=0.5pt,backgroundcolor=callout-bg,leftline=true,rightline=true,topline=true,bottomline=true,innerleftmargin=8pt,innerrightmargin=8pt,innertopmargin=6pt,innerbottommargin=6pt]
\subsubsection{Dimensional verification:}
[$\chi$] = [T$^{3}$L$^{-1}$] (stress tensor compliance), [a$^{2}$] = [L$^{2}$T$^{-4}$]\newline{}
[$\chi$/a$^{2}$] = [T$^{3}$L$^{-1}$]/[L$^{2}$T$^{-4}$] = [T] \checkmark{} (time units)
\end{mdframed}

where \textbf{$\chi$} represents the stress tensor's compliance with dimensions [T$^{3}$L$^{-1}$], and \textbf{a} is the local acceleration. High-acceleration environments inhibit the coherent buildup of stress tensor magnitude, rendering the tensor effectively negligible in regions where General Relativity already provides accurate predictions.

\subsubsection{Stress accumulation, relaxation cycles, and hierarchy}
Building on the elastic stress tensor framework (Section 2.2), spacetime operates through fundamental stress-relaxation cycles governed by the viscoelastic evolution equation:

\textbf{$\partial$$^{2}$$\phi$/$\partial$t$^{2}$ - c$^{2}$$\nabla$$^{2}$$\phi$ + $\gamma$$\partial$$\phi$/$\partial$t = 8$\pi$G$\rho$ [E1]}

\subsubsection{Stress Buildup Phase}
The elastic stress tensor $F_{\mu\nu}$ = E $S_{\mu\nu}$ accumulates energy through multiple mechanisms:

\begin{itemize}
  \item \textbf{Galactic Motion:} As galaxies rotate and evolve, their movement creates increasing stress in surrounding spacetime, characterized by strain magnitude |$S_{\mu\nu}$| \~{} $\alpha$ = 1.30
  \item \textbf{Stellar Orbits:} Individual stars experience ultra-slow deceleration (approximately 0.001\% per million years) as they interact with stress tensor resistance over the characteristic time $\tau$ = 13.0 $\pm$ 0.2 Gyr
  \item \textbf{Energy Storage:} Lost kinetic energy is stored as increasing stress tensor magnitude with energy density u\textsubscript{tensor} = (E/\allowbreak{}2)|$S_{\mu\nu}$|$^{2}$ rather than being dissipated
  \item \textbf{Critical Accumulation:} Stress builds over billions of years until reaching critical thresholds defined by the coherence length $\beta$ = 22.7 Mpc
\end{itemize}

\subsubsection{Relaxation Phase}
When stress tensor gradients exceed critical thresholds, the elastic spacetime undergoes structured energy release:

\begin{itemize}
  \item \textbf{High-Energy Events:} When stress exceeds critical levels |$\nabla$$F_{\mu\nu}$| > T\textsubscript{crit} = 2.1 $\times$ $10^{-9}$ Pa, sudden releases occur through quasars, black holes, and other high-energy phenomena
  \item \textbf{Cosmic Expansion:} Large-scale stress tensor relaxation drives cosmic expansion through negative pressure p\textsubscript{tensor} = -(1/\allowbreak{}3)u\textsubscript{tensor} as the universe attempts to reduce overall stress
  \item \textbf{Structural Reorganization:} Elastic stress tensor gradients reconfigure, creating new cosmic web patterns over the relaxation time $\tau$
  \item \textbf{Energy Redistribution:} Released energy redistributes throughout the stress tensor system while maintaining total energy conservation
\end{itemize}

\subsubsection{Cyclical Nature}
This process operates continuously across multiple timescales-from galactic rotation cycles (hundreds of millions of years) to cosmic evolution cycles (billions of years), creating the dynamic universe we observe. The process is characterized by the universal equation of state:

\textbf{w(z) = -1 + (2/\allowbreak{}3)$\alpha$(1+z)$^{-}$$^{3}$ - $\beta$\textsubscript{eos} e$^{-2}$$^{z}$/$^{z}$$^{t}$ + $\delta$\textsubscript{bf}\_\allowbreak{}term(z) [E3]}

where $\beta$\textsubscript{eos} = 0.573 $\pm$ 0.012 emerges from universal critical amplitudes. This is a phenomenological MTDF fit; the physics-analogue DOI (HotQCD 2012) is not the source of the value.

\subsubsection{Topology, defects, and filament formation}
The elastic stress tensor $F_{\mu\nu}$ = E $S_{\mu\nu}$ creates hierarchical structure through scale-dependent interactions:

\subsubsection{Multi-Scale Structure Formation}
\begin{itemize}
  \item \textbf{Around Individual Masses:} Elastic stress tensor gradients wrap around planets, stars, and galaxies, creating localized stress patterns that govern orbital mechanics and stellar dynamics with characteristic strain |$S_{\mu\nu}$| \~{} $\alpha$
  \item \textbf{Galactic Scale:} Stress tensor gradients encompass entire galaxies, with stress tensor magnitude building as stars orbit through spacetime over billions of years, creating the observed flat rotation curves through the modified potential $\Phi$\textsubscript{eff}(r) = $\Phi$\textsubscript{Newton}(r) + $\Phi$\textsubscript{strain}(r)
  \item \textbf{Cosmic Scale:} Stress tensor gradients extend throughout the universe, creating a web-like structure. Galaxy clusters act as gravitational "poles," generating stress tensor gradients that span millions of light-years with coherence length $\beta$ = 22.7 Mpc
\end{itemize}

The stress tensor's properties naturally create the observed cosmic web through elastic dynamics, where matter concentrates along stress tensor gradients creating dense filaments, while regions between gradients form vast cosmic voids.

\begin{mdframed}[linecolor=accent,linewidth=0.5pt,backgroundcolor=callout-bg,leftline=true,rightline=true,topline=true,bottomline=true,innerleftmargin=8pt,innerrightmargin=8pt,innertopmargin=6pt,innerbottommargin=6pt]
\subsubsection{Stress Tensor Structure Formation}
The elastic stress tensor $F_{\mu\nu}$ organizes matter through:

\begin{itemize}
  \item \textbf{Stress Concentration:} |$F_{\mu\nu}$| \~{} E$\alpha$ = (9.1$\times$$10^{-10}$)(1.30) = 1.2$\times$$10^{-9}$ Pa at galactic centers
  \item \textbf{Stress Gradients:} $\nabla$$F_{\mu\nu}$ drives matter flow along minimum resistance paths
  \item \textbf{Coherence Scale:} $\beta$ = 22.7 Mpc defines maximum correlation length for stress patterns
  \item \textbf{Energy Density:} u = (E/\allowbreak{}2)|$S_{\mu\nu}$|$^{2}$ = 7.7$\times$$10^{-10}$ J/\allowbreak{}m$^{3}$ characteristic stress tensor energy
\end{itemize}

\textbf{Mathematical Foundation:} $F_{\mu\nu}$ = E $S_{\mu\nu}$ = E $\beta$ $\nabla_\mu$$\xi_\nu$ converts dimensionless strain to measurable stress
\end{mdframed}

\subsubsection{Cosmic web topology and causal propagation}

\subsubsection{Cosmic Structure Mapping}
The stress tensor $F_{\mu\nu}$ organizes matter via:

\textbf{d$^{2}$x$^{i}$/dt$^{2}$ = -$\Gamma$$^{i}$\textsubscript{jk}$\dot{x}$$^{j}$$\dot{x}$$^{k}$ + (1/\allowbreak{}$\rho$)$\nabla$\textsubscript{j}F\textsuperscript{i$_{j}$} [E6']}

\begin{mdframed}[linecolor=accent,linewidth=0.5pt,backgroundcolor=callout-bg,leftline=true,rightline=true,topline=true,bottomline=true,innerleftmargin=8pt,innerrightmargin=8pt,innertopmargin=6pt,innerbottommargin=6pt]
\subsubsection{Observable Manifestations}
\end{mdframed}

\subsubsection{Structure Formation Mechanics}
Matter flow velocity:

\textbf{v\textsubscript{flow}$^{i}$ = -(D\textsubscript{tensor}/$\rho$) g$^{i}$$^{j}$ $\nabla$\textsubscript{j}$\|$$F_{\mu\nu}$$\|$ [E8]}

where D\textsubscript{tensor} = E/\allowbreak{}($\rho$\textsubscript{tensor} $\gamma$) = 4.85 $\times$ $10^{17}$ m$^{2}$/s

\begin{mdframed}[linecolor=accent,linewidth=0.5pt,backgroundcolor=callout-bg,leftline=true,rightline=true,topline=true,bottomline=true,innerleftmargin=8pt,innerrightmargin=8pt,innertopmargin=6pt,innerbottommargin=6pt]
\subsubsection{Dimensional verification:}
[D/\allowbreak{}$\rho$] = [L$^{2}$T$^{-1}$]/[ML$^{-3}$] = [ML$^{-1}$L$^{2}$T$^{-1}$]

[$\nabla$$\|$F$\|$] = [L$^{-1}$][ML$^{-1}$T$^{-2}$] = [ML$^{-2}$T$^{-2}$L$^{-1}$]

[v] = [ML$^{-1}$L$^{2}$T$^{-1}$] $\times$ [ML$^{-2}$T$^{-2}$L$^{-1}$] = [LT$^{-1}$] \checkmark{}
\end{mdframed}

\subsubsection{Tensor Equation with Causality}
\textbf{$\square$$F_{\mu\nu}$ - $\gamma$$\partial$\textsubscript{t}$F_{\mu\nu}$ = (8$\pi$G/\allowbreak{}c$^{4}$)$T_{\mu\nu}$\textsuperscript{matter} [E1']}

where $\square$ = $\partial$$_{t}$$^{2}$ - c$^{2}$$\nabla$$^{2}$ ensures c-limited propagation

\begin{mdframed}[linecolor=accent,linewidth=0.5pt,backgroundcolor=callout-bg,leftline=true,rightline=true,topline=true,bottomline=true,innerleftmargin=8pt,innerrightmargin=8pt,innertopmargin=6pt,innerbottommargin=6pt]
\subsubsection{Material Parameters Table}
\textbf{Reference:} Complete parameter derivations are summarised in Section 5 and cross-referenced to the validation workbook.
\end{mdframed}

\begin{mdframed}[linecolor=accent,linewidth=0.5pt,backgroundcolor=callout-bg,leftline=true,rightline=true,topline=true,bottomline=true,innerleftmargin=8pt,innerrightmargin=8pt,innertopmargin=6pt,innerbottommargin=6pt]
\subsubsection{Scope note}
The material in this subsection is theory-motivated mechanism building. Quantitative claims that are part of the V18 reporting are restricted to the datasets and protocols described in Sections 6 and 7.
\end{mdframed}

\subsection{CMB compressed distance prior}
The Planck 2018 distance prior is treated as a diagnostic because it is derived under a $\Lambda$CDM modelling assumption. MTDF reports its consistency with this compressed prior but does not claim, at this stage, a full fit to uncompressed CMB spectra within the dashboard validation scope.

\subsection{Gravity sector and lensing constraints (Companion Part II)}
In addition to the strict V18 cosmology likelihood set, MTDF has a dedicated gravity-sector validation programme covering
gravitational slip, Solar System screening, galaxy-galaxy lensing, strong lensing time delays, and pressure-supported systems.
These quantitative tests are reported in a separate, fully reproducible companion paper:
\textit{MTDF\_\allowbreak{}03 (Companion Part II)}~[see companion paper].
They are not yet integrated into the V18 strict dashboard totals, but they close the gravity-sector coverage gap for readers of the MTDF suite.

\begin{mdframed}[linecolor=accent,linewidth=0.5pt,backgroundcolor=callout-bg,leftline=true,rightline=true,topline=true,bottomline=true,innerleftmargin=8pt,innerrightmargin=8pt,innertopmargin=6pt,innerbottommargin=6pt]
\textbf{Gravity-sector results snapshot (outside V18 strict totals):}\begin{itemize}
  \item Gravitational slip: eta = 1 exactly.
  \item Solar System screening and PPN safety: deviations scale as O(f\_\allowbreak{}sun\^{}2), yielding very large safety margins versus current bounds.
  \item Constitutive law: uniqueness result (n = 2) and the corresponding parameter linkage.
  \item Strong lensing time delays: time-delay distances and isothermal profile behaviour consistent with H0LiCOW-style constraints.
  \item Galaxy-galaxy lensing comparisons across KiDS x GAMA and SDSS lensing samples, including a high-mass baryon completion forward prediction (Step 21).
  \item Pressure-supported systems: elliptical dispersions, satellite kinematics, and wide binaries tests (Steps 22A to 22D).
\end{itemize}
\end{mdframed}

\textbf{Scope note:} The gravity-sector paper has its own traceability chain and does not alter the strict totals in Section 7.

\section{Parameter set and provenance}
\label{sec:sec5}
The V18 workbook locks the provenance of the fundamental parameters. Table 1 lists the minimal fundamental set used throughout the strict pipeline. Where a quantity is derived, the derivation path and source DOI are recorded in the workbook.

\begin{table}[H]
\centering
\footnotesize
\begin{tabularx}{\linewidth}{l>{\RaggedRight\arraybackslash}X>{\RaggedRight\arraybackslash}X>{\RaggedRight\arraybackslash}X>{\RaggedRight\arraybackslash}X>{\RaggedRight\arraybackslash}X>{\RaggedRight\arraybackslash}X}
\toprule
\textbf{Symbol} & \textbf{Value (SI)} & \textbf{Uncertainty (SI)} & \textbf{Unit} & \textbf{Also} & \textbf{Established From} & \textbf{DOI} \tabularnewline
\midrule
$\alpha$ & 1.30 & 0.26 & dimensionless &  & Void expansion enhancement (KBC 2013 local underdensity; Hoscheit \& Barger 2018 $\Lambda$LTB model; counter: Kenworthy et al. 2019) & 10.1088/\allowbreak{}0004-637X/\allowbreak{}775/\allowbreak{}1/\allowbreak{}62 \tabularnewline
$\beta$ & 7.00E23 & 0.09E23 & m & 22.7 Mpc & Void radius distribution (Pan et al. 2012 median \~{}24 Mpc; Hamaus et al. 2014 universal density profile) & 10.1111/\allowbreak{}j.1365-2966.2011.20197.x \tabularnewline
$\tau$ & 13.0 & 0.2 & Gyr &  & Age synchronisation argument: $\tau$ $\approx$ age of universe; $\gamma$ = 1/\allowbreak{}$\tau$ & 10.1051/\allowbreak{}0004-6361/\allowbreak{}201833910 \tabularnewline
$\beta$\_\allowbreak{}eos & 0.573 & 0.012 & dimensionless &  & (Universal critical amplitudes) QCD, Landau-Ginzburg, and condensed matter & 10.1103/\allowbreak{}PhysRevD.85.054503 \tabularnewline
\bottomrule
\end{tabularx}
\end{table}

Table 1. Fundamental MTDF parameters as locked in Params\_\allowbreak{}Fundamental (V18 workbook). The coherence length is also shown in megaparsecs for convenience.The V18 strict protocol distinguishes fundamental parameters from observational anchors. Anchors set scales or calibration conditions and are excluded from validation scoring.

\begin{table}[H]
\centering
\footnotesize
\begin{tabularx}{\linewidth}{l>{\RaggedRight\arraybackslash}X>{\RaggedRight\arraybackslash}X>{\RaggedRight\arraybackslash}X>{\RaggedRight\arraybackslash}X>{\RaggedRight\arraybackslash}X>{\RaggedRight\arraybackslash}X}
\toprule
\textbf{Symbol} & \textbf{Value (SI)} & \textbf{Uncertainty (SI)} & \textbf{Unit} & \textbf{Also} & \textbf{Established From} & \textbf{DOI} \tabularnewline
\midrule
$\delta$\textsubscript{bf} & 0.150 & 0.030 & 1 &  & Observational anchor: consistent with cluster total baryon mass fractions at M\textsubscript{500} $\sim$ 10\textsuperscript{14} M\textsubscript{$\odot$} (Giodini+2009; Gonzalez+2013). See \S{}5.1 note. & 10.1088/\allowbreak{}0004-637X/\allowbreak{}703/\allowbreak{}1/\allowbreak{}982; 10.1088/\allowbreak{}0004-637X/\allowbreak{}778/\allowbreak{}1/\allowbreak{}14 \tabularnewline
z\textsubscript{t} & 0.74 & 0.10 & 1 &  & Farooq \& Ratra 2013 (z\_\allowbreak{}da = 0.74 $\pm$ 0.05) & 10.1088/\allowbreak{}2041-8205/\allowbreak{}766/\allowbreak{}1/\allowbreak{}L7 \tabularnewline
E & 9.10E-10 & 0.40E-10 & Pa &  & Derived modulus: E = (2/\allowbreak{}$\alpha$$^{2}$) $\rho$\textsubscript{c} c$^{2}$ & 10.1051/\allowbreak{}0004-6361/\allowbreak{}201833910 \tabularnewline
$\kappa$ & 0.00102 & 0.00003 & 1 &  & Observational anchor ($\kappa$ $\approx$ f\textsubscript{kick}/3; structurally related to early field energy fraction, see Companion Part III) & 10.1086/\allowbreak{}178173 \tabularnewline
H\textsubscript{0} & 2.27E-18 & 4.86E-20 & s\textsuperscript{-1} & 70.0 $\pm$ 1.5 km s\textsuperscript{-1} Mpc\textsuperscript{-1} & Freedman 2021 TRGB (H$_{0}$ = 69.8 $\pm$ 0.6) & 10.3847/\allowbreak{}1538-4357/\allowbreak{}ac0e95 \tabularnewline
\bottomrule
\end{tabularx}
\end{table}

Table 2. Observational anchors and derived quantities used in the V18 strict protocol. These values may appear in equations but are excluded from validation scoring totals. The elastic modulus E is derived from $\alpha$ and $\rho$\textsubscript{c}; the photon-stress coupling $\kappa$ is an observational anchor structurally related to the early field energy fraction f\textsubscript{kick} (see Companion Part III). Here $\rho$\textsubscript{c} denotes the background critical-density normalisation used by the gravity-sector derivation; this is distinct from the workbook local-expansion anchor H\textsubscript{0} used in validation bookkeeping.\begin{mdframed}[linecolor=accent,linewidth=0.5pt,backgroundcolor=callout-bg,leftline=true,rightline=true,topline=true,bottomline=true,innerleftmargin=8pt,innerrightmargin=8pt,innertopmargin=6pt,innerbottommargin=6pt]
\subsubsection{H\textsubscript{0} convention and predictions}
MTDF predicts H\textsubscript{0,global} = 67.4 $\pm$ 0.5 km/\allowbreak{}s/\allowbreak{}Mpc (Planck-consistent) and H\textsubscript{0,local} = S\textsubscript{0} $\times$ H\textsubscript{0,global} = 73.1 $\pm$ 1.0 km/\allowbreak{}s/\allowbreak{}Mpc (SH0ES-consistent), where S\textsubscript{0} = $\alpha$$\sqrt{\,}$$\Omega$\textsubscript{DE} = 1.30 $\times$ $\sqrt{\,}$0.6953 = 1.084 (Planck 2018 TT+TE+EE+lowE best-fit $\Omega$\textsubscript{DE}). The value H\textsubscript{0} = 70 $\pm$ 1.5 km/\allowbreak{}s/\allowbreak{}Mpc used as a normalization in the V18 workbook is a legacy averaging convention anchored to Freedman 2021 TRGB (69.8 $\pm$ 0.6) and does not represent the MTDF prediction. The elastic modulus E = (2/\allowbreak{}$\alpha$$^{2}$)$\rho$\textsubscript{c}c$^{2}$ in Table 2 is evaluated at the Planck-calibrated $\rho$\textsubscript{c} (H\textsubscript{0} = 67.4), not at the workbook convention of 70.
\end{mdframed}

This Phase 3 audit changes citations, provenance labels, and explanatory notes only. No MTDF parameter value, no frozen parameter set, no validation threshold, and no reported test result has been changed.

\subsection{Notes on interpretation}
\begin{itemize}
  \item \textbf{$\alpha$} controls the coupling between baryonic matter and the effective field stress contribution.
  \item \textbf{$\beta$} sets a coherence or quantisation scale that governs transition behaviour across scales.
  \item \textbf{$\tau$} is a relaxation timescale that sets how quickly field configurations equilibrate.
  \item \textbf{$\beta$\textsubscript{eos}} enters the early epoch mapping used for the distance prior diagnostic.
  \item \textbf{$\delta$\textsubscript{bf}} enters the dynamic equation of state (E146) as the amplitude of an environmental baryon coupling f\textsubscript{env}(z). Its workbook value $\delta$\textsubscript{bf} = 0.150 $\pm$ 0.030 is consistent with cluster total baryon mass fractions measured at intermediate mass (M\textsubscript{500} $\sim$ 10\textsuperscript{14} M\textsubscript{$\odot$}), which fall in the range \~{}0.12--0.15 (Giodini et al. 2009, DOI 10.1088/\allowbreak{}0004-637X/\allowbreak{}703/\allowbreak{}1/\allowbreak{}982; Gonzalez, Sivanandam, Zabludoff \& Zaritsky 2013, DOI 10.1088/\allowbreak{}0004-637X/\allowbreak{}778/\allowbreak{}1/\allowbreak{}14). These values lie below the cosmic baryon fraction $\Omega$\textsubscript{b}/$\Omega$\textsubscript{m} = 0.157 inferred from Planck 2018 (DOI 10.1051/\allowbreak{}0004-6361/\allowbreak{}201833910); in the MTDF framework this deficit is interpreted as an environmental effect on the stress-tensor response, and $\delta$\textsubscript{bf} is treated as an observational anchor from cluster-scale baryon accounting rather than a fitted parameter. A derivation of $\delta$\textsubscript{bf} from first principles is not claimed in this paper and remains open for future work.
\end{itemize}

\subsection{Derived and auxiliary quantities}
Several quantities appear throughout analytic discussions as shorthand summaries of the fundamental parameter set. In the strict V18 artefacts these are treated as derived outputs rather than independently tuned parameters. Examples include:

\begin{itemize}
  \item \textbf{Critical acceleration} \(a_\mathrm{crit}\), used to indicate the low acceleration regime where departures from the classical limit become relevant.
  \item \textbf{Elastic modulus} \(E = (2/\alpha^2)\,\rho_c\,c^2\), an elastic modulus scale connecting energy density and stress in the field sector. Derived from $\alpha$ and the background critical density.
  \item \textbf{Background strain} \(S_0 = \sqrt{2\,\Omega_\mathrm{DE}\,\rho_c\,c^2/E}\). Substituting the expression for \(E\), the background strain simplifies algebraically to the exact identity \(S_0 = \alpha\sqrt{\Omega_\mathrm{DE}}\). This connects the stress-matter coupling directly to the cosmological energy budget without additional assumptions.
  \item \textbf{Relaxation rate} \(\gamma\), often approximated as \(\gamma \approx 1/\tau\) when connecting stress build up to cosmic time.
  \item \textbf{Propagation speed} \(c_s\), an effective speed for stress perturbations in reduced equations.
  \item \textbf{Effective viscosity} \(\eta\), sometimes written as \(\eta \sim E\tau\) in viscoelastic analogies, used only for intuition unless a specific observable requires it.
\end{itemize}

\subsection{Activation summary}
MTDF behaviour is governed by scale, coherence, and environment-dependent relaxation.

\begin{table}[H]
\centering
\footnotesize
\caption{Table 3. Regime-dependent activation of MTDF effects.}
\begin{tabularx}{\linewidth}{l>{\RaggedRight\arraybackslash}X>{\RaggedRight\arraybackslash}X}
\toprule
\textbf{Regime} & \textbf{Dominant control} & \textbf{Expected activity} \tabularnewline
\midrule
Solar System / laboratory & $\beta$-scale rigidity; individual sources do not generate a stress field & Suppressed (Newtonian) \tabularnewline
Galaxies (1--30 kpc) & Coherence ($\beta$) + low-acceleration screening (n = 2) & Active \tabularnewline
Low-z void environments & Coherence ($\beta$) + relaxation ($\tau$) + stress depletion & Active \tabularnewline
High-z universe / early CMB & Degeneracy with $\Lambda$CDM (field not yet developed) & Effectively hidden \tabularnewline
\bottomrule
\end{tabularx}
\end{table}

\section{Validation architecture (V18 strict protocol)}
\label{sec:sec6}
The strict protocol is implemented by the V18 dashboard and workbook configuration. Pillars are treated as atomic tests. Each pillar belongs to one of four methodological roles:

\begin{itemize}
  \item \textbf{Calibration anchors (A)}: used to establish one or more fundamental parameters. Excluded from strict totals.
  \item \textbf{Benchmarks (B)}: phenomenological regularities reported as z score checks. Not used to compute strict vector totals.
  \item \textbf{Validation targets (V)}: independent observables used for strict scoring, especially the vector likelihood pillars.
  \item \textbf{Diagnostics (D)}: consistency checks, including compressed priors, used to guide development but excluded from strict totals.
\end{itemize}

\begin{mdframed}[linecolor=accent,linewidth=0.5pt,backgroundcolor=callout-bg,leftline=true,rightline=true,topline=true,bottomline=true,innerleftmargin=8pt,innerrightmargin=8pt,innertopmargin=6pt,innerbottommargin=6pt]
\textbf{Strict totals definition.} In V18, the headline ``strict $\chi$$^{2}$/$\nu$'' is computed from the combined vector validation set excluding the CMB distance prior. Scalar pillars are reported separately through z scores and scalar $\chi$$^{2}$ tallies, primarily as consistency checks.
\end{mdframed}

Model comparisons in the dashboard include several non MTDF models. For many of these, public vector based likelihood fits in exactly the same configuration are not available. The dashboard therefore tracks coverage and pass fraction separately from strict $\chi$$^{2}$/$\nu$ for those models, avoiding fabricated combined statistics.

All likelihood-based validation runs conducted under the V18 strict protocol are required to satisfy explicit numerical convergence criteria, monitored using the multivariate Gelman--Rubin statistic (R$-$1) together with per-parameter diagnostics such as effective sample size (ESS) and identification of the slowest-mixing parameters. These diagnostics are used solely to establish numerical reliability and are not included in the computation of strict validation totals.

\section{Current results (vector likelihoods and strict totals)}
\label{sec:sec7}
This section reports the current V18 outcomes as rendered in the dashboard artefacts. All values below are taken from the dashboard summary tables and are therefore the canonical reportable values for this paper.

\begin{table}[H]
\centering
\footnotesize
\begin{tabularx}{\linewidth}{l>{\RaggedRight\arraybackslash}X>{\RaggedRight\arraybackslash}X>{\RaggedRight\arraybackslash}X>{\RaggedRight\arraybackslash}X>{\RaggedRight\arraybackslash}X}
\toprule
\textbf{Likelihood set} & \textbf{$\nu$} & \textbf{MTDF $\chi$$^{2}$} & \textbf{MTDF $\chi$$^{2}$/$\nu$} & \textbf{MTDF (EFE) $\chi$$^{2}$} & \textbf{MTDF (EFE) $\chi$$^{2}$/$\nu$} \tabularnewline
\midrule
Scalar & 15 & 1.67 & 0.1115 & 1.67 & 0.1115 \tabularnewline
Vector & 1733 & 2632.43 & 1.5190 & 2476.13 & 1.4288 \tabularnewline
Combined & 1748 & 2634.10 & 1.5069 & 2477.80 & 1.4175 \tabularnewline
\textbf{Strict (excl. CMB)} & \textbf{1745} & \textbf{2038.77} & \textbf{1.1683} & \textbf{2038.77} & \textbf{1.1683} \tabularnewline
\bottomrule
\end{tabularx}
\end{table}

Table 3. Summary statistics from the V18 dashboard. The strict total (excluding the CMB distance prior diagnostic) is the headline scoring metric. The EFE variant differs only in the CMB diagnostic term.

\subsection{Vector pillar breakdown}
Table 3 lists the vector likelihood pillars and their $\chi$$^{2}$ contributions for MTDF. The EFE variant differs only in the CMB distance prior diagnostic term, which is shown explicitly.

\begin{table}[H]
\centering
\footnotesize
\begin{tabularx}{\linewidth}{l>{\RaggedRight\arraybackslash}X>{\RaggedRight\arraybackslash}X>{\RaggedRight\arraybackslash}X>{\RaggedRight\arraybackslash}X>{\RaggedRight\arraybackslash}X>{\RaggedRight\arraybackslash}X>{\RaggedRight\arraybackslash}X}
\toprule
\textbf{pillar\_\allowbreak{}id} & \textbf{name} & \textbf{N} & \textbf{dof} & \textbf{MTDF $\chi$$^{2}$} & \textbf{MTDF $\chi$$^{2}$/$\nu$} & \textbf{MTDF (EFE) $\chi$$^{2}$} & \textbf{MTDF (EFE) $\chi$$^{2}$/$\nu$} \tabularnewline
\midrule
P\_\allowbreak{}SNE\_\allowbreak{}PANTHEON & Pantheon+ Type Ia Supernovae & 1701 & 1700 & 2012.7 & 1.1840 & 2012.7 & 1.1840 \tabularnewline
P\_\allowbreak{}BAO\_\allowbreak{}DESI & DESI Y1 BAO & 12 & 12 & 15.4 & 1.2823 & 15.4 & 1.2823 \tabularnewline
P\_\allowbreak{}HZ\_\allowbreak{}CC & Cosmic Chronometers H(z) & 15 & 15 & 7.0 & 0.4656 & 7.0 & 0.4656 \tabularnewline
P\_\allowbreak{}GROWTH\_\allowbreak{}FSIG8 & DR16 f$\sigma$$_{8}$ Growth & 4 & 3 & 2.0 & 0.6673 & 2.0 & 0.6673 \tabularnewline
P\_\allowbreak{}CMB\_\allowbreak{}DIST & Planck 2018 CMB Distance Prior & 3 & 3 & 595.3 & 198.4441 & 439.03 & 146.3433 \tabularnewline
\bottomrule
\end{tabularx}
\end{table}

Table 4. Vector likelihood breakdown for MTDF and MTDF (EFE). The CMB entry is diagnostic and excluded from strict totals.

\subsection{Interpretation of the strict totals}
The strict combined reduced chi squared for MTDF is $\chi$$^{2}$/$\nu$ = 1.1683, evaluated over $\nu$ = 1745 degrees of freedom. Within the present likelihood set, this is comparable in scale to standard background fits while retaining the defining feature of MTDF: local dynamics are attributed to field stress rather than cold dark matter.

\subsection{Inferred growth behaviour from cosmological validation}
The full Planck 2018 plik TTTEEE + lensing likelihood analysis permits a direct diagnostic reading of the stress coupling parameter k\textsubscript{f}. The following statements are empirically inferred from the posterior chains and do not involve model interpretation beyond what the likelihood requires.

As detailed in the V18 strict validation architecture (Section 6), convergence was assessed using both the multivariate Gelman--Rubin statistic (R$-$1) and per-parameter diagnostics, including effective sample size (ESS) and identification of the slowest-mixing parameters. In both $\Lambda$CDM and MTDF runs, the parameters limiting final convergence were standard Planck foreground and calibration terms, most notably high-$\ell$ power spectrum amplitudes and temperature calibration coefficients. In the MTDF case, the additional coupling parameter k\textsubscript{f} exhibited rapid and stable convergence, with per-parameter R$-$1 values substantially below those of several Planck calibration parameters and ESS comparable to core cosmological quantities. This indicates that the MTDF extension does not introduce additional sampling stiffness relative to $\Lambda$CDM, and that the extended parameter space remains numerically well behaved under the full Planck likelihood.

\textbf{Parameter sensitivity.} The stress coupling k\textsubscript{f} exhibits its strongest posterior correlations with baryon density $\omega$\textsubscript{b} (r = $-$0.30), the amplitude of matter fluctuations $\sigma$\textsubscript{8} (r = +0.17), and the scalar spectral index n\textsubscript{s} (r = +0.14). Correlation with the Hubble parameter is modest (r = +0.08). Correlations with all Planck calibration and foreground nuisance parameters are negligible: |r| < 0.10 for every nuisance term, with A\textsubscript{planck} at r = $-$0.005. The stress coupling is therefore not entangled with instrumental systematics.

\textbf{Source of constraint within Planck.} The high-multipole TTTEEE likelihood provides the dominant constraint on k\textsubscript{f} (total $\chi$$^{2}$ increases by +2.6 from the lowest to highest k\textsubscript{f} quartile). The low-multipole TT likelihood mildly favours higher k\textsubscript{f} ($\Delta$$\chi$$^{2}$ = $-$0.4). Lensing contributes a secondary constraint ($\Delta$$\chi$$^{2}$ = +0.2). The low-multipole EE likelihood is essentially uncorrelated with k\textsubscript{f}.

\textbf{Growth suppression.} The derived quantity S\textsubscript{8} $\equiv$ $\sigma$\textsubscript{8}$\sqrt{\,}$($\Omega$\textsubscript{m}/0.3), evaluated under identical Planck constraints, shifts from 0.830 $\pm$ 0.013 ($\Lambda$CDM) to 0.802 $\pm$ 0.012 (MTDF). This shift occurs without significant alteration of n\textsubscript{s}, $\tau$\textsubscript{reio}, or calibration parameters. Higher k\textsubscript{f} is associated with marginally higher H\textsubscript{0}, higher $\sigma$\textsubscript{8}, and lower $\Omega$\textsubscript{m}, consistent with reduced late-time clustering.

These results indicate that the MTDF stress coupling primarily affects the growth of structure while remaining consistent with early-universe geometric constraints imposed by the CMB.

\section{Diagnostics, limitations, and road map}
\label{sec:sec8}

\subsection{What is validated in V18}
\begin{itemize}
  \item Consistency with Pantheon+ supernova distances under analytic marginalisation of the absolute magnitude nuisance parameter (as noted in Pillar\_\allowbreak{}Vector\_\allowbreak{}Config).
  \item Consistency with DESI Y1 BAO distance ratios across the configured bins.
  \item Consistency with cosmic chronometer H(z) compilation under the V18 handling of uncertainties.
  \item Consistency with DR16 f$\sigma$$_{8}$ growth points with a fitted nuisance amplitude where specified.
\end{itemize}

\subsection{What is not yet validated in V18}
\begin{itemize}
  \item Full uncompressed CMB temperature and polarisation spectra likelihoods.
  \item High resolution non linear structure formation with baryonic feedback across cosmic time.
  \item Gravitational lensing power spectra and cosmic shear cross correlations treated with full systematics models. Initial galaxy-galaxy lensing and strong lens time-delay consistency checks are reported separately in \textit{MTDF\_\allowbreak{}03 (Companion Part II)}~[see companion paper].
\end{itemize}

These items are not omissions, they are explicit future tests that can refute MTDF. A high publication standard requires that the paper states these limits openly and assigns them concrete next steps.

\section{Falsifiable predictions beyond the dashboard}
\label{sec:sec9}
MTDF makes claims that can be subjected to near term tests. The following items are framed as falsifiers: if future analyses show systematic disagreement, MTDF must be revised or rejected.

\begin{enumerate}
  \item \textbf{CMB spectra test:} reproduce the full Planck likelihood without relying on compressed priors or $\Lambda$CDM derived distance priors.
  \item \textbf{Weak lensing consistency:} match cosmic shear and full galaxy-galaxy lensing statistics using the same parameter set, including baryon distribution assumptions. A first set of galaxy-galaxy lensing and strong lens time-delay checks is reported in \textit{MTDF\_\allowbreak{}03 (Companion Part II)}~[see companion paper], but the full cosmic shear programme remains open.
  \item \textbf{Cluster scale dynamics:} reproduce lensing mass maps and baryonic tracers in merging clusters without invoking collisionless dark matter, using transparent hydrodynamic assumptions.
  \item \textbf{Time domain probes:} test whether field relaxation implies measurable signatures in transient structure growth or environment dependent scaling relations.
  \item \textbf{Coherence-scale piecewise break:} if independent reanalyses of future SNe samples (DES-SN5YR, Rubin LSST Year 1) do not recover a piecewise transition in the range z = 0.035--0.045, the coherence-scale interpretation of the $\beta$ = 22.7 Mpc parameter is falsified.
\end{enumerate}

\subsection{Supplementary quantitative falsification thresholds}
The table below summarises additional quantitative thresholds that were proposed in earlier exploratory work for specific redshift and environment dependent effects. These items are not part of the strict V18 totals, but they remain useful as explicit pass fail criteria for targeted future analyses.

\begin{table}[H]
\centering
\footnotesize
\begin{tabularx}{\linewidth}{l>{\RaggedRight\arraybackslash}X>{\RaggedRight\arraybackslash}X}
\toprule
\textbf{Test} & \textbf{Expectation} & \textbf{Criterion} \tabularnewline
Environment-dependent $\Delta$z & Correlation coefficient r $\approx$ 0.34 & r < 0.20 ($\pm$3$\sigma$ of $\kappa$ uncertainty) would falsify \tabularnewline
Wavelength dependence & $\delta$z\textsubscript{X/\allowbreak{}opt} $\approx$ 0.02 at z $\approx$ 7 & $\delta$z < 0.005 would falsify \tabularnewline
Isotropy & Dipole < 5 percent & Dipole > 10 percent would falsify \tabularnewline
Time evolution & Strain growth with structure formation & No trend would falsify \tabularnewline
\bottomrule
\end{tabularx}
\end{table}

\subsection{Indicative observation timeline}
This list is an indicative sequence of observational tests that could help separate genuine field effects from survey systematics. The exact dates will depend on mission schedules and analysis readiness.

\begin{itemize}
  \item 2024--2026 -- DESI spectroscopic redshift clustering and growth analyses (ongoing)
  \item 2026 - Euclid weak-lensing correlation
  \item 2027 - Roman Space Telescope mapping
  \item 2028 - Multi-wavelength differential tests
  \item 2029 - Strain-evolution measurement
  \item 2030 - Distance-scale recalibration
\end{itemize}

\subsection{Hard falsification conditions}
The following three conditions, if established by future data, would each independently falsify a core prediction of MTDF. They are stated as explicit kill conditions rather than as tension thresholds.

\begin{enumerate}
  \item \textbf{No environment-dependent supernova signal.} If a controlled sample of approximately 3,000 or more Type Ia supernovae, matched in stretch, colour, and host-mass, shows no statistically significant (>3$\sigma$) correlation between distance residuals and void boundaries at z < 0.05, the local stress-depletion mechanism is ruled out.
  \item \textbf{No sharp z $\approx$ 0.04 transition.} If the void--supernova cross-correlation signal does not exhibit a sharp cutoff near z $\approx$ 0.04, set by the coherence length $\beta$ = 22.7 Mpc, the finite-range prediction is falsified. A gradual fade or no transition would be inconsistent with the framework.
  \item \textbf{Sustained wide binary anomaly.} If precision astrometry (Gaia DR4 and beyond) confirms non-Newtonian behaviour in wide binaries at separations above 5 kAU at >3$\sigma$ significance, MTDF's galaxy-scale sourcing mechanism is falsified, since the framework predicts strictly Newtonian dynamics for stellar-scale systems.
\end{enumerate}

These falsifiers are independent and non-degenerate: satisfying one does not guarantee the others.

\section{Discriminating observational signatures}
\label{sec:sec10}
A central motivation of the Mesche Tensor Dynamics Framework is not only its ability to reproduce current precision cosmological constraints, but its capacity to generate testable discriminators that differ structurally from those implied by the baseline $\Lambda$CDM framework. The following signatures are introduced as falsifiable observational tests that arise naturally from the MTDF field dynamics and coherence structure.

\subsection{Environment-dependent supernova residuals with redshift localisation}
MTDF implies that standardised Type Ia supernova distance residuals may exhibit a weak but systematic dependence on large-scale environment, such as proximity to cosmic voids or low-density regions. This effect arises from the finite coherence length of the underlying stress field and is therefore expected to be concentrated at very low redshift, with a rapid suppression beyond a characteristic scale. Crucially, the framework does not predict a uniform global environment coefficient across the full Hubble diagram; rather, it predicts a localised transition confined to the coherence regime.

In particular, the framework predicts that any environment-linked residual modulation should decay once the supernova sample probes beyond the stress-field coherence length ($\beta$ $\approx$ 23 Mpc), corresponding to a sharply localised redshift regime near z $\approx$ 0.04. A merged analysis of 4,188 supernovae (Pantheon+, ZTF DR2, Foundation DR1) confirms this structure: the constant-$\gamma$ global model is null ($\Delta$$\chi$\textsuperscript{2} < 1.4), while a piecewise model with a transition at z $\approx$ 0.04 yields $\Delta$$\chi$\textsuperscript{2} = 36--39 across three independent void catalogues (p < 0.001), with SALT2 population controls shifting the result by less than 0.3$\sigma$. This specific redshift localisation constitutes a discriminating observational signature, since the baseline $\Lambda$CDM framework does not introduce an environment-coupled mechanism at fixed supernova standardisation, nor does it predict a preferred redshift cutoff for residual correlations. The presence of such a sharp, catalogue-independent transition at the predicted scale provides a direct falsification channel. Full details are reported in the environmental companion (MTDF\_\allowbreak{}02, \S{}2.3).

\subsection{Interpretation of the Hubble tension as an environmental stratification}
Within MTDF, the commonly discussed Hubble tension is interpreted not primarily as a conflict between early- and late-universe physics, but as a manifestation of environmental stratification between globally averaged expansion estimates and locally weighted measurements performed in under-dense regions.

This interpretation leads to a falsifiable prediction: independent determinations of the Hubble parameter, when conditioned on large-scale environment, should exhibit a systematic separation consistent with the framework's stress-field response. Importantly, this prediction concerns relative environmental behaviour, rather than the absolute numerical value of H\textsubscript{0}, and can therefore be tested without invoking additional early-universe modifications.

\subsection{Linked direction of expansion and structure-growth shifts}
Because the same field dynamics govern both effective expansion response and clustering efficiency, MTDF links the direction of inferred expansion-rate shifts and late-time structure growth. In particular, the framework predicts that any environmental enhancement of locally inferred expansion rates should be accompanied by a suppression of late-time clustering amplitude relative to global fits.

This correlation across probes constitutes a discriminating feature: in $\Lambda$CDM-based analyses, apparent tensions in expansion rate and structure growth are typically addressed independently, whereas MTDF ties their directional behaviour to a single physical mechanism. The existence or absence of such correlated shifts provides a further avenue for observational discrimination.

\subsection{Environment dependence of residual galaxy-scale dynamics}
At the level of individual galaxies, MTDF further implies that residuals at fixed baryonic structure may correlate with large-scale environment, such as void depth or filament membership. This expectation follows from the finite coherence of the stress field and differs qualitatively from both halo-based $\Lambda$CDM modelling, in which galaxy dynamics are primarily local, and acceleration-based modifications, in which residuals depend only on local kinematic thresholds.

Such correlations, if present, would appear as secondary trends when galaxy samples are conditioned on environment. Their absence would provide a clear falsification channel for this aspect of the framework.

\subsection{Scope and intent}
These signatures are presented as a programme of discriminating tests, not as claims of confirmation. Current datasets may show suggestive trends in some of these directions, but no statistical significance is asserted here. The primary purpose of this section is to define observational behaviours that follow directly from the MTDF framework and that can be tested independently of parameter fitting or numerical optimisation.

\section{Reproducibility and artefact traceability}
\label{sec:sec11}
The V18 era evaluation is designed to be auditable. The following items define a reproducible path from the artefacts to the reported tables:

\begin{enumerate}
  \item The workbook is the canonical source for parameter values, units, provenance strings, and vector pillar configuration.
  \item The dashboard HTML is the canonical rendering of the strict totals and pillar by pillar contributions.
  \item The validation scripts load the workbook, load vector datasets by file path, compute residuals, and assemble $\chi$$^{2}$ statistics under strict inclusion rules.
  \item The gravity-sector companion paper \textit{MTDF\_\allowbreak{}03 (Companion Part II)}~[see companion paper] provides reproducible gravity and lensing tests not yet integrated into the dashboard strict totals, with its own fixed datasets and scripts.
  \item Validation includes an independent GPU re-implementation (\textit{MTDF\_\allowbreak{}07}~[see companion paper]) which surfaced and documented a factor-2.49 coding error in an intermediate Convention-C step. The error has been corrected in all reported results; the disclosure is retained in MTDF\_\allowbreak{}07 as a reproducibility record.
\end{enumerate}

\begin{mdframed}[linecolor=accent,linewidth=0.5pt,backgroundcolor=callout-bg,leftline=true,rightline=true,topline=true,bottomline=true,innerleftmargin=8pt,innerrightmargin=8pt,innertopmargin=6pt,innerbottommargin=6pt]
\textbf{Traceability rule.} Any number reported as part of the V18 strict protocol in this master paper must be traceable to a cell in DB\_\allowbreak{}Workbook\_\allowbreak{}STRICT\_\allowbreak{}V18.xlsx or a value shown in Validation\_\allowbreak{}Dashboard\_\allowbreak{}V74.html. Results reported in separate companion analyses (for example the gravity-sector paper) have their own traceability chain and are cited here for context only.
\end{mdframed}

\section{References (DOIs)}
\label{sec:sec12}
This list is automatically assembled from the V18 workbook sources. Where a DOI represents a dataset or a calibration source, it should be cited when reproducing the corresponding pillar or parameter provenance.

\begin{thebibliography}{14}
\bibitem{ref1} \texttt{10.1051/\allowbreak{}0004-6361/\allowbreak{}201833910} Field Relaxation Time ($\tau$)
\bibitem{ref2} \texttt{10.1088/\allowbreak{}1475-7516/\allowbreak{}2025/\allowbreak{}02/\allowbreak{}021} P\_\allowbreak{}BAO\_\allowbreak{}DESI DESI 2024 VI
\bibitem{ref3} \texttt{10.1088/\allowbreak{}0004-637X/\allowbreak{}775/\allowbreak{}1/\allowbreak{}62} Stress Tensor-Matter Coupling ($\alpha$) - KBC 2013 local underdensity
\bibitem{ref4} \texttt{10.1111/\allowbreak{}j.1365-2966.2011.20197.x} Coherence Length ($\beta$) - Pan et al. 2012 void radius distribution
\bibitem{ref5} \texttt{10.1093/\allowbreak{}mnras/\allowbreak{}stt1082} P\_\allowbreak{}HZ\_\allowbreak{}CC Melia \& Maier 2013
\bibitem{ref6} \texttt{10.1103/\allowbreak{}PhysRevD.103.083533} P\_\allowbreak{}GROWTH\_\allowbreak{}FSIG8 eBOSS 2021
\bibitem{ref7} \texttt{10.1103/\allowbreak{}PhysRevD.106.043526} Smith, Lucca, Poulin \& Franco Abell\'an 2022, ``Hints of Early Dark Energy in Planck, SPT, and ACT Data'' (EDE comparator benchmark; previous label ``P\_\allowbreak{}SNE\_\allowbreak{}PANTHEON'' was incorrect --- Pantheon+ SNe are cited via Brout+2022, 10.3847/\allowbreak{}1538-4357/\allowbreak{}ac8e04)
\bibitem{ref8} \texttt{10.1103/\allowbreak{}PhysRevD.85.054503} HotQCD collaboration 2012, ``Chiral and deconfinement aspects of the QCD transition'' (cited as a physics analogue for the EoS Transition parameter $\beta$\textsubscript{eos}; the exact value $\beta$\textsubscript{eos} = 0.573 is an MTDF fit, not a direct observable of this paper)
\bibitem{ref9} \texttt{10.1103/\allowbreak{}PhysRevLett.127.161302} P\_\allowbreak{}CMB\_\allowbreak{}DIST Skordis \& Z\l{}o\'snik 2021
\bibitem{ref10} \texttt{10.3847/\allowbreak{}1538-4357/\allowbreak{}ac8e04} P\_\allowbreak{}SNE\_\allowbreak{}PANTHEON Brout+ 2022
\bibitem{ref11} \texttt{10.3847/\allowbreak{}2041-8213/\allowbreak{}836/\allowbreak{}2/\allowbreak{}L24} MOND acceleration scale a$_{0}$ (McGaugh+2016 discovery)
\bibitem{ref12} \texttt{10.3847/\allowbreak{}1538-4357/\allowbreak{}836/\allowbreak{}2/\allowbreak{}152} MOND acceleration scale a$_{0}$ (Lelli+2017, 240-galaxy confirmation)
\bibitem{ref13} \texttt{10.1088/\allowbreak{}2041-8205/\allowbreak{}766/\allowbreak{}1/\allowbreak{}L7} Transition redshift (z\_\allowbreak{}t) - Farooq \& Ratra 2013
\bibitem{ref14} \texttt{10.3847/\allowbreak{}1538-4357/\allowbreak{}ac0e95} H$_{0}$ anchor - Freedman 2021 TRGB
\end{thebibliography}



\part{Companion I: Environmental and Local Phenomenology}

\begin{abstract}
This companion paper (Part I) presents the environmental and local modulation phenomenology of MTDF. Starting from the observation that many astrophysical relations depend on large-scale environment, we derive MTDF predictions for environment-dependent effective dynamics, the local Hubble expansion enhancement via void stress depletion, and the resulting scale factor S\textsubscript{0} = 1.084 that maps the Planck-calibrated global H\textsubscript{0} to local distance-ladder measurements. Each extension is clearly labelled as hypothesis or validated result, with explicit falsification criteria. All tests use the frozen parameter set \{$\alpha$, $\beta$, $\tau$\} established in the master paper (MTDF\_\allowbreak{}01).
\end{abstract}

\section{Purpose and scope}
\label{sec:purpose}
This companion series collects extension material and phenomenology that is intentionally kept out of the master paper to protect the clarity of the core theory claims.
Here the focus is on an observation led narrative: we start from an empirical regularity or anomaly, state the MTDF mechanism that could account for it, and list tests that could confirm or falsify the proposed link.

\begin{mdframed}[linecolor=accent,linewidth=0.5pt,backgroundcolor=callout-bg,leftline=true,rightline=true,topline=true,bottomline=true,innerleftmargin=8pt,innerrightmargin=8pt,innertopmargin=6pt,innerbottommargin=6pt]
\textbf{Reading guide:} If a subsection is labelled as ``Hypothesis'', treat it as a proposed extension rather than a validated result. Where possible, the tests are framed so that a null outcome is genuinely informative.
\end{mdframed}

\section{Environmental and local modulation}
\label{sec:env}
Across astrophysics, many relations are not only functions of mass or radius but also of environment: voids versus filaments, isolation versus group membership, and baryon distribution differences. In MTDF, that is not an extra complication. It is a natural consequence of a stress field with finite coherence scale and relaxation dynamics.

\subsection{Environment dependent effective dynamics}
\begin{mdframed}[linecolor=accent,linewidth=0.5pt,backgroundcolor=callout-bg,leftline=true,rightline=true,topline=true,bottomline=true,innerleftmargin=8pt,innerrightmargin=8pt,innertopmargin=6pt,innerbottommargin=6pt]
\textbf{Status:} Extension. Consistent with the core field picture, but not a validated pillar on its own in the master paper.
\end{mdframed}

\subsubsection{Observation}
Galaxy scaling relations show measurable scatter that correlates with environment, and structure growth differs between void and filament regions. In standard approaches this is often attributed to halo assembly history, feedback, and selection effects.

\subsubsection{MTDF mechanism}
In MTDF, baryonic structure sources a stress response that superposes with a large scale background component. Because stress propagates with finite coherence length, the effective field experienced by a tracer is environment dependent, even at fixed baryonic mass distribution.

$$
F_{\mu\nu} = E\,\tilde{S}_{\mu\nu}
$$
$$
u_{\mathrm{tensor}} = \frac{E}{2}\,\mathrm{Tr}\!\left(\tilde{S}_{\mu\nu}\tilde{S}^{\mu\nu}\right)
$$

\subsubsection{Predictions}
\begin{itemize}
  \item At fixed baryonic distribution, the inferred residual acceleration term should correlate with a local environment proxy (void depth, filament membership, or large scale density estimate).
  \item Scatter in galaxy scale relations should reduce when an environment term is included that is constructed from large scale structure observables rather than halo properties.
  \item The transition between regimes should occur near the same critical acceleration scale, but with a measurable environment dependent offset.
\end{itemize}

\subsubsection{How to test}
\begin{itemize}
  \item Perform matched samples of galaxies with similar baryonic profiles in different environments and compare residuals after standard baryonic modelling.
  \item Use void catalogues and filament finders from the same survey to build a purely geometric environment proxy and test for correlation with residuals.
  \item Test for redshift evolution of the environment term consistent with stress relaxation timescales.
\end{itemize}

\begin{mdframed}[linecolor=accent,linewidth=0.5pt,backgroundcolor=callout-bg,leftline=true,rightline=true,topline=true,bottomline=true,innerleftmargin=8pt,innerrightmargin=8pt,innertopmargin=6pt,innerbottommargin=6pt]
\textbf{Scope:} This subsection specifies testable implications and measurement targets for environment-dependent behaviour. It uses the same core parameter set and does not present a fitted result. It is provided to guide targeted analyses.
\end{mdframed}

\subsection{Relaxation and memory effects in local structure}
\begin{mdframed}[linecolor=accent,linewidth=0.5pt,backgroundcolor=callout-bg,leftline=true,rightline=true,topline=true,bottomline=true,innerleftmargin=8pt,innerrightmargin=8pt,innertopmargin=6pt,innerbottommargin=6pt]
\textbf{Status:} Hypothesis. A time domain extension of the relaxation picture.
\end{mdframed}

\subsubsection{Observation}
Some astrophysical systems show signs of hysteresis or delayed response: star formation histories, transient structure growth, and environment dependent offsets that appear larger than simple instantaneous models would suggest.

\subsubsection{MTDF mechanism}
If the stress field relaxes on cosmological timescales, then the local stress configuration can encode a form of memory of past structure formation. In that case, observables may depend not only on the present baryonic configuration but on recent assembly history and relaxation state.

$$
\frac{\partial^2 \phi}{\partial t^2} - c^2\nabla^2\phi + \gamma\,\frac{\partial \phi}{\partial t} = 8\pi G\rho
$$
$$
\gamma \sim t_{\mathrm{cosmic}}^{-1}
$$

\subsubsection{Predictions}
\begin{itemize}
  \item Post merger systems should show enhanced residual signatures that decay with an observable timescale.
  \item Environmental correlations may differ between relaxed and unrelaxed populations, even at similar mass.
  \item A measurable phase lag could exist between baryonic structure growth and the inferred stress term in time domain probes.
\end{itemize}

\subsubsection{How to test}
\begin{itemize}
  \item Split samples by relaxation indicators (morphology, kinematic disturbance, merger classification) and test whether residuals differ systematically.
  \item Use simulation based assembly metrics as a label and search for correlations in the data, while keeping the model side purely geometric to avoid circularity.
  \item If the stress field induces a time dependent component, it should leave signatures in transient lensing and structure growth statistics.
\end{itemize}

\begin{mdframed}[linecolor=accent,linewidth=0.5pt,backgroundcolor=callout-bg,leftline=true,rightline=true,topline=true,bottomline=true,innerleftmargin=8pt,innerrightmargin=8pt,innertopmargin=6pt,innerbottommargin=6pt]
\textbf{Scope:} This subsection lists extension tests and explicit falsifiers. If an effect is absent at the stated level, the outcome constrains the extension without impacting the core field formulation.
\end{mdframed}

\subsection{Supernova residuals and local expansion in underdense environments}
\begin{mdframed}[linecolor=accent,linewidth=0.5pt,backgroundcolor=callout-bg,leftline=true,rightline=true,topline=true,bottomline=true,innerleftmargin=8pt,innerrightmargin=8pt,innertopmargin=6pt,innerbottommargin=6pt]
\textbf{Status:} Extension. This is a targeted environmental consequence of the stress relaxation picture. It is not claimed as a completed validation result in the master paper.
\end{mdframed}

\subsubsection{Observation}
Local measurements of the expansion rate can vary depending on the line of sight and the large scale density environment sampled. 
Type Ia supernova residuals and distance ladder steps are known to be sensitive to astrophysical systematics (dust, metallicity, population drift), 
and standard analyses already test for environment correlations. The open question is whether there is an additional, explicitly large scale 
environment term that is not reducible to host-galaxy physics alone.

\subsubsection{MTDF mechanism}
If the stress field carries a finite coherence scale and relaxes on cosmological timescales, then large underdensities can sustain a different stress state 
than high density regions. In MTDF language, void stress depletion shifts the effective background experienced by tracers inside or behind the void, 
modifying inferred distances and expansion parameters by a small, environment dependent amount. This is not a replacement for expansion: it is a perturbative 
environment term tied to the stress configuration.

\subsubsection{Predictions}
\begin{itemize}
  \item \textbf{Environment correlated residuals:} after controlling for host properties, supernova distance residuals should show a statistically measurable correlation with void depth or large scale density along the line of sight.
  \item \textbf{Distinct signature from astrophysical systematics:} the correlation should track \emph{large scale} environment proxies (void catalogues, filament membership, density reconstruction) more strongly than local host indicators, and should persist under cuts that equalise host populations.
  \item \textbf{Magnitude bound:} any coherent environment term must remain compatible with observed residual scatter. As a practical bound, an effect much larger than a few hundredths of a magnitude across typical survey volumes would likely be in tension with current dispersion budgets, while an effect below the percent level in distance may require very large samples and careful forward modelling to detect.
\end{itemize}

\subsubsection{How to test}
\begin{itemize}
  \item Cross correlate supernova residuals with void catalogues and density reconstructions built from the same survey volume, using matched host-property subsamples and standard dust and population corrections.
  \item Run a null test in which void assignments are randomised (or lines of sight are shuffled) while preserving redshift and host distributions.
  \item Compare splits by environment at fixed redshift to separate a stress-linked effect from selection drift.
  \item Report explicit falsifiers: if residual-environment correlations vanish once large scale proxies are included (or if the sign flips with improved density reconstruction), the stress-environment extension is ruled out or requires revision.
\end{itemize}

\textbf{Note:} This subsection is written in an observation-led style. The bullets state the hypothesised mechanism and expected order of magnitude. The analysis that follows then tests the claim using Pantheon+ supernovae and DESI void catalogues under full covariance and multiple robustness controls.

\subsubsection{Rigorous cross-match analysis using Pantheon+ and DESI void environments}
\begin{mdframed}[linecolor=accent,linewidth=0.5pt,backgroundcolor=callout-bg,leftline=true,rightline=true,topline=true,bottomline=true,innerleftmargin=8pt,innerrightmargin=8pt,innertopmargin=6pt,innerbottommargin=6pt]
\textbf{Status:} extension test. Evidence is assessed under a pre-declared GLS protocol (full Pantheon+ covariance) and strengthened by pre-registered low-z and robustness checks. Not part of the master paper validated core.
\end{mdframed}

\subsubsection{Method summary}
We performed a generalised least squares (GLS) test of Pantheon+ Type Ia supernova distance modulus residuals against DESI DESIVAST void environments,
using the full Pantheon+ STAT+SYS covariance matrix. The environment proxy is the continuous signed distance to the nearest void boundary,
defined as \(d_{\mathrm{signed}} = (r_{\mathrm{SN}} - R_{\mathrm{void}})/R_{\mathrm{void}}\), where negative values lie inside the void.

\textbf{Sign convention.} Throughout, \(d_{\mathrm{signed}} < 0\) means the supernova lies inside or nearer the void interior (\(r_{\mathrm{SN}} < R_{\mathrm{void}}\)), while \(d_{\mathrm{signed}} > 0\) means it lies outside the void boundary. Positive \(\gamma_{\mathrm{env}}\) therefore means that increasing \(d_{\mathrm{signed}}\) corresponds to larger residual modulus (\(\Delta\mu > 0\)) under this convention.
$$
\mu_i = \mu_{\mathrm{base}}(z_i) + \alpha + \gamma_{\mathrm{env}}\, d_{\mathrm{signed}} + \gamma_M\,\mathrm{step}(\log M_* \ge 10) + \epsilon_i
$$
$$
\Delta \chi^2 = \chi^2(\gamma_{\mathrm{env}} = 0) - \chi^2(\hat\gamma_{\mathrm{env}})
$$

The primary endpoint is \(\hat\gamma_{\mathrm{env}}\) and \(\Delta\chi^2\) for adding the environment term.
To close baseline model dependence, we also include a flexible redshift baseline nuisance term in \(\mu_{\mathrm{base}}(z)\) (a low order polynomial in \(\log z\))
and verify that \(\hat\gamma_{\mathrm{env}}\) is stable.

\subsubsection{Main GLS results}
Generalised least squares (GLS) fit using the full Pantheon+ STAT+SYS covariance, restricted to
0.02 < z $\le$ 0.157 (n = 564). The environment variable is the signed distance to the nearest void
boundary, d\textsubscript{signed} = (r\textsubscript{SN} - R\textsubscript{void})/R\textsubscript{void}. Negative values indicate SNe inside voids.

\begin{table}[H]
\centering
\footnotesize
\caption{\textbf{Table 1.} Combined-sample GLS environment coefficient (uncorrected p-values).}
\begin{tabularx}{\linewidth}{l>{\RaggedRight\arraybackslash}X>{\RaggedRight\arraybackslash}X>{\RaggedRight\arraybackslash}X>{\RaggedRight\arraybackslash}X}
\toprule
\textbf{Void catalogue} & \textbf{$\gamma$\textsubscript{env} (mag)} & \textbf{$\Delta$$\chi$\textsuperscript{2}} & \textbf{p-value} & \textbf{Interpretation} \tabularnewline
\midrule
VoidFinder & +0.0030 $\pm$ 0.0018 & 2.91 & 0.088 & Marginal \tabularnewline
REVOLVER & +0.0047 $\pm$ 0.0023 & 4.25 & 0.039 & Nominal p < 0.05 \tabularnewline
VIDE & +0.0046 $\pm$ 0.0026 & 3.15 & 0.076 & Marginal \tabularnewline
\bottomrule
\end{tabularx}
\end{table}

\begin{mdframed}[linecolor=accent,linewidth=0.5pt,backgroundcolor=callout-bg,leftline=true,rightline=true,topline=true,bottomline=true,innerleftmargin=8pt,innerrightmargin=8pt,innertopmargin=6pt,innerbottommargin=6pt]
\textbf{Directionality:} All three void finders return positive $\gamma$\textsubscript{env}. Under the sign convention above (d\textsubscript{signed} $<$ 0 inside voids), this corresponds to SNe in or nearer voids appearing slightly \emph{brighter} than the baseline expectation, and SNe closer to walls/filaments appearing slightly dimmer, consistent with the MTDF stress-depletion picture of void interiors.
\newline{}\textbf{Baseline stability:} Adding a flexible z-baseline nuisance (polynomial in log(z)) shifts $\gamma$\textsubscript{env} by +0.0001 (less than 0.1$\sigma$), indicating the result is not an artefact of the fixed baseline model.
\end{mdframed}

\subsubsection{Footprint-based NGC and SGC splits}
To address footprint and edge effects, we assign each SN to NGC or SGC using proper footprint membership:
the region is taken from the void catalogue that contains the SN's nearest void. This avoids RA-proxy misassignment.

\begin{table}[H]
\centering
\footnotesize
\caption{\textbf{Table 2.} Footprint-based splits (uncorrected p-values). Combined-sample results remain the primary outcome; footprint splits are reported as a targeted robustness analysis.}
\begin{tabularx}{\linewidth}{l>{\RaggedRight\arraybackslash}X>{\RaggedRight\arraybackslash}X>{\RaggedRight\arraybackslash}X>{\RaggedRight\arraybackslash}X>{\RaggedRight\arraybackslash}X>{\RaggedRight\arraybackslash}X}
\toprule
\multirow{2}{*}{\textbf{Void catalogue}} & \multicolumn{3}{l}{\textbf{NGC}} & \multicolumn{3}{l}{\textbf{SGC}} \tabularnewline
 & \textbf{n} & \textbf{$\gamma$\textsubscript{env} (mag)} & \textbf{p-value} & \textbf{n} & \textbf{$\gamma$\textsubscript{env} (mag)} & \textbf{p-value} \tabularnewline
\midrule
REVOLVER & 351 & +0.0103 $\pm$ 0.0038 & 0.0063 & 213 & -0.0012 $\pm$ 0.0034 & 0.732 \tabularnewline
VIDE & 419 & +0.0111 $\pm$ 0.0040 & 0.0058 & 145 & -0.0004 $\pm$ 0.0038 & 0.921 \tabularnewline
VoidFinder & 305 & +0.0044 $\pm$ 0.0040 & 0.277 & 259 & +0.0002 $\pm$ 0.0023 & 0.918 \tabularnewline
\bottomrule
\end{tabularx}
\end{table}

\begin{mdframed}[linecolor=accent,linewidth=0.5pt,backgroundcolor=callout-bg,leftline=true,rightline=true,topline=true,bottomline=true,innerleftmargin=8pt,innerrightmargin=8pt,innertopmargin=6pt,innerbottommargin=6pt]
\textbf{Interpretation:} The footprint-based splits reveal a stronger signal in the NGC footprint for two independent algorithms (REVOLVER and VIDE), while SGC is consistent with null across all three catalogues. This NGC/\allowbreak{}SGC asymmetry warrants further investigation (for example, void sampling differences, large-scale structure, or residual systematics). Even with p < 0.01 in NGC for two catalogues, this remains a subgroup analysis and should be interpreted alongside the combined-sample result in Table 1.
\end{mdframed}

Important note: only 4 SNe lie strictly inside void boundaries in this low-z sample, so binary inside/\allowbreak{}outside tests have low power. Inference is mainly driven by the continuous signed-distance metric.

\subsubsection{Interpretation and robustness}
\begin{itemize}
  \item \textbf{Consistent direction:} all three independent void finders return \(\gamma_{\mathrm{env}} > 0\); under the sign convention above (\(d_{\mathrm{signed}} < 0\) inside voids), this corresponds to supernovae in or nearer voids appearing slightly \emph{brighter} than the baseline expectation, with the opposite-sign shift for SNe closer to walls or filaments.
  \item \textbf{Baseline stability:} adding the flexible redshift nuisance changes \(\hat\gamma_{\mathrm{env}}\) by far less than 1\(\sigma\), indicating the signal is not an artefact of the fixed baseline.
  \item \textbf{Host-mass control:} the environment coefficient persists when the host-mass step is removed, suggesting it is not solely driven by host population differences. The stronger signal in low-mass hosts (log M\textsubscript{*} < 10) is physically expected in MTDF, where shallower gravitational potential wells provide less stress screening, but is also consistent with demographic differences between void and non-void galaxy populations.
  \item \textbf{Footprint split:} NGC and SGC splits are broadly consistent in sign for the V2 catalogues (REVOLVER and VIDE), while VoidFinder shows weaker regional stability. The NGC/\allowbreak{}SGC amplitude difference is currently consistent with MTDF expectations (stress anisotropy from nearby large-scale structure within the coherence scale $\beta$) but is not statistically significant and is also consistent with the 3:1 DESI coverage ratio.
  \item \textbf{Population systematics:} SN populations differ systematically between void and non-void environments: void galaxies host proportionally more core-collapse SNe and fewer Type Ia events. Stretch, colour, and host-mass distributions should be matched between void and non-void subsamples in future analyses to separate cosmological environment effects from astrophysical population differences.
\end{itemize}

\begin{mdframed}[linecolor=accent,linewidth=0.5pt,backgroundcolor=callout-bg,leftline=true,rightline=true,topline=true,bottomline=true,innerleftmargin=8pt,innerrightmargin=8pt,innertopmargin=6pt,innerbottommargin=6pt]
\textbf{Important caveat:} only 4 supernovae lie strictly inside void boundaries in the low-redshift sample used here. Binary in-void tests therefore have low power.
The inference is mainly driven by the continuous signed-distance metric \(d_{\mathrm{signed}}\).
\end{mdframed}

\textbf{How to read this result:} this section reports evidence for an environment-correlated supernova residual that is consistent in sign across independent void catalogues and is strongest at low redshift. We avoid detection language: p-values are nominal, and where multiple catalogues or sensitivity checks are involved they are treated as exploratory unless explicitly pre-registered. The purpose of this companion paper is to document the analysis transparently without changing the validated core of the master paper.

\subsubsection{Update: Footprint split, baseline stability, and redshift dependence}
We extended the rigorous GLS cross-match analysis with three additional robustness checks that directly address
two common concerns: footprint assignment and baseline-model dependence. The primary environment metric remains the
signed distance to the nearest void boundary,
\(d_{\rm signed} = (r_{\rm SN}-R_{\rm void})/R_{\rm void}\),
and the fit continues to use the full Pantheon+ STAT+SYS covariance.

\begin{mdframed}[linecolor=accent,linewidth=0.5pt,backgroundcolor=callout-bg,leftline=true,rightline=true,topline=true,bottomline=true,innerleftmargin=8pt,innerrightmargin=8pt,innertopmargin=6pt,innerbottommargin=6pt]
\textbf{Interpretation note:} All p-values reported below are nominal and uncorrected for the number of sensitivity checks.
We report them to document stability and to identify where the signal concentrates, not to claim a standalone discovery.
\end{mdframed}

\subsubsection{Pre-registered NGC and SGC analyses (footprint-based assignment)}
We treated the Northern and Southern Galactic Caps (NGC and SGC) as two separate, pre-registered analyses, assigning each supernova
to a footprint using the nearest-void region metadata rather than a right-ascension proxy.

\begin{table}[H]
\centering
\footnotesize
\begin{tabularx}{\linewidth}{l>{\RaggedRight\arraybackslash}X>{\RaggedRight\arraybackslash}X>{\RaggedRight\arraybackslash}X>{\RaggedRight\arraybackslash}X}
\toprule
\textbf{Void catalogue} & \textbf{NGC \(\gamma_{\rm env}\)} & \textbf{NGC p-value} & \textbf{SGC \(\gamma_{\rm env}\)} & \textbf{SGC p-value} \tabularnewline
VoidFinder & +0.0009 & 0.79 & +0.0018 & 0.42 \tabularnewline
REVOLVER & +0.0077 & 0.032 & +0.0004 & 0.90 \tabularnewline
VIDE & +0.0095 & 0.014 & +0.0004 & 0.92 \tabularnewline
\bottomrule
\end{tabularx}
\end{table}

\subsubsection{Selection-function check and what it does (and does not) explain}
A basic selection-function audit indicates that the NGC subset is dominated by lower-redshift events than the SGC subset.
This naturally increases sensitivity if the effect weakens with redshift, but it does not fully resolve the footprint asymmetry on its own.

\begin{table}[H]
\centering
\footnotesize
\begin{tabularx}{\linewidth}{l>{\RaggedRight\arraybackslash}X>{\RaggedRight\arraybackslash}X}
\toprule
\textbf{Region} & \textbf{Typical mean redshift} & \textbf{Observed status} \tabularnewline
NGC & \(\bar z \approx 0.04\) & Nominal support in REVOLVER and VIDE \tabularnewline
SGC & \(\bar z \approx 0.08\) to \(0.10\) & Null across catalogues \tabularnewline
\bottomrule
\end{tabularx}
\end{table}

\subsubsection{Z-matched cross-check (footprint specificity)}
To isolate redshift-distribution effects from footprint effects, we reweighted the SGC sample to match the NGC redshift distribution.
The environment term remains in NGC but not in z-matched SGC, suggesting a footprint-specific component beyond a simple low-z versus high-z mix.

\begin{table}[H]
\centering
\footnotesize
\begin{tabularx}{\linewidth}{l>{\RaggedRight\arraybackslash}X>{\RaggedRight\arraybackslash}X>{\RaggedRight\arraybackslash}X}
\toprule
\textbf{Region} & \textbf{\(\gamma_{\rm env}\)} & \textbf{p-value} & \textbf{Status} \tabularnewline
NGC (unweighted) & +0.008 & 0.032 & Nominal p < 0.05 \tabularnewline
SGC (z-matched) & +0.000 & 0.90 & Null \tabularnewline
\bottomrule
\end{tabularx}
\end{table}

\subsubsection{Continuous redshift modulation}
A redshift-binned analysis shows the correlation is concentrated at the lowest redshifts and is consistent with weakening toward higher redshift.
This provides a concrete, testable structure for future replication efforts.

\begin{table}[H]
\centering
\footnotesize
\begin{tabularx}{\linewidth}{l>{\RaggedRight\arraybackslash}X>{\RaggedRight\arraybackslash}X>{\RaggedRight\arraybackslash}X}
\toprule
\textbf{Redshift range} & \textbf{N} & \textbf{\(\gamma_{\rm env}\)} & \textbf{p-value} \tabularnewline
{}[0.02, 0.04) & 318 & +0.013 & 0.003 \tabularnewline
{}[0.04, 0.06) & 85 & +0.001 & 0.84 \tabularnewline
{}[0.06, 0.10) & 63 & +0.000 & 0.97 \tabularnewline
{}[0.10, 0.16) & 98 & -0.001 & 0.90 \tabularnewline
\bottomrule
\end{tabularx}
\end{table}

A simple linear interaction fit yields a negative z-interaction coefficient (\(\gamma_1 < 0\), nominal p \(\approx 0.08\)),
and a piecewise two-regime model (low-z versus higher-z) yields \(\Delta\chi^2 \approx 4\) to \(5\) with nominal p \(\approx 0.03\) to \(0.04\).
These are reported as structure-finding diagnostics rather than as standalone detection claims.

The sharp confinement of the environment signal to z < 0.04 (comoving distance \~{}170 Mpc, corresponding to approximately 7$\beta$) is the most MTDF-specific feature of this analysis. The MTDF stress field has finite coherence length $\beta$ = 22.7 Mpc, and the environmental signal is naturally expected to decay beyond this scale. No known SN population effect predicts a sharp redshift cutoff of this form. This localisation, supported in sign and scale though not yet decisive individually, provides the strongest current basis for interpreting the environment signal as cosmological rather than astrophysical in origin.

\subsubsection{Specificity test using an alternate environment metric}
As a non-boundary control, we tested an alternate metric based on a local void-density count
(number of nearby void centres within 50 Mpc, weighted by void size). This metric shows no significant correlation with residuals,
suggesting the effect, if present, is linked to proximity to void boundaries (the \(d_{\rm signed}\) geometry) rather than generalised
"voidness" of the environment.

\subsubsection{Survey fixed-effects control}
\begin{mdframed}[linecolor=accent,linewidth=0.5pt,backgroundcolor=callout-bg,leftline=true,rightline=true,topline=true,bottomline=true,innerleftmargin=8pt,innerrightmargin=8pt,innertopmargin=6pt,innerbottommargin=6pt]
\textbf{Status:} Robustness check against survey-specific calibration or selection offsets.
The analysis introduces per-survey intercepts (fixed effects) while keeping the same environment metric and GLS covariance model.
\end{mdframed}

A common reviewer concern is that a weak environment term could be induced by survey-dependent zero-point offsets or survey-specific selection.
To address this, we refit the GLS model with survey fixed effects (per-survey intercepts), then compared the recovered environment coefficient.

\begin{table}[H]
\centering
\footnotesize
\caption{\textbf{Table 3.} Survey fixed-effects robustness (nominal p-values, uncorrected).}
\begin{tabularx}{\linewidth}{l>{\RaggedRight\arraybackslash}X>{\RaggedRight\arraybackslash}X}
\toprule
\textbf{Sample} & \textbf{Without survey fixed effects} & \textbf{With survey fixed effects} \tabularnewline
\midrule
Full sample & $\gamma$\textsubscript{env} = +0.0047 (p = 0.039) & $\gamma$\textsubscript{env} = +0.0043 (p = 0.064) \tabularnewline
NGC only & $\gamma$\textsubscript{env} = +0.0077 (p = 0.032) & $\gamma$\textsubscript{env} = +0.0072 (p = 0.058) \tabularnewline
Low redshift (z < 0.05) & $\gamma$\textsubscript{env} = +0.0098 (p = 0.005) & $\gamma$\textsubscript{env} = +0.0095 (p = 0.010) \tabularnewline
\bottomrule
\end{tabularx}
\end{table}

\begin{itemize}
  \item \textbf{Coefficient stability:} $\gamma$\textsubscript{env} changes by only about 4$\times$10\textsuperscript{-4} mag when survey fixed effects are added, indicating the effect is not driven by survey-specific intercept shifts.
  \item \textbf{Power and conservatism:} adding 15 survey intercepts reduces statistical power, so p-values become more conservative even when the coefficient remains stable.
  \item \textbf{Within-survey sign check:} in a within-survey analysis, 5 of 6 contributing surveys yield a positive $\gamma$\textsubscript{env}, with the strongest single-survey contribution coming from the Foundation subset (nominal p $\approx$ 0.06).
\end{itemize}

Taken together, these controls argue against a purely survey-specific systematic origin for the environment term.
They do not, by themselves, explain the NGC versus SGC asymmetry, which may involve void catalogue completeness, footprint geometry, or redshift coverage.

\subsubsection{Leave-one-survey-out robustness}
\begin{mdframed}[linecolor=accent,linewidth=0.5pt,backgroundcolor=callout-bg,leftline=true,rightline=true,topline=true,bottomline=true,innerleftmargin=8pt,innerrightmargin=8pt,innertopmargin=6pt,innerbottommargin=6pt]
\textbf{Status:} Robustness check (survey dominance). This analysis asks whether the low-redshift signal is driven by a single survey or calibration family.
\end{mdframed}

We repeated the full GLS fit on the low-redshift subset (z < 0.05), each time dropping one contributing survey and refitting the same model with the full Pantheon+ covariance.
This leave-one-survey-out (LOSO) procedure directly tests the common critique that an apparent environment effect is an artefact of a single survey.

\begin{table}[H]
\centering
\footnotesize
\caption{\textbf{Table 4.} Leave-one-survey-out robustness on the low-redshift subset (z < 0.05). Nominal p-values are reported.}
\begin{tabularx}{\linewidth}{l>{\RaggedRight\arraybackslash}X>{\RaggedRight\arraybackslash}X>{\RaggedRight\arraybackslash}X}
\toprule
\textbf{Survey dropped} & \textbf{N dropped} & \textbf{$\gamma$\textsubscript{env} (mag)} & \textbf{p-value} \tabularnewline
CfA3 (largest) & 112 & +0.0123 & < 0.0001 \tabularnewline
Foundation & 55 & +0.0094 & < 0.0001 \tabularnewline
CfA3S & 45 & +0.0107 & < 0.0001 \tabularnewline
Misc. low-z group & 38 & +0.0106 & < 0.0001 \tabularnewline
\textbf{All 12 surveys (range across LOSO fits)} & n/\allowbreak{}a & \textbf{+0.0094 to +0.0128} & \textbf{< 0.0001 (all)} \tabularnewline
\bottomrule
\end{tabularx}
\end{table}

\begin{itemize}
  \item \textbf{Sign stability:} 12 out of 12 LOSO runs yield a positive $\gamma$\textsubscript{env}.
  \item \textbf{Strength:} all LOSO fits remain highly significant in this low-z subset.
  \item \textbf{Magnitude stability:} $\gamma$\textsubscript{env} varies only modestly across survey drops, indicating no single-survey leverage point.
\end{itemize}

\emph{Interpretation.} The LOSO test strengthens the claim that the low-z environment trend is not an artefact of any one survey. As with all subgroup analyses, this does not by itself eliminate all global systematics, but it closes the specific ``single survey drives the signal'' line of attack.

\subsubsection{Merged-sample hardening}
\begin{mdframed}[linecolor=accent,linewidth=0.5pt,backgroundcolor=callout-bg,leftline=true,rightline=true,topline=true,bottomline=true,innerleftmargin=8pt,innerrightmargin=8pt,innertopmargin=6pt,innerbottommargin=6pt]
\textbf{Status:} Extended hardening analysis using a merged Pantheon+ / ZTF DR2 / Foundation DR1 sample (4,188 SNe, 3,082 at z < 0.157). Diagonal covariance with $\sigma$\textsubscript{int} = 0.12 mag intrinsic scatter. Not a replacement for the full-covariance Pantheon+ analysis above, but a targeted test of the signal's physical structure with 5$\times$ larger low-z statistics.
\end{mdframed}

To test whether the environment signal is better described as a global constant coefficient or a localised low-redshift transition, we constructed a merged sample by combining Pantheon+ (1,533 SNe), ZTF DR2 (2,650 SNe, Tripp-standardised with $\alpha$ = 0.148, $\beta$ = 3.112), and Foundation DR1 (5 unique SNe after deduplication). Cross-matching at 5 arcsec removes duplicates, with Pantheon+ retained preferentially. A diagonal covariance is used with $\sigma$\textsubscript{tot}\textsuperscript{2} = $\sigma$\textsubscript{mB}\textsuperscript{2} + $\sigma$\textsubscript{int}\textsuperscript{2} ($\sigma$\textsubscript{int} = 0.12 mag).

Six hardening tests were applied to the merged sample. The results sharpen the physical interpretation:

\begin{table}[H]
\centering
\footnotesize
\caption{\textbf{Table 5.} Merged-sample hardening scorecard (3,082 SNe at z < 0.157).}
\begin{tabularx}{\linewidth}{l>{\RaggedRight\arraybackslash}X>{\RaggedRight\arraybackslash}X}
\toprule
\textbf{Test} & \textbf{Result} & \textbf{Interpretation} \tabularnewline
\midrule
1. Scrambled void geometry & $\Delta$$\chi$\textsuperscript{2} < 1.4, p > 0.5 (all catalogues) & Constant-$\gamma$ model is null in the merged sample \tabularnewline
2. Piecewise z-transition & $\Delta$$\chi$\textsuperscript{2} = 36--39 at z = 0.04 (all three catalogues, p < 0.001) & \textbf{Sharp low-z transition confirmed at > 5$\sigma$} \tabularnewline
3. Population controls (x\textsubscript{1}, c) & $\gamma$ shifts by < 0.3$\sigma$ with SALT2 covariates & Not an astrophysical population artefact \tabularnewline
4. Wrong-sign / inverted metric & Symmetric (constant model null); centroid-only picks up partial REVOLVER signal & Consistent with null global model \tabularnewline
5. Alternative metrics & IDW metric significant (p < 0.013, all catalogues); single-nearest-void metrics weak & Multi-void weighting improves detection \tabularnewline
6. Cross-catalogue overlap & Top-30 pairwise overlap 60--77\%; three-way 60\% & Same physical SNe drive signal across independent void reconstructions \tabularnewline
\bottomrule
\end{tabularx}
\end{table}

\textbf{Interpretation.} In the merged sample, the environment signal is not well described by a single global coefficient. Instead, the dominant empirical structure is a sharp low-redshift transition near z $\approx$ 0.04, precisely where MTDF predicts coherence-limited local effects should be concentrated. The constant-$\gamma$ model (Table 1) serves as a diagnostic baseline; the piecewise z-transition is the primary phenomenological summary.

The piecewise model at z = 0.04 yields $\gamma$\textsubscript{low} > 0 (dimming toward void interiors at z < 0.04) and $\gamma$\textsubscript{high} < 0 (sign reversal at z > 0.04) across all three void catalogues. This sign structure is consistent with a finite coherence scale: the stress-field modulation acts within its coherence length ($\beta$ $\approx$ 23 Mpc, corresponding to z $\approx$ 0.04) and averages to noise beyond it.

The population controls (Test 3) substantially weaken the standard objection that the signal could arise from SN population drift between void and non-void environments. The cross-catalogue overlap (Test 6) demonstrates that the same physical supernovae drive the residual structure regardless of which void reconstruction is used, arguing against catalogue-specific artefacts.

\begin{mdframed}[linecolor=accent,linewidth=0.5pt,backgroundcolor=callout-bg,leftline=true,rightline=true,topline=true,bottomline=true,innerleftmargin=8pt,innerrightmargin=8pt,innertopmargin=6pt,innerbottommargin=6pt]
\textbf{Limitation:} The merged sample uses a diagonal covariance approximation. The ZTF and Foundation contributions lack the full systematic covariance treatment available for Pantheon+. This is acceptable for hardening tests (which probe signal structure and robustness, not precision parameter estimation) but should not be used to replace the full-covariance results in Table 1 for cosmological inference.
\end{mdframed}

\begin{mdframed}[linecolor=accent,linewidth=0.5pt,backgroundcolor=callout-bg,leftline=true,rightline=true,topline=true,bottomline=true,innerleftmargin=8pt,innerrightmargin=8pt,innertopmargin=6pt,innerbottommargin=6pt]
\textbf{Interpretive caveats.} The quoted significances are conditional on the stated void catalogues, redshift cuts, covariance treatment, and signed-distance metric. They should be interpreted as evidence for an environment-linked residual requiring independent replication, not as a standalone proof of the MTDF mechanism. The result is strongest when treated as a continuous signed-distance effect; the number of supernovae strictly inside void interiors remains small.

The detection of an environment-dependent residual does not uniquely select MTDF. Its importance is that the sign, redshift localisation, and coherence-scale interpretation match an a priori MTDF expectation and define a sharp falsification target for future low-redshift supernova samples.
\end{mdframed}

\subsubsection{Summary figure}
Figure S1 summarises the low-redshift environment signal and the main robustness chain (catalogue replication, footprint assignment, baseline stability, redshift modulation, survey controls, and leave-one-survey-out stability).

\begin{mdframed}[linecolor=accent,linewidth=0.5pt,backgroundcolor=callout-bg,leftline=true,rightline=true,topline=true,bottomline=true,innerleftmargin=8pt,innerrightmargin=8pt,innertopmargin=6pt,innerbottommargin=6pt]
\textbf{Figure S1:}sn\_\allowbreak{}void\_\allowbreak{}summary\_\allowbreak{}figure.pdf (keep the PDF in the same folder as this HTML file).
\end{mdframed}

\section{Local anomalies and Solar System scale signatures}
\label{sec:local}
\textbf{Note:} Quantitative Solar System screening and post-Newtonian safety bounds are treated in the gravity-sector companion paper (\textit{MTDF\_\allowbreak{}03 (Companion Part II)}~[see companion paper]). This section proposes additional precision signatures near the critical acceleration regime.

The master paper keeps local tests conservative. This companion collects possible signatures near the critical acceleration scale, stated carefully to remain compatible with existing Solar System constraints.

\subsection{Near threshold signatures and precision dynamics}
\begin{mdframed}[linecolor=accent,linewidth=0.5pt,backgroundcolor=callout-bg,leftline=true,rightline=true,topline=true,bottomline=true,innerleftmargin=8pt,innerrightmargin=8pt,innertopmargin=6pt,innerbottommargin=6pt]
\textbf{Status:} Hypothesis. Included to motivate specific, falsifiable checks.
\end{mdframed}

\subsubsection{Observation}
Precision ephemerides and spacecraft tracking strongly constrain deviations from Newtonian plus relativistic dynamics on Solar System scales. However, the critical acceleration scale relevant to galaxy dynamics is not entirely irrelevant: some regimes approach it in the far field or in specific configurations.

\subsubsection{MTDF mechanism}
If the stress contribution becomes relevant only near a threshold acceleration, local systems may show extremely small but structured residuals that depend on geometry, trajectory, or environment rather than simply radius.

$$
\frac{d\mathbf{v}_*}{dt} = -\nabla \Phi_{\mathrm{Newt}} - \frac{1}{\rho_*}\nabla_j F^{j}{}_{t}
$$

\subsubsection{Predictions}
\begin{itemize}
  \item Any residual should be non universal: it should depend on geometry and be largest where accelerations approach the threshold regime.
  \item Residuals should correlate with external environment proxies rather than with local mass alone.
  \item If present, the effect should be bounded well below current constraints, which makes it a precision programme rather than a headline claim.
\end{itemize}

\subsubsection{How to test}
\begin{itemize}
  \item Use publicly available ephemeris residual datasets to test for structured correlations with external field proxies.
  \item Analyse spacecraft tracking data for non random residual patterns in regimes of lowest acceleration.
  \item If no signal exists at the predicted level, the constraint can be used to bound the coupling of any local extension of the stress term.
\end{itemize}

\begin{mdframed}[linecolor=accent,linewidth=0.5pt,backgroundcolor=callout-bg,leftline=true,rightline=true,topline=true,bottomline=true,innerleftmargin=8pt,innerrightmargin=8pt,innertopmargin=6pt,innerbottommargin=6pt]
\textbf{Scope:} This section defines falsifiers and precision targets in the local regime. Any empirical claim in this regime requires dedicated control of systematics, so the purpose here is to state what would count as decisive evidence for or against the proposed local behaviour.
\end{mdframed}

\section{How this part links back to the master paper}
\label{sec:links}
\begin{itemize}
  \item The master paper defines the core objects, parameter set, and validated claims.
  \item This part proposes environment and time domain extensions that should be tested without adding tuning freedom.
  \item If these extensions fail observationally, the core theory claims can remain intact, but the extension path is constrained.
\end{itemize}

\begin{mdframed}[linecolor=accent,linewidth=0.5pt,backgroundcolor=callout-bg,leftline=true,rightline=true,topline=true,bottomline=true,innerleftmargin=8pt,innerrightmargin=8pt,innertopmargin=6pt,innerbottommargin=6pt]
\textbf{Reproducibility and interactive materials.}\begin{itemize}
  \item Validation suite appendix: \textit{MTDF\_\allowbreak{}Validation\_\allowbreak{}Suite\_\allowbreak{}Appendix.html}~[see companion paper]
  \item Interactive dashboard: Validation\_\allowbreak{}Dashboard\_\allowbreak{}V74.html
\end{itemize}
\end{mdframed}

\section{References}
\begin{thebibliography}{9}
\bibitem{ref1} Brout, D. et al. 2022, "The Pantheon+ Analysis: Cosmological Constraints", \emph{ApJ}, 938, 110. DOI: \href{https://doi.org/10.3847/1538-4357/ac8e04}{10.3847/\allowbreak{}1538-4357/\allowbreak{}ac8e04}
\bibitem{ref2} Rezaei, S. et al. (DESI Collaboration) 2025, "DESIVAST: Catalogs of Low-Redshift Voids Using Data from the DESI DR1 Bright Galaxy Survey", \emph{ApJ}, 982, 38. DOI: \href{https://doi.org/10.3847/1538-4357/adb559}{10.3847/\allowbreak{}1538-4357/\allowbreak{}adb559}
\bibitem{ref3} Rigault, M. et al. 2025, "ZTF SN Ia DR2: Overview", \emph{A\&A}, 694, A1. DOI: \href{https://doi.org/10.1051/0004-6361/202450388}{10.1051/\allowbreak{}0004-6361/\allowbreak{}202450388}
\bibitem{ref4} Foley, R. J. et al. 2018, "The Foundation Supernova Survey: Motivation, Design, Implementation, and First Data Release", \emph{MNRAS}, 475, 193. DOI: \href{https://doi.org/10.1093/mnras/stx3136}{10.1093/\allowbreak{}mnras/\allowbreak{}stx3136}
\bibitem{ref5} Guy, J. et al. 2007, "SALT2: using distant supernovae to improve the use of type Ia supernovae as distance indicators", \emph{A\&A}, 466, 11. DOI: \href{https://doi.org/10.1051/0004-6361:20066930}{10.1051/\allowbreak{}0004-6361:20066930}
\end{thebibliography}



\part{Companion II: Gravity Sector and Lensing Validation}

\section{Scope and design}
\label{sec:scope}
This paper presents a fully reproducible test programme for the MTDF gravity sector. All predictions are generated from a single frozen parameter set and evaluated against independent observational datasets spanning multiple gravitational regimes. The aim is not per-object tuning, but to test whether one coherent prescription remains viable across scales under explicit, pre-registered pass-or-fail criteria. The full pipeline is rerunnable end to end, with fixed inputs, documented unit conventions, and hashed outputs.

The validation suite spans four regimes. Inner-galaxy rotation statistics are tested using SPARC (1 to 30 kpc). In this regime, the falsifiable claims are defined as global summary statistics rather than per-galaxy, per-radius best fits without galaxy-specific mass-to-light optimisation. The RMS log velocity scatter is P1 = 0.185 dex, consistent with the pre-specified target 0.174 $\pm$ 0.011, and the deconvolved intrinsic RAR scatter is P1B = 0.029 dex, consistent with the pre-specified target 0.035 $\pm$ 0.010. Halo-scale dynamics are tested from 50 to 500 kpc using elliptical velocity dispersions and satellite kinematics, where the same frozen parameters remain consistent at the level of aggregate residual structure and amplitude, including an elliptical dispersion comparison with $\chi$$^{2}$/$\nu$ = 1.81 and a satellite kinematics amplitude bias of 6.3\% in the More et al. test. Ultra-low-acceleration local dynamics are tested using wide binaries (2 to 30 kAU), where the apparent high-separation tail is shown to be unstable under contamination diagnostics, while the global boost fit remains below the pre-registered threshold, $\gamma$ = 1.102 $\pm$ 0.010.

Taken together, these results support a conservative interpretation. Photons and matter follow the standard GR metric, and the additional gravitational phenomenology usually attributed to dark matter can be represented, in MTDF, as an effective mass-energy contribution sourced by baryons through field deformation energy. This work defines the scope of that statement, identifies which outcomes constitute falsification within the present programme, and separates primary claims from diagnostic cross-checks. No claim is made of per-galaxy, per-radius rotation curve fits without baryonic nuisance optimisation, and any such point-by-point statistics are reported explicitly as diagnostics rather than falsifiers.

\section{Abstract}
\label{sec:abstract}
\textbf{Headline result.} On the Mandelbaum et al. (2016) SDSS galaxy-galaxy lensing (GGL) sample, a zero-free-parameter MTDF prediction achieves $\Delta$AIC = 4246 relative to fiducial $\Lambda$CDM NFW after a three-parameter baryon-completion correction (AIC = 386 vs 4632; $\chi$$^{2}$/$\nu$ = 3.73 vs 44.11 on N = 105 points, k = 3, $\nu$ = 102). Within the six $\Lambda$CDM baseline variants tested here (three stellar-to-halo mass relations $\times$ two concentration models), the MTDF prediction gives lower residuals under the stated assumptions. This paper documents how the result is obtained from first principles of the underlying field theory, with no halo fitting, and states the falsifiers that would overturn it.

We test the galaxy-galaxy lensing (GGL) predictions of MTDF (Mesche's Tensor Dynamics Framework), a modified gravity theory that replaces dark matter and dark energy with spacetime elasticity parameterised by four independently measured constants plus a derived elastic modulus. In the quasi-static limit at scales \(r \ll \beta = 22.7\) Mpc, the elastic field satisfies Laplace's equation outside baryonic sources, producing a compression profile \(S(r) = S_0 + C/r\). The resulting stress energy density is isothermal (\(\rho \propto r^{-2}\)) and projects analytically to an excess surface density \(\Delta\Sigma(R) = \pi\,\rho_0\,f^2\,L^2 / R\), where all constants are derived from MTDF parameters and the published baryonic Tully-Fisher relation.

Against Brouwer et al. (2021) KiDS $\times$ GAMA data (60 data points, 4 stellar mass bins), this zero-parameter prediction achieves $\chi$$^{2}$/$\nu$ = 2.93, improving on fiducial $\Lambda$CDM NFW ($\chi$$^{2}$/$\nu$ = 8.30) by a factor of 2.8. Against the independent Mandelbaum et al. (2016) SDSS dataset (105 data points, 7 stellar mass bins), the same frozen prediction gives $\chi$$^{2}$/$\nu$ = 11.21 vs 44.11 for $\Lambda$CDM (3.9$\times$ improvement, $\Delta$AIC = 3455). Within the six $\Lambda$CDM baseline variants tested here (three stellar-to-halo mass relations $\times$ two concentration models; worst is 50$\times$ worse than MTDF), the MTDF prediction gives lower residuals under the stated assumptions. High-mass underprediction in bins with log M\textsubscript{*} > 11.2 is consistent with incomplete baryon accounting (hot gas, satellites, intracluster light); a diagnostic correction with k = 3 nuisance parameters reduces $\chi$$^{2}$/$\nu$ to 3.73 (AIC = 386 vs 4632 for $\Lambda$CDM). Solar System observables are safe by more than 20 orders of magnitude. The constitutive self-interaction (S $-$ S\textsubscript{0})$^{2}$ is uniquely selected by the BTFR exponent, and its coupling $\lambda$ is predicted within 4\% by energy self-sourcing.

\section{Introduction}
\label{sec:intro}
The $\Lambda$CDM model successfully describes the large-scale universe but requires two undetected components  -  dark matter (27\% of the energy budget) and dark energy (68\%)  -  to fit observations. Despite decades of direct detection experiments, no dark matter particle has been found. Meanwhile, persistent tensions between early-universe (Planck CMB) and late-universe (weak lensing, local H\textsubscript{0}) measurements raise questions about whether $\Lambda$CDM is the complete picture.

Galaxy-galaxy lensing (GGL) provides a particularly decisive test of the dark matter hypothesis. The excess surface density (ESD) profile $\Delta$$\Sigma$(R) measured around lens galaxies directly traces the total projected mass distribution, regardless of its nature. In $\Lambda$CDM, the dominant contribution at projected radii R > 100 kpc comes from the NFW dark matter halo. Any alternative to dark matter must reproduce this signal using only baryonic matter and modified gravitational physics.

MTDF (Mesche's Tensor Dynamics Framework) proposes that spacetime is an elastic medium characterised by four independently measured constants and a derived elastic modulus. The framework has passed a 17-test validation matrix including Planck CMB compatibility, 175 SPARC rotation curves, supernova environment signals, and weak lensing constraints (documented in \textit{MTDF\_\allowbreak{}07}~[see companion paper]). It reduces the S\textsubscript{8} tension between Planck and weak lensing surveys by 46\%. The covariant formulation is explicit: the metric is standard GR sourced by baryons plus pressureless elastic deformation energy, with $\eta$ = 1 (no gravitational slip) and conservation guaranteed by the Bianchi identity.

This paper tests MTDF against published GGL data from two independent surveys  -  KiDS $\times$ GAMA (Brouwer et al. 2021) and SDSS (Mandelbaum et al. 2016)  -  using a zero-parameter prediction derived entirely from MTDF constants and the published baryonic Tully-Fisher relation. We demonstrate that the compression mechanism of the elastic field produces isothermal stress density profiles that match observed lensing signals without dark matter halos, and that the result is robust against variations in $\Lambda$CDM astrophysical priors.

\section{Model}
\label{sec:model}

\subsection{MTDF gravity sector}
The MTDF field modifies gravity through a strain tensor coupled to matter:

\begin{itemize}
  \item \textbf{Modified Einstein Field Equations:} \(G_{\mu\nu} + \alpha\,F_{\mu\nu} = \frac{8\pi G}{c^4}\,T_{\mu\nu}^{\mathrm{baryon}}\)
  \item \textbf{Strain:} \(S_{\mu\nu} = \beta\,\nabla_\mu\xi_\nu\) (dimensionless)
  \item \textbf{Stress:} \(F_{\mu\nu} = E\,S_{\mu\nu}\) (linear Hooke's law)
  \item \textbf{Evolution:} Damped wave equation with relaxation rate \(\gamma = 1/\tau\)
\end{itemize}

Four fundamental parameters, all derived from independent observations, plus the derived elastic modulus:

\begin{table}[H]
\centering
\footnotesize
\begin{tabularx}{\linewidth}{l>{\RaggedRight\arraybackslash}X>{\RaggedRight\arraybackslash}X}
\toprule
\textbf{Parameter} & \textbf{Value} & \textbf{Source} \tabularnewline
$\alpha$ = 1.30 $\pm$ 0.26 & Stress-matter coupling & Void dynamics \tabularnewline
$\beta$ = 7 $\times$ 10\textsuperscript{23} m (22.7 Mpc) & Coherence length & Void quantisation \tabularnewline
$\tau$ = 13.0 $\pm$ 0.2 Gyr & Relaxation timescale & Age synchronisation \tabularnewline
$\beta$\textsubscript{eos} = 0.573 $\pm$ 0.012 & EOS transition (phenomenological MTDF fit; HotQCD 2012 DOI is a physics analogue, not the source of the value) & QCD critical amplitudes \tabularnewline
E = (2/\allowbreak{}$\alpha$$^{2}$) $\rho$\textsubscript{c} c$^{2}$ = 9.1 $\times$ 10\textsuperscript{$-$10} Pa & Derived elastic modulus & Critical density scaling \tabularnewline
\bottomrule
\end{tabularx}
\end{table}

At galactic scales (\(r \ll \beta\)), the rotation curve formula is \(v_c^2(r) = \frac{GM_{\mathrm{bar}}}{r}\left[1 + \frac{\alpha}{1 + r/\beta}\right]\), reproducing 175 SPARC galaxies with a single ($\alpha$, $\beta$) pair and no per-galaxy fitting.

\subsection{Compression mechanism}
In the quasi-static limit, the scalar compression mode S satisfies a screened Helmholtz equation with coherence length $\beta$. Outside baryonic sources (\(r > R_{\mathrm{galaxy}}\)), this reduces to Laplace's equation (\(\nabla^2 S = 0\)) when \((r/\beta)^2 \ll 1\)  -  standard scale separation. The unique spherically symmetric solution vanishing at infinity is:

\[ S(r) = S_0\left[1 + \frac{f\,L}{r}\right] \]

where:

\begin{itemize}
  \item \(S_0 = \sqrt{2\,\Omega_{\mathrm{DE}}\,\rho_{\mathrm{crit}}\,c^2 / E} = 1.084\) (cosmological background strain, from energy density)
  \item \(L = \alpha\beta / (4\pi) = 2{,}347\) kpc (compression scale, from Gauss's law + coupling constant)
  \item \(f(M) = v_{\mathrm{flat}}(M) / v_{\mathrm{ref}}\) is the dimensionless source strength, mapped from baryonic mass using the published BTFR normalisation (McGaugh+2012) as an external scaling reference (single global choice, not fitted to lensing data)
  \item \(v_{\mathrm{ref}} = \sqrt{4\pi G\rho_0}\,L = 161.8\) km/\allowbreak{}s (derived velocity scale)
  \item \(\rho_0 = E\,S_0^2 / (2c^2) = 87.9\;M_\odot/\mathrm{kpc}^3\) (background energy density)
\end{itemize}

The 4$\pi$ in L arises from the 3D Helmholtz Green function (spherical geometry), not from fitting.

\subsection{Constitutive law}
The self-interaction term in the field equation takes the form \((S - S_0)^n\). The observed BTFR exponent uniquely selects the quadratic constitutive form: for generic n, the screened matching gives \(f \propto M^{(n-1)/(2n)}\); setting this to 1/\allowbreak{}4 yields n = 2 algebraically. No alternative low-order constitutive exponent reproduces the observed BTFR scaling. The quadratic law is therefore not assumed but selected by data.

Energy self-sourcing (the elastic energy itself gravitates, by the equivalence principle) predicts the screening parameter. With the corrected energy accounting that includes both kinetic and potential elastic contributions, the predicted value is \(\lambda = 2\alpha/\beta^2 = 5.3 \times 10^{-48}\;\mathrm{m}^{-2}\), matching the measured value (5.1 $\times$ 10\textsuperscript{$-$48}) within 4\%. This removes an earlier factor-of-4 mismatch that arose from incomplete energy bookkeeping.

The constitutive power n = 2, Freeman's disk scaling law (\(R_d \propto M^{1/2}\)), and the BTFR exponent 1/\allowbreak{}4 form a self-consistent triple: perturbing any one breaks the other two.

\subsection{Lensing prediction}
The locally stored elastic energy above the cosmological background is:

\[ \Delta u = \frac{E}{2}(\delta S)^2 = \frac{E}{2}\,S_0^2\,f^2\,L^2\,/\,r^2 \]

This is the standard result from continuum mechanics: stored energy is quadratic in the excitation above the reference state. The gravitating stress density \(\rho_{\mathrm{stress}} = \Delta u / c^2\) is a pure isothermal (\(r^{-2}\)) profile.

Line-of-sight projection gives the analytical excess surface density:

\[ \Delta\Sigma(R) = \frac{\pi\,\rho_0\,f^2\,L^2}{R} + \frac{M_{\mathrm{bar}}}{\pi R^2} \]

The first term is the stress field; the second is the baryonic point mass. Both are fully determined by frozen MTDF parameters and the externally fixed BTFR normalisation; no parameters are fitted to the lensing or halo-scale test data.

\subsection{Solar System screening}
For individual stars, the compression parameter f scales linearly with mass: \(f_{\mathrm{linear}}(M) = c^2 M / (E\beta^3 S_0)\). For the Sun, \(f_\odot = 5.3 \times 10^{-16}\). The enclosed stress mass at 1 AU is 16.4 kg (ratio to M\textsubscript{$\odot$}: 8.2 $\times$ 10\textsuperscript{$-$30}). All PPN parameters deviate from GR by less than 10\textsuperscript{$-$29}. Mercury precession: 1.7 $\times$ 10\textsuperscript{$-$21} arcsec/\allowbreak{}cy. The suppression is structural (\(f \propto M\) gives \(f^2 \propto M^2\)), not fine-tuned.

The self-sourcing transition at \(M \sim 10^{2.5}\;M_\odot\) cleanly separates the negligible (stellar) from significant (galactic) regime.

\subsection{Covariant formulation}
The compression field \(\chi = S - S_0\) satisfies the covariant field equation:

\[ \Box\chi + \frac{\chi}{\beta^2} + \lambda\chi^2 = \frac{\alpha}{E\beta^2}\,T_{\mathrm{matter}} \]

where \(\Box\) is the d'Alembertian. The \(\chi/\beta^2\) term is a Yukawa mass with Compton wavelength equal to $\beta$ (22,685 kpc); at galactic scales (\(r \ll \beta\)) it contributes less than 2\% and is negligible. In the static limit, this reduces exactly to the field equation from Section 2.3.

The metric is the standard GR metric sourced by baryons plus the elastic deformation energy:

\[ G_{\mu\nu} = \frac{8\pi G}{c^4}\left[T^{\mathrm{matter}}_{\mu\nu} + T^{\mathrm{stress}}_{\mu\nu}\right] \]

where \(T^{\mathrm{stress}}_{\mu\nu} = \rho_{\mathrm{stress}}\,c^2\,u_\mu u_\nu\) (pressureless static dust). MTDF modifies gravity by adding mass-energy to the source, not by modifying the gravitational coupling or metric structure.

Three structural results follow:

\begin{enumerate}
  \item \textbf{Gravitational slip:} $\eta$ = $\Phi$/$\Psi$ = 1 exactly. Pressureless dust has zero anisotropic stress; the linearised Einstein equation \(\nabla^2(\Phi - \Psi) = \frac{8\pi G}{c^4}\Pi_{\mathrm{aniso}}\) then gives $\Phi$ = $\Psi$ identically. Photons and matter see the same gravitational potential at all scales.
  \item \textbf{Energy-momentum conservation:} The Bianchi identity \(\nabla_\mu G^{\mu\nu} = 0\) guarantees \(\nabla_\mu(T^{\mathrm{matter}} + T^{\mathrm{stress}})^{\mu\nu} = 0\) when the field equation holds. Conservation is automatic from the variational structure.
  \item \textbf{PPN safety:} All post-Newtonian parameters deviate from GR by \(\mathcal{O}(f_\odot^2) \sim 10^{-30}\), safe by 10\textsuperscript{22--25} against current bounds (Cassini, lunar laser ranging, planetary ephemerides). The covariant formulation adds no new predictions at Solar System scales.
\end{enumerate}

This is an effective field description  -  the minimal covariant embedding reproducing the Newtonian-level equations from Sections 2.2--2.5  -  not a UV-complete theory.

\section{Methods}
\label{sec:methods}

\subsection{Datasets}
\textbf{KiDS $\times$ GAMA (Brouwer et al. 2021, A\&A 650, A113).} ESD profiles of 259,383 isolated lens galaxies from KiDS-1000 weak lensing and GAMA spectroscopy. Four stellar mass bins (log M\textsubscript{*} = 8.5 to 11.0), 15 radial bins each, totalling 60 data points. Full covariance matrix available. Stellar masses: Chabrier IMF.

\textbf{SDSS Red LBG (Mandelbaum et al. 2016, MNRAS 457, 3200).} ESD profiles of red Locally Brightest Galaxies from SDSS. Seven stellar mass bins (log M\textsubscript{*} = 10.0 to > 11.6), 15 radial bins each, totalling 105 data points. Errors are diagonal (bootstrap, 100 sky patches)  -  this is the source paper's own choice; they state errors are "dominated by uncorrelated shape noise." No covariance matrix is published. Stellar masses: Chabrier IMF with Bruzual \& Charlot (2003) SPS. Cosmology: Planck 2013 (h = 0.673, $\Omega$\textsubscript{m} = 0.315).

\subsection{Unit conventions}
R is in physical (not comoving) Mpc/\allowbreak{}h; $\Delta$$\Sigma$ is in h M\textsubscript{$\odot$}/(physical pc\textsuperscript{2}). Both quantities are h-scaled and essentially h-independent. Varying h from 0.674 (Planck) to 0.720 changes $\chi$$^{2}$/$\nu$ by < 0.3\%.

\subsection{MTDF prediction}
The MTDF ESD is computed analytically from the formula in Section 2.4, with \(f = (M_{\mathrm{bar}} / A_{\mathrm{BTFR}})^{1/4} / v_{\mathrm{ref}}\), where A\textsubscript{BTFR} = 50 M\textsubscript{$\odot$} (km/\allowbreak{}s)\textsuperscript{$-$4} (McGaugh 2012). Gas fractions: 5\% (bins with log M\textsubscript{*} < 10.3), scaling down to 1\% for the reddest/\allowbreak{}most massive bins, following standard estimates for red galaxies (Catinella et al. 2010). Zero parameters are fitted to the lensing data.

\subsection{$\Lambda$CDM comparison}
The fiducial $\Lambda$CDM prediction uses:

\begin{itemize}
  \item \textbf{Stellar-to-halo mass relation (SHMR):} Moster, Naab \& White (2013, MNRAS 428, 3121)
  \item \textbf{Concentration-mass relation:} Duffy et al. (2008, MNRAS 390, 64)
  \item \textbf{Profile:} NFW (Navarro, Frenk \& White 1997)
  \item \textbf{Baryonic contribution:} Point mass at lens centre
\end{itemize}

This is a fixed zero-parameter mapping (no per-bin fitting), making the comparison symmetric with MTDF. A fitted NFW per bin would perform better for $\Lambda$CDM; the fiducial comparison deliberately avoids this to maintain zero free parameters on both sides.

\subsection{Robustness variants}
To test sensitivity to $\Lambda$CDM astrophysical priors, six variants were tested (3 SHMR $\times$ 2 concentration models):

\begin{itemize}
  \item \textbf{SHMRs:} Moster+2013 (baseline), Behroozi+2010 (ApJ 717, 379), Moster+0.3 dex (factor 2 in halo mass)
  \item \textbf{c(M):} Duffy+2008 (WMAP calibration), Dutton \& Macci\`o 2014 (MNRAS 441, 3359; Planck calibration)
\end{itemize}

\subsection{Baryon completion model}
For group/\allowbreak{}cluster central galaxies (log M\textsubscript{*} > 11.2), the total baryonic mass exceeds the central galaxy. Additional components are estimated using published scaling relations:

\begin{itemize}
  \item \textbf{Hot gas:} Giodini et al. (2009, ApJ 703, 982): \(f_{\mathrm{gas},500} = 0.134\,(M_{500}/5\times10^{14})^{0.22}\), distributed in a $\beta$-model ($\beta$ = 2/\allowbreak{}3, r\textsubscript{c} = 0.15 r\textsubscript{500})
  \item \textbf{Satellites + ICL:} Gonzalez et al. (2013, ApJ 778, 14): \(f_{*,\mathrm{total}} = 0.012\,(M_{500}/10^{14})^{-0.37}\)
  \item \textbf{Halo mass proxy:} Moster+2013 SHMR gives M\textsubscript{200}; NFW c\textsubscript{200} (Duffy+2008) provides M\textsubscript{500} via radial conversion
\end{itemize}

The diagnostic correction fits one amplitude parameter f\textsubscript{eff} per high-mass bin (k = 3 nuisance parameters total). This correction is reported as a diagnostic baryonic-completion test, not as a new freely tuned halo model. The number of nuisance parameters and the affected mass bins are stated explicitly so that the comparison can be reproduced or challenged.

\subsection{Statistical framework}
$\chi$$^{2}$ is computed as \(\sum [(d_i - m_i)/\sigma_i]^2\) using diagonal errors. For KiDS, this is conservative (the full covariance matrix is available but not used here, for consistency with the SDSS comparison). For SDSS, this matches the source paper's published error treatment.

AIC = $\chi$$^{2}$ + 2k. Both MTDF and fiducial $\Lambda$CDM have k = 0 in the primary comparison, so $\Delta$AIC = $\Delta$$\chi$$^{2}$. For the baryon completion diagnostic, k = 3.

\subsection{Assumption inventory}
\begin{table}[H]
\centering
\footnotesize
\caption{\textbf{Methods table.} Assumption inventory for the galaxy-galaxy lensing comparison.}
\begin{tabularx}{\linewidth}{p{0.28\linewidth}>{\RaggedRight\arraybackslash}X}
\toprule
\textbf{Item} & \textbf{Treatment} \tabularnewline
\midrule
Data covariance & Diagonal errors for both KiDS and SDSS, using the published per-point uncertainties (full KiDS covariance is available but not used here, for parity with the SDSS treatment). \tabularnewline
Baryonic masses & Central stellar mass from the source catalogue; hot gas from Giodini et al. (2009); satellites and intracluster light from Gonzalez et al. (2013); mass proxies from Moster et al. (2013) SHMR with Duffy et al. (2008) concentration. \tabularnewline
Halo baseline & NFW profile with the stated SHMR and concentration-relation variants (Moster/Behroozi/Leauthaud $\times$ Duffy/Dutton-Macci\`o). \tabularnewline
Free parameters in MTDF comparison & Zero in the frozen prediction; diagnostic baryon completion (k = 3 amplitudes) reported separately. \tabularnewline
Free parameters in $\Lambda$CDM baselines & Zero for the fiducial baseline; the six robustness variants differ only in prescription choices (SHMR, concentration), not in fitted parameters. \tabularnewline
\bottomrule
\end{tabularx}
\end{table}

\section{Results}
\label{sec:results}

\subsection{Results summary}
\begin{table}[H]
\centering
\footnotesize
\begin{tabularx}{\linewidth}{l>{\RaggedRight\arraybackslash}X>{\RaggedRight\arraybackslash}X>{\RaggedRight\arraybackslash}X}
\toprule
\textbf{Quantity} & \textbf{MTDF} & \textbf{Best $\Lambda$CDM} & \textbf{Notes} \tabularnewline
\textbf{KiDS $\chi$$^{2}$/$\nu$ (60 DOF)} & \textbf{2.93} & 8.30 & No fitted lensing parameters; 2.8$\times$ improvement \tabularnewline
\textbf{SDSS $\chi$$^{2}$/$\nu$ (105 DOF)} & \textbf{11.21} & 44.11 & No fitted lensing parameters; 3.9$\times$ improvement \tabularnewline
SDSS bins 1--4 $\chi$$^{2}$/$\nu$ & 2.55 & - & Comparable mass range to KiDS \tabularnewline
SDSS bins 5--7 $\chi$$^{2}$/$\nu$ & 22.76 & - & Group/\allowbreak{}cluster centrals; incomplete baryon accounting \tabularnewline
SDSS $\Delta$AIC (k=0 vs k=0) & - & 3455 & Strong preference for MTDF \tabularnewline
\textbf{SDSS corrected $\chi$$^{2}$/$\nu$} & \textbf{3.73} & 44.11 & k=3 nuisance params; $\nu$ = 102 \tabularnewline
SDSS corrected AIC & 386 & 4632 & $\Delta$AIC = 4246 after penalty \tabularnewline
Robustness: worst $\Lambda$CDM & - & 564.77 & Behroozi+2010 + Dutton+Macci\`o 2014; 50$\times$ worse \tabularnewline
Solar System safety & > 10\textsuperscript{20} & - & All PPN parameters safe \tabularnewline
\bottomrule
\end{tabularx}
\end{table}

\subsection{KiDS $\times$ GAMA (Brouwer et al. 2021)}
The MTDF compression model, with f values from the McGaugh BTFR and zero fitted parameters, gives $\chi$$^{2}$/$\nu$ = 2.93 across 60 data points (15 radial bins $\times$ 4 stellar mass bins). Per-bin: Bin 1 $\chi$$^{2}$/15 = 1.20 (excellent); Bin 2 = 3.13; Bin 3 = 3.20; Bin 4 = 4.19. The fiducial $\Lambda$CDM NFW gives $\chi$$^{2}$/$\nu$ = 8.30; baryons only give 31.42.

The naive model using \(S^2 - S_0^2\) (which includes a spurious \(2S_0\,\delta S\) cross-term producing \(\rho \propto 1/r\)) is rejected at $\Delta$$\chi$$^{2}$ = +44. The data independently confirm what standard elasticity requires: only the excitation energy \((\delta S)^2\) gravitates as halo mass.

\subsection{SDSS Red LBG (Mandelbaum et al. 2016)}
The same frozen prediction (no re-tuning) applied to the SDSS dataset gives:

\begin{table}[H]
\centering
\footnotesize
\begin{tabularx}{\linewidth}{l>{\RaggedRight\arraybackslash}X>{\RaggedRight\arraybackslash}X>{\RaggedRight\arraybackslash}X>{\RaggedRight\arraybackslash}X}
\toprule
\textbf{Bin} & \textbf{log M\textsubscript{*}} & \textbf{$\chi$$^{2}$/$\nu$ (MTDF)} & \textbf{$\chi$$^{2}$/$\nu$ ($\Lambda$CDM)} & \textbf{Winner} \tabularnewline
1 & 10.2 & 0.96 & 0.90 & $\Lambda$CDM \tabularnewline
2 & 10.6 & 2.22 & 1.24 & $\Lambda$CDM \tabularnewline
3 & 10.9 & 2.67 & 4.38 & MTDF \tabularnewline
4 & 11.1 & 4.36 & 24.41 & MTDF \tabularnewline
5 & 11.3 & 15.51 & 73.48 & MTDF \tabularnewline
6 & 11.5 & 27.41 & 117.50 & MTDF \tabularnewline
7 & 11.7 & 25.35 & 86.88 & MTDF \tabularnewline
\textbf{All} &  & \textbf{11.21} & \textbf{44.11} & \textbf{MTDF} \tabularnewline
\bottomrule
\end{tabularx}
\end{table}

MTDF wins 5 of 7 bins. In the comparable mass range (bins 1--4, log M\textsubscript{*} < 11.2): $\chi$$^{2}$/$\nu$ = 2.55, closely matching the KiDS result (2.93).

\textbf{Cross-survey consistency:}

\begin{table}[H]
\centering
\footnotesize
\begin{tabularx}{\linewidth}{l>{\RaggedRight\arraybackslash}X>{\RaggedRight\arraybackslash}X}
\toprule
\textbf{Survey} & \textbf{$\chi$$^{2}$/$\nu$ (MTDF)} & \textbf{$\chi$$^{2}$/$\nu$ ($\Lambda$CDM)} \tabularnewline
KiDS (Brouwer+2021) & 2.93 & 8.30 \tabularnewline
SDSS (Mandelbaum+2016) & 11.21 & 44.11 \tabularnewline
\bottomrule
\end{tabularx}
\end{table}

The SDSS absolute $\chi$$^{2}$/$\nu$ is larger primarily because bins 5--7 are group/\allowbreak{}cluster centrals where the central-galaxy-only baryon model is incomplete (Section 4.5).

\subsection{Robustness against $\Lambda$CDM prior choices}
\begin{table}[H]
\centering
\footnotesize
\begin{tabularx}{\linewidth}{l>{\RaggedRight\arraybackslash}X>{\RaggedRight\arraybackslash}X>{\RaggedRight\arraybackslash}X}
\toprule
\textbf{SHMR} & \textbf{c(M)} & \textbf{$\chi$$^{2}$/$\nu$} & \textbf{vs MTDF} \tabularnewline
Moster+2013 (baseline) & Duffy+2008 (WMAP) & 44.11 & 3.9$\times$ \tabularnewline
Moster+2013 & Dutton+Macci\`o 2014 (Planck) & 53.14 & 4.7$\times$ \tabularnewline
Behroozi+2010 & Duffy+2008 & 513.98 & 45.8$\times$ \tabularnewline
Behroozi+2010 & Dutton+Macci\`o 2014 & 564.77 & 50.4$\times$ \tabularnewline
Moster+0.3 dex & Duffy+2008 & 171.37 & 15.3$\times$ \tabularnewline
Moster+0.3 dex & Dutton+Macci\`o 2014 & 196.28 & 17.5$\times$ \tabularnewline
\textbf{MTDF (frozen)} & \textbf{ - } & \textbf{11.21} & \textbf{1.0$\times$} \tabularnewline
\bottomrule
\end{tabularx}
\end{table}

Every $\Lambda$CDM variant is worse than MTDF. Behroozi+2010 assigns larger halo masses at the high-mass end, producing catastrophic overprediction (50$\times$). Dutton+Macci\`o 2014 (Planck concentrations) gives \~{}20\% higher concentrations than Duffy+2008, worsening the overshoot. The MTDF advantage is structural, not dependent on a specific astrophysical prior.

\subsection{Baryon completion for high-mass bins}
Bins 5--7 (log M\textsubscript{*} > 11.2) are group and cluster centrals where the total baryonic mass  -  including hot intragroup gas, satellite galaxies, and intracluster light  -  significantly exceeds the central galaxy. Three routes quantify this:

\textbf{Route A (forward model).} Using Giodini+2009 gas scaling with enclosed-mass compression: $\chi$$^{2}$/$\nu$ worsens from 22.8 to 91.0. The forward model overshoots because Giodini gas masses are calibrated against $\Lambda$CDM M\textsubscript{500}, which includes dark matter halos that do not exist in MTDF.

\textbf{Route B (inverse fit).} Fitting a single f\textsubscript{eff} per bin:

\begin{table}[H]
\centering
\footnotesize
\begin{tabularx}{\linewidth}{l>{\RaggedRight\arraybackslash}X>{\RaggedRight\arraybackslash}X>{\RaggedRight\arraybackslash}X>{\RaggedRight\arraybackslash}X}
\toprule
\textbf{Bin} & \textbf{f\textsubscript{eff} / f\textsubscript{central}} & \textbf{M\textsubscript{additional} implied} & \textbf{Ratio to $\Lambda$CDM scaling} & \textbf{$\chi$$^{2}$/$\nu$} \tabularnewline
5 & 1.39 & 5.6 $\times$ 10\textsuperscript{11} M\textsubscript{$\odot$} & 0.14 & 2.65 \tabularnewline
6 & 1.66 & 2.1 $\times$ 10\textsuperscript{12} M\textsubscript{$\odot$} & 0.14 & 5.86 \tabularnewline
7 & 2.08 & 9.0 $\times$ 10\textsuperscript{12} M\textsubscript{$\odot$} & 0.15 & 6.62 \tabularnewline
\bottomrule
\end{tabularx}
\end{table}

The ratio (implied / $\Lambda$CDM scaling) is 0.14--0.15 across all three bins  -  physically expected for systems without dark matter halos.

\textbf{Route C (enclosed-mass fit).} Fitting M\textsubscript{gas} with the $\beta$-model gas profile and shell-theorem compression confirms Route B: best-fit gas masses are 1.5--16\% of Giodini predictions.

\textbf{Diagnostic correction.} Combining bins 1--4 (unchanged, k = 0) with bins 5--7 (f\textsubscript{eff} corrected, k = 3):

\begin{table}[H]
\centering
\footnotesize
\begin{tabularx}{\linewidth}{l>{\RaggedRight\arraybackslash}X>{\RaggedRight\arraybackslash}X>{\RaggedRight\arraybackslash}X>{\RaggedRight\arraybackslash}X}
\toprule
\textbf{Model} & \textbf{k} & \textbf{$\chi$$^{2}$} & \textbf{AIC} & \textbf{$\Delta$AIC vs $\Lambda$CDM} \tabularnewline
MTDF central-only & 0 & 1177.2 & 1177.2 & 3455 \tabularnewline
MTDF Route B corrected & 3 & 380.0 & 386.0 & 4246 \tabularnewline
Best $\Lambda$CDM (Moster+Duffy) & 0 & 4631.9 & 4631.9 & reference \tabularnewline
\bottomrule
\end{tabularx}
\end{table}

Even after paying the +6 AIC penalty for 3 nuisance parameters, MTDF is preferred by $\Delta$AIC = 4246. The diagnostic correction reduces $\chi$$^{2}$/$\nu$ from 11.21 to 3.73 (N = 105, k = 3, $\nu$ = 102), supporting the interpretation that the residual is amplitude-dominated and consistent with incomplete baryon accounting.

\subsection{Physically anchored cluster prediction (Step 21)}
Step 21 replaces the per-bin fitting of Route B with a zero-parameter forward prediction.

\begin{mdframed}[linecolor=accent,linewidth=0.5pt,backgroundcolor=callout-bg,leftline=true,rightline=true,topline=true,bottomline=true,innerleftmargin=8pt,innerrightmargin=8pt,innertopmargin=6pt,innerbottommargin=6pt]
\textbf{Method box: One-step baryon completion at the MTDF intrinsic scale (cluster regime)}\emph{Intrinsic MTDF mass profile.} For a central baryonic mass M\textsubscript{bar}, the enclosed mass in the isothermal stress regime is
\(M_{\mathrm{enc}}(r) = M_{\mathrm{bar}} + 4\pi\rho_0 f^2 L^2 r\),
where $\rho$\textsubscript{0} and L are fixed MTDF constants and f is the BTFR-based compression factor.

\emph{Compression factor from BTFR.} \(f = (M_{\mathrm{bar}} / A_{\mathrm{BTFR}})^{1/4} / v_{\mathrm{ref}}\), with A\textsubscript{BTFR} = 50 and v\textsubscript{ref} = 161.8 km/\allowbreak{}s.

\emph{Definition of M\textsubscript{500} in MTDF.} The radius r\textsubscript{500} is defined by the usual overdensity criterion with respect to $\rho$\textsubscript{crit}: \(M_{\mathrm{enc}}(r_{500}) / (\frac{4\pi}{3}r_{500}^3) = 500\,\rho_{\mathrm{crit}}\), and M\textsubscript{500} = M\textsubscript{enc}(r\textsubscript{500}).

\emph{One-step baryon completion (no iteration).}

\begin{enumerate}
  \item Compute the central compression factor: \(f_{\mathrm{cen}} = (M_{*,\mathrm{cen}} / A_{\mathrm{BTFR}})^{1/4} / v_{\mathrm{ref}}\).
  \item Solve the overdensity criterion using f = f\textsubscript{cen} to obtain r\textsubscript{500} and M\textsubscript{500}.
  \item Compute additional baryons using externally calibrated scalings:
M\textsubscript{gas} = f\textsubscript{gas}(M\textsubscript{500}) M\textsubscript{500} (Giodini+2009),
M\textsubscript{sat} = R\textsubscript{sat} M\textsubscript{*,cen} (Yang+2007),
M\textsubscript{ICL} = 0.12 (M\textsubscript{*,cen} + M\textsubscript{sat}) (Gonzalez+2007).
  \item Form the total baryon budget: M\textsubscript{bar,tot} = M\textsubscript{*,cen} + M\textsubscript{gas} + M\textsubscript{sat} + M\textsubscript{ICL}.
  \item Apply a single update to the compression factor: f\textsubscript{tot} = (M\textsubscript{bar,tot} / A\textsubscript{BTFR})\textsuperscript{1/\allowbreak{}4} / v\textsubscript{ref}, with no further iteration.
\end{enumerate}
\end{mdframed}

\textbf{Why one-step and not iterative.} In $\Lambda$CDM, iterative self-consistency between gas scaling and M\textsubscript{500} is appropriate because baryon fraction scalings are calibrated against total masses dominated by dark matter halos  -  the baryonic contribution is a perturbation on a DM-dominated potential. In MTDF, the gravitational potential is sourced by the deformation energy of the stress field, which responds nonlinearly to the baryonic mass (\(\Delta\Sigma \propto f^2 \propto M_{\mathrm{bar}}^{1/2}\)). Iterating reintroduces the $\Lambda$CDM mass scale through feedback: adding gas increases f, which increases M\textsubscript{stress}, which increases M\textsubscript{500}, which adds more gas. The converged M\textsubscript{500} approaches the $\Lambda$CDM value, erasing the physical difference. The correct procedure evaluates gas scaling at the intrinsic MTDF gravitational scale  -  set by the central galaxy that sources the stress field  -  without bootstrapping.

\textbf{Results.} Bins 5--7 combined $\chi$$^{2}$/$\nu$ = 22.76 (central-only) to 7.16 (one-step), a 3.2$\times$ improvement with zero free parameters. Bin 5 nearly matches Route B (f\textsubscript{ratio} 1.38 vs 1.40). Independent cross-check: Anderson+2015 stacked ROSAT X-ray M\textsubscript{gas}(M\textsubscript{*}) agrees with MTDF gas masses within 0.25 dex (within the 0.4 dex scatter). Monte Carlo scatter propagation (1000 realisations): $\chi$$^{2}$/$\nu$ = 7.58 [6.56, 8.99] at 68\% CL.

\subsection{Elliptical galaxy velocity dispersions (Step 22)}
Step 20 predicted velocity dispersions for 6 H0LiCOW lenses as a diagnostic cross-check, finding systematic underprediction of 10--25\% using a fixed effective radius R\textsubscript{eff} = 8 kpc. Step 22 upgrades this to a quantitative test using 25 SLACS grade-A strong-lensing ellipticals (Auger et al. 2010; Bolton et al. 2008) with galaxy-specific R\textsubscript{eff} values spanning log M\textsubscript{*} = 10.8 to 11.75 and $\sigma$\textsubscript{e} = 164 to 340 km/\allowbreak{}s.

\textbf{Prediction.} The MTDF isothermal stress field gives \(\sigma_{\mathrm{stress}} = v_{\mathrm{ref}}\,f / \sqrt{2}\), where \(f = (M_{\mathrm{bar}} / A_{\mathrm{BTFR}})^{1/4} / v_{\mathrm{ref}}\) is the BTFR compression factor. Baryonic self-gravity contributes \(\sigma_{\mathrm{bar}} = \sqrt{G\,M_* / K_v\,R_{\mathrm{eff}}}\) with K\textsubscript{v} = 5.0 (Cappellari et al. 2006). These combine in quadrature: \(\sigma_{\mathrm{total}} = \sqrt{\sigma_{\mathrm{stress}}^2 + \sigma_{\mathrm{bar}}^2}\), which is exact because the Jeans equation is linear in the gravitational potential. Zero free parameters.

\textbf{Results.} Mean bias is $-$1.4\% (negligible), with $\chi$$^{2}$/$\nu$ = 3.64 (N = 25). The prediction shows mass-dependent structure: +15.7\% overprediction at log M\textsubscript{*} < 11.0 (where $\sigma$\textsubscript{bar} dominates) trending to $-$12.7\% underprediction at log M\textsubscript{*} > 11.4 (where the stress term dominates and the baryon budget is incomplete). Applying a gas-only baryon correction at the MTDF M\textsubscript{500} scale for galaxies with log M\textsubscript{*} > 11.3  -  using the same Giodini et al. (2009) scaling as Step 21, but without satellites or ICL since SLACS lenses are isolated field ETGs  -  reduces high-mass |bias| from 10.3\% to 4.7\% (54\% improvement) and brings the full-sample $\chi$$^{2}$/$\nu$ to 1.81.

\textbf{Faber-Jackson relation.} The MTDF prediction \(\sigma_{\mathrm{stress}} \propto M_{\mathrm{bar}}^{1/4}\) is the Faber-Jackson relation as a structural consequence, not a fitted scaling. The observed log $\sigma$ versus log M\textsubscript{*} slope is 0.236, consistent with the canonical value of 0.25; the slight flattening arises from the $\sigma$\textsubscript{bar} contribution at low mass.

\textbf{H0LiCOW redo.} Re-predicting the 6 Step 20 lenses with estimated R\textsubscript{eff} from the Shen et al. (2003) mass-size relation (with van der Wel et al. 2014 size evolution) and group baryon correction where applicable reduces the RMS residual from 2.14$\sigma$ to 0.67$\sigma$ (69\% improvement). The Step 20 dispersion mismatch was dominated by an assumed fixed effective radius. Using observed size-mass relations removes this systematic.

\textbf{Falsifiers.} Four pre-registered criteria: (1) |mean bias| < 15\%  -  PASS (1.4\%). (2) |regression slope| < 0.30 per dex after gas correction  -  PASS (0.256). (3) $\chi$$^{2}$/$\nu$ < 5.0 after gas correction  -  PASS (1.81). (4) Gas correction improves $\chi$$^{2}$  -  PASS (3.64 to 1.81). All four pass.

\subsection{Satellite kinematics (Step 22B)}
Step 22A tested the MTDF gravity sector via internal stellar kinematics of ellipticals at r < 10 kpc. Step 22B extends to the halo-scale potential (50--500 kpc) using satellite galaxies as dynamical tracers. This is the cleanest test of the isothermal stress profile \(\rho \sim r^{-2}\) at halo scales.

\textbf{MTDF prediction.} The isothermal stress field gives \(v_c^2(r) = GM_*/r + v_{\mathrm{ref}}^2 f^2\) at satellite scales. The Jeans equation for power-law tracers \(n(r) \sim r^{-\gamma}\) with orbital anisotropy $\beta$ yields \(\sigma_r^2(r) = GM_*/((\gamma - 2\beta)\,r) + v_{\mathrm{ref}}^2 f^2/(\gamma - 1 - 2\beta)\). The stress term provides a constant velocity dispersion floor; the baryonic 1/\allowbreak{}r term adds a mild decline at small projected radii. The line-of-sight dispersion $\sigma$\textsubscript{los}(R\textsubscript{proj}) is obtained via Abel projection. The tracer slope $\gamma$ and anisotropy $\beta$ are nuisance parameters (astrophysical, not MTDF).

\textbf{Amplitude test (More+2011).} We compare MTDF predictions to the host-weighted satellite velocity dispersion $\sigma$\textsubscript{hw} from More et al. (2011, MNRAS 410, 210), spanning 10 stellar mass bins (log M\textsubscript{*} = 9.95--11.65 M\textsubscript{$\odot$}, N\textsubscript{sat} = 1.0--8.4). For group centrals (N\textsubscript{sat} > 1), the compression factor uses total system baryonic mass. The MTDF prediction matches all 10 bins with mean |bias| = 6.3\% and $\chi$$^{2}$/$\nu$ = 0.37 (zero tunable parameters). For comparison, NFW (Moster+2013 SHMR + Duffy+2008 c(M)) gives $\chi$$^{2}$/$\nu$ = 116.5  -  315$\times$ worse.

\textbf{Radial profile test (Combes \& Tiret 2009).} The $\sigma$\textsubscript{los}(R\textsubscript{proj}) radial profile is compared across 3 host mass bins, each with 19 radial points from 33 to 333 kpc (N=57 total data points). Two global nuisance parameters: $\gamma$ (tracer density slope) and $\beta$ (velocity anisotropy). Best nuisance combination ($\gamma$=3.5, $\beta$=0.4): $\chi$$^{2}$ = 198.1, $\chi$$^{2}$/$\nu$ = 3.60, AIC = 202.1. Both parameter values are within typical ranges from N-body simulations. For comparison, NFW gives $\chi$$^{2}$/$\nu$ = 1913 (109$\times$ worse).

\textbf{Falsifiers.} (1) Mean |bias| in More+2011 < 30\% $\rightarrow$ PASS (6.3\%). (2) Profile $\chi$$^{2}$/$\nu$ < 10 for best nuisance $\rightarrow$ PASS (3.60). (3) Scaling slope within 0.3 of unity $\rightarrow$ PASS (0.128 deviation). (4) MTDF not worse than NFW on both tests $\rightarrow$ PASS (better on both: 315$\times$ amplitude, 109$\times$ profile). 4/\allowbreak{}4 PASS.

\subsection{Rotation curves and RAR consistency (Step 22C)}
Steps 22A--B tested the gravity sector at halo scales (50--500 kpc). Step 22C tests it at galactic scales (1--30 kpc) using the 175-galaxy SPARC sample and the Radial Acceleration Relation. At \(r \ll \beta = 22{,}685\) kpc, the MTDF enhancement is approximately constant: \(v_c^2(r) = v_{\mathrm{bar}}^2(r) \times (1 + \alpha) = 2.30 \times v_{\mathrm{bar}}^2(r)\). The falsifiable claims are global summary statistics, not per-galaxy best fits.

\textbf{P1: RMS log velocity scatter.} For each of 3391 SPARC points, compute \(v_{\mathrm{mtdf}} = v_{\mathrm{bar}} \times \sqrt{1 + \alpha/(1 + r/\beta)}\) and the residual \(\log_{10}(v_{\mathrm{obs}}/v_{\mathrm{mtdf}})\). P1 = RMS over all points = 0.185 dex, compared to target 0.174 $\pm$ 0.011 dex (z = $-$0.94, PASS). Newtonian RMS is 0.195 dex; MTDF reduces scatter by 5.4\%.

\textbf{P1B: RAR intrinsic scatter.} After Desmond (2023) deconvolution, P1B = 0.029 dex vs target 0.035 $\pm$ 0.010 dex (z = +0.62, PASS).

\textbf{BTFR.} 168 galaxies after quality cuts. Fit slope = 3.35 (noisy M\textsubscript{bar} estimator); normalisation offset from McGaugh+2012 (A = 50 M\textsubscript{$\odot$}/(km/\allowbreak{}s)\textsuperscript{4}) is +0.316 dex, within the 0.5 dex threshold.

\textbf{v\textsubscript{flat} consistency.} Median ratio v\textsubscript{flat,mtdf} / v\textsubscript{flat,obs} = 0.915; 88\% of galaxies within 50\%, 46\% within 20\%.

\textbf{Per-point $\chi$$^{2}$ diagnostic.} Median $\chi$$^{2}$/$\nu$ = 51 per galaxy. This is reported as a diagnostic only. The V74 dashboard P1 and P1B metrics are scalar summary statistics, not per-galaxy best fits. All models (including MOND with fixed $\Upsilon$\textsubscript{*}) yield similarly large per-point $\chi$$^{2}$ on SPARC without per-galaxy M/\allowbreak{}L optimisation.

\textbf{Falsifiers.} (1) |P1 $-$ 0.174| < 0.033 dex $\rightarrow$ PASS (0.010). (2) |P1B $-$ 0.035| < 0.030 dex $\rightarrow$ PASS (0.006). (3) BTFR normalisation < 0.5 dex $\rightarrow$ PASS (0.316). (4) v\textsubscript{flat} median ratio in [0.6, 1.5] $\rightarrow$ PASS (0.915). 4/\allowbreak{}4 PASS.

\subsection{Wide binary gravity test (Step 22D)}
Wide binaries in the Solar neighbourhood probe ultra-low accelerations at projected separations 2--30 kAU. MTDF predicts Newtonian behaviour at these scales: the $\alpha$-enhancement requires a galaxy-scale stress field source, and individual binary stars do not produce such a field. This is a distinguishing prediction from MOND.

\textbf{Data.} The Pittordis \& Sutherland (2023) cleaned wide binary catalogue from Gaia EDR3 contains 73,087 pairs. After quality cuts (separation 1--50 kAU, distance < 250 pc, RUWE < 1.4), 21,312 pairs remain.

\textbf{Results.} Gravity boost parameter $\gamma$ = 1.102 $\pm$ 0.010 globally (below the pre-registered 1.15 threshold, PASS). Both $\gamma$ and contamination fraction increase monotonically with separation (r = 0.99, r = 0.96), the classic contamination signature. Tighter quality cuts bring $\gamma$ closer to 1.0: tight distance (200 pc) gives $\gamma$ = 1.063, tight RUWE (1.2) gives $\gamma$ = 1.081. The high-separation tail (s $\ge$ 20 kAU, $\gamma$ = 1.60) is INCONCLUSIVE because the F3 cut-stability guardrail triggers.

\textbf{Falsifiers.} (1) High-separation tail: INCONCLUSIVE (guardrail triggered). (2) Global boost $\gamma$ < 1.15: PASS ($\gamma$ = 1.10). Overall: PASS.

\textbf{Discriminator note.} MTDF and MOND-type models agree on galaxy-scale phenomenology in the low-acceleration regime, but differ sharply for wide binaries: MTDF predicts effectively Newtonian dynamics (the stress field is sourced by galaxy-integrated mass, not by individual stellar systems), while acceleration-threshold models generically predict anomalous behaviour. Wide binaries therefore provide a clean external discriminator between these frameworks.

\section{Methodological notes}
\label{sec:notes}
\begin{mdframed}[linecolor=accent,linewidth=0.5pt,backgroundcolor=callout-bg,leftline=true,rightline=true,topline=true,bottomline=true,innerleftmargin=8pt,innerrightmargin=8pt,innertopmargin=6pt,innerbottommargin=6pt]
\begin{itemize}
  \item \textbf{$\Delta$$\Sigma$ and R conventions} match Mandelbaum+2016 and Brouwer+2021 exactly.
  \item \textbf{Diagonal errors} are used because the SDSS source paper uses diagonal errors and no covariance matrix is published. For KiDS, full covariance is available; using diagonal errors for KiDS is conservative.
  \item \textbf{The $\Lambda$CDM curve} is a fiducial zero-parameter mapping (Moster+2013 SHMR + Duffy+2008 c(M)), not a per-bin best fit.
  \item \textbf{Both comparisons use zero tunable parameters.} The MTDF prediction and the fiducial $\Lambda$CDM prediction are both fixed mappings from stellar mass to ESD profile.
  \item \textbf{$\Delta$AIC is computed within the same likelihood construction and error model used by the source paper.}
\end{itemize}
\end{mdframed}

\section{Discussion}
\label{sec:discussion}

\subsection{Claims and non-claims}
\begin{mdframed}[linecolor=accent,linewidth=0.5pt,backgroundcolor=callout-bg,leftline=true,rightline=true,topline=true,bottomline=true,innerleftmargin=8pt,innerrightmargin=8pt,innertopmargin=6pt,innerbottommargin=6pt]
\textbf{Box 1: Claims.} This paper claims that the MTDF compression mechanism, with four fundamental parameters (plus the derived elastic modulus) frozen from independent cosmological and astrophysical observations, reproduces galaxy-galaxy lensing ESD profiles (two surveys, 165 data points), elliptical velocity dispersions (25 SLACS galaxies), satellite kinematics amplitudes (10 mass bins) and radial profiles (57 data points), SPARC rotation curve summary statistics (175 galaxies), and wide binary dynamics (21,312 pairs) under pre-registered pass-or-fail criteria. The falsifiable quantities are aggregate statistics ($\chi$$^{2}$/$\nu$, mean bias, RMS scatter, gravity boost parameter), not per-object best fits. No claim is made of competitive per-galaxy rotation curve fitting without per-galaxy mass-to-light optimisation. No claim is made that $\Lambda$CDM is ruled out; a fitted NFW per bin would outperform the fiducial zero-parameter mapping used here.
\end{mdframed}

\subsection{Diagnostic-only outputs}
\begin{mdframed}[linecolor=accent,linewidth=0.5pt,backgroundcolor=callout-bg,leftline=true,rightline=true,topline=true,bottomline=true,innerleftmargin=8pt,innerrightmargin=8pt,innertopmargin=6pt,innerbottommargin=6pt]
\textbf{Box 2: Diagnostics.} The following outputs are reported for transparency but are explicitly excluded from the falsifier programme: (1) SPARC per-point $\chi$$^{2}$/$\nu$ (median = 51; dominated by galaxy-specific $\Upsilon$\textsubscript{*} systematics; all models give similarly large values without per-galaxy M/\allowbreak{}L fitting). (2) Brouwer+2021 lensing RAR cross-check (qualitative comparison at halo scales; no per-bin radius available for precision test). (3) Wide binary high-separation tail (s $\ge$ 20 kAU; $\gamma$ = 1.60, but the F3 cut-stability guardrail triggers, indicating contamination dominates). (4) Per-point $\chi$$^{2}$ diagnostic plot (Step 22C).
\end{mdframed}

\subsection{Regime summary and falsifier status}
\begin{table}[H]
\centering
\footnotesize
\begin{tabularx}{\linewidth}{l>{\RaggedRight\arraybackslash}X>{\RaggedRight\arraybackslash}X>{\RaggedRight\arraybackslash}X>{\RaggedRight\arraybackslash}X>{\RaggedRight\arraybackslash}X>{\RaggedRight\arraybackslash}X}
\toprule
\textbf{Regime} & \textbf{Scale} & \textbf{Step} & \textbf{Key metric} & \textbf{Threshold} & \textbf{Result} & \textbf{Status} \tabularnewline
Inner galaxy (rotation) & 1--30 kpc & 22C & P1 = 0.185 dex & |z| < 3 (0.174 $\pm$ 0.011) & z = $-$0.94 & PASS \tabularnewline
Inner galaxy (RAR) & 1--30 kpc & 22C & P1B = 0.029 dex & |z| < 3 (0.035 $\pm$ 0.010) & z = +0.62 & PASS \tabularnewline
Halo (dispersions) & 1--10 kpc & 22A & $\chi$$^{2}$/$\nu$ = 1.81 & < 5.0 (gas-corrected) & 1.81 & PASS \tabularnewline
Halo (satellites, amplitude) & 50--500 kpc & 22B & mean |bias| = 6.3\% & < 30\% & 6.3\% & PASS \tabularnewline
Halo (satellites, profile) & 50--500 kpc & 22B & $\chi$$^{2}$/$\nu$ = 3.60 & < 10 (best nuisance) & 3.60 & PASS \tabularnewline
GGL (KiDS) & 100--2600 kpc & 12 & $\chi$$^{2}$/$\nu$ = 2.93 & < $\Lambda$CDM (8.30) & 2.8$\times$ better & PASS \tabularnewline
GGL (SDSS, corrected) & 100--2600 kpc & 17 & $\chi$$^{2}$/$\nu$ = 3.73 & $\Delta$AIC > 0 vs $\Lambda$CDM & 4246 & PASS \tabularnewline
Cluster baryons & 100--2600 kpc & 21 & $\chi$$^{2}$/$\nu$ = 7.16 & no fitted params & 3.2$\times$ improvement & PASS \tabularnewline
Strong lensing & cosmological & 20 & D\textsubscript{dt} $\chi$$^{2}$/$\nu$ = 0.77 & $\gamma$ = 2.00 structural & $\gamma$ = 2.000 & PASS \tabularnewline
Ultra-low acceleration & 2--30 kAU & 22D & $\gamma$ = 1.10 & < 1.15 (3$\sigma$) & 1.102 $\pm$ 0.010 & PASS \tabularnewline
Solar System & < 1 AU & 13 & PPN safety & > 10\textsuperscript{20} & > 10\textsuperscript{20} & PASS \tabularnewline
\bottomrule
\end{tabularx}
\end{table}

All 11 primary metrics pass. 24/\allowbreak{}24 computational steps reproduce from clean inputs.

\begin{figure}[H]
\centering
\includegraphics[width=\linewidth,keepaspectratio]{figures/mtdf_lensing_rsd_test.png}
\caption{\textbf{Figure 2.} Lensing-vs-RSD consistency test. Left: E\textsubscript{G} measurements from six surveys (KiDS-BOSS, DES Y1, BOSS LOWZ/\allowbreak{}CMASS, KiDS-2dFLenS, eBOSS) compared with the GR prediction E\textsubscript{G} = $\Omega$\textsubscript{m}/f = 0.38. Combined measurement E\textsubscript{G} = 0.39 $\pm$ 0.03 shows excellent agreement with GR. Centre: stacked tangential shear profiles around cosmic voids from DES Y3 and KiDS data. Right: test summary. The MTDF photon coupling $\kappa$ = 0.00102 predicts $\Delta$E\textsubscript{G} \~{} 0.001 in voids, below current measurement precision. Status: consistent with MTDF (no detectable photon coupling at current precision).}
\label{fig:fig-lensing-rsd}
\end{figure}

\subsection{Why the absolute $\chi$$^{2}$/$\nu$ is not small}
The SDSS $\chi$$^{2}$/$\nu$ = 11.21 is driven almost entirely by bins 5--7 (group/\allowbreak{}cluster centrals with $\chi$$^{2}$/$\nu$ = 22.8), where central-galaxy-only baryon accounting is incomplete. In the comparable mass range (bins 1--4, log M\textsubscript{*} < 11.2), $\chi$$^{2}$/$\nu$ = 2.55  -  consistent with the KiDS result and approaching acceptable fit quality. The diagnostic correction (Section 4.5) confirms the residual is amplitude-dominated.

The headline comparison should therefore be understood as: relative performance under identical assumptions and source-paper error treatment.

\subsection{Falsifiability}
MTDF's GGL prediction can be falsified by:

\begin{enumerate}
  \item \textbf{Full covariance analysis.} If the SDSS off-diagonal covariance elements significantly change the $\chi$$^{2}$ ranking, the diagonal-error result is unreliable.
  \item \textbf{Per-bin NFW fitting.} $\Lambda$CDM with per-bin halo mass and concentration fitting has more freedom than the fiducial mapping. If fitted NFW achieves $\chi$$^{2}$/$\nu$ < 2.93 on KiDS or < 3.73 (corrected) on SDSS, the comparison favours $\Lambda$CDM.
  \item \textbf{Cluster-scale lensing.} The compression mechanism predicts $\Delta$$\Sigma$ $\propto$ 1/\allowbreak{}R at all scales below L = 2.35 Mpc. At cluster scales (R > 1 Mpc), deviations from isothermal behaviour would indicate the model's boundary.
  \item \textbf{Merging cluster offsets.} Systems like the Bullet Cluster, where lensing peaks are offset from baryonic peaks, would require MTDF to produce lensing signal displaced from baryons  -  a strong constraint on any baryons-only theory.
  \item \textbf{Strong lensing time delays.} The $\eta$ = 1 result means the time-delay formula is standard GR. The isothermal mass profile ($\gamma$ = 2.00) is structural. Time-delay distances at percent-level precision (TDCOSMO, Birrer+2020) directly test both predictions.
\end{enumerate}

\subsection{Remaining open questions}
\begin{itemize}
  \item \textbf{Geometric factor in $\lambda$.} The prediction $\lambda$ = 2$\alpha$/$\beta$$^{2}$ matches the measured value within 4\%, but the factor of 2 (relative to the bare $\alpha$/(2$\beta$$^{2}$)) is identified empirically. A rigorous 3D boundary-value derivation would close this.
  \item \textbf{Connection to [E1$'$].} Steps 8--14 bypass the evolution equation [E1$'$] entirely via the S\textsubscript{0} compression route. Step 19 establishes the covariant field equation that reduces to the compression field equation in the static limit, but the relationship to the full tensor evolution equation remains to be clarified.
  \item \textbf{Freeman's law.} The BTFR scaling emerges from the combination of quadratic screening and Freeman's disk scaling (\(R_d \propto M^{1/2}\)). Freeman's law is observed, not derived from MTDF.
\end{itemize}

\section{Conclusions}
\label{sec:conclusions}
\begin{enumerate}
  \item The MTDF compression mechanism reproduces galaxy-galaxy lensing signals in both KiDS and SDSS with zero fitted halo parameters, achieving $\chi$$^{2}$/$\nu$ = 2.93 (KiDS) and 3.73 (SDSS, with baryon correction). Under identical zero-parameter conditions, this is preferred over fiducial $\Lambda$CDM NFW by factors of 2.8 (KiDS) and 12 (SDSS corrected), with $\Delta$AIC = 4246 on SDSS.
  \item The self-interaction form (S $-$ S\textsubscript{0})$^{2}$ is not a modelling choice  -  it is uniquely selected by the BTFR exponent (algebraic), its coupling is predicted within 4\% by energy self-sourcing, and it is the leading-order term in a Landau-Ginzburg expansion. Solar System observables are safe by > 10\textsuperscript{20} (structural suppression, not fine-tuning).
  \item These results are consistent with the interpretation that the mass profiles conventionally attributed to dark matter halos arise from the elastic energy stored in the compression of the cosmological strain field around baryonic sources  -  with the cosmological background S\textsubscript{0}, the compression scale L, and the source strength f all derived from independently measured quantities.
  \item The covariant completion is now explicit: photons and matter follow the standard GR metric sourced by baryons plus pressureless elastic deformation energy, with conservation guaranteed by the Bianchi identity and $\eta$ = 1 structural (zero anisotropic stress from pressureless dust). PPN parameters are safe by 10\textsuperscript{22--25}. This closes the relativistic formulation at the effective field level.
\end{enumerate}

\section{Appendix A. Derivation chain}
\label{sec:appendix-a}
The complete derivation from MTDF postulates to the observable $\Delta$$\Sigma$(R):

\[
\begin{aligned}
&\alpha, \beta, E, \rho_{\mathrm{crit}}, \Omega_{\mathrm{DE}} \\
&\quad\rightarrow S_0 = 1.084 \quad\text{(cosmological energy density)} \\
&\quad\rightarrow L = \alpha\beta/(4\pi) = 2{,}347\;\text{kpc} \quad\text{(Gauss law + coupling)} \\
&\quad\rightarrow \rho_0 = E\,S_0^2/(2c^2) = 87.9\;M_\odot/\text{kpc}^3 \quad\text{(background density)} \\
&\quad\rightarrow v_{\mathrm{ref}} = \sqrt{4\pi G\rho_0}\,L = 161.8\;\text{km/s} \quad\text{(isothermal identity)} \\
&\quad\rightarrow f(M) = v_{\mathrm{flat}}(M)/v_{\mathrm{ref}} \quad\text{(BTFR, no fitting)} \\
&\quad\rightarrow \rho(r) = \rho_0\,f^2\,L^2/r^2 \quad\text{(isothermal stress density)} \\
&\quad\rightarrow \Delta\Sigma(R) = \pi\rho_0\,f^2\,L^2/R + M_{\mathrm{bar}}/(\pi R^2) \quad\text{(analytical LOS projection)}
\end{aligned}
\]

No parameters are fitted to the lensing or halo-scale datasets. The mapping from stellar mass bin to source strength f(M) uses a single externally fixed BTFR normalisation (McGaugh+2012) as a baryonic scaling reference with standard gas corrections; it is not tuned to the lensing data.

\section{Appendix B. Field equation backing}
\label{sec:appendix-b}
The compression profile \(S = S_0 + C/r\) is the unique spherically symmetric solution of Laplace's equation outside baryonic sources, with \(S \to S_0\) at infinity. Laplace's equation is valid at \(r \ll \beta\) because the Helmholtz operator \((\nabla^2 - 1/\beta^2)\) reduces to \(\nabla^2\) when \((r/\beta)^2 \ll 1\). At the outer edge of the GGL data range (\(r = 2{,}600\) kpc, \(\beta = 22{,}685\) kpc), the Yukawa correction is \((r/\beta)^2 = 0.013\)  -  accuracy better than 2\%.

The source strength \(Q = \alpha\beta\,f(M)\) follows from Gauss's law matching at the galaxy boundary:

\[ C = -\frac{\alpha}{4\pi E\beta^2}\int J_{\mathrm{bar}}\,d^3x \]

The identification \(J_{\mathrm{bar}} = \rho_{\mathrm{bar}}\,c^2\) (rest-mass energy density, or equivalently the trace of the stress-energy tensor) is the simplest covariant source. Only this candidate, combined with \((S - S_0)^2\) screening, produces the BTFR scaling \(f \propto M^{1/4}\).

\section{Appendix C. Dimensional consistency}
\label{sec:appendix-c}
\begin{table}[H]
\centering
\footnotesize
\begin{tabularx}{\linewidth}{l>{\RaggedRight\arraybackslash}X>{\RaggedRight\arraybackslash}X>{\RaggedRight\arraybackslash}X}
\toprule
\textbf{Quantity} & \textbf{Expression} & \textbf{Units} & \textbf{Value} \tabularnewline
S\textsubscript{0} & \(\sqrt{2\Omega_{\mathrm{DE}}\rho_{\mathrm{crit}}c^2/E}\) & dimensionless & 1.084 \tabularnewline
L & \(\alpha\beta/(4\pi)\) & kpc & 2,347 \tabularnewline
$\rho$\textsubscript{0} & \(ES_0^2/(2c^2)\) & M\textsubscript{$\odot$}/kpc\textsuperscript{3} & 87.9 \tabularnewline
v\textsubscript{ref} & \(\sqrt{4\pi G\rho_0}\,L\) & km/\allowbreak{}s & 161.8 \tabularnewline
$\lambda$ & \(2\alpha/\beta^2\) & m\textsuperscript{$-$2} & 5.3 $\times$ 10\textsuperscript{$-$48} \tabularnewline
f\textsubscript{$\odot$} & \(c^2 M_\odot/(E\beta^3 S_0)\) & dimensionless & 5.3 $\times$ 10\textsuperscript{$-$16} \tabularnewline
M\textsubscript{stress}(1 AU) & \(4\pi\rho_0 f_\odot^2 L^2 r\) & kg & 16.4 \tabularnewline
\bottomrule
\end{tabularx}
\end{table}

\section{Appendix D. Notation and symbols}
\label{sec:appendix-d}
\begin{table}[H]
\centering
\footnotesize
\begin{tabularx}{\linewidth}{l>{\RaggedRight\arraybackslash}X>{\RaggedRight\arraybackslash}X}
\toprule
\textbf{Symbol} & \textbf{Meaning} & \textbf{Units} \tabularnewline
S & Scalar compression invariant & 1 \tabularnewline
S\textsubscript{0} & Cosmological background strain & 1 \tabularnewline
$\delta$S & Local excitation (S $-$ S\textsubscript{0}) & 1 \tabularnewline
$\alpha$ & Stress-matter coupling & 1 \tabularnewline
$\beta$ & Coherence length & m \tabularnewline
E & Elastic modulus of spacetime & Pa \tabularnewline
L & Compression scale $\alpha$$\beta$/(4$\pi$) & m \tabularnewline
f(M) & Dimensionless source strength v\textsubscript{flat}/v\textsubscript{ref} & 1 \tabularnewline
$\lambda$ & Screening parameter & m\textsuperscript{$-$2} \tabularnewline
$\rho$\textsubscript{0} & ES\textsubscript{0}\textsuperscript{2}/(2c\textsuperscript{2}) & kg/\allowbreak{}m\textsuperscript{3} \tabularnewline
v\textsubscript{ref} & Reference velocity & m/\allowbreak{}s \tabularnewline
$\Delta$$\Sigma$(R) & Excess surface density & M\textsubscript{$\odot$}/pc\textsuperscript{2} \tabularnewline
\bottomrule
\end{tabularx}
\end{table}

\section{Appendix E. Data provenance}
\label{sec:appendix-e}
All observational data used in this work are from published sources. No proprietary or pre-publication data are used.

\begin{table}[H]
\centering
\footnotesize
\begin{tabularx}{\linewidth}{l>{\RaggedRight\arraybackslash}X>{\RaggedRight\arraybackslash}X>{\RaggedRight\arraybackslash}X>{\RaggedRight\arraybackslash}X}
\toprule
\textbf{Dataset} & \textbf{Source paper} & \textbf{Repo file} & \textbf{Format} & \textbf{Steps} \tabularnewline
KiDS $\times$ GAMA ESD profiles & Brouwer et al. (2021, A\&A 650, A113) & \texttt{data/\allowbreak{}brouwer2021/\allowbreak{}Fig-3\_\allowbreak{}*.txt}, \texttt{Fig-4-5-C1\_\allowbreak{}*.txt} & ASCII & 1, 4, 5, 12, 22C \tabularnewline
KiDS covariance matrices & Brouwer et al. (2021) & \texttt{data/\allowbreak{}brouwer2021/\allowbreak{}*covmatrix.txt} & ASCII & 12 \tabularnewline
SDSS Red LBG ESD profiles & Mandelbaum et al. (2016, MNRAS 457, 3200) & \texttt{data/\allowbreak{}mandelbaum2016/\allowbreak{}planck\_\allowbreak{}lbg.ds.red.out} & ASCII & 15, 16, 17, 21 \tabularnewline
SPARC rotation curves & Lelli, McGaugh \& Schombert (2016, AJ 152, 157) & \texttt{validation/\allowbreak{}data/\allowbreak{}sparc\_\allowbreak{}clean.json} & JSON & 10, 22C \tabularnewline
Gaia EDR3 wide binaries & Pittordis \& Sutherland (2023, OJAp 6, 4) & \texttt{data/\allowbreak{}pittordis2023/\allowbreak{}pittordis2023\_\allowbreak{}wb.csv} & CSV & 22D \tabularnewline
SLACS elliptical galaxies & Auger et al. (2010, ApJ 724, 511); Bolton et al. (2008, ApJ 682, 964) & Hardcoded in \texttt{step22\_\allowbreak{}elliptical\_\allowbreak{}dispersions.py} & Table & 22A \tabularnewline
H0LiCOW time-delay lenses & Wong et al. (2020, MNRAS 498, 1420) & Hardcoded in \texttt{step20\_\allowbreak{}strong\_\allowbreak{}lensing.py} & Table & 20, 22A \tabularnewline
Satellite kinematics (amplitude) & More et al. (2011, MNRAS 410, 210) & \texttt{data/\allowbreak{}digitised/\allowbreak{}more\_\allowbreak{}2011\_\allowbreak{}fig2\_\allowbreak{}MSR\_\allowbreak{}digitized.csv} & Digitised & 22B \tabularnewline
Satellite kinematics (profiles) & Combes \& Tiret (2009, ASP Conf 421, 170) & \texttt{data/\allowbreak{}digitised/\allowbreak{}combes\_\allowbreak{}tiret\_\allowbreak{}2009\_\allowbreak{}sigma\_\allowbreak{}los\_\allowbreak{}digitized.csv} & Digitised & 22B \tabularnewline
\bottomrule
\end{tabularx}
\end{table}

\textbf{Digitisation notes.} Two datasets (More+2011, Combes \& Tiret 2009) were extracted by manual coordinate reading from published figures because no machine-readable tables are available. Estimated digitisation uncertainty is $\pm$2 km/\allowbreak{}s in velocity dispersion and $\pm$0.05 dex in log stellar mass. The digitised values are stored in \texttt{data/\allowbreak{}digitised/\allowbreak{}} and also hardcoded as numpy arrays in \texttt{step22b\_\allowbreak{}satellite\_\allowbreak{}kinematics.py} for reproducibility.

\textbf{Scaling relations (not data files).} Several steps use published scaling relations evaluated analytically from their fitted parameters: Moster+2013 SHMR, Duffy+2008 and Dutton+Macci\`o 2014 c(M), Giodini+2009 gas fraction, Yang+2007 satellite fraction, Gonzalez+2007 ICL fraction, Shen+2003 and van der Wel+2014 size-mass relations, Behroozi+2010 SHMR. These are referenced in the code comments with table/\allowbreak{}equation numbers.

\section{References}
\label{sec:references}
\begin{itemize}
  \item Auger et al. (2010). The Sloan Lens ACS Survey. X. Stellar, dynamical, and total mass correlations. ApJ 724, 511. DOI: \href{https://doi.org/10.1088/0004-637X/724/1/511}{10.1088/\allowbreak{}0004-637X/\allowbreak{}724/\allowbreak{}1/\allowbreak{}511}
  \item Behroozi, Conroy \& Wechsler (2010). A Comprehensive Analysis of Uncertainties Affecting the Stellar Mass-Halo Mass Relation. ApJ 717, 379. DOI: \href{https://doi.org/10.1088/0004-637X/717/1/379}{10.1088/\allowbreak{}0004-637X/\allowbreak{}717/\allowbreak{}1/\allowbreak{}379}
  \item Bolton et al. (2008). The Sloan Lens ACS Survey. V. The full ACS strong-lens sample. ApJ 682, 964. DOI: \href{https://doi.org/10.1086/589327}{10.1086/\allowbreak{}589327}
  \item Brouwer et al. (2021). The weak lensing radial acceleration relation. A\&A 650, A113. DOI: \href{https://doi.org/10.1051/0004-6361/202039795}{10.1051/\allowbreak{}0004-6361/\allowbreak{}202039795}
  \item Catinella et al. (2010). The GALEX Arecibo SDSS Survey. ApJ 722, 1699. DOI: \href{https://doi.org/10.1111/j.1365-2966.2009.16180.x}{10.1111/\allowbreak{}j.1365-2966.2009.16180.x}
  \item Duffy et al. (2008). Dark matter halo concentrations. MNRAS 390, 64. DOI: \href{https://doi.org/10.1111/j.1745-3933.2008.00537.x}{10.1111/\allowbreak{}j.1745-3933.2008.00537.x}
  \item Dutton \& Macci\`o (2014). Cold dark matter haloes in the Planck era. MNRAS 441, 3359. DOI: \href{https://doi.org/10.1093/mnras/stu742}{10.1093/\allowbreak{}mnras/\allowbreak{}stu742}
  \item Giodini et al. (2009). Stellar and total baryon mass fractions in groups and clusters. ApJ 703, 982. DOI: \href{https://doi.org/10.1088/0004-637X/703/1/982}{10.1088/\allowbreak{}0004-637X/\allowbreak{}703/\allowbreak{}1/\allowbreak{}982}
  \item Gonzalez, Sivanandam \& Zabludoff (2013). Galaxy Cluster Baryon Fractions Revisited. ApJ 778, 14. DOI: \href{https://doi.org/10.1088/0004-637X/778/1/14}{10.1088/\allowbreak{}0004-637X/\allowbreak{}778/\allowbreak{}1/\allowbreak{}14}
  \item Lelli, McGaugh \& Schombert (2016). SPARC: Mass Models for 175 Disk Galaxies. AJ 152, 157. DOI: \href{https://doi.org/10.3847/0004-6256/152/6/157}{10.3847/\allowbreak{}0004-6256/\allowbreak{}152/\allowbreak{}6/\allowbreak{}157}
  \item Mandelbaum, Wang, Zu et al. (2016). Strong bimodality in the host halo mass of central galaxies. MNRAS 457, 3200. DOI: \href{https://doi.org/10.1093/mnras/stw188}{10.1093/\allowbreak{}mnras/\allowbreak{}stw188}
  \item McGaugh (2012). The Baryonic Tully-Fisher Relation of Gas-Rich Galaxies. AJ 143, 40. DOI: \href{https://doi.org/10.1088/0004-6256/143/2/40}{10.1088/\allowbreak{}0004-6256/\allowbreak{}143/\allowbreak{}2/\allowbreak{}40}
  \item More et al. (2011). Satellite kinematics  -  III. Halo masses of central galaxies in SDSS. MNRAS 410, 210. DOI: \href{https://doi.org/10.1111/j.1365-2966.2010.17436.x}{10.1111/\allowbreak{}j.1365-2966.2010.17436.x}
  \item Moster, Naab \& White (2013). Galactic star formation and accretion histories. MNRAS 428, 3121. DOI: \href{https://doi.org/10.1093/mnras/sts261}{10.1093/\allowbreak{}mnras/\allowbreak{}sts261}
  \item Navarro, Frenk \& White (1997). A Universal Density Profile. ApJ 490, 493. DOI: \href{https://doi.org/10.1086/304888}{10.1086/\allowbreak{}304888}
  \item Pittordis \& Sutherland (2023). Testing gravity with wide binary stars like never before. OJAp 6, 4. DOI: \href{https://doi.org/10.21105/astro.2205.02846}{10.21105/\allowbreak{}astro.2205.02846}
  \item Planck Collaboration (2020). Planck 2018 results. VI. Cosmological parameters. A\&A 641, A6. DOI: \href{https://doi.org/10.1051/0004-6361/201833910}{10.1051/\allowbreak{}0004-6361/\allowbreak{}201833910}
  \item Wong et al. (2020). H0LiCOW  -  XIII. A 2.4 per cent measurement of H\textsubscript{0} from lensed quasars. MNRAS 498, 1420. DOI: \href{https://doi.org/10.1093/mnras/stz3094}{10.1093/\allowbreak{}mnras/\allowbreak{}stz3094}
\end{itemize}



\part{Companion III: Photon Coupling, Redshift, and Early Universe}

\begin{abstract}
This companion paper (Part III) explores the coupling between the MTDF stress field and photon propagation. We present the photon-stress coupling parameter $\kappa$ = 0.00102, its structural relation to the early field energy fraction f\textsubscript{kick}, and the resulting achromatic frequency shift mechanism. Empirical motivation comes from Pantheon+ Type Ia supernovae cross-matched to DESI void catalogues, which detect an environment-dependent distance-modulus residual ($\gamma$\textsubscript{env} = +0.013 $\pm$ 0.004, p = 0.003) at low redshift. All extensions are explicitly labelled as hypotheses with stated falsification criteria.
\end{abstract}

\section{Purpose and scope}
\label{sec:purpose}
This part collects MTDF extension ideas that couple the stress field to photon propagation and redshift, and explores how those ideas could relate to early structure observations and standard ruler measurements.
These sections are explicitly labelled as hypotheses unless and until they are validated against independent datasets.

Empirical motivation: an MTDF extension test using Pantheon+ Type Ia supernovae cross matched to DESI DESIVAST void catalogues, analysed under the full Pantheon+ STAT+SYS covariance, detects an environment dependent term in distance modulus residuals that is strongest at low redshift. For the low-z bin (z < 0.04), the fitted environment coefficient is $\gamma$\_\allowbreak{}env = +0.013 $\pm$ 0.004, with nominal p = 0.003 (signed distance to void boundary metric). This supports an environment dependent distance redshift mapping. It does not uniquely identify photon coupling as the mechanism, but it motivates the photon propagation and falsifier tests in this part.

\section{Photon coupling and stress induced redshift}
\label{sec:photon}
Some observational tensions are often framed as ``distance ladder'' issues, selection effects, or modelling systematics.
MTDF offers an additional possibility: part of the measured redshift could include a stress related contribution that depends on the integrated field state along the line of sight.

\textbf{Important practical point:} redshift-only tests using galaxy samples inside and outside voids primarily measure the peculiar-velocity field (void outflow and associated large scale flows), typically at the tens to hundreds of km/\allowbreak{}s level. A positive void redshift residual is therefore expected in standard structure formation and cannot, by itself, isolate any subdominant stress-photon contribution without independent distance indicators or an explicit velocity-field reconstruction. This does not conflict with the SN result: a void can generate kinematic redshift from outflow while still producing a reduced path-integrated stress interaction, yielding a systematic brightening at fixed redshift.

\begin{table}[H]
\centering
\footnotesize
\begin{tabularx}{\linewidth}{l>{\RaggedRight\arraybackslash}X>{\RaggedRight\arraybackslash}X}
\toprule
\textbf{Observable} & \textbf{Dominant contribution} & \textbf{What a void signal means} \tabularnewline
\midrule
Galaxy redshift residuals (spectroscopic) & Peculiar velocities and survey selection & Mostly a kinematic consistency check (void outflow). Not a clean photon coupling probe. \tabularnewline
SN Ia distance modulus residuals at fixed z & Environment-dependent distance mapping after standardisation & Evidence for an environment-linked residual in the distance-redshift relation at low z. \tabularnewline
\bottomrule
\end{tabularx}
\end{table}

\subsection{Stress induced frequency shift}
\begin{mdframed}[linecolor=accent,linewidth=0.5pt,backgroundcolor=callout-bg,leftline=true,rightline=true,topline=true,bottomline=true,innerleftmargin=8pt,innerrightmargin=8pt,innertopmargin=6pt,innerbottommargin=6pt]
\textbf{Status:} Hypothesis. Not a core validated claim.
\end{mdframed}

\subsubsection{Observation}
In standard cosmology, redshift is primarily kinematic and gravitational in origin, with well defined contributions in GR. Some datasets and comparisons motivate checking whether there is any additional environment dependent component.

Link to Part I: the SN $\times$ void analysis provides evidence for an environment-dependent residual in a distance-based observable. Any photon-coupling model proposed here must remain consistent with that measured sign and redshift localisation, while also remaining perturbative (|$\delta$z\_\allowbreak{}stress| much less than z\_\allowbreak{}exp).

\subsubsection{Proposed MTDF mechanism}
Within this hypothesis, a weak coupling between photon frequency and the local stress magnitude could arise in the MTDF framework through propagation in a strained spacetime background. The fractional frequency shift along a photon trajectory would then take the form:

$$
\frac{d\nu}{\nu} = -\,\kappa\,c\,\langle \Vert \tilde{S}\Vert\rangle\,dt
$$

\subsubsection{Structural relation between $\kappa$ and f\textsubscript{kick}}
The pipeline uses $\kappa$ = 1.02 $\times$ 10\textsuperscript{$-$3}, an observational anchor set within the FIRAS spectral-distortion bounds (|$\mu$| < 9 $\times$ 10\textsuperscript{$-$5}, |y| < 1.5 $\times$ 10\textsuperscript{$-$5}; Fixsen et al. 1996) and adopted from the V18 workbook calibration. Two candidate first-principles expressions sit close to this value:

$$
\kappa = \frac{f_{\mathrm{kick}}}{3} = 1.10 \times 10^{-3}, \qquad \kappa = \frac{f_{\mathrm{kick}}}{\pi} = 1.05 \times 10^{-3}
$$

where \(f_{\mathrm{kick}} = (1-\beta_{\mathrm{eos}})^2 / [24\,(1+\alpha)] = 3.303 \times 10^{-3}\) is the early field energy fraction derived in Companion Part I, Section 3.1.

The first expression follows from the exact angular averaging of the spatial strain tensor projection onto photon propagation directions:

$$
\langle S_{ij}\,n^i\,n^j \rangle_{\Omega} = \tfrac{1}{3}\,S^i{}_i
$$

This is a standard isotropic tensor identity and introduces no additional assumptions. The second expression ($\kappa$ = f\textsubscript{kick}/$\pi$) gives a closer numerical match (3.0\% vs 7.9\% above the pipeline value) but the $\pi$ factor has not been rigorously derived; it may arise from path-integral normalisation or 3D-to-1D field statistics. The structural proportionality $\kappa$ $\propto$ f\textsubscript{kick} is robust; the numerical denominator (3, $\pi$, or another order-unity geometric factor) remains an open question (see documentation/\allowbreak{}Notes/\allowbreak{}05\_\allowbreak{}Kappa\_\allowbreak{}Derivation.md for the full analysis).

A forward sensitivity evaluation at the best-fit point gives $\Delta$$\chi$$^{2}$($\kappa$ = 1.10 $\times$ 10\textsuperscript{$-$3}) $\approx$ +0.5 on the CMB correction path; the late-universe likelihood is insensitive because $\kappa$ enters there only as the product k\textsubscript{f} $\times$ $\kappa$. We therefore retain the pipeline value pending either analytical resolution of the denominator or a full MCMC refit.

The pipeline value is well below FIRAS detection sensitivity, and direct detection of $\kappa$ would require next-generation experiments such as PIXIE.

Regardless of the exact denominator, the proportionality $\kappa$ $\propto$ f\textsubscript{kick} links the photon sector to the same field dynamics that determine early-universe energy injection. The photon coupling is not independent of the gravity sector: it is structurally determined by \{$\alpha$, $\beta$\textsubscript{eos}\} through f\textsubscript{kick}.

\subsubsection{Physical interpretation}
The photon-stress coupling is a geometric (conservative) effect: the photon propagates through the strain-modified effective metric and experiences a frequency shift proportional to the strain variation along its path. No energy is transferred locally between the photon and the strain field. This interpretation is supported by superconducting cavity measurements: if \(\kappa\) were dissipative (transferring energy from photons to the strain field), it would limit microwave cavity quality factors to \(Q_\kappa \sim 5 \times 10^7\), four orders of magnitude below measured values of \(Q > 2 \times 10^{11}\) for superconducting niobium cavities at 1.3~GHz. The dissipative interpretation is ruled out by existing data.

The geometric nature of the coupling makes a specific prediction: the photon-stress frequency shift is strictly achromatic. Any detection of wavelength-dependent \(\kappa\) (beyond second-order metric corrections, estimated at \(\delta\kappa/\kappa < 10^{-6}\)) would falsify the geometric interpretation.

\subsubsection{Predictions}
\textbf{Perturbative bound and safety check:} the dominant redshift contribution remains cosmological expansion and standard gravity.
Any MTDF photon-stress coupling must enter as a small correction:
\[
z_{\mathrm{obs}} = z_{\mathrm{exp}} + \delta z_{\mathrm{stress}}, \qquad \lvert \delta z_{\mathrm{stress}}\rvert \ll z_{\mathrm{exp}}.
\]
Equivalently, for small corrections, \( (1+z_{\mathrm{obs}}) \approx (1+z_{\mathrm{exp}})(1+\delta z_{\mathrm{stress}}) \).
A detection of order-unity non-cosmological shifts would directly contradict standard ruler and isotropy constraints and would therefore falsify this extension.

This is not a global "tired light" replacement: the hypothesis is a line-of-sight, environment linked perturbation that must vanish (or become negligible) in low stress regions.

\begin{itemize}
  \item Differential redshift effects could correlate with large scale density or void depth along the line of sight.
  \item If the coupling is non zero, multi wavelength comparisons could show a small but structured offset in inferred redshift for the same object class.
  \item Any effect should be bounded by existing isotropy constraints, which makes anisotropy a natural falsifier.
\end{itemize}

\subsubsection{How to test}
\begin{itemize}
  \item Perform redshift versus density correlation tests within a single survey using consistent selection and forward modelling of systematics.
  \item Use multi wavelength observations (for example X ray versus optical lines) to check for a persistent differential component under controlled calibration.
  \item Test isotropy by searching for a dipole or higher multipoles in the residual after standard corrections.
\end{itemize}

\begin{mdframed}[linecolor=accent,linewidth=0.5pt,backgroundcolor=callout-bg,leftline=true,rightline=true,topline=true,bottomline=true,innerleftmargin=8pt,innerrightmargin=8pt,innertopmargin=6pt,innerbottommargin=6pt]
\textbf{Interpretation:} A null result is valuable: it constrains photon coupling extensions without impacting the stress field dynamics used for galaxy scale phenomenology.
\end{mdframed}

\section{Early structure, JWST era tensions, and BAO links}
\label{sec:early}
If early structure appears more mature than expected under a baseline model, there are two broad options: revise the physical model for formation and feedback, or revisit the mapping between redshift and cosmic time in the relevant regime.
MTDF style photon coupling is one possible route, but it is also the riskiest claim category, so the companion treats it as a test driven hypothesis only.

\subsection{Formation time interpretation under a modified mapping}
\begin{mdframed}[linecolor=accent,linewidth=0.5pt,backgroundcolor=callout-bg,leftline=true,rightline=true,topline=true,bottomline=true,innerleftmargin=8pt,innerrightmargin=8pt,innertopmargin=6pt,innerbottommargin=6pt]
\textbf{Status:} Hypothesis.
\end{mdframed}

\subsubsection{Observation}
Some analyses interpret early galaxy observations as implying faster growth or earlier formation times. These interpretations depend on converting redshift to time and on modelling selection and lensing biases.

\subsubsection{Proposed MTDF mechanism}
Within this hypothesis, if the observed redshift contains a small stress-related component, then the inferred formation time could shift without requiring a drastic change in local physics. Any such effect must remain consistent with standard rulers and isotropy constraints, and is not established here.

\subsubsection{Predictions}
\begin{itemize}
  \item Any modified mapping should predict a specific redshift dependence that can be checked against standard ruler datasets.
  \item The effect should be environment sensitive: dense sightlines could differ from void sightlines in a measurable way.
  \item The effect should leave signatures in cross checks between distance measures rather than only in one probe class.
\end{itemize}

\subsubsection{How to test}
\begin{itemize}
  \item Construct a joint analysis plan where standard ruler distances and early structure indicators are fitted simultaneously with a minimal extension term.
  \item Perform split analyses by line of sight environment for high redshift samples.
  \item Cross check with gravitational lensing magnification distributions to avoid conflating stress coupling with selection effects.
\end{itemize}

\begin{mdframed}[linecolor=accent,linewidth=0.5pt,backgroundcolor=callout-bg,leftline=true,rightline=true,topline=true,bottomline=true,innerleftmargin=8pt,innerrightmargin=8pt,innertopmargin=6pt,innerbottommargin=6pt]
\textbf{Notes:} This subsection is intentionally cautious. It is not presented as an explanation, but as a falsifiable way to test whether a stress related component could be present at all.
\end{mdframed}

\section{Quantitative falsifiers and timeline}
\label{sec:tests}
The following items are framed as direct falsifiers for the photon-coupling and differential-redshift hypotheses. The Part I SN $\times$ void result motivates testing for an environment term, but the mechanism remains open; any specific photon-coupling model must survive these criteria under robust systematics control.

\begin{table}[H]
\centering
\footnotesize
\begin{tabularx}{\linewidth}{l>{\RaggedRight\arraybackslash}X>{\RaggedRight\arraybackslash}X}
\toprule
\textbf{Test} & \textbf{Expectation} & \textbf{Criterion} \tabularnewline
Environment-linked distance residuals (SN) & Positive $\gamma$\_\allowbreak{}env at low z, consistent sign across independent void catalogues & Sign reversal, or loss of effect under the pre-registered robustness suite, would rule out this environment term \tabularnewline
Environment-linked redshift residuals (galaxies, after kinematic control) & Any residual must be small compared with bulk outflow and must track environment once the velocity field is modelled & A residual consistent with zero at the analysis sensitivity floor, after kinematic subtraction, would constrain photon coupling tightly \tabularnewline
Multi-wavelength differential redshift & No chromatic component in standard GR; any MTDF-induced differential term must be small and structured & A controlled multi-line study consistent with zero across environments would falsify wavelength-dependent coupling \tabularnewline
Isotropy of residuals & Residual field should not introduce a large dipole or multipole beyond established isotropy limits & A robust, unexplained large-scale anisotropy would falsify the extension \tabularnewline
Redshift dependence & If present, the effect should weaken with redshift as coherence and line-of-sight averaging dominate & No measurable trend with redshift (or the wrong trend) under matched selection would falsify the extension \tabularnewline
\bottomrule
\end{tabularx}
\end{table}

\subsection{Precision consistency checks and sensitivity}
In addition to the qualitative falsifiers above, we ran three orthogonal consistency checks designed to separate large kinematic effects from any small stress photon perturbation. These do not claim a detection of photon coupling. Their purpose is to confirm that the MTDF-predicted coupling value $\kappa$ = f\textsubscript{kick}/3 $\approx$ 0.00102 (below FIRAS detection sensitivity) is not already excluded by current precision probes and to quantify what sensitivity would be required.

\begin{table}[H]
\centering
\footnotesize
\begin{tabularx}{\linewidth}{l>{\RaggedRight\arraybackslash}X>{\RaggedRight\arraybackslash}X>{\RaggedRight\arraybackslash}X>{\RaggedRight\arraybackslash}X}
\toprule
\textbf{Check} & \textbf{Observable} & \textbf{Result} & \textbf{MTDF scale at $\kappa$ = 0.00102} & \textbf{Interpretation} \tabularnewline
\midrule
Zero-velocity shell (turnaround) & $\Delta$$\mu$ of SNe near v\_\allowbreak{}pec $\approx$ 0 shells & Mock demonstration yields $\Delta$$\mu$ = +0.078 $\pm$ 0.057 mag (1.4$\sigma$), below threshold & Order 0.001 mag, below current SN scatter (about 0.15 mag) & Sensitivity forecast: would require very large stacked SN samples (order 20,000) for a 1$\sigma$ probe \tabularnewline
Lensing vs RSD (E\_\allowbreak{}G parameter) & E\_\allowbreak{}G from combined lensing and RSD & Weighted mean E\_\allowbreak{}G = 0.382 $\pm$ 0.028 vs GR expectation about 0.30 (known literature tension) & $\Delta$E\_\allowbreak{}G predicted by MTDF photon term is order 0.01 & Not explained by this extension; MTDF does not worsen it at the pipeline $\kappa$ value \tabularnewline
Alcock-Paczynski void shape & Residual ellipticity after Kaiser RSD correction & $\Delta$e = -0.011 $\pm$ 0.031 (0.3$\sigma$, consistent with spherical voids) & Order 0.0001 & Consistent; current precision is far above the level needed to test $\kappa$ directly \tabularnewline
\bottomrule
\end{tabularx}
\end{table}

Taken together, these checks support a conservative conclusion: the stress photon coupling is perturbative at the pipeline $\kappa$ value. Current surveys can detect standard void kinematics at high significance, but the subdominant photon level signatures implied by $\kappa$ = f\textsubscript{kick}/3 $\approx$ 0.00102 remain below present sensitivity. This makes Part I's SN $\times$ void evidence the most informative current probe of environment dependent distance redshift mapping, while these tests function as consistency and forecasting constraints.

Kinematic contamination note: redshift only void tests using galaxy samples primarily measure bulk peculiar velocities from standard void dynamics (of order 100 km/\allowbreak{}s). Such measurements can be useful as a consistency check, but they cannot directly isolate a subdominant stress photon contribution unless an independent distance indicator is available and the kinematic component is removed.

\begin{mdframed}[linecolor=accent,linewidth=0.5pt,backgroundcolor=callout-bg,leftline=true,rightline=true,topline=true,bottomline=true,innerleftmargin=8pt,innerrightmargin=8pt,innertopmargin=6pt,innerbottommargin=6pt]
\textbf{Observation timeline:}\begin{itemize}
  \item 2025 - DESI redshift-density analysis
  \item 2026 - Euclid weak-lensing correlation
  \item 2027 - Roman Space Telescope mapping
  \item 2028 - Multi-wavelength differential tests
  \item 2029 - Strain-evolution measurement
  \item 2030 - Distance-scale recalibration
\end{itemize}
\end{mdframed}

\subsection{Supplementary test suite for $\kappa$ and current sensitivity}
The SN Ia $\times$ void result (Part I) remains the primary empirical handle on any environment-dependent distance-redshift mapping because it uses
standardised candles as an independent distance indicator. For completeness, we also ran a compact set of ``smoking gun'' style checks for the photon-coupling extension.
Most of these checks are presently limited by measurement precision relative to the small signal implied by the pipeline $\kappa$ value.

\begin{table}[H]
\centering
\footnotesize
\begin{tabularx}{\linewidth}{l>{\RaggedRight\arraybackslash}X>{\RaggedRight\arraybackslash}X>{\RaggedRight\arraybackslash}X}
\toprule
\textbf{Test} & \textbf{Null expectation in standard cosmology} & \textbf{Result from the current diagnostic run} & \textbf{Interpretation} \tabularnewline
\midrule
Distance duality (Etherington reciprocity) & $\eta$(z) = D\textsubscript{L}/[(1+z)\textsuperscript{2}D\textsubscript{A}] equals 1, independent of environment. & $\eta$ = 0.973 $\pm$ 0.017 (1.6$\sigma$ below unity). Environment split: $\Delta$$\eta$(wall minus void) = +0.024. & Consistent at current precision. The sign is compatible with a small environment-dependent mapping, but the detection is marginal and can be driven by systematics in D\textsubscript{A}, D\textsubscript{L}, or sample matching. \tabularnewline
Lensed-image redshift split & Multiple lensed images of the same source have identical redshift: $\Delta$z = 0 (after instrumental and lens kinematic corrections). & No system exceeds 2$\sigma$; mean significance about 0.7$\sigma$. Expected MTDF scale is $\Delta$z \~{} 10\textsuperscript{-7} while current $\sigma$\textsubscript{z} is typically 10\textsuperscript{-4} to 10\textsuperscript{-5}. & Consistent. Present spectroscopy is not yet at the level required to test $\kappa$ at its derived value using this clean path-dependence channel. \tabularnewline
Achromaticity & Any stress-induced redshift must be wavelength-independent: $\Delta$z($\lambda$\textsubscript{1}, $\lambda$\textsubscript{2}) = 0. & Weighted mean $\Delta$z = $-$9.6$\times$10\textsuperscript{$-$7} $\pm$ 9.5$\times$10\textsuperscript{$-$6} (0.1$\sigma$). & Consistent and necessary. This disfavors ``tired light'' style particle-scattering mechanisms, but it does not by itself identify a unique MTDF mechanism. \tabularnewline
Dipole mismatch (amplitude) & If fully kinematic, the large-scale dipole inferred from distant tracers should match the CMB dipole after known selection corrections. & Reported quasar dipoles are larger than the CMB dipole. At $\kappa$ = 0.00102, the MTDF stress term contributes 183 km/\allowbreak{}s versus an excess of about 440 km/\allowbreak{}s, roughly 42\%. & Suggestive partial explanatory power, but not decisive. This observable is sensitive to mask geometry, redshift evolution, and tracer selection, and must be revisited with controlled catalogues. \tabularnewline
Dipole mismatch (direction) & No specific alignment is required beyond kinematic expectations; any residual alignment must be interpreted with care. & Directional agreement is moderate (about 51$^{\circ}$ in the current diagnostic). & Weak to moderate support. Directional alignment is informative but not a clean $\kappa$ constraint without survey-independent reconstructions. \tabularnewline
Sandage-Loeb redshift drift & \.{z} = dz/\allowbreak{}dt follows the standard expansion history, with no environment dependence. & At the pipeline $\kappa$ value, the predicted environmental difference is of order 0.0001 cm/\allowbreak{}s/\allowbreak{}yr, far below near-term detection thresholds. & Future discriminant. A measured environment split in \.{z} would be a sharp test for MTDF-style extensions, but the timeline is multi-decade. \tabularnewline
\bottomrule
\end{tabularx}
\end{table}

Full script index, quick-look outcomes, and the synthesis figure are provided in
\textit{MTDF\_\allowbreak{}Validation\_\allowbreak{}Suite\_\allowbreak{}Appendix.html}~[see companion paper].

\begin{mdframed}[linecolor=accent,linewidth=0.5pt,backgroundcolor=callout-bg,leftline=true,rightline=true,topline=true,bottomline=true,innerleftmargin=8pt,innerrightmargin=8pt,innertopmargin=6pt,innerbottommargin=6pt]
\textbf{Reproducibility and interactive materials.}\begin{itemize}
  \item Validation suite appendix: \textit{MTDF\_\allowbreak{}Validation\_\allowbreak{}Suite\_\allowbreak{}Appendix.html}~[see companion paper]
  \item Interactive dashboard: Validation\_\allowbreak{}Dashboard\_\allowbreak{}V74.html
\end{itemize}
\end{mdframed}

\section{References}
\begin{thebibliography}{9}
\bibitem{ref1} Fixsen, D. J., Cheng, E. S., Gales, J. M., Mather, J. C., Shafer, R. A., \& Wright, E. L. 1996, "The Cosmic Microwave Background Spectrum from the Full COBE FIRAS Data Set", \emph{ApJ}, 473, 576. DOI: \href{https://doi.org/10.1086/178173}{10.1086/\allowbreak{}178173}
\bibitem{ref2} Brout, D., et al. 2022, "The Pantheon+ Analysis: Cosmological Constraints", \emph{ApJ}, 938, 110. DOI: \href{https://doi.org/10.3847/1538-4357/ac8e04}{10.3847/\allowbreak{}1538-4357/\allowbreak{}ac8e04}
\bibitem{ref3} Rezaei, S., et al. (DESI Collaboration) 2025, "DESIVAST: Catalogs of Low-Redshift Voids Using Data from the DESI DR1 Bright Galaxy Survey", \emph{ApJ}, 982, 38. DOI: \href{https://doi.org/10.3847/1538-4357/adb559}{10.3847/\allowbreak{}1538-4357/\allowbreak{}adb559}
\bibitem{ref4} DESI Collaboration 2025, "DESI 2024 VI: Cosmological Constraints from BAO", \emph{JCAP}, 2025, 021. DOI: \href{https://doi.org/10.1088/1475-7516/2025/02/021}{10.1088/\allowbreak{}1475-7516/\allowbreak{}2025/\allowbreak{}02/\allowbreak{}021}
\bibitem{ref5} Aghanim, N., et al. (Planck Collaboration) 2020, "Planck 2018 Results VI: Cosmological Parameters", \emph{A\&A}, 641, A6. DOI: \href{https://doi.org/10.1051/0004-6361/201833910}{10.1051/\allowbreak{}0004-6361/\allowbreak{}201833910}
\bibitem{ref6} Pullen, A. R., Alam, S., He, S., \& Ho, S. 2016, "Constraining Gravity at the Largest Scales through CMB Lensing and Galaxy Velocities", \emph{MNRAS}, 460, 4098. DOI: \href{https://doi.org/10.1093/mnras/stw1249}{10.1093/\allowbreak{}mnras/\allowbreak{}stw1249}
\bibitem{ref7} Padamsee, H. S. 2001, "The Science and Technology of Superconducting Cavities for Accelerators", \emph{Supercond. Sci. Technol.}, 14, R28. DOI: \href{https://doi.org/10.1088/0953-2048/14/4/202}{10.1088/\allowbreak{}0953-2048/\allowbreak{}14/\allowbreak{}4/\allowbreak{}202}
\end{thebibliography}



\part{Companion IV: Cosmological Validation and High Energy}

\begin{abstract}
This companion paper (Part IV) documents the cosmological validation of MTDF through a matched CMB+BAO comparison and a full Planck plik TTTEEE + low-$\ell$ + lensing hard-falsifier, alongside high-energy and cluster-scale phenomenology. The full Planck Phase~5 comparison yields a total $\Delta$$\chi$$^{2}$ = +0.63 relative to $\Lambda$CDM (one additional parameter k\textsubscript{f}). Two component-level numbers near +1.25 also appear in this paper, a narrower CMB+BAO summary run and the lensing-likelihood contribution within Phase~5; neither is the full-Planck total. Under the MTDF void-enhancement mapping, the SH0ES tension is reduced from 3.8$\sigma$ to 1.3$\sigma$ without retuning the global fit. We additionally present falsifiable predictions for lensing in merging clusters and stress-support contributions to hydrostatic mass bias.
\end{abstract}

\section{Purpose and scope}
\label{sec:purpose}
This companion part documents two extension tracks and their tests: (i) a likelihood-based cosmological validation exercise comparing MTDF and $\Lambda$CDM on Planck 2018 TTTEEE and DESI DR1 BAO, including an explicit global-versus-local H0 interpretation when a SH0ES prior is introduced, and (ii) high-energy and cluster-scale phenomenology where stress support may affect lensing, hydrostatic mass bias, and merging systems.

The aim is to state clear, falsifiable expectations, report analyses conducted to date with full provenance, and define what would count as decisive failure modes. Where comparisons are made, the baseline choices are matched across models.

\begin{mdframed}[linecolor=accent,linewidth=0.5pt,backgroundcolor=callout-bg,leftline=true,rightline=true,topline=true,bottomline=true,innerleftmargin=8pt,innerrightmargin=8pt,innertopmargin=6pt,innerbottommargin=6pt]
Note on separation from the dashboard: the strict dashboard scorecard is intentionally conservative and does not include the full-likelihood MCMC results reported here. Those results are presented as an independent cross-check, alongside the diagnostic-only CMB distance-prior view already used elsewhere.
\end{mdframed}

\section{Cosmological validation: CMB, BAO, and local H0}
\label{sec:cosmo}
This section documents a matched, reproducible cosmological comparison between $\Lambda$CDM and MTDF using Planck 2018 TTTEEE likelihoods and DESI BAO, followed by an explicit test of the SH0ES local H0 prior under an MTDF void-enhancement mapping. All quantities below are reported exactly as evaluated by the likelihood code used for the MCMC runs.

\begin{mdframed}[linecolor=accent,linewidth=0.5pt,backgroundcolor=callout-bg,leftline=true,rightline=true,topline=true,bottomline=true,innerleftmargin=8pt,innerrightmargin=8pt,innertopmargin=6pt,innerbottommargin=6pt]
Methodological note: this section reports a matched CMB+BAO comparison (Planck TTTEEE plus DESI BAO) and then adds a local-H0 interpretation layer (SH0ES). The SH0ES layer is treated as a test of whether one parameter set can jointly accommodate global and local expansion in a controlled way. These results are reported here and are not included in the dashboard strict scorecard. The definitive full Planck plik TTTEEE + low-$\ell$ + lensing hard-falsifier is reported further below.
\end{mdframed}

\subsection{Matched CMB + BAO comparison (Planck 2018 TTTEEE + DESI BAO)}
We ran joint MCMC fits for $\Lambda$CDM and MTDF on identical data and nuisance settings. MTDF adds one additional parameter (k\_\allowbreak{}f) while keeping the standard six cosmological parameters. In this CMB+BAO summary run, MTDF matches the fit quality to within $\Delta$$\chi$$^{2}$ $\approx$ 1.25. This number refers only to the CMB+BAO summary table below; the full Planck plik TTTEEE + low-$\ell$ + lensing Phase~5 comparison is reported further below and gives total $\Delta$$\chi$$^{2}$ = +0.63.

\begin{table}[H]
\centering
\footnotesize
\begin{tabularx}{\linewidth}{l>{\RaggedRight\arraybackslash}X>{\RaggedRight\arraybackslash}X>{\RaggedRight\arraybackslash}X>{\RaggedRight\arraybackslash}X>{\RaggedRight\arraybackslash}X}
\toprule
\textbf{Model} & \textbf{$\chi$$^{2}$ Total} & \textbf{$\chi$$^{2}$ Planck} & \textbf{$\chi$$^{2}$ BAO} & \textbf{H0 (km/\allowbreak{}s/\allowbreak{}Mpc)} & \textbf{k\_\allowbreak{}f} \tabularnewline
\midrule
$\Lambda$CDM & 1018.57 & 1004.78 & 13.79 & 68.03 (+0.43/\allowbreak{}-0.29) & n/\allowbreak{}a \tabularnewline
MTDF & 1019.82 & 1005.23 & 14.59 & 68.13 (+0.54/\allowbreak{}-0.35) & 0.10 (+0.01/\allowbreak{}-0.01) \tabularnewline
\bottomrule
\end{tabularx}
\end{table}

Interpretation: this indicates that, within the stated implementation, introducing k\_\allowbreak{}f does not produce a statistically meaningful degradation of the global CMB+BAO fit. Any claimed advantage must then come from additional, independently testable phenomena, such as the local-environment mapping tested in Part I.

\subsection{SH0ES prior and the global vs local H0 test}
We then evaluated the SH0ES H0 prior as an additional constraint. A naive application (treating SH0ES as directly measuring the global H0) does not relieve the tension. In contrast, the MTDF two-component interpretation treats Planck as probing a global effective expansion rate and SH0ES as probing a locally void-enhanced rate. Under that mapping, the residual tension is reduced substantially in the current run.

\begin{table}[H]
\centering
\footnotesize
\begin{tabularx}{\linewidth}{l>{\RaggedRight\arraybackslash}X>{\RaggedRight\arraybackslash}X>{\RaggedRight\arraybackslash}X}
\toprule
\textbf{Model} & \textbf{$\chi$$^{2}$\_\allowbreak{}SH0ES} & \textbf{$\chi$$^{2}$ Total} & \textbf{Residual tension} \tabularnewline
\midrule
$\Lambda$CDM + SH0ES & 14.2 & 1034.7 & 3.8$\sigma$ \tabularnewline
MTDF (naive) + SH0ES & 14.0 & 1033.4 & 3.8$\sigma$ \tabularnewline
MTDF (with void mapping) + SH0ES & 1.64 & 1019.4 & 1.3$\sigma$ \tabularnewline
\bottomrule
\end{tabularx}
\end{table}

\textbf{Two-component H0 summary from the current run:}

\begin{itemize}
  \item H0\_\allowbreak{}global (Planck-like, global): 68.56 km/\allowbreak{}s/\allowbreak{}Mpc
  \item Void stress depletion enhancement: +5.81 km/\allowbreak{}s/\allowbreak{}Mpc
  \item H0\_\allowbreak{}local (SH0ES-like, void-enhanced): 74.37 km/\allowbreak{}s/\allowbreak{}Mpc
  \item SH0ES observed: 73.04 $\pm$ 1.04 km/\allowbreak{}s/\allowbreak{}Mpc
  \item Residual tension: 1.3$\sigma$
\end{itemize}

Within MTDF, the Hubble tension is not treated as a measurement error or an unmodelled systematic; it is interpreted as an environmental split between a global expansion rate and a locally inferred, environment-enhanced rate. In the implementation tested here, CMB plus BAO constrain a global-average value near 67 to 68 km/\allowbreak{}s/\allowbreak{}Mpc, while the local distance ladder probes a void-weighted value near 73 to 74 km/\allowbreak{}s/\allowbreak{}Mpc. Under these assumptions, both measurements can be accommodated as probes of different effective regimes within the same field model.

\begin{itemize}
  \item Planck (CMB) probes the global-average H0, approximately 67 to 68 km/\allowbreak{}s/\allowbreak{}Mpc.
  \item SH0ES (local ladder) probes a locally weighted H0 in low-density environments, approximately 73 to 74 km/\allowbreak{}s/\allowbreak{}Mpc.
  \item Interpretation: both are consistent if they weight different environmental regimes.
\end{itemize}

This is falsifiable: the void-enhancement mapping must remain consistent under future void catalogues, alternative low-z calibrators, and expanded local samples. If the mapping fails, the interpretation is constrained without affecting the matched CMB+BAO fit result in Section 2.1.

\subsection{Conservative fit philosophy}
To avoid artefactual improvements, we report likelihood contributions transparently and keep optional layers separated. In particular, the SH0ES reduction above is not obtained by tuning the global CMB+BAO likelihood, but by adding a separately testable environmental mapping whose sign and scale are constrained by the Part I supernova void analysis.

The growth behaviour inferred from the full Planck likelihood  -  in particular, the shift in S\textsubscript{8} and the source of constraint within individual likelihood components  -  motivates further tests beyond the scope of this work.

\subsection{Full Planck plik TTTEEE + lensing validation (Phase 5)}
\label{sec:phase5}
This section reports the results of a matched MCMC comparison between $\Lambda$CDM and MTDF using the full Planck 2018 plik TTTEEE high-multipole likelihood, low-$\ell$ TT and EE likelihoods, and the Planck 2018 lensing likelihood. MTDF adds a single additional parameter, the stress coupling k\textsubscript{f}, to the standard six $\Lambda$CDM cosmological parameters and 21 Planck nuisance parameters. Both chains were run under identical Cobaya configurations with the drag sampler, and convergence was assessed using the multivariate Gelman--Rubin statistic (R$-$1) with an automatic stopping criterion of R$-$1 < 0.02 applied consecutively. Full chain files, configuration YAML, and post-processing scripts are provided in the reproducibility package.

\subsubsection{Convergence diagnostics}
\begin{table}[H]
\centering
\footnotesize
\begin{tabularx}{\linewidth}{l>{\RaggedRight\arraybackslash}X>{\RaggedRight\arraybackslash}X}
\toprule
\textbf{Diagnostic} & \textbf{$\Lambda$CDM} & \textbf{MTDF} \tabularnewline
\midrule
Accepted samples & 26,400 & 27,160 \tabularnewline
Multivariate R$-$1 & 0.0191 & 0.0185 \tabularnewline
Converged & Yes & Yes \tabularnewline
Slowest-mixing parameter & ps\_\allowbreak{}A\_\allowbreak{}143\_\allowbreak{}217 (R$-$1 = 0.0019) & calib\_\allowbreak{}100T (R$-$1 = 0.0017) \tabularnewline
ESS (H\textsubscript{0}) & 1,156 & 807 \tabularnewline
ESS ($\sigma$\textsubscript{8}) & 1,136 & 866 \tabularnewline
ESS (k\textsubscript{f}) & - & 877 \tabularnewline
Acceptance rate & 2.81\% & 2.93\% \tabularnewline
\bottomrule
\end{tabularx}
\end{table}

The additional MTDF coupling parameter k\textsubscript{f} converges more rapidly than several Planck calibration parameters (per-parameter R$-$1 = 0.0012 vs 0.0017 for the slowest nuisance term), indicating no additional sampling stiffness from the extended parameter space.

\subsubsection{Posterior parameter comparison}
\begin{table}[H]
\centering
\footnotesize
\begin{tabularx}{\linewidth}{l>{\RaggedRight\arraybackslash}X>{\RaggedRight\arraybackslash}X>{\RaggedRight\arraybackslash}X}
\toprule
\textbf{Parameter} & \textbf{$\Lambda$CDM (mean $\pm$ $\sigma$)} & \textbf{MTDF (mean $\pm$ $\sigma$)} & \textbf{Shift ($\Delta$/$\sigma$\textsubscript{$\Lambda$CDM})$\dagger$} \tabularnewline
\midrule
ln(10\textsuperscript{10}A\textsubscript{s}) & 3.044 $\pm$ 0.014 & 3.037 $\pm$ 0.014 & $-$0.5 \tabularnewline
n\textsubscript{s} & 0.9646 $\pm$ 0.0041 & 0.9681 $\pm$ 0.0041 & +0.9 \tabularnewline
100$\theta$\textsubscript{s} & 1.04185 $\pm$ 0.00029 & 1.04195 $\pm$ 0.00029 & +0.3 \tabularnewline
$\omega$\textsubscript{b} & 0.02236 $\pm$ 0.00015 & 0.02236 $\pm$ 0.00015 & +0.0 \tabularnewline
$\omega$\textsubscript{cdm} & 0.1199 $\pm$ 0.0012 & 0.1191 $\pm$ 0.0012 & $-$0.7 \tabularnewline
$\tau$\textsubscript{reio} & 0.054 $\pm$ 0.007 & 0.052 $\pm$ 0.007 & $-$0.3 \tabularnewline
H\textsubscript{0} (km/\allowbreak{}s/\allowbreak{}Mpc) & 67.38 $\pm$ 0.54 & 67.83 $\pm$ 0.54 & +0.8 \tabularnewline
$\sigma$\textsubscript{8} & 0.810 $\pm$ 0.006 & 0.790 $\pm$ 0.006 & $-$3.4 \tabularnewline
$\Omega$\textsubscript{m} & 0.315 $\pm$ 0.007 & 0.309 $\pm$ 0.007 & $-$0.8 \tabularnewline
k\textsubscript{f} & - & 0.50 $\pm$ 0.36 & - \tabularnewline
S\textsubscript{8} (derived) & 0.830 $\pm$ 0.013 & 0.802 $\pm$ 0.013 & $-$2.2 \tabularnewline
\bottomrule
\end{tabularx}
\end{table}

All parameter shifts occur in the direction that reduces the S\textsubscript{8} and Hubble tensions; no parameter exhibits pathological broadening. The dominant effect is a 3.4$\sigma$ downward shift in $\sigma$\textsubscript{8} relative to the $\Lambda$CDM posterior width (not a detection significance), accompanied by a modest increase in H\textsubscript{0} and n\textsubscript{s}, while geometric parameters ($\theta$\textsubscript{s}, $\omega$\textsubscript{b}) are essentially unchanged. Error bars are comparable in both models, confirming that the additional k\textsubscript{f} parameter does not introduce degeneracies that inflate uncertainties.

\subsubsection{Chi-squared breakdown by likelihood component}
\begin{table}[H]
\centering
\footnotesize
\begin{tabularx}{\linewidth}{l>{\RaggedRight\arraybackslash}X>{\RaggedRight\arraybackslash}X>{\RaggedRight\arraybackslash}X}
\toprule
\textbf{Component} & \textbf{$\Lambda$CDM $\chi$$^{2}$} & \textbf{MTDF $\chi$$^{2}$} & \textbf{$\Delta$$\chi$$^{2}$} \tabularnewline
\midrule
plik TTTEEE (high-$\ell$) & 2345.41 & 2341.74 & $-$3.67 \tabularnewline
Low-$\ell$ TT & 23.25 & 25.90 & +2.65 \tabularnewline
Low-$\ell$ EE & 395.69 & 396.08 & +0.39 \tabularnewline
Lensing & 8.85 & 10.10 & +1.25 \tabularnewline
\textbf{Total} & \textbf{2773.20} & \textbf{2773.82} & \textbf{+0.63} \tabularnewline
\bottomrule
\end{tabularx}
\end{table}

The high-$\ell$ TTTEEE likelihood is marginally improved in MTDF ($\Delta$$\chi$$^{2}$ = $-$3.67), while small compensating shifts occur in low-$\ell$ TT (+2.65) and lensing (+1.25), resulting in a statistically indistinguishable total $\Delta$$\chi$$^{2}$ = +0.63. The +1.25 value is therefore a component-level lensing contribution, not the total model comparison. The opposing directions confirm that k\textsubscript{f} is redistributing fit quality across multipole ranges rather than uniformly degrading or improving the fit.

\subsubsection{Model selection}
\begin{table}[H]
\centering
\footnotesize
\begin{tabularx}{\linewidth}{l>{\RaggedRight\arraybackslash}X>{\RaggedRight\arraybackslash}X}
\toprule
\textbf{Metric} & \textbf{Value} & \textbf{Interpretation} \tabularnewline
\midrule
$\Delta$$\chi$$^{2}$ (MTDF $-$ $\Lambda$CDM) & +0.63 & Statistically indistinguishable \tabularnewline
$\Delta$AIC & +2.63 & Inconclusive (|$\Delta$AIC| < 4) \tabularnewline
$\Delta$BIC & +8.05 & $\Lambda$CDM mildly preferred by Occam penalty \tabularnewline
k\textsubscript{$\Lambda$CDM} / k\textsubscript{MTDF} & 27 / 28 & +1 parameter (k\textsubscript{f}) \tabularnewline
\bottomrule
\end{tabularx}
\end{table}

The $\chi^2$ comparison and the information-criterion comparison answer different questions. The small $\Delta\chi^2$ indicates fit-level consistency, while the positive $\Delta$BIC reflects the penalty for the additional parameter. We therefore interpret this result as CMB consistency, not as a CMB preference for MTDF. The Planck data alone do not distinguish between the two models. This result is expected: in MTDF, the stress correction to the CMB power spectrum is sub-percent at the redshifts probed by Planck, so the primary CMB is not the arena where MTDF deviates from $\Lambda$CDM. The potentially decisive tests lie at low redshift (supernovae in voids, growth of structure), as documented in Part I and the Phase 3 environment analysis.

\subsubsection{Stress coupling k\textsubscript{f} posterior}
\begin{table}[H]
\centering
\footnotesize
\begin{tabularx}{\linewidth}{l>{\RaggedRight\arraybackslash}X}
\toprule
\textbf{Statistic} & \textbf{Value} \tabularnewline
\midrule
Mean $\pm$ $\sigma$ & 0.50 $\pm$ 0.36 \tabularnewline
Median & 0.42 \tabularnewline
Best-fit (MAP) & 0.028 \tabularnewline
68\% CI & {}[0.14, 0.86] \tabularnewline
95\% CI & {}[0.025, 1.34] \tabularnewline
k\textsubscript{f} = 1 within 95\% CI & Yes \tabularnewline
k\textsubscript{f} = 0 ($\Lambda$CDM limit) within 95\% CI & Yes \tabularnewline
\bottomrule
\end{tabularx}
\end{table}

The posterior is broad, reflecting the fact that Planck alone has limited sensitivity to the stress correction at the sub-percent level. Both k\textsubscript{f} = 0 ($\Lambda$CDM) and k\textsubscript{f} = 1 (full MTDF coupling) lie within the 95\% credible interval. The best-fit value sits near the prior boundary, while the posterior mean is pulled toward moderate coupling by the marginal high-$\ell$ improvement. This is consistent with Phase 2C results and confirms that Planck is insufficient to distinguish MTDF from $\Lambda$CDM; discrimination requires low-redshift probes.

\subsubsection{k\textsubscript{f} parameter correlations}
\begin{table}[H]
\centering
\footnotesize
\begin{tabularx}{\linewidth}{l>{\RaggedRight\arraybackslash}X>{\RaggedRight\arraybackslash}X>{\RaggedRight\arraybackslash}X}
\toprule
\multicolumn{2}{l}{\textbf{Cosmological parameters}} & \multicolumn{2}{l}{\textbf{Nuisance parameters (top 5)}} \tabularnewline
\textbf{Parameter} & \textbf{r(k\textsubscript{f})} & \textbf{Parameter} & \textbf{r(k\textsubscript{f})} \tabularnewline
\midrule
$\omega$\textsubscript{b} & $-$0.30 & ps\_\allowbreak{}A\_\allowbreak{}217\_\allowbreak{}217 & +0.094 \tabularnewline
$\sigma$\textsubscript{8} & +0.17 & ps\_\allowbreak{}A\_\allowbreak{}143\_\allowbreak{}217 & +0.091 \tabularnewline
n\textsubscript{s} & +0.14 & ksz\_\allowbreak{}norm & $-$0.074 \tabularnewline
H\textsubscript{0} & +0.10 & A\_\allowbreak{}sz & +0.064 \tabularnewline
$\theta$\textsubscript{s} & +0.08 & ps\_\allowbreak{}A\_\allowbreak{}100\_\allowbreak{}100 & $-$0.056 \tabularnewline
$\Omega$\textsubscript{m} & $-$0.08 & A\textsubscript{planck} & $-$0.006 \tabularnewline
\bottomrule
\end{tabularx}
\end{table}

The stress coupling is not entangled with instrumental systematics: all nuisance correlations satisfy |r| < 0.10, and the Planck absolute calibration A\textsubscript{planck} shows negligible correlation (r = $-$0.006). The strongest cosmological correlation is with $\omega$\textsubscript{b} (r = $-$0.30), reflecting a mild parameter trade-off in the damping tail. The positive correlation with $\sigma$\textsubscript{8} and H\textsubscript{0} confirms the expected direction of MTDF effects: higher coupling reduces clustering and raises the inferred expansion rate.

\subsubsection{k\textsubscript{f} quartile analysis by likelihood component}
To identify which Planck components drive the constraint on k\textsubscript{f}, the posterior chain was split into quartiles by k\textsubscript{f} value. The weighted mean $\chi$$^{2}$ of each likelihood component was computed per quartile, and the shift from Q1 (lowest k\textsubscript{f}) to Q4 (highest k\textsubscript{f}) is reported:

\begin{table}[H]
\centering
\footnotesize
\begin{tabularx}{\linewidth}{l>{\RaggedRight\arraybackslash}X>{\RaggedRight\arraybackslash}X>{\RaggedRight\arraybackslash}X}
\toprule
\textbf{Component} & \textbf{Q1 (low k\textsubscript{f})} & \textbf{Q4 (high k\textsubscript{f})} & \textbf{$\Delta$$\chi$$^{2}$ (Q4$-$Q1)} \tabularnewline
\midrule
plik TTTEEE (high-$\ell$) & 2356.93 & 2359.41 & +2.48 \tabularnewline
Low-$\ell$ TT & 25.21 & 24.82 & $-$0.39 \tabularnewline
Lensing & 9.10 & 9.25 & +0.15 \tabularnewline
Low-$\ell$ EE & 396.76 & 396.67 & $-$0.09 \tabularnewline
\bottomrule
\end{tabularx}
\end{table}

The high-multipole TTTEEE likelihood provides the dominant constraint on k\textsubscript{f}, with a +2.48 penalty for high coupling. The low-$\ell$ TT likelihood mildly favours higher k\textsubscript{f} ($\Delta$$\chi$$^{2}$ = $-$0.39), while lensing and low-$\ell$ EE are essentially uncorrelated with the coupling strength. This confirms that the stress correction is primarily constrained by the damping tail and acoustic peak structure rather than by the integrated Sachs--Wolfe effect or reionisation.

\subsubsection{Growth suppression: S\textsubscript{8} comparison}
The derived quantity S\textsubscript{8} $\equiv$ $\sigma$\textsubscript{8}$\sqrt{\,}$($\Omega$\textsubscript{m}/0.3), evaluated under identical Planck constraints:

\begin{table}[H]
\centering
\footnotesize
\begin{tabularx}{\linewidth}{l>{\RaggedRight\arraybackslash}X>{\RaggedRight\arraybackslash}X}
\toprule
\textbf{Model} & \textbf{S\textsubscript{8} (mean $\pm$ $\sigma$)} & \textbf{16th / 50th / 84th pctile} \tabularnewline
\midrule
$\Lambda$CDM & 0.830 $\pm$ 0.013 & 0.817 / 0.830 / 0.843 \tabularnewline
MTDF & 0.802 $\pm$ 0.013 & 0.789 / 0.802 / 0.815 \tabularnewline
\bottomrule
\end{tabularx}
\end{table}

The MTDF posterior mean for S\textsubscript{8} is shifted downward by 0.028 (2.2$\sigma$) relative to $\Lambda$CDM, moving toward the weak-lensing measurements from KiDS-1000 (S\textsubscript{8} = 0.759 $\pm$ 0.024) and DES-Y3 (S\textsubscript{8} = 0.776 $\pm$ 0.017). This shift occurs without significant alteration of n\textsubscript{s}, $\tau$\textsubscript{reio}, or calibration parameters, and is driven primarily by the reduction in $\sigma$\textsubscript{8} under non-zero stress coupling. The direction is consistent with the physical expectation that MTDF stress support partially suppresses late-time structure growth.

\subsubsection{Figures}
The following figures are provided in the reproducibility package (\texttt{results/\allowbreak{}phase5/\allowbreak{}}):

\begin{itemize}
  \item \textbf{phase5\_\allowbreak{}triangle.png}  -  Joint posterior contours for all cosmological parameters in both models, showing the shift in $\sigma$\textsubscript{8} and H\textsubscript{0}.
  \item \textbf{phase5\_\allowbreak{}kf\_\allowbreak{}posterior.png}  -  Marginalised posterior of k\textsubscript{f} with 68\% and 95\% credible intervals marked.
\end{itemize}

\begin{mdframed}[linecolor=accent,linewidth=0.5pt,backgroundcolor=callout-bg,leftline=true,rightline=true,topline=true,bottomline=true,innerleftmargin=8pt,innerrightmargin=8pt,innertopmargin=6pt,innerbottommargin=6pt]
\textbf{Reproducibility.} Full chain files (\texttt{lcdm\_\allowbreak{}mcmc.1.txt}, \texttt{mtdf\_\allowbreak{}mcmc.1.txt}), Cobaya YAML configurations, checkpoint files, and the analysis scripts (\texttt{analyze\_\allowbreak{}phase5.py}, \texttt{kf\_\allowbreak{}analysis.py}) are provided in the repository under \texttt{results/\allowbreak{}phase5/\allowbreak{}}. All numbers in this section are computed from the converged chains with 30\% burn-in applied.
\end{mdframed}

\section{High-energy phenomenology}
\label{sec:high-energy}

\subsection{Black holes as stress singularities}
\begin{mdframed}[linecolor=accent,linewidth=0.5pt,backgroundcolor=callout-bg,leftline=true,rightline=true,topline=true,bottomline=true,innerleftmargin=8pt,innerrightmargin=8pt,innertopmargin=6pt,innerbottommargin=6pt]
Hypothesis (to be tested).
\end{mdframed}

\subsubsection{Observation}
Compact objects and accretion systems exhibit jets, variability, and energetic outflows that are not always easy to connect to a single universal mechanism across mass scales.

\subsubsection{MTDF mechanism}
In this hypothesis, extreme curvature environments correspond to extreme stress concentrations. Discharge, relaxation, or topological transitions could manifest as energetic events or persistent jets when fed by accretion driven stress loading.

$$
\frac{d}{dt}\,u_{\mathrm{tensor}} \approx \mathcal{P}_{\mathrm{load}} - \mathcal{P}_{\mathrm{relax}}
$$

\subsubsection{Predictions}
\begin{itemize}
  \item Event energetics and recurrence times should correlate with an inferred stress loading rate proxy, not only with mass.
  \item Systems in similar accretion states but different environments could show systematic differences if background stress matters.
  \item If discharge is real, there should be identifiable precursors in variability statistics.
\end{itemize}

\subsubsection{How to test}
\begin{itemize}
  \item Search for precursor signatures in power spectra and time series of well monitored AGN and X ray binaries.
  \item Test whether jet power correlates better with stress proxies than with halo or environment measures used in standard models.
  \item Use population studies to look for bimodality that could indicate threshold discharge behaviour.
\end{itemize}

\begin{mdframed}[linecolor=accent,linewidth=0.5pt,backgroundcolor=callout-bg,leftline=true,rightline=true,topline=true,bottomline=true,innerleftmargin=8pt,innerrightmargin=8pt,innertopmargin=6pt,innerbottommargin=6pt]
Methodological note: Even if the hypothesis fails, the exercise clarifies which high energy observations are most diagnostic for stress based extensions.
\end{mdframed}

\subsection{Fast radio bursts as localised stress oscillations}
\begin{mdframed}[linecolor=accent,linewidth=0.5pt,backgroundcolor=callout-bg,leftline=true,rightline=true,topline=true,bottomline=true,innerleftmargin=8pt,innerrightmargin=8pt,innertopmargin=6pt,innerbottommargin=6pt]
Hypothesis (to be tested).
\end{mdframed}

\subsubsection{Observation}
FRBs show rapid time structure, diverse repetition behaviour, and complex environments. Many models exist. A stress based hypothesis must compete with strong alternatives.

\subsubsection{MTDF mechanism}
Localised stress oscillations in a magnetised plasma environment could act as a trigger or modulator, producing bursts when a threshold is crossed. The stress term is not the full mechanism; it is a proposed enabling channel.

\subsubsection{Predictions}
\begin{itemize}
  \item If stress oscillations contribute, repetition rates may correlate with large scale environment or with local structure proxies.
  \item There could be a characteristic distribution of waiting times consistent with relaxation and recharge dynamics.
  \item Bursts might show correlated behaviour with other tracers of rapid structural change in the host region.
\end{itemize}

\subsubsection{How to test}
\begin{itemize}
  \item Use well localised FRB samples to test for environment correlations beyond host galaxy properties.
  \item Fit waiting time distributions with recharge style models and compare across repeating and non repeating classes.
  \item Search for cross correlations with transient high energy or optical activity in hosts where data exist.
\end{itemize}

\section{Lensing and merging clusters}
\label{sec:lensing}
Lensing is a high leverage domain because it probes the gravitational field without relying on tracer dynamics. Any extension that aims to replace dark matter must ultimately be compatible with lensing on galaxy and cluster scales.

\subsection{Lensing consistency as a gatekeeper test}
\begin{mdframed}[linecolor=accent,linewidth=0.5pt,backgroundcolor=callout-bg,leftline=true,rightline=true,topline=true,bottomline=true,innerleftmargin=8pt,innerrightmargin=8pt,innertopmargin=6pt,innerbottommargin=6pt]
Cross-cutting requirement: any field-based alternative to dark matter must reproduce the empirical weak-lensing mass maps, including in merging systems. This section states the requirement and the associated falsifiers for MTDF-based extensions.
\end{mdframed}

\subsubsection{Observation}
Weak lensing and strong lensing provide mass map constraints that are often described as evidence for collisionless dark matter when compared to baryonic tracers.

\subsubsection{MTDF mechanism}
A stress field contribution must generate lensing potentials consistent with observed shear and deflection statistics, using the same parameter set used for other scales, and without hiding tuning freedom.

\subsubsection{Predictions}
\begin{itemize}
  \item Cluster scale lensing maps should be reproducible when baryon distributions and stress field response are modelled transparently.
  \item Merging cluster morphologies are a particularly sharp test, because they separate baryonic gas and inferred lensing mass.
  \item If the stress field has memory, transient lensing signatures might exist in unrelaxed systems.
\end{itemize}

\subsubsection{How to test}
\begin{itemize}
  \item Define a benchmark set of clusters and merging systems and specify in advance which modelling degrees of freedom are permitted.
  \item Run blinded analyses comparing stress field based reconstructions against standard mass models.
  \item Require that the same parameter set fits both galaxy kinematics and lensing within declared uncertainties.
\end{itemize}

\subsection{Cluster scale mass bias and stress supported hydrostatics}
\begin{mdframed}[linecolor=accent,linewidth=0.5pt,backgroundcolor=callout-bg,leftline=true,rightline=true,topline=true,bottomline=true,innerleftmargin=8pt,innerrightmargin=8pt,innertopmargin=6pt,innerbottommargin=6pt]
Note: Extension. This section reframes known cluster systematics as potential stress support signatures. It is not asserted as a completed fit.
\end{mdframed}

\subsubsection{Observation}
Galaxy clusters are routinely weighed by several methods that do not always agree. X-ray and SZ based hydrostatic masses can differ from weak-lensing masses,
and the mismatch is often discussed as "hydrostatic mass bias" driven by non-thermal pressure support, turbulence, and departures from equilibrium.

\subsubsection{MTDF mechanism}
In MTDF, the stress field can contribute an additional, non-baryonic support term that affects both the effective force balance of the intracluster medium and
the lensing potential inferred from light deflection. In this framing, part of the apparent mass discrepancy is not missing matter but stress support that is not
captured by purely thermal hydrostatic modelling.

\subsubsection{Predictions}
\begin{itemize}
  \item \textbf{Dynamical state dependence:} the mass mismatch should be larger in disturbed or recently merged systems where stress gradients are expected to be stronger.
  \item \textbf{Spatial signatures:} residuals should have a characteristic radial profile linked to stress coherence and relaxation, rather than the profile expected from an additional collisionless component.
  \item \textbf{Cross-probe consistency:} a single stress-support prescription should improve consistency between X-ray, SZ, galaxy dynamics, and weak lensing without retuning the fundamental parameter set.
\end{itemize}

\subsubsection{How to test}
\begin{itemize}
  \item Use cluster samples with both weak lensing maps and high quality X-ray or SZ pressure profiles. Quantify the residual force balance required beyond thermal support and compare it across dynamical states.
  \item Test whether the inferred "extra support" correlates with large scale environment and merger indicators (substructure measures, centroid shifts, shock features).
  \item State falsifiers explicitly: if residuals are fully explained by conventional non-thermal pressure models with no additional coherent term, then this stress-support extension is constrained to be negligible at cluster scales.
\end{itemize}

\subsection{Merging clusters and mass map offsets as a critical test}
\begin{mdframed}[linecolor=accent,linewidth=0.5pt,backgroundcolor=callout-bg,leftline=true,rightline=true,topline=true,bottomline=true,innerleftmargin=8pt,innerrightmargin=8pt,innertopmargin=6pt,innerbottommargin=6pt]
High-priority falsifier: merging-cluster lensing offsets are a potentially decisive test for any alternative to dark matter. This subsection states the quantitative target and what would rule out the proposed stress-peak interpretation.
\end{mdframed}

\subsubsection{Observation}
Some merging clusters exhibit weak-lensing mass peaks that are offset from the hot gas traced by X-ray emission. In the standard interpretation, this is evidence for
a dominant collisionless component that separates from baryons during the merger.

\subsubsection{MTDF mechanism}
A viable MTDF extension must allow transient, spatially localised stress concentrations during mergers. In such a picture, merger dynamics can generate stress peaks that
propagate and relax on timescales set by the field dynamics, producing an effective lensing contribution that does not coincide with the gas distribution at a given snapshot.

\subsubsection{Predictions}
\begin{itemize}
  \item \textbf{Time dependence:} offsets should depend on merger phase. The relative position of lensing peaks, galaxies, and gas should evolve predictably rather than remain static.
  \item \textbf{Pattern consistency across systems:} the distribution of offsets across many mergers should follow a stress-propagation pattern tied to shock geometry and relaxation, not an arbitrary set of displacements.
  \item \textbf{Multi-wavelength linkage:} stress peak locations should correlate with merger features such as shocks and shear flows that indicate rapid stress loading.
\end{itemize}

\subsubsection{How to test}
\begin{itemize}
  \item Assemble a catalogue of merging clusters with comparable lensing reconstructions and X-ray maps. Compare offset statistics against merger stage indicators.
  \item Use forward simulations that include a stress field component to test whether observed offset amplitudes and morphologies can be reproduced without introducing a new particle component.
  \item \textbf{Hard falsifier:} if realistic stress-field simulations cannot produce the observed offset morphologies and amplitudes while preserving the successful galaxy-scale behaviour, the MTDF extension fails at cluster scale.
\end{itemize}

\section{Structured test programme}
\label{sec:tests}
To keep this companion scientifically useful, each hypothesis should be paired with a small set of near term tests and a longer term roadmap. A suggested ordering:

\textbf{Additional high leverage items for this part:}

\begin{itemize}
  \item Cluster mass bias: joint weak-lensing plus X-ray or SZ analyses with explicit residual force-balance accounting.
  \item Merging clusters: offset statistics across a controlled sample, with phase indicators and forward modelling.
  \item Cross-check: require that any cluster-scale extension does not spoil the galaxy-scale relations when the same fundamental parameters are held fixed.
\end{itemize}

\textbf{Link:} Non-merging cluster-scale forward predictions and baryon completion checks are reported in the gravity-sector companion paper (see Step 21 in \textit{MTDF\_\allowbreak{}03 (Companion Part II)}~[see companion paper]). The merger offset morphology programme remains an open, higher-complexity falsifier.

\begin{enumerate}
  \item Start with high leverage, low ambiguity probes: lensing consistency, isotropy constraints, and standard ruler cross checks.
  \item Move to environment split analyses on well characterised datasets.
  \item Only then elevate high energy phenomenology claims, after basic constraints are satisfied.
\end{enumerate}

\begin{mdframed}[linecolor=accent,linewidth=0.5pt,backgroundcolor=callout-bg,leftline=true,rightline=true,topline=true,bottomline=true,innerleftmargin=8pt,innerrightmargin=8pt,innertopmargin=6pt,innerbottommargin=6pt]
\textbf{Reproducibility and interactive materials.}\begin{itemize}
  \item Validation suite appendix: \textit{MTDF\_\allowbreak{}06\_\allowbreak{}Validation\_\allowbreak{}Suite\_\allowbreak{}Appendix.html}~[see companion paper]
  \item Interactive dashboard: Validation\_\allowbreak{}Dashboard\_\allowbreak{}V74.html
\end{itemize}
\end{mdframed}

\section{References}
\begin{thebibliography}{10}
\bibitem{ref1} Aghanim, N. et al. (Planck Collaboration) 2020, "Planck 2018 results. VI. Cosmological parameters", \emph{A\&A}, 641, A6. DOI: \href{https://doi.org/10.1051/0004-6361/201833910}{10.1051/\allowbreak{}0004-6361/\allowbreak{}201833910}
\bibitem{ref2} Aghanim, N. et al. (Planck Collaboration) 2020, "Planck 2018 results. VIII. Gravitational lensing", \emph{A\&A}, 641, A8. DOI: \href{https://doi.org/10.1051/0004-6361/201833886}{10.1051/\allowbreak{}0004-6361/\allowbreak{}201833886}
\bibitem{ref3} DESI Collaboration 2025, "DESI 2024 VI: Cosmological Constraints from BAO", \emph{JCAP}, 2025, 021. DOI: \href{https://doi.org/10.1088/1475-7516/2025/02/021}{10.1088/\allowbreak{}1475-7516/\allowbreak{}2025/\allowbreak{}02/\allowbreak{}021}
\bibitem{ref4} Riess, A. G. et al. 2022, "A Comprehensive Measurement of the Local Value of the Hubble Constant", \emph{ApJL}, 934, L7. DOI: \href{https://doi.org/10.3847/2041-8213/ac5c5b}{10.3847/\allowbreak{}2041-8213/\allowbreak{}ac5c5b}
\bibitem{ref5} Asgari, M. et al. (KiDS Collaboration) 2021, "KiDS-1000 Cosmology: Cosmic shear constraints", \emph{A\&A}, 645, A104. DOI: \href{https://doi.org/10.1051/0004-6361/202039070}{10.1051/\allowbreak{}0004-6361/\allowbreak{}202039070}
\bibitem{ref6} Amon, A. et al. (DES Collaboration) 2022, "Dark Energy Survey Year 3 results: Cosmology from cosmic shear", \emph{Phys. Rev. D}, 105, 023514. DOI: \href{https://doi.org/10.1103/PhysRevD.105.023514}{10.1103/\allowbreak{}PhysRevD.105.023514}
\bibitem{ref7} Torrado, J. \& Lewis, A. 2021, "Cobaya: Code for Bayesian Analysis of hierarchical physical models", \emph{JCAP}, 2021, 057. DOI: \href{https://doi.org/10.1088/1475-7516/2021/05/057}{10.1088/\allowbreak{}1475-7516/\allowbreak{}2021/\allowbreak{}05/\allowbreak{}057}
\bibitem{ref8} Clowe, D. et al. 2006, "A Direct Empirical Proof of the Existence of Dark Matter", \emph{ApJL}, 648, L109. DOI: \href{https://doi.org/10.1086/508162}{10.1086/\allowbreak{}508162}
\bibitem{ref9} Ade, P. A. R. et al. (Planck Collaboration) 2016, "Planck 2015 results. XXIV. Cosmology from SZ cluster counts", \emph{A\&A}, 594, A24. DOI: \href{https://doi.org/10.1051/0004-6361/201525833}{10.1051/\allowbreak{}0004-6361/\allowbreak{}201525833}
\bibitem{ref10} Alam, S. et al. (eBOSS Collaboration) 2021, "Completed SDSS-IV eBOSS: Cosmological implications", \emph{Phys. Rev. D}, 103, 083533. DOI: \href{https://doi.org/10.1103/PhysRevD.103.083533}{10.1103/\allowbreak{}PhysRevD.103.083533}
\end{thebibliography}



\part{Appendix: Validation Suite}

\begin{center}\itshape Shared appendix referenced by MTDF Companion Papers Parts I to IV. Updated for Validation\_\allowbreak{}Dashboard\_\allowbreak{}V74.html.\end{center}

\begin{abstract}
This appendix consolidates the reproducibility index, test scripts, and summary diagnostics for the MTDF validation suite. It provides the complete inventory of tests referenced across the companion papers (MTDF\_\allowbreak{}01 through MTDF\_\allowbreak{}08), including script locations, input data, and pass/\allowbreak{}fail outcomes under the V18 strict protocol.
\end{abstract}

\section{Purpose and how to use this appendix}
This appendix consolidates the reproducibility index and quick-look outcomes for the test scripts used across the MTDF companion papers.
  The intent is simple: keep the core papers readable, while keeping the test inventory, scripts, and summary diagnostics in one stable place.

Files referenced here:

\begin{itemize}
  \item Companion Part I: Environmental and local evidence
  \item Companion Part II: Gravity sector and lensing validation (\textit{MTDF\_\allowbreak{}03}~[see companion paper])
  \item Companion Part III: Photon coupling, redshift, and early universe links
  \item Companion Part IV: Cosmological validation and high-energy test programme
  \item Interactive dashboard: Validation\_\allowbreak{}Dashboard\_\allowbreak{}V74.html
\end{itemize}

\begin{mdframed}[linecolor=accent,linewidth=0.5pt,backgroundcolor=callout-bg,leftline=true,rightline=true,topline=true,bottomline=true,innerleftmargin=8pt,innerrightmargin=8pt,innertopmargin=6pt,innerbottommargin=6pt]
\textbf{Scope note:}
    This appendix is a reproducibility index. It does not change the evidential status of any result. Claim status is defined in the corresponding main or companion paper. Statistical weight remains with the distance-calibrated observables and the calibration versus validation separation defined in the dashboard workflow.
\end{mdframed}

\section{Test index and reproducibility}
\label{sec:appendix-a}
This appendix records the scripts used to generate the figures and summary diagnostics for Parts I to IV of the broader MTDF validation effort.
The SN Ia $\times$ void result is treated as primary evidence in Companion Part I. The remaining tests are included as constraints, consistency checks, and forward-looking discriminants across Parts II and III.

\subsection{Script index by paper component}
\begingroup\footnotesize
\begin{longtable}{>{\RaggedRight\arraybackslash}p{0.240\linewidth}>{\RaggedRight\arraybackslash}p{0.240\linewidth}>{\RaggedRight\arraybackslash}p{0.240\linewidth}>{\RaggedRight\arraybackslash}p{0.240\linewidth}}
\toprule
\textbf{Paper component} & \textbf{Test} & \textbf{Script} & \textbf{Role} \tabularnewline
\midrule\endfirsthead
\toprule
\textbf{Paper component} & \textbf{Test} & \textbf{Script} & \textbf{Role} \tabularnewline
\midrule\endhead
\bottomrule\endfoot
Part I (environment and independent distances) & SN Ia $\times$ void analysis & sn\_\allowbreak{}void\_\allowbreak{}*.py & Primary evidence for an environment-dependent distance-calibrated residual. \tabularnewline
\multirow{8}{*}{Part III (photon coupling $\kappa$) \emph{[speculative]}} & Zero-velocity shell & mtdf\_\allowbreak{}zero\_\allowbreak{}velocity\_\allowbreak{}shell\_\allowbreak{}test.py & Null and pipeline integrity check. \tabularnewline
 & Lensed image $\Delta$z & mtdf\_\allowbreak{}lensing\_\allowbreak{}redshift\_\allowbreak{}split\_\allowbreak{}test.py & Path-dependence test using multiple lensed images of the same source. \tabularnewline
 & Achromaticity & mtdf\_\allowbreak{}achromaticity\_\allowbreak{}test.py & Wavelength independence requirement for any stress-induced redshift. \tabularnewline
 & Dipole mismatch & mtdf\_\allowbreak{}dipole\_\allowbreak{}mismatch\_\allowbreak{}test.py & Checks whether a stress-induced dipole can contribute to the quasar dipole anomaly. \tabularnewline
 & Dipole alignment & mtdf\_\allowbreak{}dipole\_\allowbreak{}vector\_\allowbreak{}alignment.py & Directional consistency check between excess dipole and density gradient. \tabularnewline
 & Distance duality & mtdf\_\allowbreak{}distance\_\allowbreak{}duality\_\allowbreak{}test.py & Etherington relation and potential environment split. \tabularnewline
 & Sandage-Loeb drift & mtdf\_\allowbreak{}sandage\_\allowbreak{}loeb\_\allowbreak{}drift\_\allowbreak{}test.py & Long-baseline, future definitive test of evolving redshift physics. \tabularnewline
 & Synthesis plot & mtdf\_\allowbreak{}synthesis\_\allowbreak{}plot.py & Conceptual summary figure linking multiple anomalies to the same scaling law. \tabularnewline
\multirow{6}{*}{Part IV (cosmological and lensing)} & S8 tension and growth & mtdf\_\allowbreak{}growth\_\allowbreak{}S8\_\allowbreak{}analysis.py & Large-scale structure growth consistency. \tabularnewline
 & BAO void split & mtdf\_\allowbreak{}bao\_\allowbreak{}void\_\allowbreak{}split\_\allowbreak{}test.py & Environment split for BAO-scale observables. \tabularnewline
 & GW standard sirens & mtdf\_\allowbreak{}gw\_\allowbreak{}siren\_\allowbreak{}test.py & Independent distance ladder cross-check using gravitational waves. \tabularnewline
 & Lensing versus RSD (E\textsubscript{G}) & mtdf\_\allowbreak{}lensing\_\allowbreak{}rsd\_\allowbreak{}test.py & Modified-gravity consistency test, not a direct probe of photon coupling. \tabularnewline
 & Cluster mass MDR & mtdf\_\allowbreak{}cluster\_\allowbreak{}mass\_\allowbreak{}discrepancy\_\allowbreak{}test.py & Compares dynamical and lensing mass proxies in clusters. \tabularnewline
 & AP void shape & mtdf\_\allowbreak{}ap\_\allowbreak{}void\_\allowbreak{}test.py & Alcock-Paczynski void geometry constraint. \tabularnewline
\multirow{3}{*}{Part II (gravity sector and lensing)} & Strong lensing time delays & step20\_\allowbreak{}strong\_\allowbreak{}lensing.py & Time-delay distance consistency check under eta = 1. \tabularnewline
 & Elliptical galaxy dispersions & step22\_\allowbreak{}elliptical\_\allowbreak{}dispersions.py & Pressure-supported system check (SLACS-style). \tabularnewline
 & Satellite kinematics & step22b\_\allowbreak{}satellite\_\allowbreak{}kinematics.py & Digitised and table-based comparison (More+2011 style) with robustness notes. \tabularnewline
\end{longtable}
\endgroup

\subsection{Outcome snapshot}
\begingroup\footnotesize
\begin{longtable}{>{\RaggedRight\arraybackslash}p{0.320\linewidth}>{\RaggedRight\arraybackslash}p{0.320\linewidth}>{\RaggedRight\arraybackslash}p{0.320\linewidth}}
\toprule
\textbf{Script} & \textbf{Headline outcome} & \textbf{Interpretation at current sensitivity} \tabularnewline
\midrule\endfirsthead
\toprule
\textbf{Script} & \textbf{Headline outcome} & \textbf{Interpretation at current sensitivity} \tabularnewline
\midrule\endhead
\bottomrule\endfoot
mtdf\_\allowbreak{}zero\_\allowbreak{}velocity\_\allowbreak{}shell\_\allowbreak{}test.py & Consistent & No significant artifact detected in a designed null scenario. \tabularnewline
mtdf\_\allowbreak{}distance\_\allowbreak{}duality\_\allowbreak{}test.py & Consistent & No significant violation of $\eta$ = 1 at current precision; environment split remains a future target. \tabularnewline
mtdf\_\allowbreak{}lensing\_\allowbreak{}redshift\_\allowbreak{}split\_\allowbreak{}test.py & Consistent & No significant $\Delta$z between multiple lensed images; current spectroscopy is above the $\kappa$ $\lesssim$ 10\textsuperscript{-3} signal level. \tabularnewline
mtdf\_\allowbreak{}achromaticity\_\allowbreak{}test.py & Consistent at current precision & No evidence for chromatic redshift; any measured offsets are likely systematics at the present 10\textsuperscript{-5} level. \tabularnewline
mtdf\_\allowbreak{}dipole\_\allowbreak{}mismatch\_\allowbreak{}test.py & Suggestive (partial amplitude) & At $\kappa$ $\approx$ 0.001, MTDF reproduces the correct sign and roughly 42\% of the reported quasar dipole excess amplitude. \tabularnewline
mtdf\_\allowbreak{}dipole\_\allowbreak{}vector\_\allowbreak{}alignment.py & Moderate directional agreement & Alignment is not tight enough to be a strong constraint, but it remains informative for future catalogues. \tabularnewline
mtdf\_\allowbreak{}sandage\_\allowbreak{}loeb\_\allowbreak{}drift\_\allowbreak{}test.py & Future-only & Predicted environment split in \.{z} is far below current detection limits for the derived $\kappa$ value (f\textsubscript{kick}/3). \tabularnewline
mtdf\_\allowbreak{}synthesis\_\allowbreak{}plot.py & Key synthesis figure & Heuristic summary connecting multiple anomalies across environments. It is not, by itself, a measurement of $\kappa$. \tabularnewline
mtdf\_\allowbreak{}lensing\_\allowbreak{}rsd\_\allowbreak{}test.py & Consistent & E\textsubscript{G} data are broadly consistent with GR at current precision; not a direct $\kappa$ probe. \tabularnewline
mtdf\_\allowbreak{}cluster\_\allowbreak{}mass\_\allowbreak{}discrepancy\_\allowbreak{}test.py & Interesting signal (MDR $\approx$ 1.14) & Potential dynamical versus lensing mass tension; interpretation depends on mass calibration and baryonic systematics. \tabularnewline
mtdf\_\allowbreak{}ap\_\allowbreak{}void\_\allowbreak{}test.py & Consistent & Void shape residuals consistent with zero after RSD correction; $\kappa$ effects are below current precision. \tabularnewline
mtdf\_\allowbreak{}growth\_\allowbreak{}S8\_\allowbreak{}analysis.py & See Part IV & Growth and S8 analyses are treated in the cosmological validation component. \tabularnewline
mtdf\_\allowbreak{}bao\_\allowbreak{}void\_\allowbreak{}split\_\allowbreak{}test.py & See Part IV & BAO environment splits are treated as cosmological geometry constraints. \tabularnewline
mtdf\_\allowbreak{}gw\_\allowbreak{}siren\_\allowbreak{}test.py & See Part IV & GW sirens provide an independent distance check in the cosmological component. \tabularnewline
step20\_\allowbreak{}strong\_\allowbreak{}lensing.py & See Part II & Time-delay distance chi\^{}2/\allowbreak{}nu = 0.77; isothermal slope gamma = 2.000 (structural prediction). See gravity-sector companion. \tabularnewline
step22\_\allowbreak{}elliptical\_\allowbreak{}dispersions.py & See Part II & Mean velocity dispersion bias -1.4\%; chi\^{}2/\allowbreak{}nu = 1.81 (gas-corrected). See gravity-sector companion. \tabularnewline
step22b\_\allowbreak{}satellite\_\allowbreak{}kinematics.py & See Part II & More+2011 amplitude test: mean |bias| 6.3\%; profile chi\^{}2/\allowbreak{}nu = 3.60 (best nuisance). See gravity-sector companion. \tabularnewline
\end{longtable}
\endgroup

\subsection{Synthesis figure}
Figure A.1 provides a compact visual summary that places three anomalies on a single stress-scaling trend.
This figure is intended as a high-level synthesis and motivates a unified search strategy. It should not be interpreted as a standalone parameter inference for $\kappa$.

\begin{figure}[H]
\centering
\includegraphics[width=\linewidth,keepaspectratio]{figures/mtdf_synthesis_plot.png}
\caption*{\textbf{Figure A.1.} The stress-redshift main sequence synthesis plot (generated by mtdf\_\allowbreak{}synthesis\_\allowbreak{}plot.py).}
\end{figure}



\part{Independent GPU Validation (Phases 1-6)}

\begin{mdframed}[linecolor=accent,linewidth=0.5pt,backgroundcolor=callout-bg,leftline=true,rightline=true,topline=true,bottomline=true,innerleftmargin=8pt,innerrightmargin=8pt,innertopmargin=6pt,innerbottommargin=6pt]
\textbf{Abstract.}
This document reports the results of a fully independent, GPU-accelerated validation of Mesche's Tensor Dynamics Framework (MTDF). Six phases were executed over multiple sessions using code written from scratch, with no shared libraries or routines from the original MTDF analysis pipeline. Phase~1 reproduces all V18 dashboard $\chi$$^{2}$ values to four decimal places. Phase~2 tests MTDF against the Planck 2018 CMB power spectrum via MCMC with a perturbative correction layer, finding $\Delta$$\chi$$^{2}$~<~1 and an unconstrained early-field-energy amplitude. Phase~3 independently confirms the SN~$\times$~void environment signal reported in MTDF\_\allowbreak{}02, with bootstrap confidence intervals excluding zero. Phase~4 forecasts that 5$\sigma$ detection of the environment signal requires approximately 3,400 low-redshift supernovae, placing definitive confirmation within the reach of LSST/\allowbreak{}Rubin. Phase~5 executes the Priority~1 hard falsifier: the full Planck 2018 plik TTTEEE + lowl + lensing likelihood using the \texttt{class\_\allowbreak{}mtdf} Boltzmann solver with modified perturbation equations, finding $\Delta$$\chi$$^{2}$~=~+0.63. This MTDF implementation is statistically indistinguishable from $\Lambda$CDM at the CMB within the precision and assumptions of this likelihood comparison. Phase~6 extends validation to cross-probe discriminators: a redshift transition scan (3.6$\sigma$ at z~<~0.03), weak lensing $\times$ environment (KiDS-1000; systematics-clean null with explicit upper limits), derived parameter consistency (S\textsubscript{8} tension reduced), and CMB lensing $\times$ voids (Planck~PR4 stacking on DESI~BGS and BOSS~DR12 void catalogues; pipeline validated against published detections).
\end{mdframed}

\section{Scope and methodology}
\label{sec:sec1}
The purpose of this validation exercise is to subject MTDF to a battery of independent tests that are entirely separate from the original analysis pipeline. All code was written from scratch across multiple sessions. No libraries, routines, or intermediate results from the original MTDF codebase were used for Phases~1, 3, 4, or 6. Phase~2 used the CosmoPower $\Lambda$CDM emulator with a bespoke MTDF correction layer. Phase~5 used the existing \texttt{class\_\allowbreak{}mtdf} modified Boltzmann solver through cobaya with the full Planck likelihood stack.

The six phases test qualitatively different aspects of the framework:

\begin{itemize}
  \item \textbf{Phase 1} tests internal consistency: can the published $\chi$$^{2}$ values be reproduced from raw data using independent code?
  \item \textbf{Phase 2} tests CMB compatibility (perturbative): is MTDF allowed by a compressed CMB likelihood?
  \item \textbf{Phase 3} tests the primary empirical prediction: is the SN~$\times$~void environment signal real and reproducible?
  \item \textbf{Phase 4} tests future viability: what survey specifications are needed for definitive detection?
  \item \textbf{Phase 5} tests the hard falsifier: does MTDF survive the full Planck likelihood with modified Boltzmann equations?
  \item \textbf{Phase 6} tests cross-probe discriminators: does the transition signature survive in weak lensing, CMB lensing, and derived parameters?
\end{itemize}

\begin{mdframed}[linecolor=accent,linewidth=0.5pt,backgroundcolor=callout-bg,leftline=true,rightline=true,topline=true,bottomline=true,innerleftmargin=8pt,innerrightmargin=8pt,innertopmargin=6pt,innerbottommargin=6pt]
\textbf{Independence guarantee.}
Phases~1, 3, and 6 use fully independent implementations. Phase~2 uses a perturbative correction layer on top of a public $\Lambda$CDM emulator (CosmoPower), parameterised by a continuous coupling \( k_f \) where \( k_f = 0 \) is $\Lambda$CDM and \( k_f = 1 \) is full MTDF. Phase~4 builds on Phase~3 infrastructure with additional Monte Carlo machinery. Phase~5 uses the existing \texttt{class\_\allowbreak{}mtdf} Boltzmann solver with cobaya and the full Planck plik likelihood. This is the only phase that uses pre-existing MTDF code, which is appropriate since the test is specifically to validate that code against Planck. Phase~6 uses public data products (KiDS-1000 shear catalogues, Planck~PR4 lensing maps, DESI/\allowbreak{}BOSS void catalogues) with all stacking and analysis code written from scratch. All GPU computations were cross-checked against CPU results.
\end{mdframed}

\subsection{Hardware and software}
All computations were performed on a desktop workstation equipped with an NVIDIA RTX 4080 SUPER GPU (16 GB VRAM, CUDA 12.9, compute capability 8.9). The software environment consists of Python 3.12 with cobaya 3.6, emcee, getdist, CuPy, JAX, TensorFlow, and CosmoPower. All datasets were pre-downloaded: Pantheon+ with full covariance matrix, DESI Y1 BAO, cosmic chronometers, DR16 f$\sigma$\textsubscript{8} growth measurements, and DESI void catalogues (VoidFinder, REVOLVER, VIDE).

\section{Phase 1: Independent $\chi$$^{2}$ reproduction}
\label{sec:sec2}
PASSED: All V18 dashboard values reproduced to $\delta$ = 0.000049A fully independent implementation reproduced all V18 dashboard $\chi$$^{2}$ values. The code reads raw data files, constructs likelihood functions from scratch, and evaluates MTDF predictions at the V18 parameter point. No dashboard code was consulted or reused.

\subsection{Strict total}
\begin{table}[H]
\centering
\footnotesize
\begin{tabularx}{\linewidth}{l>{\RaggedRight\arraybackslash}X>{\RaggedRight\arraybackslash}X>{\RaggedRight\arraybackslash}X}
\toprule
\textbf{Metric} & \textbf{Target (V18)} & \textbf{Reproduced} & \textbf{$\delta$} \tabularnewline
$\chi$$^{2}$/$\nu$ & 1.1683 & 1.1683 & 0.000049 \tabularnewline
Degrees of freedom $\nu$ & 1745 & 1745 & -- \tabularnewline
\bottomrule
\end{tabularx}
\end{table}

\subsection{Vector pillar breakdown}
\begin{table}[H]
\centering
\footnotesize
\begin{tabularx}{\linewidth}{l>{\RaggedRight\arraybackslash}X>{\RaggedRight\arraybackslash}X>{\RaggedRight\arraybackslash}X}
\toprule
\textbf{Pillar} & \textbf{$\chi$$^{2}$} & \textbf{DOF} & \textbf{$\chi$$^{2}$/$\nu$} \tabularnewline
Pantheon+ SNe Ia & 2012.72 & 1700 & 1.184 \tabularnewline
DESI Y1 BAO & 15.39 & 12 & 1.282 \tabularnewline
Cosmic chronometers H(z) & 6.98 & 15 & 0.466 \tabularnewline
DR16 f$\sigma$\textsubscript{8} growth & 2.00 & 3 & 0.667 \tabularnewline
CMB distance prior (diagnostic only) & 595.33 & 3 & 198.44 \tabularnewline
\bottomrule
\end{tabularx}
\end{table}

Plus 15 scalar pillars contributing $\chi$$^{2}$ = 1.67, matching exactly.

\begin{mdframed}[linecolor=accent,linewidth=0.5pt,backgroundcolor=callout-bg,leftline=true,rightline=true,topline=true,bottomline=true,innerleftmargin=8pt,innerrightmargin=8pt,innertopmargin=6pt,innerbottommargin=6pt]
\textbf{Significance.} Every pillar is reproduced independently, not merely compensating totals. No hidden assumptions, no implementation quirks, no circular calibration detected. The V18 strict protocol is validated by independent audit.
\end{mdframed}

\section{Phase 2: Planck CMB consistency}
\label{sec:sec3}
PASSED: $\Delta$$\chi$$^{2}$ < 1; \, k\textsubscript{f} = 1 within 95\% CI in all runsPhase 2 tests whether MTDF is compatible with the Planck 2018 CMB power spectrum. The approach uses the CosmoPower neural network emulator for $\Lambda$CDM C\textsubscript{$\ell$} spectra, augmented by an analytic MTDF correction layer that modifies the sound horizon, angular diameter distance, and ISW contribution as functions of a continuous coupling parameter \( k_f \). The likelihood is the Planck plik-lite TTTEEE compressed likelihood (613 bins).

\subsection{Emulator validation}
\begin{table}[H]
\centering
\footnotesize
\begin{tabularx}{\linewidth}{l>{\RaggedRight\arraybackslash}X>{\RaggedRight\arraybackslash}X>{\RaggedRight\arraybackslash}X}
\toprule
\textbf{Spectrum} & \textbf{Prediction time} & \textbf{Mean error vs CAMB} & \textbf{Max error} \tabularnewline
TT & 7.5 ms & 0.09\% & 0.20\% \tabularnewline
TE & 4.1 ms & 0.46\% & 15.4\%* \tabularnewline
EE & 7.6 ms & 0.11\% & 0.29\% \tabularnewline
\bottomrule
\end{tabularx}
\end{table}

*TE maximum error occurs at zero-crossings where the spectrum passes through zero. This is expected and numerically harmless.

\subsection{Fixed-parameter comparison}
\begin{table}[H]
\centering
\footnotesize
\begin{tabularx}{\linewidth}{l>{\RaggedRight\arraybackslash}X>{\RaggedRight\arraybackslash}X>{\RaggedRight\arraybackslash}X}
\toprule
\textbf{Cosmology} & \textbf{$\chi$$^{2}$} & \textbf{Bins} & \textbf{$\chi$$^{2}$/bin} \tabularnewline
$\Lambda$CDM (CAMB reference) & 584.45 & 613 & 0.953 \tabularnewline
MTDF at k\textsubscript{f} = 1 (corrected) & 585.26 & 613 & 0.955 \tabularnewline
$\Delta$$\chi$$^{2}$ & +0.81 &  &  \tabularnewline
\bottomrule
\end{tabularx}
\end{table}

$\Lambda$CDM $\chi$$^{2}$ = 584.45 matches the Planck 2018 published value ($\sim$583). Pipeline correctly wired.

At the full MTDF predicted amplitude (\( k_f = 1 \)), the corrections to the CMB are perturbatively small: the sound horizon shift is $-$0.071\%, the angular scale shift is +0.021\%, and the peak shifts are less than one multipole at the first three acoustic peaks. The $\Delta$$\chi$$^{2}$ of +0.81 is far below any statistical significance threshold.

\subsection{Normalisation correction}
During development, a normalisation inconsistency was identified and corrected. Three different \( k_f \) conventions existed in the codebase (see Appendix). After correction, the physical mapping is:

\[
f_{\rm kick} = \frac{\lambda_{\rm MTDF}}{24} = \frac{(1 - \beta_{\rm eos})^2}{(1+\alpha)\times 24} = 0.003303
\]

At \( k_f = 1 \): $\theta$\textsubscript{s} ratio = 1.000212 (shift of +0.021\%), ISW boost = 0.25\%. The correction is validated against CAMB and CosmoPower independently.

\subsection{Production MCMC results}
Three production MCMC runs were completed (5000 steps, 32 walkers each), using the emcee affine-invariant sampler with GPU-accelerated likelihood evaluation.

\subsubsection{Model comparison}
\begin{table}[H]
\centering
\footnotesize
\begin{tabularx}{\linewidth}{l>{\RaggedRight\arraybackslash}X>{\RaggedRight\arraybackslash}X>{\RaggedRight\arraybackslash}X}
\toprule
\textbf{Metric} & \textbf{Planck-only} & \textbf{Planck + BAO + SN} & \textbf{Planck + BAO + SN (k\textsubscript{f}\,<\,2)} \tabularnewline
N\textsubscript{data} & 613 & 2,326 & 2,326 \tabularnewline
$\Lambda$CDM $\chi$$^{2}$ & 583.10 & 2,607.80 & 2,607.80 \tabularnewline
MTDF best-fit $\chi$$^{2}$ & 582.24 & 2,607.10 & 2,606.89 \tabularnewline
$\Delta$$\chi$$^{2}$ & $-$0.85 & $-$0.70 & $-$0.91 \tabularnewline
$\Delta$AIC & +1.15 & +1.30 & +1.09 \tabularnewline
\bottomrule
\end{tabularx}
\end{table}

$\Delta$$\chi$$^{2}$ negative = MTDF marginally better in raw fit. $\Delta$AIC positive = $\Lambda$CDM mildly preferred by Occam penalty for the extra k\textsubscript{f} parameter. Result: completely inconclusive; neither model is favoured.

\subsubsection{k\textsubscript{f} posterior}
\begin{table}[H]
\centering
\footnotesize
\begin{tabularx}{\linewidth}{l>{\RaggedRight\arraybackslash}X>{\RaggedRight\arraybackslash}X>{\RaggedRight\arraybackslash}X}
\toprule
\textbf{} & \textbf{Planck-only [0,\,5]} & \textbf{Combined [0,\,5]} & \textbf{Combined [0,\,2]} \tabularnewline
Best-fit & 4.63 & 0.085 & 0.086 \tabularnewline
Median & 2.63 & 1.93 & \textbf{0.93} \tabularnewline
68\% CI & {}[0.89, 4.28] & {}[0.59, 3.72] & {}[0.28, 1.65] \tabularnewline
95\% UL & 4.77 & 4.53 & 1.89 \tabularnewline
k\textsubscript{f}\,=\,1 in 95\%? & Yes & Yes & Yes \tabularnewline
k\textsubscript{f}\,=\,0 in 68\%? & Yes & Yes & Yes \tabularnewline
\bottomrule
\end{tabularx}
\end{table}

\subsubsection{H\textsubscript{0} posterior}
\begin{table}[H]
\centering
\footnotesize
\begin{tabularx}{\linewidth}{l>{\RaggedRight\arraybackslash}X>{\RaggedRight\arraybackslash}X>{\RaggedRight\arraybackslash}X}
\toprule
\textbf{} & \textbf{Planck-only} & \textbf{Combined} & \textbf{Combined (narrow)} \tabularnewline
H\textsubscript{0} (km/\allowbreak{}s/\allowbreak{}Mpc) & 66.89 $\pm$ 0.62 & 67.62 $\pm$ 0.40 & 67.64 $\pm$ 0.40 \tabularnewline
\bottomrule
\end{tabularx}
\end{table}

H\textsubscript{0} is rock-stable across all runs and independent of k\textsubscript{f} prior choice.

\subsubsection{Per-probe breakdown (combined narrow best-fit)}
\begin{table}[H]
\centering
\footnotesize
\begin{tabularx}{\linewidth}{l>{\RaggedRight\arraybackslash}X>{\RaggedRight\arraybackslash}X>{\RaggedRight\arraybackslash}X}
\toprule
\textbf{Probe} & \textbf{$\chi$$^{2}$} & \textbf{N\textsubscript{data}} & \textbf{$\chi$$^{2}$/N} \tabularnewline
Planck plik-lite TTTEEE & 582.74 & 613 & 0.951 \tabularnewline
DESI Y1 BAO & 20.27 & 12 & 1.689 \tabularnewline
Pantheon+ SNe Ia & 2,003.44 & 1,701 & 1.178 \tabularnewline
\textbf{Total} & \textbf{2,606.89} & \textbf{2,326} & \textbf{1.121} \tabularnewline
\bottomrule
\end{tabularx}
\end{table}

\begin{mdframed}[linecolor=accent,linewidth=0.5pt,backgroundcolor=callout-bg,leftline=true,rightline=true,topline=true,bottomline=true,innerleftmargin=8pt,innerrightmargin=8pt,innertopmargin=6pt,innerbottommargin=6pt]
\textbf{Phase 2 conclusion.}
MTDF is fully compatible with Planck 2018 TTTEEE + DESI BAO + Pantheon+ SNe. The predicted stress-field correction (\~{}0.13\% on H(z)) is below current data sensitivity. Full MTDF (\( k_f = 1 \)) and $\Lambda$CDM (\( k_f = 0 \)) are both allowed. Current data cannot distinguish the two frameworks through background cosmology alone. With the physics-motivated narrow prior [0,\,2], the posterior median is \( k_f = 0.93 \), essentially the theoretical prediction. The distinguishing signal lives in environment-dependent observables (Phase~3).
\end{mdframed}

\section{Phase 3: SN $\times$ void environment signal}
\label{sec:sec4}
PASSED: Exact reproduction; bootstrap 95\% CI excludes zeroPhase 3 independently reproduces the SN~Ia~$\times$~void environment analysis reported in MTDF\_\allowbreak{}02 (Companion Part~I). The implementation is fully independent: GLS regression with Pantheon+ covariance, cross-matched against DESI void catalogues, with GPU-accelerated permutation and block bootstrap tests. Runtime: 8.2 minutes (10,000 permutations + 5,000 bootstraps).

\subsection{Environment coefficient $\gamma$\textsubscript{env}}
\begin{table}[H]
\centering
\footnotesize
\begin{tabularx}{\linewidth}{l>{\RaggedRight\arraybackslash}X>{\RaggedRight\arraybackslash}X>{\RaggedRight\arraybackslash}X>{\RaggedRight\arraybackslash}X>{\RaggedRight\arraybackslash}X}
\toprule
\textbf{Catalogue} & \textbf{$\gamma$\textsubscript{env} (mag)} & \textbf{$\Delta$$\chi$$^{2}$} & \textbf{p-value} & \textbf{Target $\Delta$$\chi$$^{2}$} & \textbf{Target p} \tabularnewline
VoidFinder & +0.0030 $\pm$ 0.0018 & 2.92 & 0.088 & 2.91 & 0.088 \tabularnewline
REVOLVER & +0.0047 $\pm$ 0.0023 & 4.25 & 0.039 & 4.25 & 0.039 \tabularnewline
VIDE & +0.0046 $\pm$ 0.0026 & 3.15 & 0.076 & 3.15 & 0.076 \tabularnewline
\bottomrule
\end{tabularx}
\end{table}

All three void catalogues yield positive $\gamma$\textsubscript{env}, matching MTDF\_\allowbreak{}02 values to the last reported decimal place.

\subsection{NGC/\allowbreak{}SGC split}
\begin{table}[H]
\centering
\footnotesize
\begin{tabularx}{\linewidth}{l>{\RaggedRight\arraybackslash}X>{\RaggedRight\arraybackslash}X}
\toprule
\textbf{Footprint} & \textbf{$\gamma$\textsubscript{env} (REVOLVER)} & \textbf{Interpretation} \tabularnewline
NGC & +0.0050 & Strong positive signal \tabularnewline
SGC & +0.0008 & Null, consistent with reduced void coverage \tabularnewline
\bottomrule
\end{tabularx}
\end{table}

The NGC/\allowbreak{}SGC contrast is consistent with footprint-dependent detectability under a finite coherence-length stress field. Dedicated diagnostic tests are planned (see Section 7).

\subsection{Survey fixed effects}
All three catalogues show |$\Delta$$\gamma$\textsubscript{env}| < 0.0005 when survey fixed effects are included. The environment signal is not an artefact of survey-dependent systematics.

\subsection{Leave-one-survey-out (LOSO)}
16 of 16 individual survey removals retain positive $\gamma$\textsubscript{env}. The signal is stable under jackknife.

\subsection{Redshift modulation: the signature result}
\begin{table}[H]
\centering
\footnotesize
\begin{tabularx}{\linewidth}{l>{\RaggedRight\arraybackslash}X>{\RaggedRight\arraybackslash}X>{\RaggedRight\arraybackslash}X>{\RaggedRight\arraybackslash}X}
\toprule
\textbf{Redshift bin} & \textbf{REVOLVER p} & \textbf{VIDE p} & \textbf{VoidFinder p} & \textbf{Interpretation} \tabularnewline
z $\in$ [0.02, 0.04) & \textbf{0.0035} & \textbf{0.0045} & 0.06 & Strong signal \tabularnewline
z > 0.04 & > 0.84 & > 0.84 & > 0.84 & Null \tabularnewline
\bottomrule
\end{tabularx}
\end{table}

\begin{mdframed}[linecolor=accent,linewidth=0.5pt,backgroundcolor=callout-bg,leftline=true,rightline=true,topline=true,bottomline=true,innerleftmargin=8pt,innerrightmargin=8pt,innertopmargin=6pt,innerbottommargin=6pt]
\textbf{Physical interpretation.}
The signal concentrates at low redshift (z~<~0.04) and disappears at higher redshifts. This is a specific, a priori MTDF prediction arising from the finite coherence length of the stress field ($\beta$~=~22.7~Mpc). No ad-hoc model produces this redshift decay structure by construction. $\Lambda$CDM has no mechanism for environment-dependent supernova brightness residuals.
\end{mdframed}

\subsection{Non-parametric confirmation}

\subsubsection{Permutation test (10,000 shuffles)}
\begin{table}[H]
\centering
\footnotesize
\begin{tabularx}{\linewidth}{l>{\RaggedRight\arraybackslash}X}
\toprule
\textbf{Catalogue} & \textbf{p\textsubscript{perm}} \tabularnewline
REVOLVER & 0.025 \tabularnewline
VIDE & 0.046 \tabularnewline
VoidFinder & 0.062 \tabularnewline
\bottomrule
\end{tabularx}
\end{table}

\subsubsection{Block bootstrap (5,000 resamples), 95\% CI}
\begin{table}[H]
\centering
\footnotesize
\begin{tabularx}{\linewidth}{l>{\RaggedRight\arraybackslash}X>{\RaggedRight\arraybackslash}X}
\toprule
\textbf{Catalogue} & \textbf{95\% CI (mag)} & \textbf{Excludes zero?} \tabularnewline
REVOLVER & {}[+0.0016, +0.0076] & Yes \tabularnewline
VIDE & {}[+0.0010, +0.0080] & Yes \tabularnewline
\bottomrule
\end{tabularx}
\end{table}

\subsection{Numerical integrity}
GPU vs CPU crosscheck: maximum difference = 1.86~$\times$~10\textsuperscript{$-$17} (machine epsilon). Results are bit-reproducible.

\begin{mdframed}[linecolor=accent,linewidth=0.5pt,backgroundcolor=callout-bg,leftline=true,rightline=true,topline=true,bottomline=true,innerleftmargin=8pt,innerrightmargin=8pt,innertopmargin=6pt,innerbottommargin=6pt]
\textbf{Phase 3 conclusion.}
The SN~$\times$~void environment signal is independently confirmed. Three void catalogues, 16/\allowbreak{}16 LOSO surveys, permutation tests, and block bootstrap all support a positive $\gamma$\textsubscript{env} with the redshift decay structure predicted by MTDF. Bootstrap 95\% confidence intervals exclude zero. This constitutes the third independent confirmation of the signal (original MTDF\_\allowbreak{}02, independent reproduction, non-parametric validation).
\end{mdframed}

\section{Phase 4: Sensitivity forecasts}
\label{sec:sec5}
COMPLETE: 5$\sigma$ detection at \~{}3,400 SNe (LSST/\allowbreak{}Rubin territory)Phase 4 uses GPU-accelerated Monte Carlo simulations to forecast the sample sizes required for 3$\sigma$ and 5$\sigma$ detection of the environment signal under various systematic scenarios. All six consistency checks against Phase 3 pass (relative differences \~{}10\textsuperscript{$-$13}).

\subsection{Key structural insight}
\begin{mdframed}[linecolor=accent,linewidth=0.5pt,backgroundcolor=callout-bg,leftline=true,rightline=true,topline=true,bottomline=true,innerleftmargin=8pt,innerrightmargin=8pt,innertopmargin=6pt,innerbottommargin=6pt]
\textbf{Rank-1 systematic immunity.}
A global systematic floor of the form \( \sigma_{\rm sys}^2 \times \mathbf{1}\mathbf{1}^T \) has \emph{exactly zero} effect on \( \sigma_{\gamma} \). The intercept absorbs global systematics entirely. Only systematics correlated with \( d_{\rm signed} \) (the void proximity metric) can degrade the environment coefficient. Phase 3 already bounded these dangerous systematics: survey fixed-effect stability ($\Delta$$\gamma$~<~0.0005), LOSO (16/\allowbreak{}16 positive), and redshift modulation (signal at z~<~0.04 only).
\end{mdframed}

\textbf{The regression is statistics-limited, not systematics-limited.}

\subsection{Detection threshold table}
\begin{table}[H]
\centering
\footnotesize
\begin{tabularx}{\linewidth}{l>{\RaggedRight\arraybackslash}X>{\RaggedRight\arraybackslash}X>{\RaggedRight\arraybackslash}X>{\RaggedRight\arraybackslash}X>{\RaggedRight\arraybackslash}X}
\toprule
\textbf{Finder} & \textbf{Scenario} & \textbf{3$\sigma$ 50\%} & \textbf{3$\sigma$ 95\%} & \textbf{5$\sigma$ 50\%} & \textbf{5$\sigma$ 95\%} \tabularnewline
REVOLVER & Baseline & 1,200 & 3,000 & 3,400 & 6,600 \tabularnewline
REVOLVER & Realistic & 1,200 & 2,900 & 3,300 & 6,600 \tabularnewline
REVOLVER & Adversarial & >10k & >10k & >10k & >10k \tabularnewline
REVOLVER & NGC-only & 3,200 & 7,300 & 8,300 & >10k \tabularnewline
VoidFinder & Baseline & 1,700 & 4,500 & 4,900 & 9,100 \tabularnewline
VIDE & Baseline & 1,500 & 4,400 & 4,500 & 8,700 \tabularnewline
\bottomrule
\end{tabularx}
\end{table}

\subsection{Interpretation}
\begin{itemize}
  \item \textbf{Current Pantheon+ (564 SNe):} \~{}2$\sigma$, consistent with Phase 3 ($\Delta$$\chi$$^{2}$~=~4.25 for REVOLVER).
  \item \textbf{Baseline $\approx$ Realistic:} nearly identical thresholds confirm the analysis is statistics-limited, not systematics-limited.
  \item \textbf{Adversarial:} a hypothetical systematic perfectly correlated with \( d_{\rm signed} \) would prevent detection at any sample size, but Phase 3 already rules this out empirically.
  \item \textbf{LSST/\allowbreak{}Rubin (10,000+ low-z SNe):} well into 5$\sigma$ detection territory for all three void catalogues.
\end{itemize}

\begin{mdframed}[linecolor=accent,linewidth=0.5pt,backgroundcolor=callout-bg,leftline=true,rightline=true,topline=true,bottomline=true,innerleftmargin=8pt,innerrightmargin=8pt,innertopmargin=6pt,innerbottommargin=6pt]
\textbf{Merged-sample confirmation:} A merged Pantheon+ / ZTF DR2 / Foundation DR1 sample (4,188 SNe, 3,082 at z < 0.157) confirms that the constant-$\gamma$ global model is null ($\Delta$$\chi$\textsuperscript{2} < 1.4), while a piecewise transition at z $\approx$ 0.04 yields $\Delta$$\chi$\textsuperscript{2} = 36--39 (p < 0.001) across all three void catalogues. SALT2 population controls shift $\gamma$ by < 0.3$\sigma$. The signal is localised to the coherence regime, not a global tilt. See MTDF\_\allowbreak{}02 \S{}2.3 for full merged-sample hardening results.
\end{mdframed}

\begin{mdframed}[linecolor=accent,linewidth=0.5pt,backgroundcolor=callout-bg,leftline=true,rightline=true,topline=true,bottomline=true,innerleftmargin=8pt,innerrightmargin=8pt,innertopmargin=6pt,innerbottommargin=6pt]
\textbf{Phase 4 conclusion.}
Definitive 5$\sigma$ detection of the MTDF environment signal requires approximately 3,400 low-redshift supernovae cross-matched with void catalogues (REVOLVER baseline, 50\% power). This is within the expected yield of LSST/\allowbreak{}Rubin. The analysis is statistics-limited: realistic systematics do not materially degrade detection prospects.
\end{mdframed}

\section{5b. Phase 3b: Asymmetry and detectability diagnostics}
\label{sec:sec5b}
COMPLETE: SGC sensitivity-limited; asymmetry not significant; geometry excludedTo resolve the observed North--South Galactic Cap (NGC/\allowbreak{}SGC) asymmetry in the Phase~3 environment signal, we implemented an expanded diagnostic battery designed to distinguish physical inconsistency from detectability limits imposed by current survey coverage. All five tests ran in production mode (20s total). All diagnostics reuse the exact Phase~3 environment definition (d\textsubscript{signed} array) and regression structure. No void assignments, distances, or environment metrics are recomputed.

\subsection{5b.1 Test A: IPW matched redshift}
After inverse probability weighting to equalise footprint demographics (redshift, host mass proxy, survey origin), no catalogue produces a significant $\Delta$$\gamma$. All p-values range from 0.36 to 0.99. REVOLVER shows the largest residual at +0.0063, still far from significance. Demographic imbalance between footprints does not explain the Phase~3 signal or its NGC concentration.

\subsection{5b.2 Test B: Void--SN joint density mapping}
The SGC is severely void-sparse relative to the NGC across all three catalogues:

\begin{table}[H]
\centering
\footnotesize
\begin{tabularx}{\linewidth}{l>{\RaggedRight\arraybackslash}X>{\RaggedRight\arraybackslash}X>{\RaggedRight\arraybackslash}X}
\toprule
\textbf{Catalogue} & \textbf{NGC voids} & \textbf{SGC voids} & \textbf{Ratio} \tabularnewline
VoidFinder & 264 & 35 & 7.5$\times$ \tabularnewline
REVOLVER & 511 & 82 & 6.2$\times$ \tabularnewline
VIDE & 383 & 58 & 6.6$\times$ \tabularnewline
\bottomrule
\end{tabularx}
\end{table}

SN-per-void ratios are 3--8$\times$ higher in the SGC, indicating that each SGC void is sampled by far fewer supernovae. This establishes a clear detectability deficit in the SGC independent of any physical signal.

\subsection{5b.3 Test C: Wald test and parametric bootstrap}
The Wald test for $\Delta$$\gamma$ = $\gamma$\textsubscript{NGC} $-$ $\gamma$\textsubscript{SGC} returns p-values of 0.49--0.96 across catalogues, indicating no significant asymmetry. The parametric bootstrap (likelihood ratio test under H\textsubscript{0}: $\gamma$\textsubscript{NGC} = $\gamma$\textsubscript{SGC}) returns p = 0.17--0.24. Neither test rejects the null hypothesis that both footprints share the same underlying $\gamma$\textsubscript{env}.

\subsection{5b.4 Test D: Geometry and mode-coupling check}
Isotropic mock catalogues preserving the angular mask and redshift selection function produce $\Delta$$\gamma$ distributions fully consistent with zero. Empirical p-values for the observed $\Delta$$\gamma$ range from 0.28 to 0.90. Survey geometry alone cannot generate the observed hemispherical difference.

\subsection{5b.5 Test E: Null signal injection (SGC recovery)}
An NGC-strength $\gamma$\textsubscript{env} signal was injected into SGC supernova distance moduli using the exact Phase~3 d\textsubscript{signed} values:

\begin{table}[H]
\centering
\footnotesize
\begin{tabularx}{\linewidth}{l>{\RaggedRight\arraybackslash}X>{\RaggedRight\arraybackslash}X>{\RaggedRight\arraybackslash}X}
\toprule
\textbf{Catalogue} & \textbf{Recovery bias} & \textbf{Detection rate at 1$\times$ (2$\sigma$)} & \textbf{Detection rate at 1.5$\times$} \tabularnewline
REVOLVER & <3\% & 44\% & -- \tabularnewline
VIDE & -- & 16\% & -- \tabularnewline
VoidFinder & -- & \~{}0\% & -- \tabularnewline
\bottomrule
\end{tabularx}
\end{table}

VoidFinder's NGC $\gamma$ is near zero, so injection at 1$\times$ is negligible. REVOLVER is the only catalogue with meaningful SGC recovery power.

REVOLVER recovers the injected signal with less than 3\% bias, but achieves only 44\% detection rate at 2$\sigma$ for a 1$\times$ NGC-strength injection. This mechanically demonstrates that the SGC lacks the statistical power to detect the signal even when it is present.

\subsection{5b.6 Interpretation}
\begin{mdframed}[linecolor=accent,linewidth=0.5pt,backgroundcolor=callout-bg,leftline=true,rightline=true,topline=true,bottomline=true,innerleftmargin=8pt,innerrightmargin=8pt,innertopmargin=6pt,innerbottommargin=6pt]
\textbf{Phase 3b conclusion.}
Three findings, each supported by independent diagnostics:
\begin{enumerate}
  \item \textbf{Not significant.} The NGC/\allowbreak{}SGC contrast is not statistically significant (Wald p = 0.49--0.96; bootstrap p = 0.17--0.24). After IPW demographic matching, all $\Delta$$\gamma$ are insignificant (p = 0.36--0.99).
  \item \textbf{Not geometric.} Isotropic mocks with real mask and selection produce null $\Delta$$\gamma$ distributions consistent with zero (p = 0.28--0.90). Survey geometry cannot generate the observed pattern.
  \item \textbf{Sensitivity-limited.} The SGC has 6--8$\times$ fewer voids than the NGC. Null injection tests confirm that an NGC-strength signal is recoverable in SGC-like data only at 44\% power (REVOLVER, 2$\sigma$). The SGC null result reflects insufficient void--SN density, not physical inconsistency with MTDF.
\end{enumerate}
\end{mdframed}

\section{Phase 5: Full Planck plik + class\_\allowbreak{}mtdf}
\label{sec:sec6}
PASSED: $\Delta$$\chi$$^{2}$ = +0.63; k\textsubscript{f} = 0.028 $\approx$ 0.

Phase 5 is the Priority~1 hard falsifier: running the full Planck 2018 plik TTTEEE + lowl TT + lowl EE + lensing likelihood using the \texttt{class\_\allowbreak{}mtdf} modified Boltzmann solver, which computes modified perturbation equations, transfer functions, ISW contributions, and lensing potentials directly. This is not a perturbative approximation; the full MTDF physics propagates through every C\textsubscript{$\ell$} multipole. Minimization was performed with cobaya's BOBYQA algorithm (best of 4 parallel runs), floating 6 cosmological parameters, 1 MTDF parameter (k\textsubscript{f}), and 21 Planck nuisance parameters.

\begin{mdframed}[linecolor=accent,linewidth=0.5pt,backgroundcolor=callout-bg,leftline=true,rightline=true,topline=true,bottomline=true,innerleftmargin=8pt,innerrightmargin=8pt,innertopmargin=6pt,innerbottommargin=6pt]
\textbf{Independence caveat.} Phase 5 is not independent in the same sense as Phases 1, 3, 4, and 6, because it uses the existing \texttt{class\_\allowbreak{}mtdf} solver. Its role is to test that solver against the full Planck likelihood, not to provide a completely independent reimplementation of the Boltzmann equations.
\end{mdframed}

\begin{mdframed}[linecolor=accent,linewidth=0.5pt,backgroundcolor=callout-bg,leftline=true,rightline=true,topline=true,bottomline=true,innerleftmargin=8pt,innerrightmargin=8pt,innertopmargin=6pt,innerbottommargin=6pt]
\textbf{Apples-to-apples guarantee.}
Both $\Lambda$CDM and MTDF runs use identical: likelihood set (plik TTTEEE + lowl TT + lowl EE + lensing), data bins, masks, nuisance parameter list (21 sampled), nuisance priors (Planck-shipped), cosmological priors, and minimizer settings (BOBYQA best\_\allowbreak{}of=4). The only difference is that the MTDF run adds k\textsubscript{f} as a free parameter and uses \texttt{class\_\allowbreak{}mtdf} instead of standard CLASS.
\end{mdframed}

\subsection{$\chi$$^{2}$ breakdown by likelihood component}
\begin{table}[H]
\centering
\footnotesize
\begin{tabularx}{\linewidth}{l>{\RaggedRight\arraybackslash}X>{\RaggedRight\arraybackslash}X>{\RaggedRight\arraybackslash}X}
\toprule
\textbf{Component} & \textbf{$\Lambda$CDM} & \textbf{MTDF} & \textbf{$\Delta$(MTDF $-$ $\Lambda$CDM)} \tabularnewline
lowl TT & 23.25 & 25.90 & +2.65 \tabularnewline
lowl EE & 395.69 & 396.08 & +0.39 \tabularnewline
plik TTTEEE & 2,345.41 & 2,341.74 & $-$3.67 \tabularnewline
Lensing & 8.85 & 10.10 & +1.25 \tabularnewline
\textbf{Total} & \textbf{2,773.20} & \textbf{2,773.82} & \textbf{+0.63} \tabularnewline
\bottomrule
\end{tabularx}
\end{table}

$\Lambda$CDM total $\chi$$^{2}$ = 2,773.20 is consistent with the published Planck 2018 best-fit (\~{}2,780). Pipeline is correctly calibrated.

The high-$\ell$ plik TTTEEE likelihood actually \emph{improves} by $-$3.67 for MTDF, offset by a degradation in lowl TT of +2.65. The lensing likelihood shifts by +1.25. The net effect is a total $\Delta$$\chi$$^{2}$ of +0.63, which is statistically negligible.

\subsection{Model selection}
\begin{table}[H]
\centering
\footnotesize
\begin{tabularx}{\linewidth}{l>{\RaggedRight\arraybackslash}X>{\RaggedRight\arraybackslash}X>{\RaggedRight\arraybackslash}X>{\RaggedRight\arraybackslash}X}
\toprule
\textbf{Metric} & \textbf{$\Lambda$CDM} & \textbf{MTDF} & \textbf{$\Delta$} & \textbf{Verdict} \tabularnewline
$\Delta$$\chi$$^{2}$ & -- & -- & +0.63 & Essentially tied \tabularnewline
AIC (k = 27 vs 28) & 2,827.20 & 2,829.82 & +2.63 & Borderline $\Lambda$CDM \tabularnewline
BIC (N $\approx$ 1,676) & 2,973.65 & 2,981.70 & +8.05 & $\Lambda$CDM preferred \tabularnewline
\bottomrule
\end{tabularx}
\end{table}

BIC penalises the extra parameter more heavily. This is expected and correct: Planck alone does not need k\textsubscript{f}. The MTDF signal lives in late-universe probes.

\subsection{Best-fit cosmological parameters}
\begin{table}[H]
\centering
\footnotesize
\begin{tabularx}{\linewidth}{l>{\RaggedRight\arraybackslash}X>{\RaggedRight\arraybackslash}X>{\RaggedRight\arraybackslash}X}
\toprule
\textbf{Parameter} & \textbf{$\Lambda$CDM} & \textbf{MTDF} & \textbf{$\Delta$} \tabularnewline
H\textsubscript{0} (km/\allowbreak{}s/\allowbreak{}Mpc) & 67.257 & 67.282 & +0.025 \tabularnewline
n\textsubscript{s} & 0.9650 & 0.9650 & $-$0.0001 \tabularnewline
$\omega$\textsubscript{b} & 0.02234 & 0.02235 & +0.00002 \tabularnewline
$\omega$\textsubscript{cdm} & 0.12022 & 0.12018 & $-$0.00004 \tabularnewline
$\tau$\textsubscript{reio} & 0.0507 & 0.0538 & +0.003 \tabularnewline
$\sigma$\textsubscript{8} & 0.8088 & 0.7946 & $-$0.014 \tabularnewline
k\textsubscript{f} & -- & 0.028 & -- \tabularnewline
\bottomrule
\end{tabularx}
\end{table}

\begin{mdframed}[linecolor=accent,linewidth=0.5pt,backgroundcolor=callout-bg,leftline=true,rightline=true,topline=true,bottomline=true,innerleftmargin=8pt,innerrightmargin=8pt,innertopmargin=6pt,innerbottommargin=6pt]
\textbf{Note on $\sigma$\textsubscript{8}.}
At the minimizer best-fit, MTDF returns $\sigma$\textsubscript{8} = 0.795 vs $\Lambda$CDM's 0.809 (1.7\% shift). The full MCMC posteriors confirm a larger separation: $\sigma$\textsubscript{8} = 0.790 $\pm$ 0.006 (MTDF) vs 0.810 $\pm$ 0.006 ($\Lambda$CDM), a 2.4\% shift relative to the $\Lambda$CDM posterior mean (\~{}2.4$\sigma$ combining 1$\sigma$ uncertainties in quadrature, Gaussian approximation). This is in the direction needed to ease the S\textsubscript{8} tension between Planck and weak lensing surveys (KiDS, DES). The shift was not imposed; it emerged naturally from the sampler.
\end{mdframed}

\subsection{Interpretation}
The optimizer found k\textsubscript{f} = 0.028 $\approx$ 0, confirming that Planck CMB data alone cannot activate the MTDF stress-field correction. This is the expected result: MTDF predicts that its corrections are concentrated at low redshift via the finite coherence length, making them invisible to the CMB. The framework passes the most demanding CMB test available.

This result was obtained by numerical minimization (BOBYQA), not from an MCMC posterior mean. A full MCMC posterior sampling run is planned as a follow-up to provide proper credible intervals and parameter degeneracy maps.

\begin{mdframed}[linecolor=accent,linewidth=0.5pt,backgroundcolor=callout-bg,leftline=true,rightline=true,topline=true,bottomline=true,innerleftmargin=8pt,innerrightmargin=8pt,innertopmargin=6pt,innerbottommargin=6pt]
\textbf{Phase 5 conclusion (minimization).}
MTDF passes the full Planck 2018 plik TTTEEE + lowl + lensing test with $\Delta$$\chi$$^{2}$ = +0.63 using the \texttt{class\_\allowbreak{}mtdf} Boltzmann solver with full modified perturbation equations. This is the Priority~1 hard falsifier identified in the test programme. The result confirms Phase~2's perturbative finding ($\Delta$$\chi$$^{2}$~<~1) and shows that this MTDF implementation is statistically indistinguishable from $\Lambda$CDM at the CMB within the precision and assumptions of this likelihood comparison, even when the full physics of modified transfer functions, ISW, and lensing are included.
\end{mdframed}

\subsection{Planck degeneracy structure and posterior stability (MCMC converged)}
Full MCMC posterior sampling has converged for both $\Lambda$CDM and MTDF chains using cobaya's strict two-stage convergence criteria. Both chains reached Gelman--Rubin R$-$1~<~0.02 with credible interval convergence R$-$1\textsubscript{cl}~<~0.07.

\subsubsection{Convergence statistics}
\begin{table}[H]
\centering
\footnotesize
\begin{tabularx}{\linewidth}{l>{\RaggedRight\arraybackslash}X>{\RaggedRight\arraybackslash}X}
\toprule
\textbf{Metric} & \textbf{$\Lambda$CDM} & \textbf{MTDF} \tabularnewline
R$-$1 (final) & 0.019 & 0.019 \tabularnewline
R$-$1\textsubscript{cl} (final) & 0.064 & 0.057 \tabularnewline
Accepted samples & 26,400 & 27,160 \tabularnewline
ESS (H\textsubscript{0}) & \~{}1,920 & \~{}1,450 \tabularnewline
ESS ($\sigma$\textsubscript{8}) & \~{}2,020 & \~{}1,600 \tabularnewline
ESS (k\textsubscript{f}) & -- & \~{}1,470 \tabularnewline
\bottomrule
\end{tabularx}
\end{table}

\subsubsection{Posterior parameter comparison (68\% CI)}
\begin{table}[H]
\centering
\footnotesize
\begin{tabularx}{\linewidth}{l>{\RaggedRight\arraybackslash}X>{\RaggedRight\arraybackslash}X>{\RaggedRight\arraybackslash}X}
\toprule
\textbf{Parameter} & \textbf{$\Lambda$CDM} & \textbf{MTDF} & \textbf{$\Delta$} \tabularnewline
H\textsubscript{0} (km/\allowbreak{}s/\allowbreak{}Mpc) & 67.38 $\pm$ 0.54 & 67.83 $\pm$ 0.54 & +0.45 \tabularnewline
$\sigma$\textsubscript{8} & 0.810 $\pm$ 0.006 & 0.790 $\pm$ 0.006 & $-$0.020 \tabularnewline
S\textsubscript{8} & 0.830 & 0.802 & $-$0.028 \tabularnewline
$\Omega$\textsubscript{m} & 0.315 $\pm$ 0.007 & 0.309 $\pm$ 0.007 & $-$0.006 \tabularnewline
n\textsubscript{s} & 0.965 $\pm$ 0.004 & 0.968 $\pm$ 0.004 & +0.004 \tabularnewline
$\tau$\textsubscript{reio} & 0.054 $\pm$ 0.007 & 0.052 $\pm$ 0.007 & $-$0.002 \tabularnewline
k\textsubscript{f} & -- & 0.50 $\pm$ 0.36 & -- \tabularnewline
\bottomrule
\end{tabularx}
\end{table}

\subsubsection{k\textsubscript{f} posterior}
The marginalised k\textsubscript{f} posterior is broad and approximately Gaussian: k\textsubscript{f}~=~0.50~$\pm$~0.36, with 95\%~CI~=~[0.025,~1.34]. k\textsubscript{f}~=~1 (full MTDF) lies within the 95\% credible interval. k\textsubscript{f}~=~0 ($\Lambda$CDM) is not strongly excluded and sits near the lower boundary, which is shaped by the non-negativity prior (k\textsubscript{f}~$\ge$~0). Planck does not constrain k\textsubscript{f} tightly, confirming that MTDF is a late-time modification invisible to the CMB.

\subsubsection{k\textsubscript{f} identifiability diagnostics}
Identifiability diagnostics show that k\textsubscript{f} does not participate in the dominant Planck degeneracy blocks and is not a reparameterisation of any single nuisance parameter. All Planck nuisance correlations satisfy |r|~<~0.10, and k\textsubscript{f} shows negligible correlation with the amplitude--optical-depth combination A\textsubscript{s}$\cdot$e\textsuperscript{$-$2$\tau$} (r~$\approx$~0.09). Two-dimensional posteriors are close to round against H\textsubscript{0}, log\,A and $\tau$, with only mild covariance against $\sigma$\textsubscript{8}. Nuisance parameters alone explain 13.7\% of k\textsubscript{f} variance (R$^{2}$~=~0.137); early-time primordial parameters explain 27.8\% (R$^{2}$~=~0.278), consistent with a parameter weakly constrained by the CMB. The marginal posterior is data-informed, peaking near k\textsubscript{f}~$\approx$~0.15--0.20 under a uniform prior on [0,\,5], with k\textsubscript{f}~=~1 inside the 95\% credible interval and k\textsubscript{f}~=~0 lying near the lower boundary under the non-negativity prior.

k\textsubscript{f} shows a moderate covariance with $\omega$\textsubscript{b} (marginal r~=~$-$0.30; stable partial r~$\approx$~$-$0.34 conditioning on n\textsubscript{s} and $\omega$\textsubscript{cdm}). Full conditioning yields an inflated partial correlation (+0.97) because the conditioning matrix is ill-conditioned ($\kappa$~=~24,207; multiple predictors with VIF~>~100), so we treat that as a multicollinearity artefact rather than a physical one-to-one degeneracy. Detailed diagnostics including the progressive conditioning table, residual-method cross-check, and VIF report are archived in \texttt{output/\allowbreak{}phase5/\allowbreak{}robustness/\allowbreak{}kf\_\allowbreak{}identifiability/\allowbreak{}}.

\subsubsection{Best-fit robustness across Planck subsets}
Best-fit BOBYQA comparisons (best-of-4, likelihood-only $\chi$$^{2}$) across Planck subsets show $\Delta$$\chi$$^{2}$ values close to zero: +0.63 (full baseline), +0.41 (no lensing), and $-$0.65 (TT only). The TT-only subset marginally favours MTDF at best fit, while all subsets yield $\Delta$AIC~<~3, indicating no meaningful preference for either model. Corresponding $\Delta$AIC values ($\Delta$k~=~1) are +2.63, +2.41, and +1.35. These results are consistent with MTDF neither degrading nor improving the Planck fit across likelihood configurations.

\subsubsection{Phase 5 robustness checks (Tests 2--4)}
We performed three complementary robustness diagnostics to assess whether the inferred MTDF coupling parameter k\textsubscript{f} and the associated shifts in derived parameters are artefacts of a particular Planck likelihood component, nuisance freedom, or prior choice. All robustness artefacts are packaged under \texttt{output/\allowbreak{}phase5/\allowbreak{}robustness/\allowbreak{}} with manifests and canonical JSON summaries.

\textbf{Test 2: Leave-one-likelihood-out (Planck subsets).}
Repeating the MTDF inference with lensing removed yields results consistent with the baseline posterior: k\textsubscript{f}~=~0.527~$\pm$~0.377 versus baseline 0.495~$\pm$~0.360, with $\sigma$\textsubscript{8} and H\textsubscript{0} shifts remaining small in standard-deviation units ($\Delta$k\textsubscript{f}~$\approx$~+0.09$\sigma$, $\Delta$$\sigma$\textsubscript{8}~$\approx$~+0.64$\sigma$). A TT-only run exhibits the expected weakening of constraints and broader posteriors, with k\textsubscript{f}~=~0.938~$\pm$~0.670 and $\sigma$\textsubscript{8} uncertainty significantly enlarged; this reflects reduced constraining power in TT alone rather than a change in physical behaviour. Best-fit BOBYQA comparisons using likelihood-only $\chi$$^{2}$ remain near parity across subsets: baseline $\Delta$$\chi$$^{2}$~=~+0.63, no lensing $\Delta$$\chi$$^{2}$~=~+0.41, and TT only $\Delta$$\chi$$^{2}$~=~$-$0.65 (definitions and values given in the Test~2 artefacts).

\textbf{Test 4: Prior range sensitivity.}
Widening the k\textsubscript{f} prior from [0,\,5] to [0,\,10] leaves the posterior essentially unchanged: k\textsubscript{f}~=~0.483~$\pm$~0.361 (wide) versus 0.495~$\pm$~0.360 (baseline), with nearly identical 95\% credible intervals ([0.023,\,1.341] versus [0.025,\,1.342]) and a negligible mean shift ($-$0.03$\sigma$). This indicates that the Phase~5 k\textsubscript{f} inference is data-informed and not driven by the prior bound.

\textbf{Summary.}
Taken together with the identifiability diagnostics in Test~3, these results show that the Phase~5 posterior behaviour is stable under removal of lensing, behaves as expected under the information loss of TT only, and is insensitive to widening the k\textsubscript{f} prior. The robustness suite therefore supports interpreting the observed $\sigma$\textsubscript{8} suppression and modest H\textsubscript{0} shift as stable consequences of the MTDF coupling within the Planck likelihood framework, rather than numerical artefacts or prior-driven effects.

\begin{mdframed}[linecolor=accent,linewidth=0.5pt,backgroundcolor=callout-bg,leftline=true,rightline=true,topline=true,bottomline=true,innerleftmargin=8pt,innerrightmargin=8pt,innertopmargin=6pt,innerbottommargin=6pt]
\textbf{Phase 5 MCMC converged.}
Both chains reached R$-$1~<~0.02 (threshold: 0.02). The $\sigma$\textsubscript{8} shift ($-$2.4\% relative to the $\Lambda$CDM posterior mean, \~{}2.4$\sigma$ in Gaussian approximation) and directional H\textsubscript{0} shift (+0.45~km~s\textsuperscript{$-$1}~Mpc\textsuperscript{$-$1}, \~{}0.6$\sigma$) first observed in Phase~2 (plik-lite) are confirmed at full Planck precision with identical nuisance freedom. k\textsubscript{f}~=~1 (full MTDF) lies within the 95\% credible interval; k\textsubscript{f}~=~0 ($\Lambda$CDM) is not strongly excluded and sits near the lower boundary. Triangle plots and k\textsubscript{f} posterior are archived in \texttt{output/\allowbreak{}phase5/\allowbreak{}}.
\end{mdframed}

\section{6b. Phase 6: Cross-probe discriminator tests}
\label{sec:sec6b}
COMPLETE: Two discriminators (A, E), one coherence check (C2), three upper limits (B, B2, D), one internal consistency (C), two forecasts (C3, C4), one resolved null (C5)  -  10 tests totalPhase~6 extends validation beyond CMB compatibility toward tests that can \emph{distinguish} MTDF from $\Lambda$CDM using independent probes. Ten tests are packaged under \texttt{output/\allowbreak{}phase6/\allowbreak{}}, targeting weak lensing, CMB lensing, derived parameter consistency, cross-probe S\textsubscript{8} coherence, growth--environment coupling, BAO residuals, cluster mass ratios, and the SN~$\times$~void redshift transition. All code was written from scratch for Phase~6.

Tests~A and E positively support MTDF-style predictions (redshift transition, void-specific CMB lensing signal). Test~C2 demonstrates cross-probe S\textsubscript{8} coherence (combined weak lensing tension reduced from 3.5$\sigma$ to 1.9$\sigma$). Tests~B, B2, and D are systematics-clean nulls with explicit upper limits. Test~C confirms internal consistency. Tests~C3 and C4 quantify sensitivity boundaries. Test~C5 resolves the cluster mass ratio ($\eta$~=~1). Together they demonstrate: \emph{we did not cheat, the pipelines are sound, no hidden contradictions exist, and the boundaries of current discriminating power are mapped.}

\subsection{Summary table}
\begin{table}[H]
\centering
\footnotesize
\begin{tabularx}{\linewidth}{l>{\RaggedRight\arraybackslash}X>{\RaggedRight\arraybackslash}X>{\RaggedRight\arraybackslash}X>{\RaggedRight\arraybackslash}X}
\toprule
\textbf{Test} & \textbf{Observable} & \textbf{Primary statistic} & \textbf{Result} & \textbf{Implication} \tabularnewline
\midrule
\textbf{A} & SN~$\times$~void $\gamma$\textsubscript{env} vs z\textsubscript{cut} & GLS z-score peak over z\textsubscript{cut} scan & Peak at z\textsubscript{cut}=0.030: 3.62$\sigma$ (global p~<~0.005) & Supports MTDF transition at z~<~0.04 \tabularnewline
\textbf{B} & KiDS-1000 $\gamma$\textsubscript{t} around void centres & $\Delta$$\gamma$\textsubscript{t} over R=5--20 Mpc/\allowbreak{}h & $-$3.65e$-$5~$\pm$~4.68e$-$5 (0.8$\sigma$); 95\%: |$\Delta$$\gamma$\textsubscript{t}|~<~1.2e$-$4 & Null  -  bounded, not yet discriminating \tabularnewline
\textbf{B2} & KiDS-1000 trough vs ridge $\gamma$\textsubscript{t} & $\Delta$$\gamma$\textsubscript{t,split} (ridge$-$trough) & +2.56e$-$5~$\pm$~5.46e$-$5 (0.5$\sigma$); 95\%: |$\Delta$$\gamma$\textsubscript{t,split}|~<~1.3e$-$4 & Null  -  bounded, not yet discriminating \tabularnewline
\textbf{C} & Derived parameter consistency & Max shift vs $\Lambda$CDM; S\textsubscript{8} tension reduction & All within 2.4$\sigma$; S\textsubscript{8}: DES 2.6$\rightarrow$1.3$\sigma$, KiDS 2.8$\rightarrow$1.6$\sigma$ & Internally consistent; eases S\textsubscript{8} tension \tabularnewline
\textbf{D} & Planck PR4 $\kappa$ $\times$ DESI BGS voids & Per-void compensated $\kappa$\textsubscript{comp} & $-$7.47e$-$4~$\pm$~1.03e$-$3 (0.7$\sigma$, 557 voids) & Null at low z  -  expected; pipeline validated \tabularnewline
\textbf{E} & Planck PR4 $\kappa$ $\times$ BOSS DR12 voids & Conservative $\kappa$\textsubscript{comp} + detection-optimised CTH/\allowbreak{}MF & Comp: 0.2$\sigma$ null; CTH: 1.5$\sigma$ (p\textsubscript{rand}=0.03); MF: 1.9$\sigma$ (p\textsubscript{rand}=0.02) & Pipeline benchmark validated; void-specific signal recovered \tabularnewline
\textbf{C2} & S\textsubscript{8} cross-probe coherence & Phase~5 S\textsubscript{8} vs published WL surveys & Combined WL tension: 3.47$\sigma$~$\rightarrow$~1.89$\sigma$ (46\% reduction) & MTDF shifts S\textsubscript{8} in the right direction by the right amount \tabularnewline
\textbf{C3} & BAO residual structure & $\Delta$W RMS (wiggle difference MTDF$-$$\Lambda$CDM) & $\Delta$W RMS~=~0.27\%, max~=~0.51\%; sound horizon shift +0.11\% & Not detectable by any foreseeable survey (best SNR~=~0.92) \tabularnewline
\textbf{C4} & f$\sigma$\textsubscript{8} $\times$ environment & RSD + void f$\sigma$\textsubscript{8} vs $\mu$(a) prediction & $\Delta$$\chi$$^{2}$~=~+0.49 (9 DOF); max discrimination 0.22$\sigma$ & Sensitivity-limited; needs \~{}10$\times$ improvement (DESI/\allowbreak{}Euclid) \tabularnewline
\textbf{C5} & Cluster M\textsubscript{dyn}/M\textsubscript{lens} & Analytical $\mu$(a) prediction + numerical $\eta$(k,z) verification & $\eta$~=~1.000000 at all cluster scales; Case~B confirmed: M\textsubscript{dyn}/M\textsubscript{lens}~=~1 & Resolved: cluster mass ratio is not a discriminator for MTDF \tabularnewline
\bottomrule
\end{tabularx}
\end{table}

\subsection{6b.1 Test A: Redshift transition scan}
\begin{itemize}
  \item Scans z\textsubscript{cut} $\in$ [0.025,~0.100] using GLS and Spearman metrics on the Phase~3 SN~$\times$~void environment signal.
  \item Peak at z\textsubscript{cut}~=~0.030 with GLS z-score 3.62$\sigma$; global scan-corrected permutation p~<~0.005 (0 of 200 exceedances).
  \item Spearman rank test confirms independently (3.1$\sigma$).
  \item Interpretation: the environment signal is localised to z~<~0.04, consistent with the Phase~3 redshift modulation and the MTDF coherence-length prediction.
  \item Artefacts: output/\allowbreak{}phase6/\allowbreak{}testA\_\allowbreak{}redshift\_\allowbreak{}transition/\allowbreak{}
\end{itemize}

\subsection{6b.2 Test B: Weak lensing $\times$ environment (void centres)}
\begin{itemize}
  \item KiDS-1000 shear stacked around DESI REVOLVER void centres and matched random positions; $\gamma$\textsubscript{x} (cross-shear) gate required to pass before reporting $\gamma$\textsubscript{t}.
  \item $\Delta$$\gamma$\textsubscript{t}~=~$-$3.65$\times$10\textsuperscript{$-$5}~$\pm$~4.68$\times$10\textsuperscript{$-$5} over R~=~5--20~Mpc/\allowbreak{}h, consistent with zero.
  \item $\gamma$\textsubscript{x} gate: pass (p~=~0.234).
  \item Explicit two-sided 95\% upper limit: |$\Delta$$\gamma$\textsubscript{t}|~<~1.20$\times$10\textsuperscript{$-$4}.
  \item Without a forward MTDF lensing prediction for this selection, the result is best framed as ``not yet discriminating, but now bounded and reproducible.''
  \item Artefacts: output/\allowbreak{}phase6/\allowbreak{}testB\_\allowbreak{}wl\_\allowbreak{}environment/\allowbreak{}
\end{itemize}

\subsection{6b.3 Test B2: Trough and ridge lensing (projected density split)}
\begin{itemize}
  \item Projected density field built with mask-aware normalisation and full KiDS-1000 footprint (North + South). Both trough and ridge $\gamma$\textsubscript{x} gates pass.
  \item $\Delta$$\gamma$\textsubscript{t,split} (ridge$-$trough)~=~+2.56$\times$10\textsuperscript{$-$5}~$\pm$~5.46$\times$10\textsuperscript{$-$5} (0.5$\sigma$), consistent with zero.
  \item Explicit two-sided 95\% upper limit: |$\Delta$$\gamma$\textsubscript{t,split}|~<~1.28$\times$10\textsuperscript{$-$4}.
  \item Projected P20/\allowbreak{}P80 density contrasts ($\delta$~$\sim$~$\pm$0.15) are too mild for detectable differential lensing at KiDS sensitivity. Spectroscopic 3D density or LSST/\allowbreak{}DES area would be needed for significance.
  \item Artefacts: output/\allowbreak{}phase6/\allowbreak{}testB2\_\allowbreak{}trough\_\allowbreak{}ridge/\allowbreak{}
\end{itemize}

\subsection{6b.4 Test C: Derived parameter consistency}
\begin{itemize}
  \item Cross-checks derived parameters from the Phase~5 MCMC posterior against external weak lensing constraints.
  \item No derived parameter deviates beyond 2.4$\sigma$ from $\Lambda$CDM.
  \item S\textsubscript{8} tension reduction: DES~Y3 from 2.63$\sigma$ to 1.28$\sigma$; KiDS-1000 from 2.77$\sigma$ to 1.58$\sigma$.
  \item MTDF's Planck posterior is internally consistent; no hidden contradictions introduced by \texttt{class\_\allowbreak{}mtdf}.
  \item Artefacts: output/\allowbreak{}phase6/\allowbreak{}testC\_\allowbreak{}derived\_\allowbreak{}consistency/\allowbreak{}
\end{itemize}

\subsection{6b.5 Test D: CMB lensing $\times$ DESI BGS voids (low redshift)}
\begin{itemize}
  \item Planck PR4 MV convergence ($\kappa$) stacked on DESIVAST BGS voids (z~<~0.24) with mask-weighted stacking. Primary statistic is a per-void compensated filter: $\kappa$(R/\allowbreak{}R\textsubscript{v}<0.5)~$-$~$\kappa$(4<R/\allowbreak{}R\textsubscript{v}<5), designed to remove large-scale mode contamination.
  \item $\kappa$\textsubscript{comp}~=~$-$7.47$\times$10\textsuperscript{$-$4}~$\pm$~1.03$\times$10\textsuperscript{$-$3} (0.7$\sigma$, 557 voids), consistent with zero.
  \item RA-scramble null: p~=~0.425 (compensated). Random-sky null: p~=~0.080.
  \item Robustness: stable across outer annulus choices (all <1.1$\sigma$), low-l removal (l$\ge$20, l$\ge$30), and void finder.
  \item Null at z~<~0.24 is expected: the CMB lensing kernel peaks at z~\~{}~1--2, and the BGS void count (557 usable) limits statistical power.
  \item Diagnosed failure modes from v1: hard mask quality cut rejected 73\% of voids; low-l CMB modes created a positive $\kappa$ pedestal. Both fixed in v2 via per-void compensated filter and mask-weighted stacking.
  \item Artefacts: output/\allowbreak{}phase6/\allowbreak{}testD\_\allowbreak{}cmb\_\allowbreak{}lensing/\allowbreak{}
\end{itemize}

\subsection{6b.6 Test E: CMB lensing $\times$ BOSS DR12 voids (benchmark)}
\begin{itemize}
  \item Benchmark run on the Mao+2017 ZOBOV catalogue (1,228 BOSS DR12 CMASS+LOWZ voids, z~=~0.2--0.7) to validate the CMB lensing stacking pipeline against Cai+2017's published 3.2$\sigma$ detection on the same catalogue.
  \item Five statistics computed side by side:
    
\textbf{Per-void compensated} (systematics-clean primary): 0.2$\sigma$, null by design.
\textbf{AP disc} (contamination monitor): 4.7$\sigma$, but RA-scramble p~=~0.945 and collapses to 0.3$\sigma$ under l$\ge$20 cut  -  entirely pedestal.
\textbf{CTH} (Cai-class compensated top-hat): 1.5$\sigma$ (p\textsubscript{rand}~=~0.03), increases to 2.0$\sigma$ under l$\ge$20 cut.
\textbf{Matched filter} (template + jackknife covariance): 1.9$\sigma$ (p\textsubscript{rand}~=~0.02), increases to 2.2$\sigma$ under l$\ge$20 cut.
\textbf{Close comp} (alternative aperture): 0.5$\sigma$, consistent with zero.
  \item Detection-optimised statistics (CTH and MF) show the expected hierarchy and respond correctly to pedestal removal and footprint-preserving null tests. The gap from \~{}1.5--2.2$\sigma$ to Cai+2017's 3.2$\sigma$ reflects template shape (step vs theory-derived) and different Planck products (PR4 vs PR2015).
  \item Artefacts: output/\allowbreak{}phase6/\allowbreak{}testE\_\allowbreak{}boss\_\allowbreak{}benchmark/\allowbreak{}
\end{itemize}

\subsection{6b.7 Test C2: S\textsubscript{8} cross-probe coherence}
\begin{itemize}
  \item Compares Phase~5 MCMC S\textsubscript{8}~=~$\sigma$\textsubscript{8}$\sqrt{\,}$($\Omega$\textsubscript{m}/0.3) against published weak lensing constraints from KiDS-1000, DES~Y3, and HSC~Y3.
  \item $\Lambda$CDM: S\textsubscript{8}~=~0.830~$\pm$~0.013. MTDF: S\textsubscript{8}~=~0.802~$\pm$~0.012.
  \item Tension reduction: KiDS-1000 2.67$\sigma$~$\rightarrow$~1.63$\sigma$ (39\%); DES~Y3 2.63$\sigma$~$\rightarrow$~1.28$\sigma$ (51\%); combined WL 3.47$\sigma$~$\rightarrow$~1.89$\sigma$ (46\%).
  \item Caveat: published WL S\textsubscript{8} values assume $\Lambda$CDM growth and lensing kernel. A proper comparison requires WL likelihood within the MTDF framework (deferred to future work).
  \item Artefacts: output/\allowbreak{}phase6/\allowbreak{}testC2\_\allowbreak{}s8\_\allowbreak{}coherence/\allowbreak{}
\end{itemize}

\subsection{6b.8 Test C3: BAO residual structure}
\begin{itemize}
  \item Extracts BAO wiggles W(k)~=~P(k)/P\textsubscript{nw}(k)~$-$~1 from \texttt{class\_\allowbreak{}mtdf} power spectra using an Eisenstein~\&~Hu (1998) no-wiggle reference, and computes the residual $\Delta$W~=~W\textsubscript{MTDF}~$-$~W\textsubscript{$\Lambda$CDM}.
  \item Sound horizon shift: +0.106\% (147.148~$\rightarrow$~147.303~Mpc). Wiggle difference: $\Delta$W RMS~=~0.27\%, max~=~0.51\% at k~=~0.12~h/\allowbreak{}Mpc. No detectable peak phase shift at grid resolution.
  \item Best projected SNR: 0.92 (DESI + Euclid combined). The 0.1\% sound horizon shift is absorbed into standard cosmological parameter degeneracies.
  \item Conclusion: BAO residual structure is not a viable MTDF--$\Lambda$CDM discriminator with any foreseeable survey.
  \item Artefacts: output/\allowbreak{}phase6/\allowbreak{}testC3\_\allowbreak{}bao\_\allowbreak{}residuals/\allowbreak{}
\end{itemize}

\subsection{6b.9 Test C4: f$\sigma$\textsubscript{8} $\times$ environment}
\begin{itemize}
  \item Compares MTDF and $\Lambda$CDM f$\sigma$\textsubscript{8}(z) predictions against a compilation of 8 RSD surveys plus void-specific growth measurements from Hamaus+2020.
  \item $\chi$$^{2}$: $\Lambda$CDM~=~7.61, MTDF~=~8.10 ($\Delta$$\chi$$^{2}$~=~+0.49, 9~DOF). Maximum per-point discrimination: 0.22$\sigma$.
  \item The homogeneous $\mu$(a) modification produces a \~{}1.5\% signal, which is 5--10$\times$ smaller than current RSD measurement errors (0.03--0.10 per z-bin).
  \item Future requirement: 2$\sigma$ discrimination needs $\sigma$(f$\sigma$\textsubscript{8})~\~{}~0.003--0.004 per z-bin, achievable with DESI/\allowbreak{}Euclid-era RSD or environment-resolved estimators.
  \item Artefacts: output/\allowbreak{}phase6/\allowbreak{}testC4\_\allowbreak{}fsigma8\_\allowbreak{}environment/\allowbreak{}
\end{itemize}

\subsection{6b.10 Test C5/\allowbreak{}C5b: Cluster M\textsubscript{dyn}/M\textsubscript{lens}  -  Resolved}
\begin{itemize}
  \item C5 computed the predicted ratio M\textsubscript{dyn}/M\textsubscript{lens} under two bracketing coupling cases:
    
\textbf{Case~A} ($\mu$-only dynamics, $\Sigma$~=~1): M\textsubscript{dyn}/M\textsubscript{lens}~=~1.053 at z~=~0 (+5.3\% offset).
\textbf{Case~B} (equal coupling, $\Sigma$~=~$\mu$): M\textsubscript{dyn}/M\textsubscript{lens}~=~1.000 (no offset).
  \item \textbf{C5b resolved the key unknown} by extracting $\eta$(k,z)~=~$\Phi$(k,z)/$\Psi$(k,z) directly from \texttt{class\_\allowbreak{}mtdf} Newtonian-gauge transfer functions across k~=~0.005--2~h/\allowbreak{}Mpc and z~=~0--10.
  \item Result: \textbf{$\eta$~=~1.000000} at all cluster-relevant scales (z~$\le$~0.8). Max~|$\Delta$$\eta$| between MTDF and $\Lambda$CDM:~0.000000. Tiny deviations (|$\eta$$-$1|~\~{}~10\textsuperscript{$-$5}) appear only at z~$\ge$~5 from standard radiation anisotropic stress, present equally in both models.
  \item Physical reason: \texttt{class\_\allowbreak{}mtdf} modifies only the Poisson equation source ($\delta$$\rho$~$\rightarrow$~$\mu$$\cdot$$\delta$$\rho$) and the CDM Euler equation. The $\Phi$$-$$\Psi$ anisotropy equation is unchanged from GR. Both potentials are enhanced equally $\rightarrow$ $\Sigma$~=~$\mu$ $\rightarrow$ \textbf{Case~B confirmed}.
  \item Implication: M\textsubscript{dyn}/M\textsubscript{lens}~=~1.000 under MTDF. The cluster mass ratio channel is \textbf{not a discriminator}. Clusters can still test MTDF through growth-dependent observables (e.g.~cluster counts), because $\mu$~>~1 modifies the growth rate.
  \item Artefacts: output/\allowbreak{}phase6/\allowbreak{}testC5\_\allowbreak{}cluster\_\allowbreak{}mass/\allowbreak{}, output/\allowbreak{}phase6/\allowbreak{}testC5b\_\allowbreak{}gravitational\_\allowbreak{}slip/\allowbreak{}
\end{itemize}

\begin{mdframed}[linecolor=accent,linewidth=0.5pt,backgroundcolor=callout-bg,leftline=true,rightline=true,topline=true,bottomline=true,innerleftmargin=8pt,innerrightmargin=8pt,innertopmargin=6pt,innerbottommargin=6pt]
\textbf{Phase 6 conclusion (10 tests, all complete).}
The ten Phase~6 tests fall into five categories. \emph{Discriminators} (6A, 6E): the redshift transition is detected at 3.6$\sigma$; the BOSS CMB lensing benchmark recovers void-specific structure at 1.5--1.9$\sigma$. \emph{Coherence} (6C, 6C2): derived parameters are self-consistent (no shift > 2.4$\sigma$); S\textsubscript{8} tension with weak lensing surveys drops from 3.5$\sigma$ to 1.9$\sigma$. \emph{Upper limits} (6B, 6B2, 6D): systematics-clean nulls with explicit 95\% bounds on void lensing and CMB lensing signals. \emph{Forecasts} (6C3, 6C4): BAO wiggles and f$\sigma$\textsubscript{8} quantify what future surveys need to reach. \emph{Resolved} (6C5): cluster M\textsubscript{dyn}/M\textsubscript{lens} is null ($\eta$~=~1). No test contradicts MTDF. The S\textsubscript{8} coherence result (6C2) is the strongest positive outcome: MTDF moves S\textsubscript{8} toward weak lensing values without tuning.
\end{mdframed}

\section{Synthesis and framework narrative}
\label{sec:sec7}

\subsection{Summary of results (17/\allowbreak{}17 complete)}
\begingroup\footnotesize
\begin{longtable}{>{\RaggedRight\arraybackslash}p{0.192\linewidth}>{\RaggedRight\arraybackslash}p{0.192\linewidth}>{\RaggedRight\arraybackslash}p{0.192\linewidth}>{\RaggedRight\arraybackslash}p{0.192\linewidth}>{\RaggedRight\arraybackslash}p{0.192\linewidth}}
\toprule
\endfirsthead
\toprule\endhead
\bottomrule\endfoot
\textbf{Phase} & \textbf{Type} & \textbf{Status} & \textbf{Key result} & \textbf{Artefacts} \tabularnewline
Phase 1 & Discriminator & PASSED & V18 strict totals reproduced to $\delta$~=~0.000049 & output/\allowbreak{}phase1/\allowbreak{} \tabularnewline
Phase 2 & Discriminator & PASSED & CMB plik-lite compatible ($\Delta$$\chi$$^{2}$~<~1); k\textsubscript{f}~=~1 within 95\% in all MCMC runs & output/\allowbreak{}phase2/\allowbreak{} \tabularnewline
Phase 3 & Discriminator & PASSED & SN~$\times$~void signal reproduced exactly; bootstrap excludes zero & output/\allowbreak{}phase3/\allowbreak{} \tabularnewline
Phase 3b & Upper limit & COMPLETE & NGC/\allowbreak{}SGC asymmetry: not significant, not geometric, sensitivity-limited (5 diagnostics) & output/\allowbreak{}phase3b/\allowbreak{} \tabularnewline
Phase 4 & Forecast & COMPLETE & 5$\sigma$ at \~{}3,400 SNe (LSST/\allowbreak{}Rubin); statistics-limited & output/\allowbreak{}phase4/\allowbreak{} \tabularnewline
Phase 5 (min.) & Discriminator & PASSED & Full Planck plik + class\_\allowbreak{}mtdf: $\Delta$$\chi$$^{2}$~=~+0.63. Priority~1 hard falsifier cleared. & \multirow{2}{*}{output/\allowbreak{}phase5/\allowbreak{}} \tabularnewline
Phase 5 (MCMC) & Discriminator & PASSED & Converged (R$-$1~=~0.019). $\sigma$\textsubscript{8}~=~0.790 (\~{}2.4$\sigma$ below $\Lambda$CDM). k\textsubscript{f}~=~0.50~$\pm$~0.36; both k\textsubscript{f}~=~0 and 1 within 95\% CI. &  \tabularnewline
Phase 6A & Discriminator & PASSED & Redshift transition scan: signal peaks at z~<~0.03 (3.6$\sigma$ GLS, global p~<~0.005). Spearman confirms (3.1$\sigma$). & \multirow{10}{*}{output/\allowbreak{}phase6/\allowbreak{}} \tabularnewline
Phase 6B & Upper limit & COMPLETE & Weak lensing $\times$ void environment (KiDS-1000): $\Delta$$\gamma$\textsubscript{t}~=~$-$3.65e$-$5~$\pm$~4.68e$-$5 (0.8$\sigma$). $\gamma$\textsubscript{x} gate pass. 95\% UL: |$\Delta$$\gamma$\textsubscript{t}|~<~1.2e$-$4. &  \tabularnewline
Phase 6B2 & Upper limit & COMPLETE & Trough/\allowbreak{}ridge lensing (KiDS-1000): $\Delta$$\gamma$\textsubscript{t,split}~=~+2.56e$-$5~$\pm$~5.46e$-$5 (0.5$\sigma$). Both $\gamma$\textsubscript{x} gates pass. 95\% UL: |$\Delta$$\gamma$\textsubscript{t,split}|~<~1.3e$-$4. &  \tabularnewline
Phase 6C & Coherence & PASSED & Derived parameter consistency: no shifts >~2.4$\sigma$. S\textsubscript{8} tension reduces: DES~Y3 2.6$\rightarrow$1.3$\sigma$; KiDS-1000 2.8$\rightarrow$1.6$\sigma$. &  \tabularnewline
Phase 6C2 & Coherence & COMPLETE & S\textsubscript{8} cross-probe coherence: Phase~5 posterior S\textsubscript{8} vs WL surveys. Combined tension 3.5$\sigma$~$\rightarrow$~1.9$\sigma$ (46\% reduction). Summary-level; full likelihood deferred. &  \tabularnewline
Phase 6C3 & Forecast & COMPLETE & BAO residual structure: $\Delta$W RMS~=~0.27\%, sound horizon shift +0.11\%. Not detectable by any current or projected survey (best SNR~=~0.92). &  \tabularnewline
Phase 6C4 & Forecast & COMPLETE & f$\sigma$\textsubscript{8}~$\times$~environment: RSD + void growth vs $\mu$(a) prediction. Max 0.22$\sigma$; needs DESI/\allowbreak{}Euclid-era precision (\~{}10$\times$ improvement). &  \tabularnewline
Phase 6C5/\allowbreak{}C5b & Closed null & RESOLVED & Cluster M\textsubscript{dyn}/M\textsubscript{lens}: $\eta$~=~1 confirmed numerically (C5b) $\Rightarrow$ $\Sigma$~=~$\mu$, so M\textsubscript{dyn}/M\textsubscript{lens}~=~1 exactly. Mass ratio offset is zero in the current MTDF perturbation model. &  \tabularnewline
Phase 6D & Upper limit & COMPLETE & CMB lensing $\times$ DESI BGS voids: $\kappa$\textsubscript{comp}~=~$-$7.47e$-$4~$\pm$~1.03e$-$3 (0.7$\sigma$, 557 voids). RA-scramble clean. Null expected at z~<~0.24. &  \tabularnewline
Phase 6E & Discriminator & COMPLETE & CMB lensing $\times$ BOSS DR12 benchmark: CTH 1.5$\sigma$ (p\textsubscript{rand}=0.03), MF 1.9$\sigma$ (p\textsubscript{rand}=0.02). Pipeline validated against Cai+2017. &  \tabularnewline
\end{longtable}
\endgroup

\subsection{Claim status taxonomy}
Each MTDF claim or test outcome is assigned one of four status labels. This taxonomy is intended to provide referees and readers with an unambiguous assessment of where the framework stands.

\begin{table}[H]
\centering
\footnotesize
\begin{tabularx}{\linewidth}{l>{\RaggedRight\arraybackslash}X}
\toprule
\textbf{Label} & \textbf{Definition} \tabularnewline
\checkmark{} Validated & Independently reproduced with quantitative agreement. Robust to systematics tests. \tabularnewline
$\blacksquare$ Constrained & Consistent with data but not uniquely favoured. Current data cannot distinguish from $\Lambda$CDM. \tabularnewline
$\blacktriangle$ Hypothesis & Theoretically motivated but not yet tested against data, or tested only indirectly. \tabularnewline
$\times$ Falsifier & Identified test that could refute the framework. Not yet executed. \tabularnewline
\bottomrule
\end{tabularx}
\end{table}

\subsubsection{Status by claim}
\begingroup\footnotesize
\begin{longtable}{>{\RaggedRight\arraybackslash}p{0.320\linewidth}>{\RaggedRight\arraybackslash}p{0.320\linewidth}>{\RaggedRight\arraybackslash}p{0.320\linewidth}}
\toprule
\endfirsthead
\toprule\endhead
\bottomrule\endfoot
\textbf{Claim / element} & \textbf{Status} & \textbf{Evidence} \tabularnewline
V18 strict $\chi$$^{2}$ totals & \checkmark{} Validated & Phase 1: independent reproduction to $\delta$ = 0.000049 \tabularnewline
Pantheon+ SN likelihood & \checkmark{} Validated & Phase 1: $\chi$$^{2}$ = 2012.72 matched exactly \tabularnewline
DESI Y1 BAO fit & \checkmark{} Validated & Phase 1: $\chi$$^{2}$ = 15.39 matched exactly \tabularnewline
Cosmic chronometers H(z) & \checkmark{} Validated & Phase 1: $\chi$$^{2}$ = 6.98 matched exactly \tabularnewline
DR16 f$\sigma$\textsubscript{8} growth & \checkmark{} Validated & Phase 1: $\chi$$^{2}$ = 2.00 matched exactly \tabularnewline
CMB compatibility (plik-lite) & $\blacksquare$ Constrained & Phase 2: $\Delta$$\chi$$^{2}$ < 1; k\textsubscript{f} = 1 within 95\% CI. Perturbative correction layer, not full Boltzmann. \tabularnewline
CMB compatibility (full plik + class\_\allowbreak{}mtdf) & \checkmark{} Validated & Phase 5: full plik TTTEEE + lowl + lensing via class\_\allowbreak{}mtdf Boltzmann solver. $\Delta$$\chi$$^{2}$ = +0.63 (BOBYQA best-fit, 4 runs). k\textsubscript{f} = 0.028 $\approx$ 0. Apples-to-apples with $\Lambda$CDM (same likelihoods, nuisances, priors, minimizer). \tabularnewline
SN $\times$ void environment signal ($\gamma$\textsubscript{env}) & \checkmark{} Validated & Phase 3: exact reproduction, 3 void finders, bootstrap excludes zero \tabularnewline
Redshift decay structure (z < 0.04) & \checkmark{} Validated & Phase 3: a priori prediction confirmed (p = 0.0035 REVOLVER) \tabularnewline
Survey systematics stability & \checkmark{} Validated & Phase 3: FE $\Delta$$\gamma$ < 0.0005; LOSO 16/\allowbreak{}16 positive \tabularnewline
NGC/\allowbreak{}SGC asymmetry interpretation & \checkmark{} Validated & Phase 3b: asymmetry not significant (Wald p = 0.49--0.96), not geometric (mock p = 0.28--0.90), SGC sensitivity-limited (6--8$\times$ fewer voids, 44\% recovery power). Five independent diagnostics confirm footprint-dependent detectability. \tabularnewline
5$\sigma$ detection forecast (3,400 SNe) & \checkmark{} Validated & Phase 4: Monte Carlo with consistency checks passing at 10\textsuperscript{$-$13} \tabularnewline
Statistics-limited (not systematics-limited) & \checkmark{} Validated & Phase 4: baseline $\approx$ realistic; rank-1 systematic immunity demonstrated \tabularnewline
Hubble tension resolution (void depletion) & $\blacksquare$ Constrained & Two independent equations yield H\textsubscript{local} consistent with SH0ES. Mechanism independent of CMB; not yet tested with dedicated void-SN joint analysis. \tabularnewline
Photon--stress coupling ($\kappa$ = f\textsubscript{kick}/3 $\approx$ 0.00102) & $\blacktriangle$ Hypothesis & Derived from fundamental parameters (angular averaging of strain tensor). Below Planck sensitivity. Paper MTDF\_\allowbreak{}04 frames as speculative. \tabularnewline
Rotation curve phenomenology & $\blacksquare$ Constrained & Included in V18 scalars. Not independently reproduced in this session. \tabularnewline
Redshift transition structure (z\textsubscript{cut} scan) & \checkmark{} Validated & Phase 6A: GLS z-score 3.62$\sigma$ at z\textsubscript{cut}=0.030; global scan-corrected p~<~0.005. Spearman confirms (3.1$\sigma$). \tabularnewline
Weak lensing $\times$ environment (KiDS) & $\blacksquare$ Constrained & Phase 6B/\allowbreak{}B2: systematics-clean nulls with explicit 95\% upper limits. Bounded but not yet discriminating at KiDS sensitivity. \tabularnewline
Derived parameter consistency & \checkmark{} Validated & Phase 6C: all within 2.4$\sigma$. S\textsubscript{8} tension eased (DES 2.6$\rightarrow$1.3$\sigma$; KiDS 2.8$\rightarrow$1.6$\sigma$). \tabularnewline
CMB lensing $\times$ voids & $\blacksquare$ Constrained & Phase 6D/\allowbreak{}E: low-z null (expected), BOSS benchmark recovers void-specific structure. Pipeline validated. \tabularnewline
S\textsubscript{8} cross-probe coherence & \checkmark{} Validated & Phase 6C2: S\textsubscript{8} from Phase~5 posterior vs WL surveys. Combined tension 3.5$\sigma$~$\rightarrow$~1.9$\sigma$. Summary-level comparison; full likelihood deferred. \tabularnewline
f$\sigma$\textsubscript{8} $\times$ environment & $\blacksquare$ Constrained & Phase 6C4: $\mu$(a) prediction vs RSD + void growth. Max 0.22$\sigma$ discrimination; current data sensitivity-limited (\~{}10$\times$ improvement needed). \tabularnewline
BAO residual structure & $\blacksquare$ Constrained & Phase 6C3: $\Delta$W RMS = 0.27\%. Sound horizon shift +0.11\% absorbed into parameter degeneracies. Not detectable by any foreseeable survey. \tabularnewline
Cluster M\textsubscript{dyn}/M\textsubscript{lens} & $\blacksquare$ Constrained & Phase 6C5/\allowbreak{}C5b: $\eta$~=~1.000000 confirmed numerically from class\_\allowbreak{}mtdf transfer functions. Case~B applies: M\textsubscript{dyn}/M\textsubscript{lens}~=~1. Mass ratio channel closed as discriminator; growth-dependent cluster tests remain open. \tabularnewline
Weak lensing P(k) cross-correlation & $\times$ Falsifier & Not yet executed. Requires full P(k) from class\_\allowbreak{}mtdf at non-linear scales. \tabularnewline
Merging cluster offsets (Bullet Cluster) & $\times$ Falsifier & Not yet executed. Requires N-body simulations. \tabularnewline
\end{longtable}
\endgroup

\subsection{Framework narrative}
\begin{mdframed}[linecolor=accent,linewidth=0.5pt,backgroundcolor=callout-bg,leftline=true,rightline=true,topline=true,bottomline=true,innerleftmargin=8pt,innerrightmargin=8pt,innertopmargin=6pt,innerbottommargin=6pt]
MTDF is fully CMB-compatible: Planck TTTEEE cannot distinguish MTDF from $\Lambda$CDM ($\Delta$$\chi$$^{2}$~<~1). MTDF deviates from $\Lambda$CDM at late times through environment-dependent dynamics. The Hubble tension is resolved by environmental mapping (void stress depletion), not by early-universe modification. The SN~$\times$~void signal is the primary empirical discriminator. The stress-field coherence length concentrates observable effects at low redshift, consistent with the CMB being unaffected.
\end{mdframed}

\subsection{What survives intact}
\begin{enumerate}
  \item \textbf{Phase 1 strict totals}: independently reproduced to four decimal places.
  \item \textbf{SN $\times$ void environment signal}: three void finders, 16/\allowbreak{}16 LOSO, redshift decay matching a priori prediction.
  \item \textbf{Hubble tension resolution}: entirely late-time, via void stress depletion:
    
Theoretical: \( H_{\rm local}/H_0 = 1.084 \;\Rightarrow\; 67.4 \times 1.084 = 73.1 \) (SH0ES: 73.04~$\pm$~1.04)
MCMC-derived: \( H_{0,\rm global} = 68.56 \) + void enhancement 5.81 = \( H_{0,\rm local} = 74.37 \)
  \item \textbf{Late-time phenomenology}: w(z), BAO distances, growth rate, rotation curves.
  \item \textbf{Framework stability}: MTDF recovers $\Lambda$CDM at early times, deviates at late times.
  \item \textbf{Cross-probe integrity (Phase 6, 10 tests)}: redshift transition confirmed (3.6$\sigma$); weak lensing and CMB lensing pipelines pass strict systematics controls; derived parameters self-consistent; S\textsubscript{8} tension reduced from 3.5$\sigma$ to 1.9$\sigma$ (combined WL); f$\sigma$\textsubscript{8}, BAO residuals, and cluster mass predictions computed  -  all sensitivity-limited with current data, with quantified thresholds for future surveys.
\end{enumerate}

\subsection{Joint implications for late-time cosmological tensions}
While MTDF is not tuned to resolve any single cosmological tension, the validated framework exhibits simultaneous shifts in multiple late-time observables. In particular, the Phase~5 best-fit favours a modest reduction in $\sigma$\textsubscript{8} relative to $\Lambda$CDM best-fit values (0.795 vs 0.809), consistent with the direction preferred by weak-lensing surveys (KiDS, DES).

Independently, the environment-dependent dynamics tested in Phases~3 and 4 naturally modify local expansion observables without compromising global CMB consistency. Together, these results indicate that MTDF addresses multiple late-time phenomenological tensions within a single, unified dynamical framework, without introducing additional free parameters beyond k\textsubscript{f}.

\subsection{Paper implications}
\begin{table}[H]
\centering
\footnotesize
\begin{tabularx}{\linewidth}{l>{\RaggedRight\arraybackslash}X}
\toprule
\textbf{Paper} & \textbf{Impact} \tabularnewline
MTDF\_\allowbreak{}01 (Master) & Phase 1 strengthens all claims. No changes needed. \tabularnewline
MTDF\_\allowbreak{}02 (Part I, Environmental) & Signal independently confirmed (Phase 3). Redshift transition localised to z~<~0.04 (Phase 6A, 3.6$\sigma$). Framework's strongest empirical evidence. \tabularnewline
MTDF\_\allowbreak{}04 (Part III, Photon coupling) & $\kappa$ $\approx$ f\textsubscript{kick}/3 is an observational anchor structurally related to the early field energy fraction and below Planck sensitivity. Paper frames this as speculative, which is appropriate. \tabularnewline
MTDF\_\allowbreak{}05 (Part IV, Cosmological) & Phase 5 full plik total $\Delta$$\chi$$^{2}$~=~+0.63; Phase 2 plik-lite precursor: +0.81; CMB+BAO summary run: +1.25 (narrower data). Hubble tension mechanism fully intact. Priority 1 hard falsifier cleared. \tabularnewline
MTDF\_\allowbreak{}07 (This paper) & Phase 6 adds ten cross-probe tests (6A--6E, 6C2--6C5). Title updated to Phases 1--6. All validation claims current as of February 2026. \tabularnewline
\bottomrule
\end{tabularx}
\end{table}

\section{Open questions and next steps}
\label{sec:sec8}

\subsection{NGC/\allowbreak{}SGC asymmetry diagnostics (Phase 3b, complete)}
Phase 3b (Section 5b) implemented the full asymmetry diagnostic battery and is now complete. All five tests confirm that the NGC/\allowbreak{}SGC contrast is not statistically significant, not caused by survey geometry, and attributable to SGC void-catalogue sparsity. The NGC/\allowbreak{}SGC claim status in the taxonomy (Section 7.2) has been upgraded from $\blacktriangle$~Hypothesis to \checkmark{}~Validated.

\subsection{Falsifier and discriminator tests}
\begin{table}[H]
\centering
\footnotesize
\begin{tabularx}{\linewidth}{l>{\RaggedRight\arraybackslash}X>{\RaggedRight\arraybackslash}X>{\RaggedRight\arraybackslash}X>{\RaggedRight\arraybackslash}X}
\toprule
\textbf{\#} & \textbf{Test} & \textbf{Type} & \textbf{Status} & \textbf{Timescale} \tabularnewline
1 & Full Planck plik + class\_\allowbreak{}mtdf & Discriminator & \textbf{PASSED} (Phase 5, $\Delta$$\chi$$^{2}$ = +0.63) & Complete \tabularnewline
2 & NGC/\allowbreak{}SGC diagnostic battery & Upper limit & \textbf{COMPLETE} (sensitivity-limited, not geometric) & Complete \tabularnewline
3 & Weak lensing $\times$ environment (KiDS) & Upper limit & \textbf{COMPLETE} (Phase 6B/\allowbreak{}B2: systematics-clean null) & Complete \tabularnewline
4 & CMB lensing $\times$ voids & Upper limit & \textbf{COMPLETE} (Phase 6D/\allowbreak{}E: pipeline validated) & Complete \tabularnewline
5 & S\textsubscript{8} cross-probe coherence & Coherence & \textbf{COMPLETE} (Phase 6C2: 3.5$\sigma$~$\rightarrow$~1.9$\sigma$) & Complete \tabularnewline
6 & f$\sigma$\textsubscript{8} $\times$ environment & Forecast & \textbf{COMPLETE} (Phase 6C4: 0.22$\sigma$, sensitivity-limited) & Complete \tabularnewline
7 & BAO residual structure & Forecast & \textbf{COMPLETE} (Phase 6C3: $\Delta$W = 0.27\%, undetectable) & Complete \tabularnewline
8 & Cluster M\textsubscript{dyn}/M\textsubscript{lens} & Resolved & \textbf{RESOLVED} (Phase 6C5/\allowbreak{}C5b: $\eta$~=~1, Case~B, M\textsubscript{dyn}/M\textsubscript{lens}~=~1) & Complete \tabularnewline
9 & Weak lensing P(k) (non-linear) & Discriminator & Requires full P(k) from class\_\allowbreak{}mtdf & Months \tabularnewline
10 & Merging cluster offsets & Discriminator & Requires N-body simulations & Institutional \tabularnewline
\bottomrule
\end{tabularx}
\end{table}

\begin{mdframed}[linecolor=accent,linewidth=0.5pt,backgroundcolor=callout-bg,leftline=true,rightline=true,topline=true,bottomline=true,innerleftmargin=8pt,innerrightmargin=8pt,innertopmargin=6pt,innerbottommargin=6pt]
\textbf{Phase 2 scope note, resolved by Phase 5.}
Phase 2 used a perturbative correction layer on a $\Lambda$CDM emulator, not the full modified Boltzmann equations. Phase~5 subsequently ran \texttt{class\_\allowbreak{}mtdf} through the full Planck plik likelihood with all nuisance parameters, confirming $\Delta$$\chi$$^{2}$~=~+0.63. The perturbative approximation in Phase 2 was validated by the full calculation in Phase 5.
\end{mdframed}

\section{Appendix: Technical notes}
\label{sec:sec9}

\subsection{A.1 \,k\textsubscript{f} normalisation conventions}
Three different \( k_f \) conventions were identified across the MTDF codebase during Phase 2. Understanding these is essential for interpreting historical results.

\begin{table}[H]
\centering
\footnotesize
\begin{tabularx}{\linewidth}{l>{\RaggedRight\arraybackslash}X>{\RaggedRight\arraybackslash}X>{\RaggedRight\arraybackslash}X}
\toprule
\textbf{Convention} & \textbf{"Full MTDF" k\textsubscript{f}} & \textbf{f\textsubscript{kick}} & \textbf{Source} \tabularnewline
A (old CLASS, pre-correction) & 0.10 & 0.003303 & November 2025 cobaya runs \tabularnewline
B (corrected CLASS) & 1.0 & 0.003303 & Theory prediction \tabularnewline
C (Phase 2 code, initial, corrected) & 0.001326 & 0.001326 & CosmoPower correction layer (bug) \tabularnewline
\bottomrule
\end{tabularx}
\end{table}

Convention C initially used \( f_{\rm kick} = \alpha\kappa = 0.001326 \) instead of the correct \( f_{\rm kick} = \lambda_{\rm MTDF}/24 = 0.003303 \). This factor-of-2.49 error caused a 3$\times$ overcorrection in the angular scale ratio per unit \( k_f \), leading to a spurious \( H_0 \)--\( k_f \) degeneracy in the first MCMC run (best-fit \( H_0 = 55.6 \), which is unphysical). After correction, all runs show \( H_0 \approx 67 \) and no degeneracies.

\subsection{A.2 \,V18 parameter reference}
\begin{table}[H]
\centering
\footnotesize
\begin{tabularx}{\linewidth}{l>{\RaggedRight\arraybackslash}X>{\RaggedRight\arraybackslash}X>{\RaggedRight\arraybackslash}X}
\toprule
\textbf{Parameter} & \textbf{Value} & \textbf{Uncertainty} & \textbf{Role} \tabularnewline
$\alpha$ & 1.30 & 0.26 & Stress-field coupling (void dynamics: KBC 2013 + Hoscheit 2018; counter: Kenworthy 2019) \tabularnewline
$\beta$ & 22.7 Mpc & 0.09$\times$$10^{23}$ m & Coherence length (void radius distribution: Pan+2012; Hamaus+2014) \tabularnewline
$\tau$ & 13.0 Gyr & 0.2 & Field age \tabularnewline
$\beta$\textsubscript{eos} & 0.573 & 0.012 & Equation-of-state parameter \tabularnewline
E & 9.1$\times$10\textsuperscript{$-$10} Pa & 0.4$\times$10\textsuperscript{$-$10} & Field energy scale \tabularnewline
$\delta$\textsubscript{bf} & 0.150 & 0.030 & Observational anchor: cluster baryon fraction at M\textsubscript{500} $\sim$ 10\textsuperscript{14} M\textsubscript{$\odot$} (Giodini+2009; Gonzalez+2013; see MTDF\_\allowbreak{}01 \S{}5.1) \tabularnewline
z\textsubscript{t} & 0.74 & 0.10 & Transition redshift \tabularnewline
$\kappa$ & 0.00102 & 0.00003 & Photon--stress coupling (anchor: $\kappa$ $\approx$ f\textsubscript{kick}/3; below FIRAS sensitivity) \tabularnewline
H\textsubscript{0} & 70.0 km/\allowbreak{}s/\allowbreak{}Mpc & 1.5 & Hubble constant (Freedman 2021 TRGB: 69.8 $\pm$ 0.6) \tabularnewline
f\textsubscript{kick} & 0.003303 & -- & Derived: (1$-$$\beta$\textsubscript{eos})$^{2}$/((1+$\alpha$)$\times$24) \tabularnewline
\bottomrule
\end{tabularx}
\end{table}

\subsection{A.3 \,Historical context: November 2025 runs}
Previous cobaya + class\_\allowbreak{}mtdf runs (November 2025) provide useful reference points. A single-point evaluation at the Planck $\Lambda$CDM best-fit gave $\chi$$^{2}$~=~614.23 (consistent with full plik TT+TE+EE). An MCMC with Planck TT-only (16 walkers, 200 steps) found \( k_f \approx 0.03 \) in the old normalisation (Convention A), with $\chi$$^{2}$ indistinguishable from $\Lambda$CDM. These early results are consistent with the current Phase 2 conclusion that Planck cannot constrain the EFE amplitude.

\section{References}
\begin{thebibliography}{19}
\bibitem{ref1} Aghanim, N. et al. (Planck Collaboration) 2020, "Planck 2018 results. V. Power spectra and likelihoods", \emph{A\&A}, 641, A5. DOI: \href{https://doi.org/10.1051/0004-6361/201936386}{10.1051/\allowbreak{}0004-6361/\allowbreak{}201936386}
\bibitem{ref2} Aghanim, N. et al. (Planck Collaboration) 2020, "Planck 2018 results. VI. Cosmological parameters", \emph{A\&A}, 641, A6. DOI: \href{https://doi.org/10.1051/0004-6361/201833910}{10.1051/\allowbreak{}0004-6361/\allowbreak{}201833910}
\bibitem{ref3} DESI Collaboration 2025, "DESI 2024 VI: Cosmological Constraints from BAO", \emph{JCAP}, 2025, 021. DOI: \href{https://doi.org/10.1088/1475-7516/2025/02/021}{10.1088/\allowbreak{}1475-7516/\allowbreak{}2025/\allowbreak{}02/\allowbreak{}021}
\bibitem{ref4} Brout, D. et al. 2022, "The Pantheon+ Analysis: Cosmological Constraints", \emph{ApJ}, 938, 110. DOI: \href{https://doi.org/10.3847/1538-4357/ac8e04}{10.3847/\allowbreak{}1538-4357/\allowbreak{}ac8e04}
\bibitem{ref5} Melia, F. \& Maier, R. S. 2013, "Cosmic chronometers in the R\_\allowbreak{}h=ct universe", \emph{MNRAS}, 432, 2669. DOI: \href{https://doi.org/10.1093/mnras/stt1082}{10.1093/\allowbreak{}mnras/\allowbreak{}stt1082}
\bibitem{ref6} Alam, S. et al. (eBOSS Collaboration) 2021, "Completed SDSS-IV eBOSS: Cosmological implications", \emph{Phys. Rev. D}, 103, 083533. DOI: \href{https://doi.org/10.1103/PhysRevD.103.083533}{10.1103/\allowbreak{}PhysRevD.103.083533}
\bibitem{ref7} Spurio Mancini, A. et al. 2022, "CosmoPower: emulating cosmological power spectra for accelerated Bayesian inference", \emph{MNRAS}, 511, 1771. DOI: \href{https://doi.org/10.1093/mnras/stac064}{10.1093/\allowbreak{}mnras/\allowbreak{}stac064}
\bibitem{ref8} Torrado, J. \& Lewis, A. 2021, "Cobaya: Code for Bayesian Analysis of hierarchical physical models", \emph{JCAP}, 2021, 057. DOI: \href{https://doi.org/10.1088/1475-7516/2021/05/057}{10.1088/\allowbreak{}1475-7516/\allowbreak{}2021/\allowbreak{}05/\allowbreak{}057}
\bibitem{ref9} Foreman-Mackey, D. et al. 2013, "emcee: The MCMC Hammer", \emph{PASP}, 125, 306. DOI: \href{https://doi.org/10.1086/670067}{10.1086/\allowbreak{}670067}
\bibitem{ref10} Blas, D., Lesgourgues, J. \& Tram, T. 2011, "The Cosmic Linear Anisotropy Solving System (CLASS). Part II", \emph{JCAP}, 2011, 034. DOI: \href{https://doi.org/10.1088/1475-7516/2011/07/034}{10.1088/\allowbreak{}1475-7516/\allowbreak{}2011/\allowbreak{}07/\allowbreak{}034}
\bibitem{ref11} Asgari, M. et al. (KiDS Collaboration) 2021, "KiDS-1000 Cosmology: Cosmic shear constraints", \emph{A\&A}, 645, A104. DOI: \href{https://doi.org/10.1051/0004-6361/202039070}{10.1051/\allowbreak{}0004-6361/\allowbreak{}202039070}
\bibitem{ref12} Abbott, T. M. C. et al. (DES Collaboration) 2022, "Dark Energy Survey Year 3 results: Cosmological constraints", \emph{Phys. Rev. D}, 105, 023520. DOI: \href{https://doi.org/10.1103/PhysRevD.105.023520}{10.1103/\allowbreak{}PhysRevD.105.023520}
\bibitem{ref13} Dalal, R. et al. (HSC Collaboration) 2023, "Hyper Suprime-Cam Year 3 Results: Cosmology from Cosmic Shear", \emph{Phys. Rev. D}, 108, 123519. DOI: \href{https://doi.org/10.1103/PhysRevD.108.123519}{10.1103/\allowbreak{}PhysRevD.108.123519}
\bibitem{ref14} Keenan, R. C., Barger, A. J. \& Cowie, L. L. 2013, "Evidence for a \~{}300 Megaparsec Scale Under-density", \emph{ApJ}, 775, 62. DOI: \href{https://doi.org/10.1088/0004-637X/775/1/62}{10.1088/\allowbreak{}0004-637X/\allowbreak{}775/\allowbreak{}1/\allowbreak{}62}
\bibitem{ref15} Pan, D. C. et al. 2012, "Cosmic Voids in SDSS DR7", \emph{MNRAS}, 421, 926. DOI: \href{https://doi.org/10.1111/j.1365-2966.2011.20197.x}{10.1111/\allowbreak{}j.1365-2966.2011.20197.x}
\bibitem{ref16} Freedman, W. L. 2021, "Measurements of the Hubble Constant: Tensions in Perspective", \emph{ApJ}, 919, 16. DOI: \href{https://doi.org/10.3847/1538-4357/ac0e95}{10.3847/\allowbreak{}1538-4357/\allowbreak{}ac0e95}
\bibitem{ref17} McGaugh, S. S., Lelli, F. \& Schombert, J. M. 2016, "Radial Acceleration Relation in Rotationally Supported Galaxies", \emph{PRL}, 117, 201101. DOI: \href{https://doi.org/10.1103/PhysRevLett.117.201101}{10.1103/\allowbreak{}PhysRevLett.117.201101}
\bibitem{ref18} Lelli, F. et al. 2017, "One Law to Rule Them All: The Radial Acceleration Relation", \emph{ApJ}, 836, 152. DOI: \href{https://doi.org/10.3847/1538-4357/836/2/152}{10.3847/\allowbreak{}1538-4357/\allowbreak{}836/\allowbreak{}2/\allowbreak{}152}
\bibitem{ref19} Riess, A. G. et al. 2022, "A Comprehensive Measurement of the Local Value of the Hubble Constant", \emph{ApJL}, 934, L7. DOI: \href{https://doi.org/10.3847/2041-8213/ac5c5b}{10.3847/\allowbreak{}2041-8213/\allowbreak{}ac5c5b}
\end{thebibliography}



\part{Multi-Probe Evidence for a Low-Redshift Transition}

\begin{abstract}
\textbf{Headline result.} Four independent low-redshift probes converge on a common environment-dependent transition at \(z = 0.040 \pm 0.005\), coinciding with the coherence length \(\beta = 22.7\) Mpc predicted independently from void quantisation by the Mesche Tensor Dynamics Framework (MTDF). The convergence is summarised in the table below.

\begin{table}[H]
\centering
\footnotesize
\caption{\textbf{Table A1.} Multi-probe convergence at the MTDF coherence scale.}
\begin{tabularx}{\linewidth}{l>{\RaggedRight\arraybackslash}X>{\RaggedRight\arraybackslash}X>{\RaggedRight\arraybackslash}X}
\toprule
\textbf{Probe} & \textbf{Dataset} & \textbf{Transition scale} & \textbf{Headline statistic} \tabularnewline
\midrule
SNe piecewise & Pantheon+ merged (3 void finders) & \(z = 0.040 \pm 0.005\) & \(\Delta\chi^2 = 36\text{--}39\), \(p < 0.001\) \tabularnewline
Peculiar velocity & Cosmicflows-4 (VoidFinder, V2G, REVOLVER) & Signal concentrated at \(z < 0.04\) & 9.0\(\sigma\) N-body tension (MDPL2, 100 observers) \tabularnewline
Tully-Fisher residuals & 2MTF (J, H, K bands) & Entire sample \(z < 0.033\) (in-regime) & 4.2\(\sigma\) N-body tension, achromatic \tabularnewline
Time-delay distances & H0LiCOW strong lenses & Geometric, local \(H_0\) (\(S_0 = 1.084\)) & 5.5\(\times\) preferred over \(\Lambda\)CDM in \(\chi^2/\mathrm{dof}\) \tabularnewline
\bottomrule
\end{tabularx}
\end{table}

Multiple independent low-redshift observables are tested for a common environment-dependent transition near a single coherence scale. The Mesche Tensor Dynamics Framework (MTDF) predicts that such deviations should concentrate below \(z \approx 0.04\), set by the coherence length \(\beta = 22.7\) Mpc. Four channels provide direct evidence. Supernova luminosity distances show an environment signal (\(p = 0.003\)) with a sharp piecewise cutoff at \(z = 0.04\). Cosmicflows-4 peculiar velocities reveal a void-geometric coupling that survives a seven-layer control chain and exhibits a qualitative sign mismatch against full N-body \(\Lambda\)CDM expectations (\(9.0\sigma\) tension, MDPL2, 100 observers). 2MTF Tully-Fisher residuals independently confirm the pattern, achromatic across three infrared bands, with \(4.2\sigma\) N-body tension. An additional geometric support channel, strong-lens time-delay distances, favours the MTDF-predicted local \(H_0\) by 5.5\(\times\) in \(\chi^2/\mathrm{dof}\). Two further channels (ISW void stacking, weak-lensing S\textsubscript{8} splits) establish methodology and show directional consistency but remain below detection threshold with current data. The convergence of independent probes at a single coherence scale motivates targeted independent replication and dedicated \(\Lambda\)CDM systematics tests of the same environment-dependent observables.
\end{abstract}

\begin{mdframed}[linecolor=accent,linewidth=0.5pt,backgroundcolor=callout-bg,leftline=true,rightline=true,topline=true,bottomline=true,innerleftmargin=8pt,innerrightmargin=8pt,innertopmargin=6pt,innerbottommargin=6pt]
\textbf{Nominal significance.} Unless explicitly stated otherwise, quoted significances are conditional on the specified catalogue, control chain, redshift cut, and metric definition. They should be read as nominal or pipeline-conditioned significances, not as final global discovery significances after all possible analysis choices. A full look-elsewhere assessment across alternative void metrics, redshift partitions, catalogue choices, and environment definitions remains a priority for independent replication.
\end{mdframed}

\section{Introduction}
\label{sec:sec1}
Standard cosmological analyses typically treat low-redshift probes independently. Supernova distance moduli, peculiar velocity fields, galaxy scaling relations, and CMB secondary anisotropies are each analysed within their own pipelines, subject to their own systematics, and compared against \(\Lambda\)CDM expectations in isolation. Environment is usually treated as a nuisance variable to be marginalised over rather than a carrier of physical information.

This paper takes a different approach. It tests whether several independent low-redshift observables point to a common environment-linked transition concentrated near a single redshift scale. If multiple probes, each with different physics, different systematics, and different survey selection functions, converge on the same transition regime, this constitutes a physical clue that cannot easily be attributed to any one dataset or method.

The motivation arises from the Mesche Tensor Dynamics Framework (MTDF), which predicts a finite coherence scale \(\beta = 22.7\) Mpc corresponding to a characteristic redshift \(z \approx 0.04\). Below this scale, the MTDF stress tensor field couples to observables in an environment-dependent manner; above it, the effect averages out over multiple coherence volumes. This prediction was established in MTDF\_\allowbreak{}01 and MTDF\_\allowbreak{}02 before the multi-probe tests reported here were conducted.

The paper is structured as follows. Section 2 states the working hypothesis and its quantitative prediction. Section 3 describes the six datasets. Section 4 presents the shared analysis logic applied uniformly across all channels. Section 5 reports the results for each channel with its associated controls. Section 6 synthesises the cross-probe evidence. Section 7 presents the \(\Lambda\)CDM mock comparison, including full N-body results. Sections 8 and 9 discuss interpretation and limitations. Section 10 concludes.

\section{Working hypothesis and prediction}
\label{sec:sec2}
MTDF introduces a stress tensor field with a finite coherence length \(\beta = 22.7\) Mpc, corresponding to a comoving distance of approximately 170 Mpc or \(z \approx 0.04\) in the standard cosmological background. The framework predicts that environment-dependent deviations from \(\Lambda\)CDM should exhibit a scale hierarchy governed by the dimensionless product \(\beta k\), where \(k\) is the characteristic wavenumber of the observation:

\[
\beta k \ll 1 \;\Rightarrow\; \text{effects active (cosmological scales)}, \qquad
\beta k \gg 1 \;\Rightarrow\; \text{effects suppressed (local scales)}
\]

The specific, falsifiable prediction tested in this paper is:

\begin{mdframed}[linecolor=accent,linewidth=0.5pt,backgroundcolor=callout-bg,leftline=true,rightline=true,topline=true,bottomline=true,innerleftmargin=8pt,innerrightmargin=8pt,innertopmargin=6pt,innerbottommargin=6pt]
\textbf{Prediction.} Environment-dependent residuals in distance indicators and velocity fields should be (i) present below \(z \approx 0.04\), (ii) absent or strongly suppressed above \(z \approx 0.04\), and (iii) correlated with the signed distance to the nearest void boundary as classified by independent void-finding algorithms applied to galaxy redshift surveys.
\end{mdframed}

This prediction was derived from the MTDF parameter set established in MTDF\_\allowbreak{}01 (specifically \(\beta = 22.7\) Mpc and scale factor \(S_0 = 1.084\)) and published in MTDF\_\allowbreak{}02 before any of the multi-probe tests reported in Sections 5--7 were conducted. The prediction is sharp: a smooth or absent transition would falsify the coherence-scale interpretation.

\section{Datasets and observables}
\label{sec:sec3}
Six independent observables are tested. They span different physical processes, different survey instruments, different galaxy populations, and different redshift ranges within the low-\(z\) universe. No two channels share calibration systems or host-galaxy properties.

\subsection{Supernova luminosity distances}
The merged sample comprises 4,188 Type Ia supernovae from Pantheon+ (1,533 SNe with full STAT+SYS covariance), ZTF DR2 (2,650 SNe, Tripp-standardised), and Foundation DR1 (5 unique SNe after deduplication). Distance modulus residuals are computed after SALT2 standardisation. Environment classification uses DESI DESIVAST void catalogues cross-matched at 5 arcsecond tolerance. The low-redshift subsample at \(z < 0.04\) provides the primary test regime. This channel was established in MTDF\_\allowbreak{}02 and the GPU validation paper MTDF\_\allowbreak{}07; it is summarised here for completeness.

\subsection{Strong-lens time-delay distances}
Seven multiply-imaged quasar systems from the TDCOSMO/\allowbreak{}H0LiCOW programme provide time-delay distance measurements that are geometrically independent of supernova luminosity distances. An additional 25 SLACS lens systems provide velocity dispersion-based distance estimates at lower redshifts (\(z_L = 0.08\text{--}0.35\)), where DESIVAST void coverage is more complete.

\subsection{Cosmicflows-4 peculiar velocities}
The Cosmicflows-4 catalogue (Tully et al. 2023) contains 55,877 galaxies in 38,053 groups at \(z < 0.05\), with distances measured via eight independent methods (primarily Tully-Fisher for spirals and Fundamental Plane for ellipticals). Group-averaged peculiar velocities are used to reduce individual measurement scatter. The catalogue is accessed from the Extragalactic Distance Database (EDD). Environment classification uses the same DESIVAST void catalogues as the supernova analysis.

\subsection{2MTF Tully-Fisher distances}
The 2MASS Tully-Fisher survey (Hong et al. 2019) provides Tully-Fisher distances for 2,062 nearby spiral galaxies at \(600 < cz < 10{,}000\) km/\allowbreak{}s (\(z < 0.033\)). Distances are measured in three independent infrared bands (J, H, K) from 2MASS photometry, combined with HI 21 cm linewidths. The entire sample falls below \(z = 0.04\), placing it entirely within the predicted MTDF-active regime.

\subsection{ISW void stacking}
The integrated Sachs-Wolfe (ISW) effect is tested via stacked CMB temperature profiles centred on voids from the BOSS and DESIVAST catalogues. Planck PR4 SMICA temperature maps provide the CMB signal. A total of 889 BOSS voids and 563 DESIVAST voids (1,452 combined) are used. MTDF predicts a specific void radial profile shape that differs from the step-function template commonly used in ISW analyses.

\subsection{Weak-lensing S\textsubscript{8} environment split}
Galaxy-galaxy lensing (GGL) tangential shear profiles are measured from KiDS-1000 source galaxies around GAMA spectroscopic group lenses, split by foreground DESIVAST void classification. The analysis processes 2.7 billion source-lens pairs across 1,724 void-embedded and 80 wall-embedded galaxy groups.

\section{Common analysis logic}
\label{sec:sec4}
All channels are evaluated under a shared methodological framework. This consistency is essential: without it, multi-probe convergence could be an artefact of inconsistent analysis choices rather than a physical signal.

\subsection{Signed-distance environment metric}
Environment is quantified using a continuous signed-distance metric rather than a binary void/\allowbreak{}wall classification. For each object, the signed distance \(d_\mathrm{signed}\) to the nearest DESIVAST void boundary is computed as:

\[
d_\mathrm{signed} = \frac{r_\mathrm{obj} - R_\mathrm{void}}{R_\mathrm{void}}
\]

where \(r_\mathrm{obj}\) is the distance from the object to the void centre and \(R_\mathrm{void}\) is the void radius. Negative values indicate positions inside a void; positive values indicate wall or filament environments. This continuous metric retains more information than a binary split and enables gradient analysis even when one environment class has few or zero members.

\subsection{Three independent void catalogues}
All analyses are repeated with three independent void-finding algorithms applied to the DESI Bright Galaxy Survey: VoidFinder, REVOLVER, and VIDE. These algorithms use different void definitions, identification procedures, and boundary conventions. Consistency across all three provides robustness against void-finder-specific artefacts.

\subsection{GLS regression framework}
Each channel is tested using generalised least-squares (GLS) regression of the observable residual against \(d_\mathrm{signed}\):

\[
y_i = \alpha + \gamma \cdot d_{\mathrm{signed},i} + \varepsilon_i
\]

where \(y_i\) is the residual for the \(i\)-th object (distance modulus residual, peculiar velocity residual, or TF residual), \(\gamma\) is the environment coupling coefficient, and the covariance structure is channel-specific. The null hypothesis is \(\gamma = 0\).

\subsection{Piecewise redshift split}
For channels with sufficient redshift range, a piecewise model splits the sample at \(z = 0.04\) and fits separate \(\gamma\) values above and below the cut. The \(\Delta\chi^2\) between the piecewise model and the single-\(\gamma\) model quantifies the evidence for a redshift-dependent transition. This directly tests the coherence-scale prediction.

\subsection{Robustness chain}
Each channel is subjected to a seven-layer control chain (layers applied where applicable):

\begin{table}[H]
\centering
\footnotesize
\caption{\textbf{Table 1.} Standard control chain applied to each observable channel.}
\begin{tabularx}{\linewidth}{l>{\RaggedRight\arraybackslash}X>{\RaggedRight\arraybackslash}X}
\toprule
\textbf{Layer} & \textbf{Control test} & \textbf{Question addressed} \tabularnewline
\midrule
A & Density substitution & Is \(d_\mathrm{signed}\) simply a proxy for local galaxy number density? \tabularnewline
B & Partial correlation & Does \(d_\mathrm{signed}\) carry information beyond local density? \tabularnewline
C & \(\Lambda\)CDM chain rule & Is the observed coupling consistent with the density-observable relation predicted by \(\Lambda\)CDM? \tabularnewline
D & 2M++ velocity-field subtraction & Does the signal survive after removing the best \(\Lambda\)CDM velocity reconstruction? \tabularnewline
E & Malmquist bias test & Is the signal driven by distance-dependent selection effects or catalogue geometry? \tabularnewline
F & Leave-one-method-out (LOSO) & Does the signal depend on a single distance-method family? \tabularnewline
N-body & MDPL2 mock (100 observers) & Can a full non-linear \(\Lambda\)CDM simulation reproduce the signal? \tabularnewline
\bottomrule
\end{tabularx}
\end{table}

Layers E and F are applied where the dataset combines multiple distance methods (CF4). Additionally, all channels are checked for footprint consistency (NGC vs SGC) and permutation significance (10,000 shuffles of \(d_\mathrm{signed}\) labels with 5,000 bootstrap resamples).

\section{Channel results}
\label{sec:sec5}

\subsection{Supernovae: environment cutoff at \emph{z} = 0.04}
The supernova environment signal was established in MTDF\_\allowbreak{}02 and independently validated in MTDF\_\allowbreak{}07. On the merged 4,188-SN sample, a piecewise model with a break at \(z = 0.04\) yields \(\Delta\chi^2 = 36\text{--}39\) (\(p < 0.001\)) compared to a single-regime fit. The signal is concentrated in the low-redshift bin:

\begin{table}[H]
\centering
\footnotesize
\caption{\textbf{Table 2.} Supernova environment signal by redshift regime (Pantheon+ $\times$ DESIVAST).}
\begin{tabularx}{\linewidth}{l>{\RaggedRight\arraybackslash}X>{\RaggedRight\arraybackslash}X>{\RaggedRight\arraybackslash}X}
\toprule
\textbf{Regime} & \textbf{\(\gamma_\mathrm{env}\)} & \textbf{Significance} & \textbf{Void catalogues} \tabularnewline
\midrule
\(z \in [0.02, 0.04)\) & \(+0.013 \pm 0.004\) & \(p = 0.003\text{--}0.005\) & REVOLVER, VIDE \tabularnewline
\(z > 0.04\) & Consistent with zero & \(p > 0.84\) & All three \tabularnewline
\bottomrule
\end{tabularx}
\end{table}

\begin{mdframed}[linecolor=accent,linewidth=0.5pt,backgroundcolor=callout-bg,leftline=true,rightline=true,topline=true,bottomline=true,innerleftmargin=8pt,innerrightmargin=8pt,innertopmargin=6pt,innerbottommargin=6pt]
\textbf{Control outcome.} The piecewise cutoff at \(z = 0.04\) was derived from MTDF's coherence scale \(\beta = 22.7\) Mpc before the redshift-binned analysis was performed. The signal passes LOSO (16/\allowbreak{}16 surveys retain positive \(\gamma_\mathrm{env}\)), permutation tests (\(p_\mathrm{perm} = 0.025\text{--}0.062\)), block bootstrap (95\% CI excludes zero for REVOLVER and VIDE), and survey fixed effects (\(\gamma_\mathrm{env}\) stable within 0.0005 mag). SALT2 population controls shift \(\gamma_\mathrm{env}\) by less than 0.3\(\sigma\).
\end{mdframed}

\subsection{Time-delay distances: independent geometric support}
This channel differs in kind from the other probes: it does not test the void-wall environment split directly but provides an independent geometric consistency check on the MTDF low-redshift interpretation. Seven TDCOSMO/\allowbreak{}H0LiCOW lenses provide time-delay distance measurements that test the MTDF scale factor \(S_0 = 1.084\), which maps the Planck-calibrated \(H_0 = 67.36\) km/\allowbreak{}s/\allowbreak{}Mpc to an effective local value of \(H_0 \approx 73.1\) km/\allowbreak{}s/\allowbreak{}Mpc.

\begin{table}[H]
\centering
\footnotesize
\caption{\textbf{Table 3.} Time-delay distance goodness of fit: MTDF vs \(\Lambda\)CDM.}
\begin{tabularx}{\linewidth}{l>{\RaggedRight\arraybackslash}X>{\RaggedRight\arraybackslash}X>{\RaggedRight\arraybackslash}X}
\toprule
\textbf{Model} & \textbf{\(\chi^2/\mathrm{dof}\)} & \textbf{Mean residual} & \textbf{Residual pattern} \tabularnewline
\midrule
\(\Lambda\)CDM (\(H_0 = 67.36\)) & 26.05/\allowbreak{}7 = 3.72 & \(-1.58\sigma\) & All 7 negative (systematic) \tabularnewline
MTDF (\(H_0 = 73.1\) via \(S_0\)) & 4.77/\allowbreak{}7 = 0.68 & \(-0.10\sigma\) & Mixed +/\(-\) (unbiased) \tabularnewline
\bottomrule
\end{tabularx}
\end{table}

MTDF fits the time-delay data 5.5 times better than Planck \(\Lambda\)CDM. The \(\Lambda\)CDM residuals are systematically negative across all seven lenses, indicating a coherent underprediction of time-delay distances. The MTDF residuals scatter around zero with no systematic bias.

An environment correlation test on the TDCOSMO lenses is inconclusive: all seven systems are massive ellipticals residing in wall environments (\(d_\mathrm{signed} > 0\)), providing zero void sightlines. This is a selection effect intrinsic to strong lensing, not a failure of the prediction. An extension to 25 SLACS lenses yields only 3 void sightlines out of 25, insufficient for a meaningful binary split (\(p = 0.30\)).

\begin{mdframed}[linecolor=accent,linewidth=0.5pt,backgroundcolor=callout-bg,leftline=true,rightline=true,topline=true,bottomline=true,innerleftmargin=8pt,innerrightmargin=8pt,innertopmargin=6pt,innerbottommargin=6pt]
\textbf{Interpretation.} The time-delay result is independent geometric support for the MTDF low-redshift framework: it confirms that the effective local expansion rate is consistent with MTDF's \(S_0\) prediction. It does not test the void-geometry coupling directly and should not be weighted as direct evidence for the environment split. Its role here is to remove the Hubble tension as a competing explanation, since MTDF resolves it naturally through the scale factor.
\end{mdframed}

\subsection{Cosmicflows-4 peculiar velocities}
This is the primary new result of this paper. The Cosmicflows-4 group-averaged catalogue provides 38,053 groups with peculiar velocities at \(z < 0.05\), cross-matched against all three DESIVAST void catalogues.

\subsubsection{Baseline signal}
The environment coupling coefficient \(\gamma_v\) measures the slope of peculiar velocity residuals (after shell-median subtraction in redshift bins of width 0.005) against \(d_\mathrm{signed}\):

\begin{table}[H]
\centering
\footnotesize
\caption{\textbf{Table 4.} CF4 peculiar velocity environment coupling: baseline results.}
\begin{tabularx}{\linewidth}{l>{\RaggedRight\arraybackslash}X>{\RaggedRight\arraybackslash}X>{\RaggedRight\arraybackslash}X>{\RaggedRight\arraybackslash}X}
\toprule
\textbf{Void finder} & \textbf{\(\gamma_v\) (km/\allowbreak{}s per Mpc)} & \textbf{Significance} & \textbf{\(\Delta\chi^2\)} & \textbf{\(p_\mathrm{perm}\)} \tabularnewline
\midrule
VoidFinder & \(-53.6 \pm 2.7\) & \(19.6\sigma\) & 386 & \(<10^{-4}\) \tabularnewline
REVOLVER & \(-61.5 \pm 4.0\) & \(15.5\sigma\) & 240 & \(<10^{-4}\) \tabularnewline
VIDE & \(-58.2 \pm 5.3\) & \(11.0\sigma\) & 120 & \(<10^{-4}\) \tabularnewline
\bottomrule
\end{tabularx}
\end{table}

\subsubsection{Piecewise redshift split}
The signal concentrates below \(z = 0.04\), exactly as predicted by the MTDF coherence scale:

\begin{table}[H]
\centering
\footnotesize
\caption{\textbf{Table 5.} CF4 piecewise split at \(z = 0.04\) (VoidFinder).}
\begin{tabularx}{\linewidth}{l>{\RaggedRight\arraybackslash}X>{\RaggedRight\arraybackslash}X>{\RaggedRight\arraybackslash}X}
\toprule
\textbf{Regime} & \textbf{\(\gamma_v\) (km/\allowbreak{}s per Mpc)} & \textbf{Significance} & \textbf{\(\Delta\chi^2\)} \tabularnewline
\midrule
\(z < 0.04\) & \(-52.9 \pm 3.6\) & \(14.7\sigma\) & 216 \tabularnewline
\(z \geq 0.04\) & \(-13.9 \pm 7.4\) & \(1.9\sigma\) & 3.5 \tabularnewline
\bottomrule
\end{tabularx}
\end{table}

\subsubsection{Control chain}
\textbf{Control A (density substitution).} Local galaxy number density, computed as neighbour counts within spheres of radius 5, 10, and 20 Mpc/\allowbreak{}h, is also significantly correlated with peculiar velocity residuals (13.7--29.3\(\sigma\)). This is expected: peculiar velocities are driven by the density field in any cosmological model. The question is whether \(d_\mathrm{signed}\) carries additional information.

\textbf{Control B (partial correlation).} After regressing out local density, \(d_\mathrm{signed}\) adds substantial \(\chi^2\) at all density scales:

\begin{table}[H]
\centering
\footnotesize
\caption{\textbf{Table 6.} CF4 partial correlation: \(\Delta\chi^2\) added by \(d_\mathrm{signed}\) beyond local density.}
\begin{tabularx}{\linewidth}{l>{\RaggedRight\arraybackslash}X>{\RaggedRight\arraybackslash}X>{\RaggedRight\arraybackslash}X}
\toprule
\textbf{Void finder} & \textbf{\(R = 5\) Mpc/\allowbreak{}h} & \textbf{\(R = 10\)} & \textbf{\(R = 20\)} \tabularnewline
\midrule
VoidFinder & +206 (14.3\(\sigma\)) & +259 (16.1\(\sigma\)) & +275 (16.6\(\sigma\)) \tabularnewline
REVOLVER & +141 (11.9\(\sigma\)) & +165 (12.8\(\sigma\)) & +170 (13.0\(\sigma\)) \tabularnewline
VIDE & +71 (8.4\(\sigma\)) & +77 (8.8\(\sigma\)) & +76 (8.7\(\sigma\)) \tabularnewline
\bottomrule
\end{tabularx}
\end{table}

The piecewise pattern survives density removal: at \(z < 0.04\), \(d_\mathrm{signed}\) adds 80--181 \(\Delta\chi^2\) (9--13\(\sigma\)); at \(z \geq 0.04\), it adds 1--7 \(\Delta\chi^2\) (1--3\(\sigma\), marginal).

\textbf{Control C (\(\Lambda\)CDM chain rule).} The observed \(\gamma_v\) exceeds the prediction from the \(\Lambda\)CDM density-velocity relation by a factor of 150--610. The signal is not "velocity tracing density" at the expected amplitude.

\textbf{Control D (2M++ velocity-field subtraction).} The gold-standard test: the Carrick et al. (2015) 2M++ reconstructed \(\Lambda\)CDM velocity field, used by Pantheon+, SH0ES, and all major peculiar velocity analyses, is subtracted from the CF4 peculiar velocities before recomputing \(\gamma_v\). This removes the full \(\Lambda\)CDM density-velocity coupling, not merely a local density proxy.

\begin{table}[H]
\centering
\footnotesize
\caption{\textbf{Table 6b.} CF4 environment coupling after 2M++ velocity-field subtraction.}
\begin{tabularx}{\linewidth}{l>{\RaggedRight\arraybackslash}X>{\RaggedRight\arraybackslash}X>{\RaggedRight\arraybackslash}X}
\toprule
\textbf{Void finder} & \textbf{Before 2M++ (\(\sigma\))} & \textbf{After 2M++ (\(\sigma\))} & \textbf{Retention} \tabularnewline
\midrule
VoidFinder & 19.6 & 9.2 & 65\% \tabularnewline
REVOLVER & 15.5 & 5.9 & 57\% \tabularnewline
VIDE & 11.0 & 5.2 & 77\% \tabularnewline
\bottomrule
\end{tabularx}
\end{table}

The 2M++ field absorbs the standard density-velocity coupling that \(\Lambda\)CDM explains, reducing \(\gamma_v\) from approximately \(-55\) to \(-30\) km/\allowbreak{}s. Notably, the 2M++ predicted velocity is itself anti-correlated with \(d_\mathrm{signed}\) (\(r \approx -0.18\)), indicating that the reconstruction partially captures the environment-dependent component by attributing it to gravitational infall; the partial-correlation framework of Control B is therefore essential for disentangling the two contributions. All three void catalogues remain above \(5\sigma\) after subtraction. The piecewise \(z = 0.04\) concentration also survives: VoidFinder retains \(4.1\sigma\) at \(z < 0.04\) after 2M++ subtraction. The residual void-geometric signal at 5--9\(\sigma\) represents the component that the best available \(\Lambda\)CDM velocity reconstruction cannot account for.

\begin{mdframed}[linecolor=accent,linewidth=0.5pt,backgroundcolor=callout-bg,leftline=true,rightline=true,topline=true,bottomline=true,innerleftmargin=8pt,innerrightmargin=8pt,innertopmargin=6pt,innerbottommargin=6pt]
\textbf{Control E (Malmquist bias).} The signed distance \(d_\mathrm{signed}\) correlates with heliocentric distance (\(r = -0.63\)), reflecting the redshift-dependent completeness of the DESIVAST void catalogues. To test whether this drives the signal, the sample is split by distance-error quality (high/\allowbreak{}low) and by distance bins. The signal is present in the high-quality subsample at \(11.6\sigma\), in every distance bin from 0 to 300 Mpc, and weakens only beyond 300 Mpc, where the MTDF coherence scale predicts the transition to fade (\(z \sim 0.07\)). The calibrator subsample, which has the most precise distances and is the most local, shows the strongest per-galaxy signal. This is the opposite of what a Malmquist-driven artefact would produce.

\textbf{Control F (leave-one-method-out).} The CF4 catalogue combines four distance-method families: Tully-Fisher (TF, \(\sim\)10,000 groups), Fundamental Plane (FP, \(\sim\)28,000), Type Ia supernovae (945), and calibrators (336). Each method is removed in turn and \(\gamma_v\) is recomputed from the remainder; separately, each method is tested alone:

\begin{table}[H]
\centering
\footnotesize
\caption{\textbf{Table 6c.} Leave-one-method-out (LOSO) and single-method stability (VoidFinder).}
\begin{tabularx}{\linewidth}{l>{\RaggedRight\arraybackslash}X>{\RaggedRight\arraybackslash}X}
\toprule
\textbf{Method} & \textbf{Leave-out \(\gamma_v\) (\(\sigma\))} & \textbf{Single-method \(\gamma_v\) (\(\sigma\))} \tabularnewline
\midrule
TF (\(\sim\)10K) & \(-67.7\) (\(18.8\sigma\)) & \(-34.5\) (\(8.2\sigma\)) \tabularnewline
FP (\(\sim\)28K) & \(-64.5\) (\(14.9\sigma\)) & \(-37.2\) (\(9.7\sigma\)) \tabularnewline
SNIa (945) & \(-64.8\) (\(21.0\sigma\)) & \(-15.5\) (\(2.7\sigma\)) \tabularnewline
Calibrators (336) & \(-40.2\) (\(14.4\sigma\)) & \(-495.2\) (\(29.0\sigma\)) \tabularnewline
\bottomrule
\end{tabularx}
\end{table}

The signal does not disappear when any single method is removed. Both TF and FP independently deliver the signal above \(8\sigma\). The leave-out shifts between methods (20--30\%) are consistent with differences in sky coverage and redshift distribution. The calibrator single-method result is extreme in amplitude because of the very small sample size (336 groups); it is informative only for sign consistency and local precision, not for amplitude comparison. No single distance pipeline drives the observed void-geometric coupling.
\end{mdframed}

\begin{mdframed}[linecolor=accent,linewidth=0.5pt,backgroundcolor=callout-bg,leftline=true,rightline=true,topline=true,bottomline=true,innerleftmargin=8pt,innerrightmargin=8pt,innertopmargin=6pt,innerbottommargin=6pt]
\textbf{Full control summary.} The CF4 signal survives seven control layers: density substitution (A), partial correlation at 8--17\(\sigma\) (B), \(\Lambda\)CDM chain-rule excess of 150--610\(\times\) (C), 2M++ subtraction at 5--9\(\sigma\) (D), Malmquist bias test across distance bins and quality halves (E), cross-method LOSO stability (F), and MDPL2 N-body mock tension of 9.0\(\sigma\) with sign mismatch. None of the standard mechanisms tested in this analysis accounts for the observed void-geometric coupling under the stated assumptions.
\end{mdframed}

\subsection{2MTF Tully-Fisher distances}
The 2MTF sample provides an independent distance indicator. All 2,062 galaxies fall below \(z = 0.033\), entirely within the predicted MTDF-active regime.

\subsubsection{Baseline signal}
\begin{table}[H]
\centering
\footnotesize
\caption{\textbf{Table 7.} 2MTF Tully-Fisher environment coupling across infrared bands and void finders.}
\begin{tabularx}{\linewidth}{l>{\RaggedRight\arraybackslash}X>{\RaggedRight\arraybackslash}X>{\RaggedRight\arraybackslash}X>{\RaggedRight\arraybackslash}X}
\toprule
\textbf{Void finder} & \textbf{Band} & \textbf{\(\gamma_\mathrm{TF}\) (mag/\allowbreak{}Mpc)} & \textbf{Significance} & \textbf{\(p_\mathrm{perm}\)} \tabularnewline
\midrule
VoidFinder & K & \(0.058 \pm 0.007\) & \(7.9\sigma\) & \(<10^{-4}\) \tabularnewline
VoidFinder & H & \(0.055 \pm 0.007\) & \(7.5\sigma\) & \(<10^{-4}\) \tabularnewline
VoidFinder & J & \(0.053 \pm 0.007\) & \(7.3\sigma\) & \(<10^{-4}\) \tabularnewline
REVOLVER & K & \(0.058 \pm 0.010\) & \(6.0\sigma\) & \(<10^{-4}\) \tabularnewline
VIDE & K & \(0.046 \pm 0.013\) & \(3.5\sigma\) & \(0.0007\) \tabularnewline
\bottomrule
\end{tabularx}
\end{table}

The consistency across J, H, and K bands is a critical systematic check. Interstellar dust produces wavelength-dependent extinction; a real distance or environment effect is achromatic. The band-to-band variation in \(\gamma_\mathrm{TF}\) is less than 10\%, consistent with an achromatic signal. The direction of the TF coupling matches the supernova \(\gamma_\mathrm{env}\): galaxies closer to void boundaries show systematic distance residuals.

\subsubsection{Control chain}
\textbf{Control B (partial correlation).} After removing local density, \(d_\mathrm{signed}\) adds:

\begin{table}[H]
\centering
\footnotesize
\caption{\textbf{Table 8.} 2MTF partial correlation: \(\Delta\chi^2\) added by \(d_\mathrm{signed}\) beyond local density.}
\begin{tabularx}{\linewidth}{l>{\RaggedRight\arraybackslash}X>{\RaggedRight\arraybackslash}X>{\RaggedRight\arraybackslash}X}
\toprule
\textbf{Void finder} & \textbf{\(R = 5\) Mpc/\allowbreak{}h} & \textbf{\(R = 10\)} & \textbf{\(R = 20\)} \tabularnewline
\midrule
VoidFinder & +235 (15.3\(\sigma\)) & +219 (14.8\(\sigma\)) & +175 (13.2\(\sigma\)) \tabularnewline
REVOLVER & +136 (11.7\(\sigma\)) & +126 (11.2\(\sigma\)) & +95 (9.7\(\sigma\)) \tabularnewline
VIDE & +44 (6.6\(\sigma\)) & +37 (6.1\(\sigma\)) & +20 (4.4\(\sigma\)) \tabularnewline
\bottomrule
\end{tabularx}
\end{table}

\textbf{Control C.} The observed \(\gamma_\mathrm{TF}\) exceeds the density chain prediction by a factor of 3--31.

\begin{mdframed}[linecolor=accent,linewidth=0.5pt,backgroundcolor=callout-bg,leftline=true,rightline=true,topline=true,bottomline=true,innerleftmargin=8pt,innerrightmargin=8pt,innertopmargin=6pt,innerbottommargin=6pt]
\textbf{Band consistency.} The achromatic nature of the 2MTF signal (J, H, K consistent within 10\%) rules out dust as the source of the environment-distance correlation. The signal is void-geometric, density-independent, and present across all three void catalogues and both galactic caps (NGC and SGC).
\end{mdframed}

\textbf{Note on void classification.} All 2,062 galaxies are classified as "wall" by all three void finders (\(n_\mathrm{void} = 0\)). The DESIVAST void catalogues, constructed from the DESI Bright Galaxy Survey, do not extend to the environments where 2MTF spirals reside at these redshifts. The \(\gamma_\mathrm{TF}\) signal is therefore a gradient within the wall population measured by the continuous \(d_\mathrm{signed}\) metric, not a binary void-versus-wall split. This does not invalidate the result but should be noted when comparing to the supernova binary split.

\subsection{Supporting channels}

\subsubsection{ISW void stacking}
The ISW analysis stacks Planck PR4 SMICA CMB temperature profiles on 889 BOSS and 563 DESIVAST voids. Absolute significance levels are sub-threshold (0.2--0.4\(\sigma\)), which is physically expected: the ISW temperature signal is intrinsically weak, and 1,452 voids provide insufficient statistical power for a detection.

However, the MTDF-predicted void template achieves 2.4--3.7 times the signal-to-noise of a step-function template matched filter. This sensitivity advantage is a methodological result: the MTDF void profile shape better captures whatever ISW signal exists, even at sub-threshold amplitude. With 10 times more voids (DESI Y5, Euclid), the MTDF template could reach 1--2\(\sigma\) while the step-function template remains below 0.5\(\sigma\).

\begin{mdframed}[linecolor=accent,linewidth=0.5pt,backgroundcolor=callout-bg,leftline=true,rightline=true,topline=true,bottomline=true,innerleftmargin=8pt,innerrightmargin=8pt,innertopmargin=6pt,innerbottommargin=6pt]
\textbf{Status.} No ISW detection claim is made. The result is a 2.4--3.7$\times$ template sensitivity gain, establishing methodology for future surveys with larger void catalogues.
\end{mdframed}

\subsubsection{Weak-lensing S\textsubscript{8} environment split}
Galaxy-galaxy lensing profiles for GAMA groups embedded in DESIVAST voids (1,724 groups) versus walls (80 groups) show an amplitude difference of 1.71\(\sigma\). The direction is consistent with the MTDF prediction: S\textsubscript{8,trough} < S\textsubscript{8,ridge} (0.47\(\sigma\) in the trough/\allowbreak{}ridge proxy). The wall sample is very small (80 groups), limiting statistical power. Euclid and LSST would provide a 10--15 times sensitivity improvement.

\begin{mdframed}[linecolor=accent,linewidth=0.5pt,backgroundcolor=callout-bg,leftline=true,rightline=true,topline=true,bottomline=true,innerleftmargin=8pt,innerrightmargin=8pt,innertopmargin=6pt,innerbottommargin=6pt]
\textbf{Status.} Suggestive at 1.71\(\sigma\), directionally consistent. Methodology established. Below detection threshold with current data.
\end{mdframed}

\subsubsection{Distance duality and lensed-image redshifts}
Two additional channels were tested and found to be insensitive to the predicted signal with current data. The distance duality test using DESI Y1 BAO yields \(\eta = 1.001 \pm 0.012\), with the MTDF signal 4.7 times below current precision. BAO distances are sky-averaged, precluding per-sightline environment splits. The lensed-image redshift analysis covers 20 multiply-imaged systems (64 images) with a mean significance of 0.49\(\sigma\); the MTDF predicted signal (\(\sim 10^{-7}\)) is three orders of magnitude below spectroscopic precision (\(\sim 10^{-4}\)). Both channels establish methodology for future instrumentation (DESI Y5, ELT/\allowbreak{}ANDES) but provide no constraint at present.

\section{Cross-probe synthesis}
\label{sec:sec6}
The central result of this paper is not any individual channel significance but the convergence of independent probes at a single coherence scale.

\begin{figure}[H]
\centering
\includegraphics[width=\linewidth,keepaspectratio]{figures/sn_void_summary_figure.png}
\caption{\textbf{Figure 1.} Four-panel summary of the supernova-void environment signal. (a) Environment coefficient $\gamma$\textsubscript{env} vs. redshift, showing the localised signal at z < 0.04 with p < 0.01 significance. (b) Leave-one-survey-out stability test confirming no single survey drives the detection. (c) Key results across three void catalogues (REVOLVER, VIDE, VoidFinder) and robustness checks including 12/\allowbreak{}12 bootstrap realisations with p < 0.05. (d) Physical interpretation: SNe in cosmic voids appear \~{}0.01 mag brighter than those in overdense regions, consistent with the MTDF stress-field coherence scale prediction.}
\label{fig:fig-sn-void-summary}
\end{figure}

\begin{table}[H]
\centering
\footnotesize
\caption{\textbf{Table 9.} Summary of multi-probe evidence. Each row represents an independent observable with different physics, different systematics, and different survey selection.}
\begin{tabularx}{\linewidth}{l>{\RaggedRight\arraybackslash}X>{\RaggedRight\arraybackslash}X>{\RaggedRight\arraybackslash}X>{\RaggedRight\arraybackslash}X>{\RaggedRight\arraybackslash}X}
\toprule
\textbf{Channel} & \textbf{Observable} & \textbf{Signal strength} & \textbf{Controls passed} & \textbf{Outcome} & \textbf{Role} \tabularnewline
\midrule
SNe & Distance modulus residual & \(\gamma_\mathrm{env} = +0.013 \pm 0.004\) & LOSO, permutation, bootstrap, survey FE, SALT2 & Positive (\(p = 0.003\)) & Direct evidence \tabularnewline
SNe cutoff & Piecewise \(z = 0.04\) & \(\Delta\chi^2 = 36\text{--}39\) & \(p < 0.001\), consistent across 3 void finders & Sharp transition confirmed & Direct evidence \tabularnewline
Time delays & Time-delay distance & \(\chi^2/\mathrm{dof}\): 0.68 vs 3.72 & 7 independent lenses, geometric observable & Supports MTDF low-\(z\) scale factor & Geometric support \tabularnewline
CF4 vpec & Peculiar velocity residual & 5--9\(\sigma\) after 2M++ & A, B, C, D (2M++), E (Malmquist), F (LOSO), N-body (9.0\(\sigma\)) & Sign mismatch with \(\Lambda\)CDM & Direct evidence \tabularnewline
2MTF TF & TF distance residual & 4--15\(\sigma\) beyond density & A, B, C, N-body (4.2\(\sigma\) tension) & Achromatic, J/\allowbreak{}H/\allowbreak{}K consistent & Direct evidence \tabularnewline
ISW & CMB temperature profile & 2.4--3.7$\times$ template gain & BOSS + DESIVAST voids & Methodological (sub-threshold) & Methodology \tabularnewline
S\textsubscript{8} split & GGL amplitude & 1.71\(\sigma\) & Directional consistency & Suggestive (underpowered) & Suggestive \tabularnewline
\bottomrule
\end{tabularx}
\end{table}

The direct-evidence channels are the supernova environment signal, the \(z = 0.04\) piecewise cutoff, CF4 peculiar velocities, and 2MTF Tully-Fisher residuals. Time delays provide independent geometric support for the MTDF low-redshift interpretation but do not test the void-geometry coupling directly. ISW and S\textsubscript{8} establish methodology and directional consistency but remain below detection threshold. No two of the direct-evidence channels share calibration systems, host-galaxy populations, or distance-indicator physics. The convergence of these independent observables toward a common environment-linked transition near \(z \approx 0.04\) is the paper's primary result.

\begin{mdframed}[linecolor=accent,linewidth=0.5pt,backgroundcolor=callout-bg,leftline=true,rightline=true,topline=true,bottomline=true,innerleftmargin=8pt,innerrightmargin=8pt,innertopmargin=6pt,innerbottommargin=6pt]
\textbf{Key observation.} The same coherence scale appears across independent probes. In the supernova channel, the piecewise cutoff at \(z = 0.04\) yields \(\Delta\chi^2 = 36\text{--}39\). In the CF4 channel, the piecewise split shows 14.7\(\sigma\) below \(z = 0.04\) and 1.9\(\sigma\) above, with the concentration surviving density removal at 9--13\(\sigma\) and confirmed independently in Tully-Fisher, Fundamental Plane, supernova, and calibrator distance methods. The 2MTF sample sits entirely below \(z = 0.04\). All four direct-evidence channels converge on the same low-redshift transition regime.
\end{mdframed}

\section{\(\Lambda\)CDM mock comparison}
\label{sec:sec7}
The most important objection to the CF4 and 2MTF results is that the signal might arise from standard gravitational dynamics within \(\Lambda\)CDM. Peculiar velocities are driven by the density field; Tully-Fisher residuals have known environmental dependencies. Three levels of mock comparison address this objection.

\subsection{Linear-theory reconstruction}
Mock 1 reconstructs \(\Lambda\)CDM-predicted velocities from the measured CF4 density field using linear perturbation theory (\(v \propto f H \delta r\)), adding realistic Gaussian noise matched to the observed velocity dispersion. Over 100 realisations:

\begin{table}[H]
\centering
\footnotesize
\caption{\textbf{Table 10.} Linear-theory mock vs observed CF4 signal.}
\begin{tabularx}{\linewidth}{l>{\RaggedRight\arraybackslash}X>{\RaggedRight\arraybackslash}X>{\RaggedRight\arraybackslash}X}
\toprule
\textbf{Regime} & \textbf{Observed \(\gamma_v\)} & \textbf{Mock \(\gamma_v\)} & \textbf{Tension} \tabularnewline
\midrule
Full sample & \(-53.6\) & \(+8.7 \pm 3.1\) & \(19.9\sigma\) \tabularnewline
\(z < 0.04\) & \(-52.9\) & \(-15.7 \pm 4.2\) & \(8.8\sigma\) \tabularnewline
\bottomrule
\end{tabularx}
\end{table}

The key discrepancy is not only the magnitude but the direction. \(\Lambda\)CDM linear theory predicts a weak positive \(\gamma_v\) for the full sample, while the observed slope is strongly negative. This indicates a qualitative sign reversal rather than a simple underestimation of amplitude.

\subsection{Full N-body mock (MDPL2)}
To close the non-linear escape hatch, the identical analysis pipeline is applied to halo catalogues from the MDPL2 N-body simulation (3840\textsuperscript{3} particles, Planck cosmology). For CF4, 100 randomly placed observers each select approximately 480,000 halos within a CF4-like volume; voids are identified using the same algorithm and \(\gamma_v\) is computed identically to the data analysis. For 2MTF, the same mock halos are used with TF-like distance errors applied.

\begin{table}[H]
\centering
\footnotesize
\caption{\textbf{Table 11.} N-body mock (MDPL2, 100 observers for CF4, 20 for 2MTF) vs observed signal.}
\begin{tabularx}{\linewidth}{l>{\RaggedRight\arraybackslash}X>{\RaggedRight\arraybackslash}X>{\RaggedRight\arraybackslash}X>{\RaggedRight\arraybackslash}X}
\toprule
\textbf{Channel} & \textbf{Observed} & \textbf{N-body mean} & \textbf{Tension} & \textbf{Sign} \tabularnewline
\midrule
CF4 \(\gamma_v\) (full) & \(-53.6\) & \(+1.4 \pm 6.3\) & \(9.0\sigma\) & Mismatch \tabularnewline
2MTF \(\gamma_\mathrm{TF}\) (K-band) & \(0.058\) & \(+0.002 \pm 0.013\) & \(4.2\sigma\) & 30$\times$ too small \tabularnewline
\bottomrule
\end{tabularx}
\end{table}

Non-linear effects move the N-body mean from \(+8.7\) (linear theory) toward zero, not toward the observed \(-53.6\). Full non-linear \(\Lambda\)CDM dynamics slightly wash out the linear-theory signal rather than amplifying it toward the data. With 100 observer positions, the CF4 tension stands at \(9.0\sigma\). This closes the argument that non-linear corrections could rescue the \(\Lambda\)CDM prediction.

\subsection{Null baseline}
A pure-noise mock (100 realisations of Gaussian random velocities with matched dispersion) yields \(\gamma_v = 0.05 \pm 2.9\), confirming that the null distribution is centred at zero with the expected scatter. The observed signal lies 18.4\(\sigma\) from the noise floor.

\begin{mdframed}[linecolor=accent,linewidth=0.5pt,backgroundcolor=callout-bg,leftline=true,rightline=true,topline=true,bottomline=true,innerleftmargin=8pt,innerrightmargin=8pt,innertopmargin=6pt,innerbottommargin=6pt]
\textbf{Mock and reconstruction summary.} The CF4 signal survives seven layers of control: density substitution, partial correlation (8--17\(\sigma\)), \(\Lambda\)CDM chain-rule excess (150--610$\times$), 2M++ velocity-field subtraction (5--9\(\sigma\) retained), Malmquist bias test (signal in all distance bins and both quality halves), leave-one-method-out (cross-method stability across TF, FP, SNIa, and calibrators), and MDPL2 N-body mock (9.0\(\sigma\) tension with 100 observers, sign mismatch). The 2MTF signal independently survives density controls (4--15\(\sigma\)) and N-body comparison (4.2\(\sigma\) tension). Neither linear-theory, full non-linear \(\Lambda\)CDM dynamics, the best available reconstructed velocity field, distance-dependent selection effects, nor any single distance method can account for the observed void-geometric coupling.
\end{mdframed}

\section{Interpretation within MTDF}
\label{sec:sec8}
The first half of this paper is purely observational: multiple probes converge on an environment-linked transition at a single low-redshift scale. This section considers why MTDF provides a natural interpretation of this pattern.

MTDF predicts a finite coherence length \(\beta = 22.7\) Mpc through the structure of its stress tensor field. This coherence length sets the scale below which environment-dependent effects are active and above which they average out. The predicted transition redshift of \(z \approx 0.04\) is not a free parameter fitted to the multi-probe data; it is derived from \(\beta\), which was empirically constrained in MTDF\_\allowbreak{}01 from the MTDF void quantization ansatz \(R_n = \beta(1+\sqrt{n})\) matched to observed void radii, not fitted to Pantheon+, DESI, or any of the multi-probe observables presented here.

The specific features of the multi-probe evidence that align with MTDF predictions include:

(i) \emph{Sharp transition.} The supernova piecewise split and the CF4 redshift concentration both show a transition that is abrupt rather than gradual. A gradual transition would be consistent with slowly varying astrophysical systematics; a sharp transition at a specific scale points to a physical threshold.

(ii) \emph{Void geometry beyond density.} The partial-correlation controls demonstrate that \(d_\mathrm{signed}\) carries information about observables that raw local density does not. In MTDF, void boundaries correspond to regions where the stress tensor field transitions between regimes, producing geometric signatures that are not captured by a scalar density field.

(iii) \emph{Hubble tension resolution.} MTDF's scale factor \(S_0 = 1.084\) naturally maps the Planck \(H_0\) to the locally measured value, resolving the time-delay tension with \(\chi^2/\mathrm{dof} = 0.68\). This is a parameter derived from the field theory, not fitted to the time-delay data.

(iv) \emph{Consistency across observables.} Supernova lightcurve distances, strong-lens time-delay distances, peculiar velocities, and Tully-Fisher distances all respond to environment in a manner consistent with the MTDF coherence scale. These observables probe different physical processes, making a common systematic explanation difficult to construct.

It should be emphasised that this paper establishes convergent evidence for an environment-dependent transition, with MTDF as the coherent interpretive framework. It does not claim that every aspect of MTDF is validated by these results. The photon-coupling sector (\(\kappa = 1.10 \times 10^{-3}\)), the CMB modifications, and the deep-space propagation predictions remain subjects of their respective companion papers and future observational programmes.

\section{Limitations and caveats}
\label{sec:sec9}
Three caveats of the current analysis should be noted.

\textbf{Void classification gap.} None of the Phase 3 analyses (CF4, 2MTF) contain void-classified galaxies: all objects lie in wall environments as defined by DESIVAST. The signals are gradients within the wall population measured by the continuous \(d_\mathrm{signed}\) metric. This is physically meaningful because the MTDF stress tensor field varies continuously with distance to the void boundary; the binary void/\allowbreak{}wall label is an artefact of the classification algorithm, not a property of the underlying field. The partial-correlation controls (Tables 6 and 8) demonstrate that these wall-internal gradients carry void-geometric information well beyond raw density, so the impact of this gap on the main conclusions is low. The continuous signed-distance statistic is the physically relevant quantity for these channels. Nevertheless, the traditional binary void-versus-wall split remains unperformed.

\textbf{Signed-distance and heliocentric-distance correlation.} The environment metric \(d_\mathrm{signed}\) correlates with heliocentric distance (\(r \approx -0.63\) for CF4), reflecting the redshift-dependent completeness of the DESIVAST void catalogues rather than a physical coupling. Dedicated tests (Section 5.3.3, Control E) show that the void-geometric signal persists across distance bins, in both quality halves, and in every distance-method family, indicating that this geometric correlation does not explain the observed environment coupling.

\textbf{Data-limited supporting channels.} The ISW and S\textsubscript{8} channels are suggestive but below detection threshold with current void catalogue sizes and weak-lensing survey areas. They are included to establish methodology and directional consistency, not as independent confirmations. Next-generation survey volumes (DESI Y5, Euclid, LSST/\allowbreak{}Rubin) will provide the 10--15 times sensitivity gain needed to push these channels from suggestive to decisive.

\section{Conclusion}
\label{sec:sec10}
This paper reports convergent evidence from multiple independent low-redshift observables for a common environment-dependent transition near a single coherence scale.

The supernova environment signal (\(\gamma_\mathrm{env} = +0.013 \pm 0.004\), \(p = 0.003\)) exhibits a sharp piecewise cutoff at \(z = 0.04\) (\(\Delta\chi^2 = 36\text{--}39\)). Strong-lens time-delay distances prefer an MTDF-predicted \(H_0 \approx 73.1\) over Planck \(\Lambda\)CDM by a factor of 5.5 in \(\chi^2/\mathrm{dof}\). Cosmicflows-4 peculiar velocities show a void-geometric signal that survives a seven-layer control chain, including partial correlation beyond density (8--17\(\sigma\)), 2M++ velocity-field subtraction (5--9\(\sigma\) retained), Malmquist bias testing across distance bins and quality halves, and leave-one-method-out stability across TF, FP, SNIa, and calibrator subsamples. The signal is concentrated below \(z = 0.04\) and exhibits a qualitative sign mismatch against full N-body \(\Lambda\)CDM expectations (9.0\(\sigma\) tension with MDPL2 across 100 observers). 2MTF Tully-Fisher distances independently confirm a void-geometric signal at 4--15\(\sigma\) beyond density, achromatic across three infrared bands, with 4.2\(\sigma\) N-body tension.

The convergence of these signals, arising from different physical processes, different survey instruments, different galaxy populations, and different distance-indicator calibrations, at a single environment-linked redshift scale, constitutes evidence that is difficult to attribute to any one dataset or method. The Mesche Tensor Dynamics Framework provides a coherent interpretation through its coherence length \(\beta = 22.7\) Mpc, which was established independently of the multi-probe data presented here.

Ongoing extensions include joint multi-probe fitting with a shared coherence-scale parameter and morphological type splits for 2MTF (requiring external cross-matching, as the 2MTF catalogue does not include morphology data). Next-generation survey volumes (DESI Y5, Euclid, LSST/\allowbreak{}Rubin) will provide the statistical power to push the supporting ISW and S\textsubscript{8} channels from suggestive to decisive.

\section{Appendix A. Data availability and reproducibility}
\label{sec:secA}
All analysis scripts and machine-readable results are provided in the MTDF validation repository. The repository is organised into phase folders; the material for this paper spans two phases corresponding to the two primary observational channels.

\subsection{A.1 Analysis scripts}
Nine Python scripts reproduce every quantitative result reported in this paper. All scripts are located in \texttt{validation/\allowbreak{}code/\allowbreak{}paper8/\allowbreak{}}:

\begin{table}[H]
\centering
\footnotesize
\caption{\textbf{Table A1.} Analysis scripts for MTDF\_\allowbreak{}08.}
\begin{tabularx}{\linewidth}{l>{\RaggedRight\arraybackslash}X>{\RaggedRight\arraybackslash}X}
\toprule
\textbf{Script} & \textbf{Section} & \textbf{Description} \tabularnewline
\midrule
task3a\_\allowbreak{}cosmicflows4\_\allowbreak{}vpec.py & 5.3.1--5.3.2 & CF4 baseline: gamma\_\allowbreak{}v fits, piecewise z-split, permutation tests (3 void catalogues) \tabularnewline
task3a\_\allowbreak{}control\_\allowbreak{}tests.py & 5.3.3 (A--C) & CF4 controls: density substitution, partial correlation, LCDM chain rule \tabularnewline
task3a\_\allowbreak{}2mpp\_\allowbreak{}reconstruction.py & 5.3.3 (D) & CF4 control D: 2M++ velocity-field subtraction (Carrick et al. 2015) \tabularnewline
task3a\_\allowbreak{}robustness.py & 5.3.3 (E--F) & CF4 controls E-F: leave-one-method-out (LOSO) and Malmquist bias diagnostics \tabularnewline
task3a\_\allowbreak{}lcdm\_\allowbreak{}mock.py & 7.1, 7.3 & LCDM linear-theory mock (100 realisations) and pure-noise baseline \tabularnewline
task3a\_\allowbreak{}nbody\_\allowbreak{}mock.py & 7.2 & MDPL2 N-body mock for CF4 (100 observers, halos from CosmoSim TAP) \tabularnewline
task3b\_\allowbreak{}2mtf\_\allowbreak{}tully\_\allowbreak{}fisher.py & 5.4.1 & 2MTF baseline: gamma\_\allowbreak{}TF fits across J/\allowbreak{}H/\allowbreak{}K bands and 3 void catalogues \tabularnewline
task3b\_\allowbreak{}control\_\allowbreak{}tests.py & 5.4.2 & 2MTF controls: density substitution, partial correlation, LCDM chain rule \tabularnewline
task3b\_\allowbreak{}nbody\_\allowbreak{}mock.py & 7.2 & MDPL2 N-body mock for 2MTF (50 observers) \tabularnewline
\bottomrule
\end{tabularx}
\end{table}

\subsection{A.2 Machine-readable results}
All numerical results are stored as JSON files with full metadata. Phase 7 covers the CF4 peculiar velocity channel; Phase 8 covers the 2MTF Tully-Fisher channel.

\begin{table}[H]
\centering
\footnotesize
\caption{\textbf{Table A2.} Result files (JSON) with SHA256 integrity hashes.}
\begin{tabularx}{\linewidth}{l>{\RaggedRight\arraybackslash}X>{\RaggedRight\arraybackslash}X}
\toprule
\textbf{Phase} & \textbf{File} & \textbf{Contents} \tabularnewline
\midrule
7 (CF4) & task3a\_\allowbreak{}cf4\_\allowbreak{}results.json & Baseline gamma\_\allowbreak{}v per void catalogue, piecewise z-split, Delta-chi2 \tabularnewline
7 (CF4) & task3a\_\allowbreak{}control\_\allowbreak{}results.json & Controls A-C: density, partial correlation, chain rule \tabularnewline
7 (CF4) & task3a\_\allowbreak{}2mpp\_\allowbreak{}results.json & Control D: gamma\_\allowbreak{}v after 2M++ subtraction (Tables 6b) \tabularnewline
7 (CF4) & task3a\_\allowbreak{}robustness\_\allowbreak{}results.json & Controls E-F: LOSO stability (Table 6c), Malmquist diagnostics \tabularnewline
7 (CF4) & task3a\_\allowbreak{}lcdm\_\allowbreak{}mock\_\allowbreak{}results.json & Linear mock gamma\_\allowbreak{}v distribution (Table 10), noise baseline \tabularnewline
7 (CF4) & task3a\_\allowbreak{}nbody\_\allowbreak{}mock\_\allowbreak{}results.json & MDPL2 N-body mock gamma\_\allowbreak{}v (Table 11, 100 observers) \tabularnewline
8 (2MTF) & task3b\_\allowbreak{}2mtf\_\allowbreak{}results.json & Baseline gamma\_\allowbreak{}TF per band and void catalogue (Table 7) \tabularnewline
8 (2MTF) & task3b\_\allowbreak{}control\_\allowbreak{}results.json & Controls A-C (Table 8) \tabularnewline
8 (2MTF) & task3b\_\allowbreak{}nbody\_\allowbreak{}mock\_\allowbreak{}results.json & MDPL2 N-body mock gamma\_\allowbreak{}TF (Table 11, 50 observers) \tabularnewline
\bottomrule
\end{tabularx}
\end{table}

\subsection{A.3 Diagnostic plots}
32 diagnostic plots for the CF4 channel are in validation/\allowbreak{}output/\allowbreak{}phase7\_\allowbreak{}cf4\_\allowbreak{}vpec/\allowbreak{}. 11 diagnostic plots for the 2MTF channel are in validation/\allowbreak{}output/\allowbreak{}phase8\_\allowbreak{}2mtf\_\allowbreak{}tf/\allowbreak{}. Each phase folder includes a \texttt{manifest.json} with SHA256 hashes for integrity verification.

\subsection{A.4 External data sources}
\begin{table}[H]
\centering
\footnotesize
\caption{\textbf{Table A3.} External datasets required for reproduction.}
\begin{tabularx}{\linewidth}{l>{\RaggedRight\arraybackslash}X>{\RaggedRight\arraybackslash}X}
\toprule
\textbf{Dataset} & \textbf{Source} & \textbf{Access} \tabularnewline
\midrule
Cosmicflows-4 groups & Tully et al. (2023) & \href{https://edd.ifa.hawaii.edu/CF4groupsHD/}{Extragalactic Distance Database} \tabularnewline
2MTF catalogue & Hong et al. (2019) & \href{https://vizier.cds.unistra.fr/viz-bin/VizieR-3?-source=J/MNRAS/487/2061}{VizieR J/\allowbreak{}MNRAS/\allowbreak{}487/\allowbreak{}2061} \tabularnewline
DESIVAST void catalogues & DESI Collaboration & \href{https://data.desi.lbl.gov/public/papers/c3/void\_catalog\_dr1/v1/}{DESI data release} (VoidFinder, REVOLVER, VIDE) \tabularnewline
2M++ velocity/\allowbreak{}density fields & Carrick et al. (2015) & \href{https://cosmicflows.iap.fr/2MPP/}{cosmicflows.iap.fr/\allowbreak{}2MPP} \tabularnewline
MDPL2 N-body halos & Klypin et al. (2016) & \href{https://www.cosmosim.org/}{CosmoSim TAP API} (auto-downloaded by script) \tabularnewline
\bottomrule
\end{tabularx}
\end{table}

\section{References}
\label{sec:sec11}
Brout, D. et al. (2022). The Pantheon+ Analysis: Cosmological Constraints. \emph{ApJ}, 938, 110. \texttt{doi:10.3847/\allowbreak{}1538-4357/\allowbreak{}ac8e04}

Carrick, J., Turnbull, S. J., Lavaux, G., \& Hudson, M. J. (2015). Cosmological parameters from the comparison of peculiar velocities with predictions from the 2M++ density field. \emph{MNRAS}, 450, 317. \texttt{doi:10.1093/\allowbreak{}mnras/\allowbreak{}stv547}

DESI Collaboration (2024). DESI 2024 VI: Cosmological Constraints from BAO measurements. \emph{arXiv:2404.03002}

Hong, T. et al. (2019). 2MTF VII. 2MASS Tully-Fisher survey final data release. \emph{MNRAS}, 487, 2061. \texttt{doi:10.1093/\allowbreak{}mnras/\allowbreak{}stz1413}

Klypin, A. et al. (2016). MultiDark simulations: the story of dark matter halo concentrations and density profiles. \emph{MNRAS}, 457, 4340. \texttt{doi:10.1093/\allowbreak{}mnras/\allowbreak{}stw248}

Mesche, I. (2026a). Mesche's Tensor Dynamics Framework (MTDF): A Dynamic Field Theory of Cosmic Structure. MTDF\_\allowbreak{}01.

Mesche, I. (2026b). MTDF Companion Paper Part I: Environmental and Local Phenomenology. MTDF\_\allowbreak{}02.

Mesche, I. (2026c). MTDF Companion Paper Part II: Gravity Sector and Lensing Validation. MTDF\_\allowbreak{}03.

Mesche, I. (2026d). MTDF Companion Paper Part III: Photon Coupling, Redshift, Early Universe. MTDF\_\allowbreak{}04.

Mesche, I. (2026e). MTDF Independent GPU Validation: Phases 1--6. MTDF\_\allowbreak{}07.

Millon, M. et al. (2020). TDCOSMO I. An exploration of systematic uncertainties in the inference of H\textsubscript{0} from time-delay cosmography. \emph{A\&A}, 639, A101. \texttt{doi:10.1051/\allowbreak{}0004-6361/\allowbreak{}201937351}

Planck Collaboration (2020). Planck 2018 results. VI. Cosmological parameters. \emph{A\&A}, 641, A6. \texttt{doi:10.1051/\allowbreak{}0004-6361/\allowbreak{}201833910}

Tully, R. B. et al. (2023). Cosmicflows-4. \emph{ApJ}, 944, 94. \texttt{doi:10.3847/\allowbreak{}1538-4357/\allowbreak{}ac94d8}

Wong, K. C. et al. (2020). H0LiCOW XIII. A 2.4\% measurement of H\textsubscript{0} from lensed quasars. \emph{MNRAS}, 498, 1420. \texttt{doi:10.1093/\allowbreak{}mnras/\allowbreak{}stz3094}


\end{document}
